\documentclass[11pt]{article}
\usepackage[a4paper,left=28mm,right=32mm,top=25mm,bottom=25mm,marginparwidth=18mm,marginparsep=4mm]{geometry}
\usepackage{fontspec,amsmath,amssymb,xcolor,needspace,array,booktabs}
\usepackage{manyfoot,perpage}
\usepackage[unicode,hidelinks]{hyperref}
\usepackage{polyglossia}
\setmainlanguage{english}
\setotherlanguages{greek,german,latin}
\setmainfont{Linux Libertine O}[Ligatures=TeX]
\setlength{\emergencystretch}{1.5em}
\newfontfamily\greekfont{FreeSerif}[Script=Greek]
\newfontfamily\SignFont{FreeSerif}
\DeclareNewFootnote{T}[alph]
\MakePerPage{footnote}\MakePerPage{footnoteT}
\makeatletter
\newcommand{\TNoteLetter}[1]{\ifnum#1>26 \@alph{\numexpr(#1-14)/26\relax}\@alph{\numexpr#1-26*((#1-14)/26)\relax}\else\@alph{#1}\fi}
\renewcommand{\thefootnoteT}{\TNoteLetter{\value{footnoteT}}}
\makeatother
\hypersetup{pdftitle={H. Suter, the German matter of his edition of al-Khwarizmi's astronomical tables: an English translation},pdfauthor={H. Suter}}
\setlength{\parindent}{0pt}\setlength{\parskip}{0.55em}
\tracinglostchars=2
\newcommand{\Pg}[1]{\marginpar{\footnotesize\color{gray}#1}}
\newcommand{\Entry}[1]{\par{\parskip=0.2em\relax\noindent\hangindent=1.5em\hangafter=1 #1\par}}
\newcommand{\PartHead}[2]{\clearpage\pdfbookmark[0]{#1}{part-#2}{\centering\LARGE #1\par}\bigskip}
\newcommand{\UnitHead}[3]{\par\needspace{8\baselineskip}\bigskip\pdfbookmark[1]{#2}{unit-#3}\hypertarget{EN-#3}{}%
 {\large\bfseries\interlinepenalty=10000 #1\par}\nobreak\medskip}
\begin{document}

\thispagestyle{empty}
\begin{center}\Large H. SUTER\par\vspace{4mm}\LARGE The Astronomical Tables of\par al-Khwārizmī\par\vspace{2mm}\large the foreword, the introduction, the commentary, the additions, the glossary and the index of his edition (Copenhagen, 1914)\par\vspace{5mm}\normalsize translated from the German\end{center}\vspace{8mm}
\noindent Heinrich Suter's \emph{Die astronomischen Tafeln des Muḥammed ibn Mūsā al-Khwārizmī} (Copenhagen, 1914), on the preliminary work of A. Bjørnbo and R. Besthorn, edits the astronomical tables of al-Khwārizmī in the recension of Maslama ibn Aḥmad al-Majrīṭī and the Latin translation of Adelard of Bath. The Latin text and tables were translated into English by O. Neugebauer (\emph{The Astronomical Tables of al-Khwārizmī}, Copenhagen, 1962), with his own commentary; Suter's German, his foreword, introduction, commentary, additions, glossary and index, has not been translated before. This volume translates it from the line-by-line transcription of the edition in this repository. Every sentence is translated, in Suter's order; nothing is summarised or left out. The words of the Latin text that Suter quotes are kept, followed by an English translation in parentheses; Greek stands as printed. Arabic names and words are given in a modern transliteration (j for Suter's dj, q for his ḳ), except in the glossary and the index, which keep Suter's transcription. Suter's footnotes are numbered; the notes of the translator are lettered: they record a misprint set right, with the reading of the print, a slip in Suter's reckoning, checked against his numbers, and the figures of the print, which are not reproduced. The pages of Suter's book are given in the margin.
\par\smallskip\noindent The translation is complete. It is AI-translated (Claude), 2026; no person has proofread it. Suter's work is in the public domain; the translation is dedicated to the public domain (CC0 1.0).
\par\clearpage
\PartHead{The Astronomical Tables of Muḥammad ibn Mūsā al-Khwārizmī}{F}\Pg{I}
\par\begin{center}in the recension of Maslama ibn Aḥmad al-Majrīṭī and the Latin translation of Adelard of Bath\end{center}
\par\begin{center}on the preliminary work of A. Bjørnbo{\SignFont\upshape †} and R. Besthorn in Copenhagen\end{center}
\par\begin{center}edited with a commentary by H. Suter in Zürich\end{center}
\par\begin{center}D. Kgl. Danske Vidensk. Selsk. Skrifter, 7. Række, historisk og filosofisk Afd. III. 1 (Memoirs of the Royal Danish Academy of Sciences and Letters, seventh series, historical and philosophical section, III. 1)\end{center}
\par\begin{center}København (Copenhagen); main agent: Andr. Fred. Høst \& Søn, royal court booksellers; printed by Bianco Luno; 1914\end{center}
\UnitHead{Contents}{Contents}{F01-6}\Pg{III}
\begin{center}\small\begin{tabular}{lr}
 & \emph{Page} \\
Foreword & V \\
Introduction & VII–XXV \\
Text of the Latin translation & 1–31 \\
Commentary & 32–106 \\
Tables & 109–231 \\
Additions and corrections & 233–237 \\
Glossary & 241–248 \\
Index of names and subjects & 249–255 \\
\end{tabular}\end{center}
\UnitHead{Chapters of the commentary}{Chapters of the commentary}{F01-8}
\par\begin{center}(corresponding to those of the Latin text)\end{center}
\begin{center}\small\begin{tabular}{lrr}
 & \emph{Page} & \emph{Table} \\
Introduction of the text & 32 &  \\
1. Description of the Arab year & 34 &  \\
2. Description of the Roman year & 34 & 1 \\
3. On the headings of the columns of the tables & 36 &  \\
4. On the beginning of the Arab months & 36 & 2–2$^{b}$ \\
5. On finding the Roman year & 38 & 3–3$^{a}$ \\
6. On the division of the sphere & 41 &  \\
7. On the mean motion of the planets & 41 & 4–20 \\
8. On finding the (true) place of the sun & 43 & 21–26 \\
9. On finding the (true) place of the moon & 45 & 21–26 \\
10. Determination of the (true) place of the three upper planets & 46 & 27–44 \\
11. Determination of the (true) place of the two lower planets & 54 & 45–56 \\
12. Determination of the longitude of the nodes of the moon & 55 & 19–20 \\
13. On the stations and the forward and backward motion of the planets & 55 & 27–56 \\
14. On the time of the direct and the retrograde motion & 56 & 27–56 \\
15. On the declination of the sun & 58 & 21–26 \\
16. On the latitude of the moon & 58 & 19–26 \\
17. On the latitude of the planets & 60 & 27–56 \\
18. On the apogee and the perigee of the planets & 64 &  \\
19. On the nodes of the planets & 64 &  \\
20 and 21. On the daily motion of the moon and of the sun & 65 &  \\
22. On the appearance of the moon on the evening of the 29th day (of the month) & 66 & 57$^{a}$ \\
\end{tabular}\end{center}
\Pg{IV}\begin{center}\small\begin{tabular}{lrr}
 & \emph{Page} & \emph{Table} \\
23. On finding the sine from the arc and conversely & 68 & 58–58$^{a}$ \\
24. On finding the latitude of any place & 71 &  \\
25. The ascensions of the signs of the zodiac in the right sphere & 72 & 59–59$^{b}$ \\
26 and 26$^{a}$. With how many degrees (of the equator) each sign of the zodiac rises at any place & 73 & 59–59$^{b}$ \\
27. Determination of the size of the seasonal hour of any day for any place & 76 &  \\
28–28$^{b}$. How the two shadows are determined from the altitude of the sun, and conversely & 76 & 60 \\
29. On \emph{elbuht} and how to find it & 78 & 61–66 \\
30 and 30$^{a}$. On the apparent diameters of the sun and the moon & 78 & 61–66 \\
The tables of the equation of days (the equation of time) & 80 & 67–68 \\
31. On the conjunction and the opposition of the sun and the moon & 81 & 69–72 \\
33. How the eclipse of the moon is determined & 86 & 73–76 \\
34. On the parallaxes of the moon in longitude and latitude & 91 & 77–77$^{a}$ \\
35. On the eclipse of the sun & 94 & 78 \\
32 and 36. On the division of the twelve astrological houses & 96 & 79–90 \\
37. On the places of the projection of the rays & 98 & 91–114 \\
On the conversion of the years of nativity & 102 & 115 \\
On the domiciles, exaltations, decans, etc. of the planets & 103 & 116 \\
\end{tabular}\end{center}
\UnitHead{Foreword}{Foreword}{F02-0}\Pg{V}
\par The edition of these astronomical tables was begun about five years ago by the eminent Danish scholar A. Bjørnbo, who unhappily died too early. After his death in October 1911 his manuscript was handed over to me by the Royal Danish Academy of Sciences, when I had declared myself ready to finish the work that A. Bjørnbo had begun. This manuscript contained the text of the work with the readings that had resulted from the collation of three manuscripts (\emph{Oxford}, \emph{Chartres} and \emph{Paris}), likewise the tables with the readings of these manuscripts, and also the rendering into Latin of most of the Arabic terms that occur in them by the meritorious Danish Arabist Dr. R. Besthorn of Copenhagen. Shortly before his death A. Bjørnbo had learned that a fourth manuscript of this translation of Adelard's existed in the National Library at \emph{Madrid}. So my task was, first, to bring this manuscript too into the collation; secondly, I had to have the manuscript of \emph{Chartres} sent once more for revision, since on several leaves of his manuscript A. Bjørnbo had noted that a further collation with it was needed; thirdly, it was necessary to work out a commentary on the short and often hardly intelligible text, and to add to the work an introduction, a glossary and an index of names and subjects. At the end of July 1913 I submitted my work to the Royal Danish Academy of Sciences for examination; in December of the same year the Academy resolved to publish it in its \emph{Skrifter}, for which I here express my sincere thanks.\par
\par \Pg{VI}I am likewise indebted to Professors Dr. D. Andersen in Copenhagen and Dr. G. Billeter in Zürich for the kind help they gave me with the proofs. Unhappily my thanks to a third scholar no longer find him among the living: Mr. C. F. Pechüle, observer at the Royal Observatory in Copenhagen, who died in May of this year, drew my attention to some points in the commentary that called for a somewhat different interpretation and presentation; I have taken his remarks into account in Addition 1).\par
\par The undersigned hopes that the publication of this work may help to advance the knowledge of the older Arabic astronomy, and in particular to make easier the understanding of the astronomical tables of the western Middle Ages, which very probably are based for the most part on this translation of Adelard's.\par
\par Küsnacht, near Zürich, July 1914.\par
\par\noindent\hfill \emph{H. Suter.}\par
\UnitHead{Introduction}{Introduction}{F02-6}\Pg{VII}
\par As early as 1846 the eminent French scholar M. Chasles, in an article in the \emph{Comptes rendus de l’Académie des Sciences}\footnote{Vol. XXIII, p. 848 ff.}, drew attention to the importance of Adelard's translation of the tables of al-Khwārizmī for the history of astronomy. But he did not study the work closely, since he did not intend to edit it, and for this reason most of the conclusions that he drew from his study no longer stand up to criticism today. For Chasles did not yet know at that time that several Arabic authors ascribe to the Spanish-Arab mathematician and astronomer Maslama al-Majrīṭī\footnote{See H. Suter, \emph{Die Mathematiker u. Astronomen der Araber}, etc., in Abhandlungen zur Geschichte der mathem. Wissensch., X (1900), pp. 76–77.} a recension, or better a reworking, of al-Khwārizmī's tables. This reworking consisted chiefly, as the Arabic biographers Ibn al-Qifṭī and Ibn Abī Uṣaybiʿa report after the Qāḍī Ibn Ṣāʿid\footnote{Ibid., pp. 106–107.}, in the conversion of the tables, which had been made for the \emph{era} of Yazdegerd, into tables based on the Muslim \emph{era}; and also, as some passages of the text and some tables prove, in the partial replacement of the meridian of \emph{Arin}, which al-Khwārizmī had taken as his basis, by that of \emph{Córdoba}, where al-Majrīṭī lived.\par
\par \Pg{VIII}From these very passages of the text and tables (see also the commentary) it follows with almost absolute certainty that in this translation of Adelard's\footnote{Also written Adhelhard, Adelhard, Adelard, etc.}\footnoteT{Suter writes Athelhard; the translation uses the usual English form of the name, Adelard.} we have before us the translation not of al-Khwārizmī's original astronomical work but of its reworking by Maslama al-Majrīṭī. As far as I know, C. H. Haskins was the first to point to this fact, in his paper on \emph{Adelard of Bath}\footnote{In \emph{The English Historical Review}, July 1911, p. 493.}; whether A. Bjørnbo recognised it too does not appear from the papers he left, but it is not likely, for otherwise he would surely have drawn attention to it in his paper \emph{Al-Chwarizmi’s trigonometriske Tavler}\footnote{In the \emph{Festskrift til H. G. Zeuthen}, København 1909.}. From this it naturally follows that we must be very cautious with regard to the conclusions that have been drawn, and may yet be drawn, from this work about the dependence of the older Arabic astronomy on the Indian; for it is very difficult, indeed in most cases impossible, to decide what in the present work still comes from al-Khwārizmī and what is to be ascribed to Maslama. But the work does not by any means lose interest on that account: for, first, the main part of the original work will probably still be present in the recension, and secondly, whatever influence this work had on the astronomy of the western Middle Ages is in any case due to this recension by the Spanish-Arab astronomer, since the original work very probably never became known in the Christian West.\par
\par Before we come to the discussion of the manuscripts on which this edition is based, it will be in place to add a few short biographical notes on the Arab scholars who are chiefly concerned here; for further material I refer to my work cited above (p. VII).\par
\par 1. Muḥammad b. Mūsā al-Khwārizmī (also written al-Khuwārizmī) Abū ʿAbdallāh\footnote{\emph{Abū Jaʿfar} is also given as his by-name (\emph{kunya}), but I take \emph{Abū ʿAbdallāh} to be the right one; the former has probably arisen from a confusion of our astronomer with Muḥ. b. Mūsā b. Shākir (d. 873), who lived somewhat later and bore this by-name.} was one of the astronomers in the service of al-Maʾmūn (who reigned as caliph from 813 to 833); he was a Persian by descent; the years of his birth and of his death are not known. At al-Maʾmūn's command he composed an abridgement of an Indian \emph{Siddhānta} (= astronomical work, system), which Ibrāhīm (according to others Muḥ. b. Ibrāhīm) al-Fazārī had translated from the Indian into Arabic and which the Arabs called “the \emph{great Sindhind}”; the abridgement was then called, in contrast to it, “the \emph{small Sindhind}”. He likewise composed astronomical tables (\emph{zīj}) based on this Indian \emph{Siddhānta}, which are said to have appeared in two editions.\footnote{The sources do not quite agree here; according to others it would also be possible that these tables were themselves the abridgement of the great \emph{Sindhind}; the question can no longer be decided.} In them he determined the mean places of the planets for the Persian \emph{era} of Yazdegerd and for the meridian of \emph{Arin} (see the commentary). These tables, then, no longer exist; but their reworking by Maslama al-Majrīṭī does, in the Latin translation of Adelard of Bath. — He further composed a book on \emph{Indian reckoning}, which no longer exists in Arabic either and was edited in a Latin translation (the translator has not been established) by B. Boncompagni\footnote{\emph{Trattati d’aritmetica}, I, Rome, 1857.}; one on \emph{algebra} (\emph{fī l-jabr wa-l-muqābala}), translated into Latin by Gerard of Cremona and Robert of Chester (\emph{Cestrensis}), edited in Arabic with an English translation by F. Rosen\footnote{\emph{Algebra of Mohammed ben Musa, edited and translated} by Frederic Rosen, London, 1831.}; a \emph{Geography} \Pg{IX}(\emph{ṣūrat al-arḍ}), a fairly independent reworking of the Geography of Ptolemy, still extant in Arabic at Strasbourg\footnote{See C. A. Nallino, \emph{Al-Ḫuwārizmī e il suo rifacimento della geografia di Tolomeo}, in the Memorie della Reale Accademia dei Lincei (Classe di scienze morali, storiche e filologiche), Serie 5, Vol. II, P. I, Rome, 1894.}; a \emph{chronicle} (\emph{kitāb al-taʾrīkh}); and three more writings on sundials and astrolabes, none of which exists any longer. The “Indian art of reckoning” and the “Algebra” are the works to which the Christian West chiefly owes its first acquaintance with these disciplines; perhaps — we cannot assert it with certainty — his astronomical tables in Maslama's recension were also the basis of the later astronomical work of the western Arab and the Christian Middle Ages.\par
\par 2. Maslama b. Aḥmad al-Majrīṭī, Abū l-Qāsim, of Córdoba, lived in the time of al-Ḥakam II (961–976) and Hishām II (976–1013) and died in 1007/08. He was the most important mathematician and astronomer of Arab Spain in his time. His chief works are: on \emph{commercial} or \emph{business arithmetic} (\emph{al-muʿāmalāt}), extant neither in Arabic nor in translation; on \emph{the construction and use of the astrolabe}, still extant in Arabic in the Escorial; a translation of the \emph{Planisphaerium} of Ptolemy with a commentary, translated into Latin by Hermannus Dalmata (Hermann of Dalmatia), printed at Basel in 1536 and at Venice in 1558; an \emph{abridgement} of the \emph{tables} of al-Battānī, no longer extant; the \emph{reworking of the tables} of al-Khwārizmī, extant only in the Latin translation of Adelard of Bath, which we here make public.\par
\par We also give here, in a literal translation, the passage that the Spanish-Arab historian Ibn Ṣāʿid (d. 1070 in Toledo)\footnote{See my work cited above (p. VII), p. 106.} has handed down to us about Maslama's reworking of the tables of al-Khwārizmī\footnote{From the edition of his work \emph{Ṭabaqāt al-umam} (= the classes of the nations) by L. Cheikho, S. J., Beirut 1912, p. 69.}: “He occupied himself with the reworking of the tables of al-Khwārizmī; he converted their Persian chronology into the Arabic one and determined the mean places of the planets for the beginning of the \emph{Hijra}; he added to the work further fine tables, in which, however, he followed him (that is, al-Khwārizmī) closely and did not draw attention to the errors in his work; I have pointed out this circumstance in my book ‘on the correction of the motion of the planets and the instruction on the errors of the astronomers’.”\par
\par This passage, then, proves to us that not all the tables in this recension come from al-Khwārizmī; but which are those that Maslama added is hard to decide. We have already drawn attention to this point above, and shall come back to it in the commentary. — By the remark that Maslama did not notice al-Khwārizmī's errors, or did not draw attention to them, he probably means in the first place the neglect of the precession of the equinoctial points.\par
\par \Pg{X}Unfortunately nothing of the Arabic texts either of al-Khwārizmī or of Maslama exists any longer, as far as our present knowledge of the Arabic manuscripts goes. We do not yet give up the hope, however, that something of Maslama's recension may yet be found in Spanish libraries, since it was only three years ago that C. H. Haskins first drew attention to the existence of a manuscript of Adelard's translation in Madrid.\footnote{\emph{English Historical Review}, July 1911, pp. 493–494.}\par
\par As for the manuscripts used by A. Bjørnbo and by me, I cannot give an exact description of them, since only that of \emph{Chartres}, which I had to have sent once more, was before me; of those of \emph{Oxford} and \emph{Madrid} I have seen only photographs of Adelard's translation, and the very incomplete manuscript of the \emph{Mazarine} Library in \emph{Paris} I have never seen. I must therefore confine myself to giving a few notes on these manuscripts from the photographs before me and from the catalogues of the \emph{Bibliothèque publique} of \emph{Chartres} and of the \emph{Bibliothèque Mazarine}\footnote{Both published at Paris in 1890.} in \emph{Paris}.\par
\par The manuscripts of \emph{Chartres}, Bibliothèque publique (C), and of \emph{Oxford}, Bodleian Library (O), are both parchment manuscripts of the twelfth century, according to the authors of the catalogues; the former has the number 214 (173)\footnote{\emph{Catalogue général des mss. des biblioth. publ. de France}, Départements, Tome XI, pp. 109–110.}, the latter Cod. Auct. F. I. 9. They are on the whole well written and easy to read, C somewhat better than O. The written area of the leaf measures 26/18 cm in C and in O (from the photographs) 15/10 cm on average. Both manuscripts are fairly complete, but C contains several larger passages that are missing in O, and the other way round (see the text). Our treatise forms the end of C, from fol. 41$^{r}$ to 102$^{r}$; in O it occupies fol. 99$^{v}$ to 159$^{v}$, but fol. 159 has been moved to the end by mistake; it should stand between 99 and 100, so that the work really ends with fol. 158.\par
\par C and O contain some other treatises as well, of which I shall name only the most important. In C there are also the first half of the treatise on the astrolabe by Hermannus Contractus, and fragments of the Geometry of Gerbert and of the \emph{Liber embadorum} of Savasorda (according to Chasles); in O there are also the \emph{Computus} of Gerlandus and the \emph{Abacus} of Bernelinus (\emph{auctore juniore Bernhelmo!}, by the younger Bernhelm!).\par
\par \Pg{XI}The manuscript of the Bibliothèque Mazarine in \emph{Paris} (M)\footnote{Bjørnbo had chosen the letter M for this manuscript; I did not want to make any more changes in this respect, and have therefore chosen the letter N for the manuscript of the \emph{Biblioteca Nacional} at Madrid.}, with the number 3642 (1258)\footnote{\emph{Catalogue général des mss. des biblioth. publ. de France}, Départements, Tom. III, pp. 151–152.}, is quite incomplete: instead of the 37 chapters of the text it contains only the first 21, which make up the smaller half of the text, and instead of c. 120 tables it has only 9. It too is a parchment manuscript, of the thirteenth century; the written area of the leaf measures 30/23 cm on average. Several hands have worked on it; Bjørnbo expressed this by the designations M$_{1}$, M$_{2}$, etc.\footnote{Thus, for example, the first table of the equation of Saturn (Table 27) appears to be present in three copies by different hands (M$_{1}$, M$_{2}$, M$_{3}$) on fol. 86 and 87, and the second table of the equation of this planet (Table 28) in two copies (M$_{1}$ and M$_{2}$) on the same leaves; the following tables are all missing.} It is very probable that M is an incomplete and very careless copy of C, since even faulty readings mostly agree. Our treatise occupies fol. 82$^{r}$–89$^{v}$ of this manuscript. Besides it, the manuscript contains among other things: Thebit, \emph{de iis quae indigent expositione antequam legatur Almagesti}, incomplete; Jordanus, \emph{de pondere}; the treatise that begins \emph{communis et consuetus} . . .\footnote{See G. Eneström, in the \emph{Biblioth. mathemat.}, 3rd series, vol. 8, p. 135 ff.}; \emph{liber alchorismi arithmeticae qui editus est a magistro Johanne}.\par
\par We must deal somewhat more fully with the manuscript of the Biblioteca Nacional in \emph{Madrid} (N), with the number 10016, which A. Bjørnbo unfortunately could no longer use, since it is of particular historical importance. According to the catalogue, whose particulars came to my knowledge through the efforts of Mr. Miguel Asín, professor of Arabic at the University of Madrid, to whom I here publicly express my most sincere thanks for his kind help, especially also in obtaining photographs of the manuscript, it is a parchment manuscript of the thirteenth century. The written area of the leaf measures, from the photographs, 16/10 cm in the text and 17/12 cm in the tables. The manuscript is very beautifully and evenly written and has few abbreviations; the initials are done in black, red and green, and the figures of the different columns of the tables and their titles also often (not throughout) alternate in black, red and green ink\footnote{This was fatal to the photographs of the tables: the titles and figures written in green and red ink came out so faint in the image that in some places even the best magnifying glass failed.}. The manuscript was originally in the cathedral of \emph{Toledo}. Prof. Asín remarked to me in a letter that it had probably come to Spain from Italy, since it had formed part of the collection of Cardinal Celada\footnoteT{Francesco Saverio de Zelada (1717–1801), whose manuscripts went to the cathedral library of Toledo.}, which had been brought together in Italy. We rather believe that it came into this collection only in Spain; the collection may well have been begun in Italy but perhaps continued in Spain.\par
\par This beautiful manuscript too is unfortunately incomplete: of the 37 chapters of the text it contains only 1–5, 22 and 23, and 31–37, even fewer than M; but of the c. 125 tables it has 110 (only the numbers 4–8, 21, 57$^{b}$, 59, 59$^{a}$, 59$^{b}$, 60, 61, 67, 68, 115 and 116 are missing).\par
\par In the first sentence (the title) of the text in Ms. N, Rodbertus Cestrensis, that is, Robert of Chester, introduces himself with the words \emph{et per rodbertum cestrensem ordine digestus} (and set in order by Robert of Chester), \Pg{XII}which are missing in the other three manuscripts, as the “arranger” of the translation of Adelard of Bath. This arranging hand can indeed be felt in the text in several respects, as we shall see below, but not in the order of the tables: there Ms. N shows a bad disorder that we do not meet in C and O; but we must certainly not lay it at the door of Robert of Chester, rather of later copyists or binders of the manuscript.\par
\par L. C. Karpinski\footnote{\emph{Biblioth. mathem.}, 3rd series, vol. 11 (1911), p. 127.} regards Robert of Chester as identical with Robertus Ketenensis, or Retinensis (Robert of Reading), who was the first to translate the Koran into Latin (1143) and then, as he himself remarks\footnote{In the introduction to his translation of the Koran.}, turned to the study and translation of astronomical and astrological works of the Arabs. A translation by him of the \emph{Algebra} of al-Khwārizmī is also known to exist\footnote{\emph{Biblioth. mathem.}, ibid., pp. 125–131.}. — It is a pity that the result of this arranging hand of Robert of Chester has come down to us once more in disorder and incomplete.\par
\par Our treatise occupies fol. 8$^{r}$–68$^{v}$ of Ms. N. Besides it, this codex contains a series of mostly anonymous texts and tables of astronomical and astrological content, of which only those that immediately follow our treatise (fol. 69–84, that is, to the end) concern us: they are tables on various astronomical subjects, especially on eclipses, composed by the Englishman Roger of Hereford (c. 1150–1220)\footnote{See my paper \emph{Die Mathematik auf den Universitäten des Mittelalters}, p. 21 (Wissenschaftl. Beilage zum Programm der Kantonsschule in Zürich, 1887).} for the latitude and longitude of that place, which are given as 52° and 24° respectively. Thus fol. 71 says, for example: \emph{Sunt autem hae tabulae compositae super Herefordiam a magistro Rogerio Herefordensi, sicut et praemissa tabula de eclipsi solis et regulae praescriptae. Est autem longitudo Herefordiae ab occidente arcus XXIV graduum, latitudo vero LII}\footnote{In another place the latitude is given as 51° 30′.} (these tables were composed for Hereford by Master Roger of Hereford, as were the preceding table of the eclipse of the sun and the rules written before it; the longitude of Hereford from the west is an arc of 24 degrees, and its latitude 52). — On fol. 73 there follows: \emph{Tabula medii motus solis super mediam noctem Herefordiae secundum annos domini} (table of the mean motion of the sun for midnight at Hereford, by the years of the Lord). — These tables may be regarded as supplements to Adelard's translation, for precisely this part of the work, on the eclipses, leaves something to be desired, since the tables of the \emph{ascensiones obliquae} for \emph{Córdoba} are missing; Roger of Hereford has indeed supplied this, but not for Córdoba: for his home town.\par
\par \Pg{XIII}How, then, are these manuscripts related to one another? As I have mentioned above, M is most probably a careless copy of C. As for N, we are of the opinion that Robert of Chester had both C and O, or their archetypes, before him: in the text he follows C almost word for word, with only a few deviations; in the figures of the tables he follows now the reading of C, now that of O, but the former more often. Robert's correcting hand shows itself above all in his having replaced, with very few exceptions, the many Arabic words that C and O leave untranslated or keep beside the translation by Latin ones: thus for \emph{dakaicae} he always has \emph{minuta}, for \emph{eltheniae} \emph{secunda}, for \emph{elwacat} \emph{medius cursus}, for \emph{tadil} \emph{coaequatio} (not \emph{aequatio}), for \emph{elgeib} \emph{sinus}, for \emph{elfelek elmustakim} \emph{circulus aequatorialis} or \emph{aequinoctialis}, etc.\par
\par The relation of C to O, and the relation of both to Adelard's autograph, is the hardest to explain. Bjørnbo, who left no written statement on the matter, was of the opinion, according to a letter to me from Prof. Heiberg in Copenhagen, “that we have only copies of two unfinished working copies of the translator, so that the restitution of an original text would be impossible; rather, in many cases the two texts would have to be printed side by side.” Yet the two texts C and O do not differ from each other so much that this would be necessary; Bjørnbo himself did it in one single chapter only (no. 7), and we have followed him in this. He is right that a restitution of the original text is impossible, and that very probably the text and the tables are incomplete in C and O too, though only slightly. But the deviations are nevertheless such that they cannot come from the mere copying of one manuscript from the other; and so the conjecture has forced itself upon us that the scribe of O or of C — which of the two we cannot decide — had another translation before him besides Adelard's, and took his text now from the one, now from the other, whichever seemed to him the better. It is indeed very probable that Hermannus Dalmata, the contemporary and friend of Robert of Chester, also translated the tables of al-Khwārizmī (as altered by Maslama). In a passage of the translation of Ptolemy's \emph{Planisphaerium} with Maslama's commentary, which was long wrongly ascribed to Rudolf of Bruges, Hermannus says: \emph{Quem locum a Ptolemaeo minus diligenter perspectum cum Albatene miratur et Alchoarismus, quorum hunc quidem opera nostra Latium habet, illius vero commodissima translatio Roberti mei industria, Latinae orationis thesaurum accumulat}\footnote{See Wüstenfeld, \emph{Die Übersetzungen arab. Werke in das Lateinische seit dem 11. Jahrh.}, p. 52 (in Abhandlungen der k. Gesellsch. d. Wissensch. zu Göttingen, histor.-philolog. Kl. XXII, 2, 1877); Steinschneider, \emph{Die europäischen Übersetzungen aus dem Arabischen}, A, p. 34 (in Sitzungsberichte der k. Akademie d. Wissensch. in Wien, philos.-histor. Kl. CXLIX, 1904); \emph{Al-Battānī}, ed. Nallino, I, Praefatio, p. XLIX.} (al-Khwārizmī, together with al-Battānī, wonders at this place, which Ptolemy examined with too little care; of these two, Latium has the former through our labour, while a most convenient translation of the latter, by the diligence of my Robert, adds to the treasure of the Latin tongue). Might the different rendering of the Arabic letter \emph{shīn} in the two manuscripts C and O, by \emph{sh} (English) and by \emph{sc} (Dalmatian-Italian?) respectively, be connected with this? I must leave it to the philologists to pursue this question further and, if possible, to decide it. There is then the further question whether Hermannus Dalmata also brought Adelard's translation into his comparison; as a friend of Robert of Chester, its “arranger”, he will surely have known this translation.\par
\par \Pg{XIV}Since I received the photographs of Ms. N from Madrid (chiefly through my own fault) only when the text and the tables had already been written out ready for the press, I have taken the following course so as not to have to make too many corrections in the footnotes: in the cases where N agrees with C, and these are the majority, I have mostly not mentioned Ms. N at all; where N is missing in the footnotes, then, its reading is the same as that of C. — Since N shows rather many deviations from C and O in the titles of the columns, I have put these together below for the two main groups of tables, so as not to have to give them again in the notes to each table:\par
\par In Tables 9 to 20, which contain the mean motion of the planets and of the head of the dragon of the lunar orbit, N has\par
\begin{center}\small\begin{tabular}{lll}
for \emph{Medialitas} & always & \emph{Medius cursus} \\
for \emph{plani} & always & \emph{residui} \\
for \emph{Heza} & always & \emph{Argumentum} \\
for \emph{partes horarum} & sometimes & \emph{minuta horarum} \\
\end{tabular}\end{center}
\par The names of the months are written in N almost throughout as follows (variants in brackets):\par
\begin{center}\small\begin{tabular}{ll}
\emph{Almuharram} & \emph{Rageb} \\
\emph{Szafar} (\emph{Szaffar, Szaphar}) & \emph{Saaben} (\emph{Szaaben}) \\
\emph{Rabbe}\footnotemark{} (\emph{Rabe}) & \emph{Ramadhan} (\emph{Ramadham}) \\
\emph{Rabbe} (\emph{Rabe}) & \emph{Sewil} \\
\emph{Jumedi} & \emph{Dulkidda} (\emph{Dulkadda}) \\
\emph{Jumedi} & \emph{Dulheigga} (\emph{Dulheiga, Dulhegga}) \\
\end{tabular}\end{center}
\addtocounter{footnote}{-1}\stepcounter{footnote}\footnotetext{Without the distinguishing numbers I and II. [\emph{Translator:} The print calls this note at both \emph{Rabbe} and at both \emph{Jumedi}.]}
\par In Tables 22 to 56 (21 is missing in N), which contain the equations of the sun, of the moon and of the five planets, N has\par
\begin{center}\small\begin{tabular}{lll}
for \emph{Examinatio} & always & \emph{Coaequatio} \\
for \emph{Catar} (\emph{Semita, Pagina, Via}) & always & \emph{Linea} \\
for \emph{Computatio minuenda} & mostly & \emph{Computatio diminuta} \\
for \emph{Computatio augenda} & mostly & \emph{Computatio aucta} \\
for \emph{Sublimatio examinata} & always & \emph{Sublimatio definita} \\
for \emph{Heza} & always & \emph{Argumentum} \\
\end{tabular}\end{center}
\par Ms. N also contains some larger additions and marginal glosses, which I could not well take into the footnotes, since almost all of them need a commentary; I give them here in the introduction, in the order in which they occur in the text and in the tables:\par
\par 1) On fol. 8$^{r}$, after the words \emph{omnia quae hic} (in our text p. 2), which should form the end of the first column, there come four more lines in the same hand which evidently do not belong to our text but probably to the text of fol. 7$^{v}$: \emph{Ab aequinoctio usque ad solstitium augmentatio sive diminutio dierum et noctium duorum minutorum et XL secundorum computatur} (from the equinox to the solstice the increase or decrease of the days and the nights is reckoned at two minutes and 40 seconds).\par
\par \Pg{XV}2) At the end of the second column of the same folio (in our text p. 2), after the words \emph{quandoque in hieme interdum in aestate accidat} (happens now in winter, now in summer), there stands the following table of the Arab months, without the number of their days:\par
\begin{center}\small\begin{tabular}{ll}
\emph{Almuharram} & . . . . . . . . . \\
\emph{Szaphar} & . . . . . . . . . \\
\emph{Rabbe} & . . . . . . . . . \\
\emph{Rabbe} & . . . . . . . . . \\
\emph{Jumedi} & . . . . . . . . . \\
\emph{Jumedi} & . . . . . . . . . \\
\emph{Rageb} & . . . . . . . . . \\
\emph{Saaben} & . . . . . . . . . \\
\emph{Ramadhan} & . . . . . . . . . \\
\emph{Sewil} & . . . . . . . . . \\
\emph{Dulkidda} & . . . . . . . . . \\
\emph{Dulheiga} & . . . . . . . . . \\
\end{tabular}\end{center}
\par 3) Table 3, which takes up one page in each of the manuscripts C and O, is spread over two pages in N (9$^{r}$ and 9$^{v}$): some columns are added here that are missing in C and O, but that occur in the \emph{Tabulae spuriae} which Nallino added to his edition of al-Battānī (II, pp. 300–307). All the spurious tables of al-Battānī can therefore be traced back to Maslama's reworking of the tables of al-Khwārizmī, and the conjecture that Nallino expressed about these tables (II, pp. 308–309: \emph{Haud improbabile est tabulas p. 302–307 ex opere Maslamah quoque, ut de p. 300–301 affirmatur, exscriptas esse, etc.}, it is not improbable that the tables of pp. 302–307 too were copied from the work of Maslama, as is asserted of pp. 300–301, etc.) has thus proved true. The columns added in Ms. N are the following:\par
\par A) On fol. 9$^{r}$, after the column \emph{Anni Alexandri collecti}, which is the second column of our Table 3, there comes a column headed \emph{Anni iazdeiardi} (\emph{Yezdegirdi}) \emph{post aeram arabum} (the years of Yazdegerd after the era of the Arabs); it agrees with column 4 of al-Battānī (II, 304), headed \emph{Anni Persarum sine intercalatione, etc.} (the years of the Persians without intercalation), except for the radix, which is given there as 9$^{a}$ 11$^{m}$ 9$^{d}$ (see also Table 1, where the incorrect radix 9$^{a}$ 11$^{m}$ 4$^{d}$ is found too).\par
\par B) On fol. 9$^{v}$ there is the column corresponding to the preceding one for the \emph{anni plani} or \emph{residui} (the single, or remaining, years); it stands after the column \emph{Anni Alexandri plani}, which is the fourth column of our Table 3. It agrees with column 4 of al-Battānī (II, 305), headed \emph{Anni singuli Persarum sine intercalatione} (the single years of the Persians without intercalation), except for the days for \emph{annus} 27, where our table has 19 but al-Battānī 18.\par
\par \Pg{XVI}C) On fol. 9$^{r}$, immediately after column A for the \emph{Anni iazdeiardi}, there is one for the \emph{Anni Diocletiani completi} (the complete years of Diocletian); on the whole it agrees with col. 3 of al-Battānī (II, 304), headed \emph{Anni Coptorum intercalationem habentes etc.} (the years of the Copts that have intercalation). With various Arab astronomers the \emph{era of Diocletian} (beginning on 29 Aug. 284 AD) was also called the \emph{era of the Copts}, a name by which various other eras were designated as well, such as the \emph{era of Nabonassar}, that of \emph{Philip Arrhidaeus}, and, in al-Battānī, also the \emph{era of Augustus} or \emph{Alexandrian era}, about whose beginning there is no full agreement\footnote{See \emph{Al-Battānī}, ed. Nallino, I, 244–245.}. The column in N shows several deviations from that of al-Battānī: under a) it should read 483 instead of 488, under b) 977 instead of 997; and the last four lines should then of course read:\par
\begin{center}\small\begin{tabular}{rrrr}
920. & 0. & 25. & $^{1}/_{4}$ \\
949. & 1. & 4. & 0 \\
978. & 2. & 13. & $^{3}/_{4}$ \\
1007. & 4. & 21. & $^{1}/_{2}$ \\
\end{tabular}\end{center}
\par Al-Battānī too has 919. 12 instead of 920. 0, and the last line is missing in his table. The numbers of the column \emph{Dies} agree with those of al-Battānī if we leave out the 25 of the second line from the top and move all the other numbers of the column up by one line, except for the two lowest numbers (13 and 21).\par
\begin{center}\scriptsize\begin{tabular}{rrrr|rrrr|rrrrr}
\multicolumn{4}{c}{\emph{A.\footnotemarkT{}}} & \multicolumn{4}{c}{\emph{B.}} & \multicolumn{5}{c}{\emph{C.}} \\
\emph{Anni collecti} & \multicolumn{3}{c}{\emph{Anni iazdeiardi post aeram arabum}} & \emph{Anni residui} & \multicolumn{3}{c}{\emph{Anni iezdeiardi}} & \emph{Anni collecti} & \multicolumn{4}{c}{\emph{Anni diocletiani completi}} \\
 & \emph{Anni} & \emph{Menses} & \emph{Dies} &  & \emph{Anni} & \emph{Menses} & \emph{Dies} &  & \emph{Anni} & \emph{Menses} & \emph{Dies} & \emph{Partes} \\
Radix & 9 & 11 & 4 & 1 & 0 & 11 & 24 & Radix & 337 & 10 & 19 & $^{3}/_{4}$ \\
30 & 19 & 2 & 12 & 2 & 1 & 11 & 14 & 30 & 366 & 11 & 25 & $^{1}/_{2}$ \\
60 & 48 & 3 & 28 & 3 & 2 & 11 & 3 & 60 & 396 & 1 & 28 & 0 \\
90 & 77 & 5 & 14 & 4 & 3 & 10 & 22 & 90 & 425 & 2 & 2 & $^{3}/_{4}$ \\
120 & 106 & 7 & 0 & 5 & 4 & 10 & 12 & 120 & 454 & 3 & 10 & $^{1}/_{2}$ \\
150 & 135 & 8 & 16 & 6 & 5 & 10 & 1 & 150 & 488$^{a}$) & 4 & 19 & $^{1}/_{4}$ \\
180 & 164 & 10 & 2 & 7 & 6 & 9 & 21 & 180 & 512 & 6 & 28 & 0 \\
210 & 193 & 11 & 18 & 8 & 7 & 9 & 10 & 210 & 541 & 7 & 7 & $^{3}/_{4}$ \\
240 & 223 & 0 & 29 & 9 & 8 & 8 & 29 & 240 & 570 & 8 & 15 & $^{1}/_{2}$ \\
270 & 252 & 2 & 15 & 10 & 9 & 8 & 19 & 270 & 599 & 10 & 24 & $^{1}/_{4}$ \\
300 & 281 & 4 & 1 & 11 & 10 & 8 & 8 & 300 & 628 & 11 & 3 & 0 \\
330 & 310 & 5 & 17 & 12 & 11 & 7 & 27 & 330 & 658 & 0 & 12 & $^{1}/_{2}$ \\
360 & 339 & 7 & 3 & 13 & 12 & 7 & 17 & 360 & 687 & 1 & 15 & $^{1}/_{4}$ \\
390 & 368 & 8 & 19 & 14 & 13 & 7 & 6 & 390 & 716 & 3 & 24 & 0 \\
420 & 397 & 10 & 5 & 15 & 14 & 6 & 25 & 420 & 745 & 4 & 3 & $^{3}/_{4}$ \\
450 & 426 & 11 & 21 & 16 & 15 & 6 & 15 & 450 & 774 & 5 & 11 & $^{1}/_{2}$ \\
480 & 456 & 1 & 2 & 17 & 16 & 6 & 4 & 480 & 803 & 6 & 20 & $^{1}/_{4}$ \\
510 & 485 & 2 & 18 & 18 & 17 & 5 & 24 & 510 & 832 & 8 & 28 & 0 \\
540 & 514 & 4 & 4 & 19 & 18 & 5 & 13 & 540 & 861 & 9 & 8 & $^{3}/_{4}$ \\
570 & 543 & 5 & 20 & 20 & 19 & 5 & 2 & 570 & 890 & 10 & 16 & $^{1}/_{2}$ \\
600 & 572 & 7 & 6 & 21 & 20 & 4 & 22 & 600 & 919 & 12 & 25 & $^{1}/_{4}$ \\
630 & 601 & 8 & 22 & 22 & 21 & 4 & 11 & 630 & 948 & 13 & 4 & 0 \\
660 & 630 & 10 & 8 & 23 & 22 & 4 & 0 & 660 & 997$^{b}$) & 14 & 13 & $^{3}/_{4}$ \\
690 & 659 & 11 & 24 & 24 & 23 & 3 & 20 & 690 & 1006 & 16 & 21 & $^{1}/_{2}$ \\
 &  &  &  & 25 & 24 & 3 & 9 &  &  &  &  &  \\
 &  &  &  & 26 & 25 & 2 & 29 &  &  &  &  &  \\
 &  &  &  & 27 & 26 & 2 & 19 &  &  &  &  &  \\
 &  &  &  & 28 & 27 & 2 & 7 &  &  &  &  &  \\
 &  &  &  & 29 & 28 & 1 & 27 &  &  &  &  &  \\
 &  &  &  & 30 & 29 & 1 & 16 &  &  &  &  &  \\
\end{tabular}\end{center}
\addtocounter{footnoteT}{-1}\stepcounter{footnoteT}\footnotetextT{The heads of the tables that Suter copies from the manuscripts stand in their Latin: \emph{anni collecti} and \emph{residui}, the collected and the remaining (single) years; \emph{anni iazdeiardi post aeram arabum}, the years of Yazdegerd after the era of the Arabs; \emph{anni diocletiani completi}, the complete years of Diocletian; \emph{anni}, \emph{menses}, \emph{dies}, \emph{partes}: years, months, days, parts.}
\par \Pg{XVII}4) At the foot of fol. 9$^{r}$ there is a small table that allows the conversion of \emph{Arab} years into \emph{Christian} ones, but only for the \emph{anni collecti} from 30 to 660; it is thus incomplete. Perhaps it is a remnant of the table mentioned in the commentary, p. 36, which is missing in the other manuscripts; to be sure, it is only a scholium by a later hand, but I give it here as well:\par
\begin{center}\small\begin{tabular}{rrrrrr}
\emph{Notae} & \emph{Anni arab.} & \emph{Anni dom. collecti} & \emph{Menses} & \emph{Dies} & \emph{Fract.} \\
5 & 30 & 650 & 7 & 23 & 3 \\
4 & 60 & 679 & 9 & 2 & 2 \\
1 & 90 & 708 & 10 & 11 & 1 \\
5 & 120 & 737 & 11 & 20 & 0 \\
4 & 150 & 767 & 0 & 28 & 2 \\
2 & 180 & 796 & 2 & 2 & 1 \\
7 & 210 & 825 & 3 & 11 & 0 \\
5 & 240 & 824$^{a}$) & 4 & 29 & 3 \\
4 & 270 & 883 & 5 & 28 & 2 \\
1 & 300 & 912 & 7 & 7 & 1 \\
5 & 330 & 941 & 8 & 16 & 0 \\
4 & 360 & 970 & 9 & 24 & 3 \\
2 & 390 & 999 & 11 & 3 & 2 \\
7 & 420 & 1029 & 0 & 7 & 0 \\
5 & 450 & 1058 & 1 & 15 & 3 \\
3 & 480 & 1087 & 2 & 24 & 2 \\
1 & 510 & 1116 & 4 & 3 & 1 \\
5 & 540 & 1145 & 5 & 12 & 0 \\
4 & 570 & 1174 & 6 & 20 & 3 \\
2 & 600 & 1203 & 7 & 29 & 2? \\
7 & 630 & 1232 & 9 & 7 & 1 \\
5 & 660 & 1261 & 10 & 17 & 0 \\
 & 690 & 1290 & 11 & 25 & 3 \\
 & 720 & 1320 & 1 & 4 & 2 \\
 & 750 & 1349 & 2 & 13 & 1 \\
\end{tabular}\end{center}
\par On this table we have to remark: under a) it should read 854; the numbers of the last column mean “quarters”; the numbers of the column \emph{Dies} (days) are often faulty, and can easily be corrected from modern tables; the \emph{Fractiones} too are not right throughout. The last three lines, for the years 690, 720 and 750, have been added by a more recent hand, and in Arabic numerals that seem to me to be those of the fourteenth or fifteenth century\footnote{The figures of the 2nd, 4th and 6th columns were illegible for me because the photographic image was too faint; for a copy of them from the original I am most grateful to Mr. Gonzalez, archivist in Madrid, a pupil of Prof. Asín.}.\par
\par 5) In the margin of fol. 9$^{r}$ there stands the following note by a second hand: \emph{Era (Aera) quoque} (?) \emph{cesaris augusti tempore posita est: dicta autem era ex quo omnis orbis es (aes) reddere professus est reipublicae. Era secundum ysidorum adcescit (sic!) — Bissextus autem} \Pg{XVIII}\emph{accescitur}\footnote{Probably for \emph{accersitur} = \emph{arcessitur} (= to fetch, to take).} \emph{a sexto. Nonas marcii usque ad diem pridie kalendarum ianuarii apponitur in cursu lunae} (?) (the era, too (?), was instituted in the time of Caesar Augustus; it was called era because the whole world then professed to pay tribute (\emph{aes}) to the state. The era, according to Isidore, \emph{adcescit} (sic!) — but the bissextile day is taken from the sixth. From the Nones of March to the day before the Kalends of January it is added in the course of the moon (?)).\par
\par The words \emph{dicta autem era ex quo . . . reipublicae} are indeed a quotation from Isidore of Seville, which L. Ideler also cites in his attempts to explain the word \emph{aera}\footnote{\emph{Handbuch der mathem. u. technischen Chronologie}, Berlin 1825–26, II, 426–430.}. The explanation of \emph{bissextus} is hardly intelligible, and was probably misunderstood by the scribe too. As is well known, the word comes from the fact that in leap years one said \emph{twice}: \emph{Sextus ante Kalendas Martii}, that is, one counted 24 February \emph{twice}. The scribe takes the \emph{Nonae Martii} instead of the \emph{Kalendae Martii}, and with what follows he probably means that the leap day was intercalated at various places, even at the end of December, which, as we shall see below, was the case in the so-called Spanish era. — What is meant by the words \emph{in cursu lunae} we do not understand.\par
\par 6) On fol. 11$^{v}$, at the words \emph{ideo quotiens duodecim menses in annum reducimus, totiens ab ipsa reductionis summa} (therefore as often as we reduce twelve months to a year, so often from the sum of the reduction itself; in our text p. 6), the following marginal note by a second hand is added: \emph{id est a mensibus et diebus residuis ex quibus fit annus qui aliis annis adiicitur; non enim complebitur annus ex XII mensibus XXX dierum nisi addatur V dies et quarta; menses enim tantum XXX dierum hic dantur} (?), \emph{unde erit eis quandoque addendum quandoque subtrahendum ut vere romani menses habeantur} (that is, from the remaining months and days, from which a year is made that is added to the other years; for a year will not be completed from 12 months of 30 days unless 5 days and a quarter are added; for here only months of 30 days are given (?), so that something will sometimes have to be added to them and sometimes taken away, that true Roman months may be had).\par
\par \Pg{XIX}7) At the lower margin of fol. 18$^{v}$ (Table 15) there is the following note: \emph{Adiectiones Veneris et Mercurii sunt adiectio solis; hae autem adiectiones omnes iuxta civitatem Londoniae in hoc libro computantur et mediis cursibus planetarum adiciuntur} (sic) (the additions of Venus and Mercury are the addition of the sun; all these additions are computed in this book for the city of London and are added to the mean courses of the planets). — By this the author of the note, and this is without doubt Robert of Chester\footnote{For according to the testimony of Hermannus Dalmata he translated the work of al-Battānī into Latin, and he is also said to have recomputed its tables for the meridian of London. See above, p. XIII, and al-Battānī, ed. Nallino, I, \emph{Praefatio}, XXXIV–XXXV, where, however, Rudolphus Brugensis still stands instead of Hermannus Dalmata.}, means to say that the mean longitudes of Venus and Mercury are the same as that of the sun, and that for these and for the three upper planets he has computed the \emph{adiectiones} for the meridian of London. By these \emph{adiectiones} the following is to be understood. The mean longitudes of the sun, the moon and the planets computed from Tables 9–20 hold for the meridian of \emph{Arin} (\emph{Ujjain}); for London a certain quantity must be added to the mean longitude of each of these heavenly bodies, namely its mean motion (\emph{motus medius}) in c. 5$^{h}$ (for Ujjain is c. 75° east of London); this quantity to be added Robert of Chester calls \emph{adiectio}. — He has given these \emph{adiectiones} in several places (each time at the lower margin of the leaf): for Saturn on fol. 24$^{v}$ (Table 9) as 18″, for Jupiter on fol. 28$^{v}$ (Table 11) as 45″, for Mars on fol. 14$^{v}$ (Table 13) as 4′ 49″, and for the dragon's head on fol. 23$^{v}$ (Table 19) as 30″; the tables for the mean motion of the sun and of the moon are missing in N. These figures, however, do not agree with our tables for a difference in time of 5$^{h}$; Robert of Chester seems to have taken it as not quite 4$^{h}$. The approximately correct figures, from Tables 10, 12, 14 and 20, are: Saturn 25″, Jupiter 1′ 2″, Mars 6′ 33″, the dragon's head 40″ (for the sun 12′ 19″, for the moon 2° 44′ 42″).\par
\par 8) On fol. 23$^{v}$ (Table 19), in the lower left corner, N adds the following figures, which differ in the seconds, and partly in the minutes too, from the data of the \emph{Medialitas draconis} (the mean motion of the dragon) that stand immediately above them:\par
\begin{center}\small\begin{tabular}{rrrrr}
\emph{collecti} & \multicolumn{4}{c}{\emph{medio die Cordubae (at noon at Córdoba)}} \\
360 & 5 & 5 & 20 & 42 \\
390 & 11 & 28 & 48 & 24 \\
420 & 6 & 22 & 16 & 50? \\
450 & 1 & 15 & 43 & 49 \\
480 & 8 & 9 & 11 & 32 \\
510 & 3 & 2 & 39 & 14 \\
540 & 9 & 26 & 6 & 57 \\
570 & 4 & 19 & 34 & 39 \\
\end{tabular}\end{center}
\par 9) On fol. 24$^{r}$ (Table 20), N has in the last column (it leaves out the thirds) the following figures for the seconds: 0, 0, 0, 0, 1, 1, 1, 1, 2, 2, 2, 3, 3, 3, 4, 4, 4, 4, 5, 5, 5, 6, 6, 6, 6, 7, 7, 7, 8, 8. One sees that several of these figures, even when the thirds that have been dropped are taken into account, do not agree with those of Mss. C and O.\par
\par 10) On fol. 42$^{r}$, after the words \emph{Sed et aequalibus idem e contrario addendo inaequales fient} (but by adding the same to the equal ones, conversely, they will become unequal; in our text p. 25), a second hand has added: \emph{Verus dies naturalis medius est XXIIII horarum, id est revolutio CCCLX graduum et LIX min. et V sec. Unde} (or \emph{una}) \emph{hora est XV grad. et II. min. et XXVIII sec. Sed quoniam gradus aequinoctialis lineae, secundum quam haec determinata} (?) \emph{sunt, et zodiaci in quo sol movetur, non similiter transeunt meridianam lineam, sed et sol aliquando citius aliquando tardius contra celum movetur, necesse est ergo naturali diei medio aliquando addi, aliquando subtrahi. Attamen haec tabula sic facta est, quod semper subtrahitur a die media, sed aliquando sequenti die maius aliquando minus quam priori. Sed quando maius tunc ille dies aequatus a meridiana linea in meridianam erit minor die media, si autem minus tunc necesse est esse maiorem; maior enim subtractio minuit de medio, minor vero subtractio de medio et aequali residuum relinquit maius quam faceret subtractio maior. Exempli gratia: si hodie a revolutione sperae}\footnote{\emph{revolutio sperae} = 360°.} \emph{et LIX min. et VIII sec. subtraherentur II min. et VIII sec., et in sequenti die III min. et VIII sec., esset tunc dies aequatus revolutio sperae et LVIII min., quod est minus die medio; item si priori die III min. et VIII sec. subtraherentur et sequenti die tantum II min. et VIII sec., erit dies aequatus revolutio sperae et LX min., quod est maius die medio} (the true mean natural day is of 24 hours, that is, a revolution of 360 degrees and 59 min. and 5 sec. Hence (or: one) hour is 15 deg. and 2 min. and 28 sec. But since the degrees of the equinoctial line, by which these things are determined (?), and of the zodiac, in which the sun moves, do not cross the meridian line alike, and the sun moreover moves against the heaven now faster, now slower, something must therefore be added to the mean natural day at one time, subtracted at another. This table, however, is so made that something is always subtracted from the mean day, but on the following day sometimes more, sometimes less than on the one before. And when more, that equated day from the meridian line to the meridian line will be less than the mean day; when less, it must be greater; for the greater subtraction diminishes the mean day, while the smaller subtraction from the mean and equal day leaves a greater remainder than the greater subtraction would. For example: if today 2 min. and 8 sec. were subtracted from a revolution of the sphere and 59 min. and 8 sec., and on the following day 3 min. and 8 sec., the equated day would then be a revolution of the sphere and 58 min., which is less than the mean day; likewise, if 3 min. and 8 sec. were subtracted on the earlier day and only 2 min. and 8 sec. on the following one, the equated day will be a revolution of the sphere and 60 min., which is greater than the mean day).\par
\par \Pg{XX}11) In the tables of the conjunction and the opposition of the sun and the moon for the \emph{anni collecti} (fol. 42$^{v}$ and 43$^{r}$, Tables 69 and 70), N also gives the data for the years 541 and 571, which are missing in C and O; the former seem to be in the same hand as the preceding ones, the latter in another hand. Since the table can always be continued by anyone, we do without reproducing these lines.\par
\par 12) On fol. 47$^{v}$ (in our text p. 29), after the words \emph{erit et ipsa meridionalis} (it too will be southern), a second hand (?) has added: \emph{Haec circa meridiem, alibi autem: Inventa vera lunae latitudine nadair}\footnote{Probably \emph{nadir}.} \emph{gradus coniunctionis ponatur in astrolabio super horam coniunctiones et videatur quot gradibus distet gradus coniunctionis a centro minimi circuli almucantarat, et si fuit vera lunae latitudo septemtrionalis, auferatur ab illis gradibus, si meridiana, addatur; et per id fiat introitus in tabulas latitudinis per aspectum} (this about noon; elsewhere, however: when the true latitude of the moon has been found, let the nadir of the degree of the conjunction be set on the astrolabe on the hour of the conjunction, and let it be seen how many degrees the degree of the conjunction is distant from the centre of the smallest almucantar; and if the true latitude of the moon was northern, let it be taken from those degrees, if southern, added; and with this let the tables of the latitude by aspect be entered).\par
\par 13) On fol. 49$^{r}$ (in our text p. 29), after the words \emph{horas coniunctionis quaesitae ponet} (will set the hours of the conjunction sought), a second hand has added: \emph{Introitus harum tabularum hic est: primo videatur coniunctionis distantia solis et lunae a capite vel cauda et si fuerit luna in longiori longitudine sui brevis circuli}\footnote{\emph{brevis circuli} = \emph{epicycli} (of the small circle = of the epicycle).} \emph{intretur in tabulas priores, aliter in posteriores et sumatur de digitis quod opponitur distantiae praedictae. Exinde sumatur diversitas aspectus lunae et videatur quot digitis ipsa aequipolleat. Et si fuerit luna tunc septemtrionalis a sole addantur illi digiti super inventos digitos in tabula, si meridiana, demantur si poterit, aliter enim non erit eclipsis; et sic inventum erit quantitas eclipsis. Et tunc cum illo invento in tabula sumantur minuta casus in eius directo et erit mora usque ad mediam eclipsim. In diversitate aspectus quot fuerint digiti sic inveniatur: Per tabulas elbut inveniatur quantitas solaris corporis et dividatur in XII aequas partes quia semper duodecima diametri solis digitus dicitur. Videatur quotiens eius duodecima in diversitate aspectus contineatur, et tot fient ex ea digiti. Si ergo inventum fuerit plus XII digitis et XLVIII min. in longitudine propiori vel plus X et XLVIII min. in longiori, et luna fuerit septemtrionalis a sole, auferantur praedicti digiti de invento et residuum clarum in sole relinquitur. Si vero luna fuerit meridiana et inventum secundum se obscurabit de sole, haec secundum eos quibus sol semper meridianus est; quia cum fuerit septemtrionalis contra operari oportebit} (the entry of these tables is this: first let the distance of the sun and the moon at the conjunction from the head or the tail be seen; and if the moon is at the greater distance of its small circle, let the first tables be entered, otherwise the second, and let the digits be taken that stand against the said distance. Then let the diversity of aspect of the moon be taken, and let it be seen to how many digits it amounts. And if the moon is then north of the sun, let those digits be added to the digits found in the table; if south, let them be taken away, if that can be done, for otherwise there will be no eclipse; and what is so found will be the magnitude of the eclipse. And then, with what has been found, let the minutes of the fall be taken in the table, in the line opposite it, and that will be the duration up to the middle of the eclipse. How many digits there are in the diversity of aspect is found thus: by the tables of \emph{elbut} let the size of the body of the sun be found, and let it be divided into 12 equal parts, since a twelfth of the diameter of the sun is always called a digit. Let it be seen how many times its twelfth is contained in the diversity of aspect, and so many digits will be made of it. If then more than 12 digits and 48 min. is found at the nearer distance, or more than 10 and 48 min. at the farther, and the moon is north of the sun, let the said digits be taken from what has been found, and the remainder is left bright in the sun. But if the moon is southern, what has been found will by itself darken so much of the sun; this for those for whom the sun is always to the south; for when it is to the north one will have to work the other way round).\par
\par Then there follows, by a third hand: \emph{Sed et in maiori elongatione ante caudam vel post caput, quia usque ad XXIII gradus fit solis eclipsis. Et tunc sciatur vera lunae latitudo et diversitas aspectus latitudinis eius; et auferatur minus de maiori, et per residuum erit quantitas scita eclipsis, cuius digiti per solis et residui quantitatem, ut praedictum est, investigentur, et per eos mora in tabula. Iungantur autem medietates}\footnote{In the text \emph{mediates}.} \emph{diametrorum solis et lunae, et si residuum praedictum fuerit eis aequale vel maius, non erit eclipsis, si nichilum fuerit residuum, totalis, si vero residuum minus fuerit, erit eclipsis partialis} (but also at a greater elongation before the tail or after the head, since an eclipse of the sun occurs up to 23 degrees. And then let the true latitude of the moon and the diversity of aspect of its latitude be known; and let the smaller be taken from the greater, and by the remainder the magnitude of the eclipse will be known, whose digits are to be found from the size of the sun and of the remainder, as was said before, and by them the duration in the table. Let the halves of the diameters of the sun and of the moon be joined; and if the said remainder is equal to them or greater, there will be no eclipse; if the remainder is nothing, it will be total; and if the remainder is less, the eclipse will be partial).\par
\par We have remarked in the commentary on chapter 35 that this chapter is very short and certainly incomplete. The author of the first of the above additions seems to have felt this too, and therefore thought he had to give a better and more intelligible account of the determination of an eclipse of the sun. But his rule is \Pg{XXI}when one compares it with that of al-Battānī, probably not correct either, or at least it too is incomplete, which led a later hand to add a second correction.\par
\par In the text and in the tables there occur some quite corrupt \emph{Arabic} words and sentences, which I should like to try to explain here, since I could not well place them in the glossary under a definite headword. How they came into the work is hard to decide, except for the Persian names of the planets; they seem to me mere games of the translator or of a copyist; it can hardly be supposed that they come from Maslama al-Majrīṭī, let alone from al-Khwārizmī. For the kind help that Professors Nallino in Palermo and Seybold in Tübingen gave me with the difficult passages, I express my sincere thanks to them here as well.\par
\par In Ms. O, on fol. 103$^{v}$ (Table 7), there stands the following note in the upper margin: \emph{kulelach. vveli. kemithl. elhileli. fikemelen. wanuczeni.} This is probably meant to be \emph{kull al-aḥwāli kamithl al-hilāli fī kamālin wa-nuqṣānin}, which means, freely translated: all changes (of life, of nature) behave like the moon; (they show themselves) in a waxing (up to fullness) and in a (following) waning.\par
\par In the same manuscript, on fol. 104$^{v}$ (Table 9), above the column title \emph{Heza Saturni supradictum}, there stands the following note: \emph{Nota quod hoc argumentum felek eltedewir servit. bifelek eltedewir lieztekim asmir.} (note that this argument serves the \emph{felek eltedewir} ...). — \emph{Felek eltedewir} = \emph{falak al-tadwīr} is the \emph{epicycle} of the planets; \emph{lieztekim asmir} might, as Nallino conjectures, mean \emph{li-yastaqīm al-sair}, or \emph{al-masīr}, that is, that it may determine, order, regulate the course (of the planet); and indeed the epicycle was introduced for this purpose.\par
\par \Pg{XXII}In the same manuscript, on fol. 127$^{v}$ (Table 55), there is the following note in the upper margin: \emph{Cullel. kavvekib. kumna. fiscerafehu. ille Otari. ahine unxi ademu.} Nallino conjectured that the whole might be a verse in honour or praise of some high personage (a prince), and should read \emph{kullu l-kawākibi qumna fī sharafihi illā ʿUṭārida}, that is, all the planets rose in his honour except Mercury; but the three words \emph{ahine unxi ademu} remained unclear to him. — On a visit to Prof. Seybold in Tübingen in the autumn of 1913 I laid this passage before him too; it at once reminded him of some passages of the Koran\footnote{See Suras II, 32; VII, 11; XX, 116.} that tell of the angels, who all fell down before \emph{Adam} except \emph{Iblīs} (Satan); and he therefore put forward the conjecture, surely well founded, that the \emph{a} of \emph{ahine} is probably corrupted from a \emph{d} and belongs at the end of \emph{Otari}, and that the \emph{x} in \emph{unxi} is to be read \emph{sh}, after the Spanish manner of transcription, so that the last three words are to be read \emph{ḥīna unshiʾa ādamu}, and the whole sentence, in half classical, half vulgar pronunciation, \emph{kull al-kawākib qumna fī sharafihi illā ʿUṭārid ḥīna unshiʾa ādamu}, that is, all the planets rose in his (Adam's) honour except Mercury, at the time when Adam was created. — That \emph{Mercury} here takes the place of \emph{Satan} (\emph{Iblīs}) probably has its reason in the fact that in the Middle Ages the Greek \emph{Hermes} was identified with the \emph{Hermes Trismegistos} or \emph{Thoth} of the Egyptians, and he was the patron god of the alchemists, astrologers and magicians. In the books of miracles and magic of the Christian Middle Ages \emph{Ophiel} is named as one of the seven princes of the underworld; this word is very probably compounded of the Greek \emph{ophis} = serpent and the Hebrew \emph{el} = god, thus serpent-god = \emph{Hermes}; joined with the \emph{megistos} of the word \emph{Trismegistos}, \emph{Ophiel} became \emph{Megistophiel}, and from this arose the corrupted, or deliberately distorted, \emph{Mephistopheles}. This is probably the best etymology of the word\footnote{See J. Goebel in the \emph{Beilage zur Allgem. Zeitung}, 1905, no. 195.}.\par
\par Again in Ms. O, in the upper margin of fol. 132$^{v}$ (in our text p. 20), there stand the four words \emph{eltabeb, helhacib, elhaazem, el-munegim}; the first, the second and the fourth are probably to be read \emph{al-ṭabīb, al-ḥāsib, al-munajjim}, and then mean: the physician, the reckoner, the astronomer; the third I cannot interpret. What the scribe means to say by these words is unclear to me; perhaps he means that the following chapter, \emph{Quot gradibus in qualibet regione quodlibet oriatur signum} (with how many degrees each sign rises in any region), is important for these professions.\par
\par In Ms. C, on fol. 83$^{r}$ (p. 27 of our text), after the words \emph{et accipies quod contra invenies} (and you will take what you find opposite), there stand the following: \emph{wakhod mebibelche}; Nallino believes, and probably rightly, that the last word must be read \emph{mekibelche}, and the sentence would then be \emph{wa-khudh mā qibālahā}\footnote{For the reader who does not know Arabic we remark that the western Arabs pronounced \emph{ā} as \emph{ē}, and in Morocco still pronounce it so (\emph{Fēs} for \emph{Fās}).}, that is, and take what stands opposite it (that is, opposite the argument); thus simply the Arabic text of the preceding translation \emph{et accipies quod contra invenies}.\par
\par In the titles of Tables 27, 33, 39 and 45 both Mss. C and O give the \emph{Persian} name of the planet after its \emph{Arabic} one each time; only in Table 51 is the Persian name of Mercury missing. Those of Saturn (\emph{Keiwen} = \emph{Kaiwān}), Jupiter (\emph{Hormi} = \emph{Hormuz}) and Mars (\emph{Baheram} = \emph{Bahrām}) I recognized, but those of Venus\footnote{The Persians had two different names for this planet.} are so corrupt that I could not succeed in interpreting them. On fol. 122$^{v}$ (Table 45) of Ms. O, after the Arabic name of Venus (\emph{elzohara}), there stand the following words: \emph{beied hakht eneid}; on fol. 65$^{r}$ of Ms. C, after the same Arabic word, there stands \emph{bei ed halt eneid}. In the last word, \emph{eneid}, I did suspect the Persian name of Venus known to me, \emph{nāhīd}, but the \emph{e} before \emph{neid} troubled me, for I did not know that the Persian name is also written \emph{anāhīd}; \emph{beied hakht} or \emph{bei ed halt}, however, was quite unintelligible to me. Here again Nallino's knowledge of these things put me on the right path: \emph{beied hakht eneid} is corrupted from \emph{bēdukht anāhīd}; these are the names of two ancient Iranian female deities, which were later transferred to the planet Venus. — It is very probable that these Persian names of the planets come from al-Khwārizmī, for Maslama will hardly have understood Persian.\par
\par \Pg{XXIII}On the principles that have guided me in reproducing the text and the tables I must remark the following.\par
\par In the text the words are reproduced as they stand in the manuscripts, with these exceptions: first, for the genitive and dative ending of the first declension singular the modern \emph{ae} has been written instead of the medieval \emph{e}; secondly, for \emph{u}, when it stands before a vowel and is therefore read \emph{v}, this letter has always been put; thirdly, both manuscripts often (not always) write \emph{c} for \emph{t} before \emph{i} and a following vowel; we have always written \emph{t}.\par
\par The text and the tables contain only Roman numerals, except for some later additions in O and N. In the text I have kept the Roman numerals, except in the addition by a later hand to chapter 21, which contains very many and rather large numbers and occurs only in C. In the tables, to make printing and reading easier, I have replaced all the numerals by Arabic ones.\par
\par For zero three signs occur, in the text as well as in the tables: C and O have now {\greekfont τ} (= \emph{teca})\footnote{See M. Curtze, \emph{Petri Philomeni de Dacia in Algorismum vulgarem Johannis de Sacrobosco Commentarius}, Hauniae 1897, pp. 2 and 26.}, now {\greekfont θ} (= \emph{theta}?), N always $\overline{0}$\footnote{{\greekfont τ} and $\overline{0}$ are also found in the \emph{Prologus N. Ocreati in Helceph ad Adelardum Batensem magistrum suum} (ed. C. Henry in the Abhandlungen z. Geschichte der mathem. Wissenschaften III, p. 131 ff.). Likewise in the manuscript of the Escorial that Nallino used for his edition of al-Battānī the sign $\overline{0}$ stands in the tables where, for example, the minutes are wanting.}. It is interesting that in N, in various places, some of which I have mentioned in the footnotes, XXX0 (or XX0) is written for 38 (a few times also for 28); so, in order not to let the number in Roman numerals become too long, VIII was replaced by 0, which is surely the initial letter of \emph{octo}.\par
\par As for the tables, all the manuscripts except N show an incredible arbitrariness and inconsistency in the titles and in the headings of the columns: sometimes both the Arabic and the Latin title stand there, sometimes only the Arabic, sometimes only the Latin, sometimes no title at all. The headings of the columns, too, are often worded differently although they are meant to designate the same thing; thus, for example, in the same manuscript one table has \emph{menses arabici} and another \emph{menses arabum}, one \emph{in annis collectis} and another only \emph{in collectis}, one \emph{via}, another \emph{semita}, a third \emph{catar} (= \emph{linea}), one \emph{recedentis} and another \emph{descendentis}, etc.; often the column titles are left out altogether, especially in C. The greatest confusion reigns in the designations of the smaller parts of the angle: one table has \emph{Min. Sec.}, another \emph{Dak.}\footnote{\emph{Dak.} = \emph{Dakaicae} (see the Arabic glossary).} \emph{Then.}\footnote{\emph{Then.} = \emph{Theniae} (see the Arabic glossary).}, a third \emph{Dak. Sec.}, a fourth (especially in C, but not throughout) \emph{P.}\footnote{\emph{P.} = \emph{Puncta} for \emph{Min.}} \emph{M.}\footnote{\emph{M.} = \emph{Minuta} for \emph{Sec.}}, and these names are now abbreviated so, now written out in full. Only in the designation of the signs of the zodiac (\emph{Signa}) and of the degrees (\emph{Gradus}) is there agreement, but they are \Pg{XXIV}sometimes more, sometimes less abbreviated. On this point I have simplified things as far as possible and always written Sig. Gr. Min. Sec.\par
\par The establishment of the text and the critical apparatus are on the whole the work of Bjørnbo. This scholar was of the opinion, as I learned from a letter of Prof. Heiberg to me, that the text of C was the best, and for this reason he generally took it as his basis. A renewed comparison of C and O\footnote{Bjørnbo too had intended this, as various notes in his manuscript show; but unfortunately he could no longer carry out a second reading, and I was therefore often obliged to make changes in his report of the readings.}, however, has shown me that in several places the reading of O is the more apt (a few times also that of N); in such cases I have therefore taken the liberty of adopting the better reading into the text. Where the reading of a manuscript is not given in the footnotes, it agrees with that of the text.\par
\par The commentary that follows the text is entirely my own work. I first asked myself whether it might be necessary to follow the text with a German translation, since the subject is often set out very briefly and obscurely; but I told myself that the reader would be only half served by it, for whoever is not acquainted with ancient and medieval astronomy could not do much with a mere translation, and a commentary would have to be added to it as well. I have therefore done without the translation and contented myself with a commentary, which I have tried to make more useful by giving, for the chief astronomical problems treated in this work, examples solved by the rules of the text. — In the commentary I have translated the titles of the individual chapters of the text fairly literally; for convenient reference the chapters are numbered, which they are not in the manuscripts.\par
\par I must confess that this commentary was a difficult and laborious work for me, and whoever reads through the individual chapters of the text will surely find this understandable. Here the excellent edition of al-Battānī, with the commentary on it, by C. A. Nallino and G. Schiaparelli\footnote{This work, which appeared in three parts from 1899 to 1907 under the title \emph{Al-Battānī sive Albatenii Opus astronomicum, edit. a} C. A. Nallino, in the \emph{Pubblicazioni del R. Osservatorio di Brera in Milano}, N. XL, I shall cite throughout as Batt., ed. Nallino, or simply as Batt.} rendered me the most valuable services; without this edition a commentary on the present work would hardly have been possible for me. In spite of this help, however, several points have remained unclear, and I hope that scholars versed in these questions will perhaps lend their support to clearing them up.\par
\par In the commentary I have allowed myself to use abbreviations for the three proper names that occur most often, since this also makes the printing considerably easier: so I always write Kh. for al-Khwārizmī, Batt. for al-Battānī and Ptol. for Ptolemy.\par
\par \Pg{XXV}The work closes with two indexes, an \emph{Arabic glossary} and a general \emph{index of names and subjects}. The basis of the former was compiled by Dr. R. O. Besthorn in Copenhagen, to whom, together with J. L. Heiberg, we owe the edition of the oldest Arabic translation of Euclid. Various things were afterwards supplemented by me, in particular the correct forms of the Egyptian and Persian names of the months. I give the Arabic names only in transcription, first as they stand in the text and then in the correct form; in this I follow the rules of the \emph{Encyclopaedia of Islam}, which is now appearing under the direction of Prof. Houtsma in Utrecht and others. For the reader I will only mention here that \emph{dj} is to be pronounced like English \emph{j}, \emph{z} like a voiced \emph{s}, \emph{kh} like the Alemannic \emph{ch}, \emph{y} like English \emph{y}, and \emph{dh} and \emph{th} like the voiced and the voiceless English \emph{th}\footnoteT{Suter gives the German equivalents: \emph{dj} as \emph{dsch}, \emph{z} as soft \emph{s}, \emph{y} as \emph{j}.}; the other signs need no explanation. The meaning of the Arabic words I have given in German\footnoteT{In this translation the meanings are given in English.}, with the agreement of Dr. Besthorn, who had originally added it in Latin. — As to the use of the glossary, I point out that in all words that begin with the article \emph{al} or \emph{el} the article is to be left out: look for \emph{elmuht} under \emph{muht}, for \emph{eltheniae} under \emph{theniae}, etc. — For words that recur very often, such as \emph{dakaicae} (= minutae), \emph{theniae} (= secundae) and a few others, I leave out the references to the pages, except those where their meaning is explained.\par
\par For both indexes the following holds: numbers without a preceding letter are the pages of the introduction, of the text and of the commentary; numbers preceded by T. are the numbers of the tables; N in brackets means “note”.\par
\PartHead{Commentary}{K}\Pg{32}
\UnitHead{Introduction of the text}{Introduction of the text}{K01-1}
\par\begin{center}(Text pp. 1–2.)\end{center}
\par \Pg{33}The Arabic word \emph{ezich} = \emph{elzīj} (pronounced \emph{ez-zīj}) is, like several other Arabic technical terms, not translated. It is originally a Persian word (\emph{zīg}) and means the warp threads of a cloth, into which the weaver or the embroiderer works the weft or the embroidery, and then, correspondingly, “astronomical tables”. Of these there was a great number in the Arabic astronomical literature: the oldest were the tables composed after the Indian \emph{Siddhāntas} (see the introduction) by Ibrāhīm al-Fazārī and Yaʿqūb b. Ṭāriq; then come the two editions of the \emph{Sindhind} (the Arabic spelling of the Indian \emph{Siddhānta}) by Muḥ. b. Mūsā al-Khwārizmī, which were probably composed in the years 800–830 AD; after them the “tested tables” of the astronomers of al-Maʾmūn (Yaḥyā b. Abī Manṣūr, Sanad b. ʿAlī, ʿAbbās b. Saʿīd al-Jawharī and others)\footnote{For these astronomers see H. Suter, \emph{Die Mathematiker und Astronomen der Araber}, etc., pp. 3–14, and our introduction to this edition.}. — It is possible that Kh.'s later edition of the \emph{Sindhind} was already somewhat influenced by the observations of al-Maʾmūn's astronomers and so departed somewhat further from the Indian models than the first edition, and that, as we should like to conclude from the content of the present work, it was probably this second edition that Maslama al-Majrīṭī (see the introduction) subjected to a reworking. But of course nothing certain can be said about this; we are only putting forward conjectures that seem to us justified. — Which of the various Indian \emph{Siddhāntas} Kh. took as the basis of his tables likewise cannot be decided with certainty. Until now it has been regarded as fairly probable that it was the \emph{Brahma-Siddhānta} (composed by Brahmagupta in the first half of the seventh century AD); but we shall see further on that other \emph{Siddhāntas} may also come into question, especially the \emph{Sūrya-Siddhānta}. This would then indicate that Kh. perhaps also consulted the Persian “Tables of the Shāh” (\emph{zīj al-shāh}, or \emph{zīj al-shahriyār})\footnote{See Batt., ed. Nallino, I, 218.}, which were translated into Arabic in the second century of the \emph{Hijra} and were probably based on the \emph{Sūrya-Siddhānta}. In particular the value of the apogee of the sun and the manner of computing the equations of the planets might have been taken from these Persian tables; for Ibn al-Qifṭī\footnote{See the edition by Lippert, Leipzig, 1903, pp. 270–271.}, in his biography of al-Fazārī, says the following after Ibn al-Ādamī\footnote{See H. Suter, \emph{Die Mathematiker und Astronomen der Araber}, etc., p. 44.}: “He (Kh.) relied in his tables on the mean places of the \emph{Sindhind}, but departed from it in the equations (of the planets) and in the obliquity (of the ecliptic); he fixed his equations by the method of the Persians, and the declination of the sun by the method of Ptolemy.” Since the latter assertion is correct, as we shall see below, the former may well be true too; we cannot prove it, since these Persian tables no longer exist.\par
\par According to the introduction of the text, the data of the tables for the mean and the true places of the sun, the moon and the planets (or, as the Arabs say, of the seven planets) refer to the meridian of \emph{Arin} (a wrong spelling for \emph{Azīn} or \emph{Uzain}, and this in turn the Arabic form of the Indian \emph{Ujjain}, in Sanskrit \emph{Ujjayini}), which lies in northern India (geogr. lat. 23° 11′, geogr. long. 75° 52′ from Greenwich). As for the time to which the \emph{radix} of the various tables refers in each case, the starting point given is noon of 1 \emph{Muḥarram} of the year 1 of the Hijra (that is, 15 July 622 AD); in Table 4, for example, the radix for the sun is given as 3$^{s}$ 23° 25′ 48″ (= 113° 25′ 48″), that is, this was the mean longitude of the sun at noon (the astronomical beginning of the day with most Arab astronomers) of 1 \emph{Muḥarram} of the year 1 of the \emph{Hijra} for the meridian of \emph{Arin}. Maslama al-Majrīṭī, then, as we already remarked in the introduction, replaced the \emph{era} of Yazdegerd (beginning on 16 June 632 AD), on which Kh. had based himself, by the Muslim era, but he kept the meridian of \emph{Arin}. This holds, however, only for the tables of the mean and the true motion of the sun, the moon and the five planets; we shall see later that Maslama referred a series of tables (especially those of the conjunctions, the oppositions and the eclipses) to the meridian of \emph{Córdoba}. The possibility is therefore not excluded that these tables were not originally contained in Kh.'s work; statements pointing to this were indeed made, as we saw in the introduction, by the Spanish scholar Ibn Ṣāʿid, and Ibn al-Qifṭī\footnote{Ed. Lippert, p. 326.} repeated them; but in that case Kh.'s work would have been all too insignificant, for precisely these tables are among the most important of an astronomical work.\par
\par In this introduction there also occurs the only proper name that these tables contain besides those of the author and the translator, the name “Ptolomaeus”. One sees, then — and this also follows from the study of the tables — that Kh., like the later astronomers after him, especially Batt., already relied on Ptolemy's \emph{Almagest} and was therefore also substantially influenced by it, as will appear in several places in the course of our discussion.\par
\UnitHead{Chapter 1. Description of the Arab year}{Chapter 1. Description of the Arab year}{K01-6}\Pg{34}
\par\begin{center}(Text p. 2.)\end{center}
\par The Arab lunar year has about 354$\tfrac{11}{30}$ days; it is divided into 12 months of alternately 30 and 29 days; to make up for the fraction of a day, the last month, which otherwise counts 29 days, is sometimes given 30 days. In which years this happens Kh. does not say: they are the 2nd, 5th, 7th, 10th, 13th, 16th, 18th, 21st, 24th, 26th and 29th of a cycle of 30 years. Such a year of 355 days is called \emph{sana kabīsa} = leap year. The Arab months are called (in the text Kh. names only the first two and the last):\par
\begin{center}\small\begin{tabular}{llll}
1) \emph{Muḥarram} & = 30 days & 7) \emph{Rajab} & = 30 days \\
2) \emph{Ṣafar} & = 29 days & 8) \emph{Shaʿbān} & = 29 days \\
3) \emph{Rabīʿ I} & = 30 days & 9) \emph{Ramaḍān} & = 30 days \\
4) \emph{Rabīʿ II} & = 29 days & 10) \emph{Shawwāl} & = 29 days \\
5) \emph{Jumādā I} & = 30 days & 11) \emph{Dhū l-Qaʿda} & = 30 days \\
6) \emph{Jumādā II} & = 29 days & 12) \emph{Dhū l-Ḥijja} & = 29 days (in leap years 30) \\
\end{tabular}\end{center}
\par These names are often badly disfigured in the text and in the tables; it should be noted in particular that \emph{j} is mostly rendered by \emph{i}, so that the last month is almost throughout written \emph{Dulheia}; \emph{sh} is often expressed by \emph{sc}, especially in Ms. O (see also the introduction).\par
\par At the end of the chapter Kh. remarks that the length of the Arab lunar year in relation to the solar year has the consequence that the beginning of the Arab year falls in different seasons, now in winter, now in summer.\par
\UnitHead{Chapter 2. Description of the Roman year}{Chapter 2. Description of the Roman year}{K01-12}
\par\begin{center}(Text pp. 2–4, Table 1.)\end{center}
\par By this name the Arab astronomers usually designate the years of the era of \emph{Dhū l-Qarnayn} (that is, of Alexander the Great), the so-called \emph{Seleucid} era, which began on 1 October 312 BC and whose years were taken as 365$\tfrac{1}{4}$ days, as in the \emph{Julian} reckoning. Interesting is the remark of our text that the leap day was added each time (thus every 4 years) at the end of December. Nallino\footnote{Batt., ed. Nallino, II, 314–316.} calls this an erroneous view; but he did not yet know this passage of our tables, and according to it it seems fairly certain that at least among the Christian Spaniards of the tenth century (in Maslama's time) this \emph{intercalation} of the leap day really took place in December instead of February. This also follows from the \emph{Notulae} or \emph{Signa} of our Table 3$^{a}$, which increase by 4 from December to January in the leap years, but by 3 in the other years. The Spaniards had, moreover, a special chronology of their own (see below), which began 38 years before ours, so they could well also have had an intercalation of their own, valid only among them \Pg{35}and was unknown in the East; for the authors of the \emph{Alfonsine} tables too say (according to Nallino): \emph{Dissentiunt de mense cui dies bissextus addere debeant; alii enim mensi Decembri eum addunt, alii Februario} (they disagree about the month to which they should add the bissextile day; for some add it to the month of December, others to February). Later western Arab authors say the same.\par
\par In the last paragraph of this chapter reference is made to a table that gives the times elapsed from important events, such as the Flood, or the accession (or the death) of eminent rulers, up to a certain point in time. In the manuscripts this table does not stand in the first place, but because it is the first to be mentioned in the text, we have put it as Table 1. This table is also mentioned in the addition enclosed in square brackets (pp. 3–4), which occurs only in Mss. CMN. That this addition did not stand in Kh.'s Arabic manuscript, and cannot have been added by Maslama either, is proved by certain sentences of it, such as: \emph{nos vero latinorum studemus utilitati, ideo et in primis tabulam suppositam praenotavimus, per quam ab eo quidem anno quo hic liber in nostrum sermonem translatus est} . . . . (but we study the advantage of the Latins, and therefore we have first of all set down the table below, by which, from the year in which this book was translated into our language ...). This addition is therefore to be regarded as a commentary or scholium of the translator or of a copyist, and it perhaps also shows that Mss. CMN may derive from a different translation from O (see my introduction).\par
\par In this addition the time elapsed from the Flood to \emph{Dhū l-Qarnayn}, that is, to the beginning of the \emph{Seleucid era}, is given as IIDCCXCIII years, II months, V days; the II before D should be written $\overline{\text{II}}$, and then means MM; the time therefore amounts to 2793 years 2 months 5 days. — Table 1 is very disordered as regards the various eras; I give the most important of them below in their direct order:\par
\begin{center}\small\begin{tabular}{lrrrl}
From Adam to the \emph{Diluvium} & 2280 years &  &  &  \\
From the \emph{Diluvium} to Nabukodonossor\footnotemark{} & 2356 years & 7 months & 23 days &  \\
From Nabukodonossor to Alexander & 436 years & 6 months & 12$\tfrac{1}{2}$ days & (?) \\
From Alexander to Christ & 311 years & 3 months &  &  \\
From Christ to the \emph{Hijra} & 621 years & 6 months & 15 days &  \\
From the \emph{Hijra} to Yazdegerd & 9 years & 11 months & 4 days &  \\
\end{tabular}\end{center}
\addtocounter{footnote}{-1}\stepcounter{footnote}\footnotetext{This is the \emph{era of Nabonassar} used by Ptol. (Batt. and others call him \emph{Bukhtanaṣṣar primus}), which begins on 26 February 747 BC; this, to be sure, does not quite agree with the figures above.}
\par \Pg{36}In this Table 1 an \emph{Aera azofar} is also named twice. This is an era that was in use in Spain among the Christian population before and after the time of \emph{Maslama}, and that began exactly 38 years before our Christian reckoning, as our Table 1 also correctly states in the fifth and the fourth lines from the end; according to the sixth and the seventh lines from the end, however, the difference between the two eras would be 38 years and 2 days, and there is probably an error here. — This Spanish era was called by the Arabs \emph{taʾrīkh al-ṣufr} (pronounced \emph{aṣ-ṣufr}); that this is connected with \emph{ṣifr} (= empty, zero), as Ideler\footnote{\emph{Handbuch der math. u. techn. Chronologie}, Berlin 1826, II, 423–425.} conjectures, is very improbable.\footnote{See also Dozy, \emph{Supplément aux dictionnaires arabes}, s. v. \emph{ṣafar} and \emph{aṣfar}.}\par
\par At the end of this addition a second table is mentioned, with which the years of the Christian reckoning (but in the text \emph{anni romani} stands as well as \emph{anni domini}) can be converted into Arab years; this table is missing in all the manuscripts (see the introduction). There is only Table 3 in CO and N, headed \emph{Pagina reducendi annos Alexandri ad elhigeram} (page for reducing the years of Alexander to the Hijra), which, however, serves in the first place to convert Arab years into years of the Seleucid era; it can, to be sure, also serve for the reverse; we shall come back to it below.\par
\UnitHead{Chapter 3. On the headings of the columns of the tables}{Chapter 3. On the headings of the columns of the tables}{K02-6}
\par\begin{center}(Text pp. 4–5.)\end{center}
\par So the title \emph{de numerorum praescriptionibus} must be translated. The author explains here what is to be understood by \emph{anni collecti}, \emph{anni plani}, \emph{gradus}, \emph{dakaicae}, etc., and how one must proceed in using the tables (as an example he takes those for determining the mean places of the planets). I have the following remarks to add to this chapter. In \emph{anni plani}, \emph{planum} is the translation of the Arabic \emph{mabsūṭ} = spread out; instead of \emph{planum} one therefore also finds (in C) \emph{expansum}; we should say \emph{anni singuli}; N has \emph{anni residui}. — The first column of the tables\footnote{After the column of the years.} for determining the mean places of the planets usually contains the path of the planet from the \emph{radix} in periods of 30 years; then follows the column for the single years from 1 to 30, then the column for the months, and so on. For the word “minute” C and O curiously often use the Arabic \emph{daqīqa}, in its plural form \emph{daqāʾiq}; this is rendered in Latin by \emph{dakaica}, which is then declined like a genuine Latin word (for example \emph{dakaicarum}). On the other hand, especially in Ms. C, the Arabic singular form \emph{thenia} = \emph{thāniya} is put for “second”, and from it the plural \emph{theniae} is formed, instead of the Arabic plural form \emph{thawānī} being used. For the rest, Ms. C mostly has \emph{puncta} instead of \emph{minutae} and \emph{minuta} instead of \emph{secundae}\footnote{I write \emph{minutae} and \emph{secundae} instead of \emph{minuta} and \emph{secunda}, taking these words as adjectives to \emph{partes}; the correct form would really be (\emph{partes}) \emph{primae}, (\emph{partes}) \emph{secundae}. [\emph{Translator:} The print calls this note after \emph{minutae} and again after \emph{secundae}.]}. — I have already drawn attention in the introduction to the arbitrariness and inconsistency that reign in these headings of the columns.\par
\UnitHead{Chapter 4. On the beginning of the Arab months}{Chapter 4. On the beginning of the Arab months}{K02-9}
\par\begin{center}(Text p. 5, Tables 2–2$^{b}$.)\end{center}
\par \Pg{37}This chapter, to which belongs Table 2 (found in all four manuscripts, but in O on the last page, 159$^{r}$), teaches how to determine the weekday (\emph{feria}) on which any Arab month begins in any year. To make clear the rules of the text, which are often somewhat hard to understand, I give the following two examples.\par
\par 1. On what weekday did \emph{Ramaḍān} of the year 529 AH begin?\par
\begin{center}\small\begin{tabular}{lr}
for 510 years the nota (that is, the number of the weekday)\footnotemark{} is & = 1 \\
for 19 years the nota is & = 6 \\
for the month of Ramaḍān the nota is & = 6 \\
 & sum = 13 \\
\end{tabular}\end{center}
\addtocounter{footnote}{-1}\stepcounter{footnote}\footnotetext{1 = Sunday, 2 = Monday, and so on . . . 7 = Saturday.}
\par From this sum 7 is subtracted (or it is divided by 7); the remainder is 6, that is, 1 \emph{Ramaḍān} was a \emph{Friday}.\par
\par 2. At its end this chapter states that 26 January of the year 1126 (in the Julian reckoning) fell on 1 \emph{Muḥarram} of the year 520 AH and was a \emph{Tuesday} (\emph{feria tertia}); let us prove the correctness of this statement too from Table 2:\par
\begin{center}\small\begin{tabular}{lr}
for 510 years the nota is & = 1 \\
for 10 years the nota is & = 1 \\
for the month of Muḥarram the nota is & = 1 \\
 & sum = 3 \\
\end{tabular}\end{center}
\par that is, 1 \emph{Muḥarram} of the year 520 was a \emph{Tuesday}. — This statement is found only in Ms. O, and probably means to say that the translation of this work into Latin was either begun or finished on that day; the conclusion would then suggest itself that Ms. O may represent, if not Adelard's autograph, at least a copy of it made soon afterwards.\par
\par In Table 2 the designation \emph{keb.} (= \emph{kebice} = leap year) stands several times in the wrong place, and a leap year is missing as well, since their number within 30 years is 11.\footnote{See above, p. 34.}\par
\par I add here a mention of two tables that are not named in the text at all, and that I have designated 2$^{a}$ and 2$^{b}$. The first table teaches with which weekday the months of the \emph{era} of Yazdegerd begin. One will easily be able to determine the weekday if we note that in the first column of the section headed \emph{Signa vel Notae} this table gives in each case the weekday of the beginning of the months of the first year of the era: thus 1 \emph{Affrordinme} of the year 1 began with a Tuesday (3), 1 \emph{Azdihestme} with a Thursday (5), and so on. For each following year the weekday moves on by one, since the Persian year had 365 days. Since every Persian month had 30 days, except \emph{Abenme}, in which the 5 days beyond 360 were intercalated, so that it came to 35 days, that is, 5 full weeks, the following month, \emph{Adarme}, has the same \emph{notae} as \emph{Abenme}.\footnote{Al-Bīrūnī gives instructions for determining the weekday of the beginning of each month for any year of this era in his \emph{Chronology} (English translation by Sachau, p. 184; Arabic text p. 203).} — Table 2$^{b}$ contains the \emph{Egyptian} (\emph{Coptic}) \Pg{38}months, which Ptol. too used for his astronomical data, with their \emph{Signa} or \emph{Notae}. — Both the Persian and the Coptic names are rather disfigured in the two Mss. M and N, which alone contain the two Tables 2$^{a}$ and 2$^{b}$; we have entered the correct ones in the glossary, after Nallino\footnote{Batt., ed. Nallino, I, 67 and 243.} and al-Bīrūnī's \emph{Chronology}\footnote{English translation by Sachau, pp. 82 and 199–217; Arabic text pp. 215–233.}.\par
\UnitHead{Chapter 5. On finding the Roman year}{Chapter 5. On finding the Roman year}{K03-1}
\par\begin{center}(Text pp. 6–7, Tables 3 and 3$^{a}$.)\end{center}
\par The title is really \emph{De anni bissextilis inventione} (on finding the bissextile year). But in fact it is a matter of converting the years of the \emph{Hijra} into years of the \emph{Seleucid era}, or of finding the \emph{Roman} date that corresponds to a given \emph{Arab} date. The years of the Roman era were also called \emph{anni bissextiles} (leap years), which seems somewhat strange, since the Arabs too had leap years; but this \emph{anni bissextiles} stands in contrast to the \emph{anni vagi} of the Egyptians, who knew no \emph{intercalation}. The corresponding Table 3 is headed \emph{Pagina reducendi annos Alexandri ad elhigeram} (page for reducing the years of Alexander to the Hijra), which is really incorrect. The table is arranged somewhat differently from the corresponding tables of Batt.\footnote{Ed. Nallino, II, 9–18.}, which give, for the beginning of each year of the \emph{Hijra} from 1 to 600, the corresponding date of the \emph{Roman era}; Kh.'s table takes up much less room, but still requires subsidiary computations. We give three examples of this reduction.\par
\par 1. To which day of the \emph{Roman era} does \emph{1 Muḥarram} of the year 520 AH correspond?\par
\begin{center}\small\begin{tabular}{lrrrl}
At the end of the 510th year AH there had elapsed & 1427 y. & 7 m. & 5$\tfrac{1}{4}$ d. & of the Roman era \\
9 Muslim years are & = 8 y. & 8 m. & 27 d. & of the Roman era \\
sum: & 1436 y. & 3 m. & 32$\tfrac{1}{4}$ d. &  \\
\end{tabular}\end{center}
\par Since a conversion of 12 months into one year has taken place once, one must subtract 5$\tfrac{1}{4}$ days from this sum (for if one takes a year for 12 months, one has taken 5$\tfrac{1}{4}$ days too many)\footnote{Since the months are reckoned at only 30 days.}; this gives 1436 y. 3 m. 27 d. Since among the 3 months from October onwards there are two months of 31 days, one must subtract 2 more from the number of days; this gives 25 days. That is, since the beginning of the Roman era (1 Oct. 312 BC) 1436 y. 3 m. 25 d. have elapsed; that is, \emph{1 Muḥarram} 520 AH corresponds to \emph{26 Kānūn II} of the year 1437, that is, to \emph{26 January} 1126 AD.\par
\par 2. To which day of the \emph{Roman era} does \emph{5 Rajab} of the year 413 AH correspond?\par
\Pg{39}\begin{center}\small\begin{tabular}{lrrrl}
At the end of the 390th year AH there had elapsed & 1311 y. & 2 m. & $\tfrac{1}{4}$ d. & of the Roman era \\
22 Muslim years are & = 21 y. & 4 m. & 5$\tfrac{3}{4}$ d. & of the Roman era \\
the first 6 Muslim months are & = & 5 m. & 27 d. & of the Roman era \\
and in Rajab a further &  &  & 4 d. &  \\
sum: & 1333 y. & 0 m. & 7 d. &  \\
\end{tabular}\end{center}
\par Since again a conversion of 12 months into one year has taken place once, one must take 5$\tfrac{1}{4}$ days from this sum; this gives 1333 y. 0 m. 2 d. That is, from the beginning of the Roman era to 5 Rajab 413 AH, 1333 y. 0 m. 2 d. have elapsed; or 5 Rajab 413 corresponds to 3 \emph{Tishrīn I} of the year 1334, that is, to \emph{3 October} 1022 AD.\footnote{If the date falls in the months October to December, one must subtract 312 from the Roman year number to obtain the Christian (Julian) year number, otherwise 311; the reader will easily find the explanation.}\par
\par 3. To which day of the \emph{Roman era} does \emph{3 Shawwāl} of the year 399 AH correspond?\par
\begin{center}\small\begin{tabular}{lrrrl}
At the end of the 390th year AH there had elapsed & 1311 y. & 2 m. & $\tfrac{1}{4}$ d. & of the Roman era \\
8 Muslim years are & = 7 y. & 9 m. & 8$\tfrac{1}{4}$ d. & of the Roman era \\
the first 9 Muslim months are & = & 8 m. & 26 d. & of the Roman era \\
and in Shawwāl a further &  &  & 2 d. &  \\
sum: & 1319 y. & 7 m. & 36$\tfrac{1}{2}$ d. &  \\
\end{tabular}\end{center}
\par Since again a conversion of 12 months into one year has taken place once, one must take 5$\tfrac{1}{4}$ days from this sum; this gives 1319 y. 7 m. 31$\tfrac{1}{4}$ d. Since among the 7 months from October onwards there are four months of 31 days, one must subtract 4 more from the number of days; this gives 1319 y. 7 m. 27$\tfrac{1}{4}$ d. Since, further, February is among these 7 months, one must add 2 to the number of days; this gives 1319 y. 7 m. 29$\tfrac{1}{4}$ d. That is, from the beginning of the Roman era to 3 Shawwāl 399 AH, 1319 y. 7 m. 29 d. have elapsed; or \emph{3 Shawwāl} 399 corresponds to \emph{30 Hiyār} (the correct form is \emph{Ayyār}) 1320, that is, to \emph{30 May} 1009.\par
\par By the modern comparative tables one obtains 27 January for the first example, 4 October for the second and 31 May for the third; Kh. or Maslama, then, like Batt., put the beginning of the Muslim era on 15 July 622 instead of the 16th.\par
\par \Pg{40}Table 3$^{a}$ serves to determine whether the Roman year found is a leap year or not, and with which weekday it and its months begin. The columns headed with the names of the months contain the \emph{Signa} (\emph{Notae}). It should be noted that, if one wants the beginning of a month of a given year in the \emph{Julian} reckoning, one has to add 311 (312) to this year number.\footnote{See the preceding note.} Further, one must take into account that in the Julian calendar the days of a month fall again on the same weekdays as before after 28 years have passed, so that in Table 3$^{a}$, under \emph{Anni romanorum plani}, only the remainders are entered that result from dividing the year number by 28; instead of the remainder 0 the number 28 stands at the end of the column. To spare this division each time, Table 3 has a column (the last) with the multiples of 28; the same multiples, in considerably greater number, are also found in Table 1. I give three examples.\par
\par 1. What weekday was 5 March (Syr. \emph{ādhār}, in our tables written adar) of the year 1363 of the \emph{Roman} era?\par
\par In Table 1 one finds the nearest number below 1363 that is divisible by 28 to be 1344; this taken from 1363 leaves 19; in Table 3$^{a}$ one finds against 19, in the column headed “\emph{Mart.}”, the \emph{nota} 1, that is, 1 March fell on a Sunday, the 5th, then, on a \emph{Thursday}. This year of the Roman era was a leap year, for in the column headed \emph{Anni bissextiles} there stands \emph{bissext.}; and so it was, for it corresponds to the year 1052 of the Julian reckoning.\par
\par 2. What weekday was 26 January 1126 of the \emph{Julian} reckoning? We add the number 311 to 1126; this gives 1437; in Table 3 one finds the nearest number divisible by 28 to be 1428; this taken from 1437 leaves 9; in Table 3$^{a}$ one finds against 9, in the column headed January, the \emph{nota} 6, that is, 1 January fell on a Friday, and so did the 8th, the 15th and the 22nd; the 26th, then, on a \emph{Tuesday}. This result agrees with the statement on p. 5 of the text (see also above, p. 37).\par
\par 3. What weekday was 10 October of the year 1912 of the \emph{Julian} reckoning?\par
\par We add the number 312 to 1912 (see p. 39); this gives 2224; the nearest number divisible by 28 is 2212 (this number no longer stands in our tables), remainder 12; to it belongs, in Table 3$^{a}$ under October, the nota 2; so 1 October fell on a Monday, and so did the 8th; the 10th, then, on a \emph{Wednesday}.\par
\par For the correct readings of the Syrian names of the months, which stand at the top of the headings of the individual columns, I refer to the glossary.\par
\par These first six tables (1, 2, 2$^{a}$, 2$^{b}$, 3, 3$^{a}$) correspond, some wholly, others only in part, to the spurious tables, that is, those not by Batt., that Nallino added at the end of his edition (II, pp. 300–305), and which certainly come from Maslama al-Majrīṭī, as is stated directly of some of them. This of course formed an essential reason for our conjecture that in this work we have before us not Kh.'s original tables but Maslama's reworking. The correspondence of these tables is as follows:\par
\begin{center}\small\begin{tabular}{lcl}
Batt. p. 301 (first half) & = & Kh. Table 2 \\
Batt. p. 301 (second half) & = & Kh. Table 2$^{a}$ \\
Batt. p. 302 & = & Kh. Table 3$^{a}$ and the left half of 1 \\
\end{tabular}\end{center}
\Pg{41}\begin{center}\small\begin{tabular}{lcl}
Batt. p. 303 & = & Kh. Table 2$^{b}$ \\
Batt. pp. 304–305 & = & \begin{tabular}[b]{@{}c@{}}Kh. Table 3 (but here only the \emph{anni Romanorum}\\are taken in, not the \emph{anni Coptorum}\\and \emph{Persarum} (see the introduction))\end{tabular} \\
\end{tabular}\end{center}
\UnitHead{Chapter 6. On the division of the sphere}{Chapter 6. On the division of the sphere}{K04-1}
\par\begin{center}(Text p. 7.)\end{center}
\par This short chapter needs no commentary.\par
\UnitHead{Chapter 7. On the mean motion of the planets}{Chapter 7. On the mean motion of the planets}{K04-4}
\par\begin{center}(Text pp. 7–9, Tables 4–20.)\end{center}
\par This chapter with its tables holds a principal place in every astronomical work of the Greeks and the Arabs. With its help one computes the mean places of the sun, the moon and the five planets: first, with the column headed \emph{anni collecti}, the places for every 30 years, then for the \emph{anni plani} or \emph{singuli}, then for the months, the days and the hours, and even for every 2 minutes; for each of these seven heavenly bodies, then, one can determine its mean longitude, that is, its distance from the vernal point, for any moment of time, down to two minutes. Since the position of the nodes of the lunar orbit was important for determining the latitude of the moon, and for astrology as well, the mean motion of the northern or ascending node of the moon (Arab. \emph{al-ʿaqd al-shamālī}) is given besides those of the sun, the moon and the five planets. In our tables it is designated by the Persian-Arabic word \emph{gavzehar} (= \emph{jawzahar}), in Latin by \emph{draco} (\emph{caput draconis} = the dragon's head would be better).\par
\par In Batt. the mean motion is called \emph{al-masīr al-wasaṭ} (= the mean course), or just \emph{al-wasaṭ} (the mean), or even just \emph{al-masīr} (course, path, journey); in our tables there stands partly the Arabic \emph{wasaṭ}, written \emph{wazat} and \emph{wacat}, partly \emph{medialitas} or \emph{mediocritas}, partly \emph{medius cursus}.\par
\par The mean place is given in each case in signs of the zodiac (\emph{Signa}), degrees, minutes and seconds; Batt. has converted the signs into degrees, and he neglects the seconds for the five planets, while he gives them for the sun and the moon.\par
\par \Pg{42}The \emph{radix} entered at the beginning of the \emph{anni collecti} gives, as was remarked above, the longitude of the planets at noon of 1 \emph{Muḥarram} of the year 1 AH for the meridian of \emph{Ujjain}; but even if one takes into account the difference in longitude between \emph{Ujjain} (c. 75° 52′ from Greenwich) and \emph{Raqqa} (c. 39° 4′ from Greenwich), the radix still does not agree with that of Batt. Our tables have, for example, the radix 113° 25′ 48″ for the sun, Batt. 113° 58′ 4″; at the beginning of the year 511 AH the difference amounts to c. 9° 15′; these are surely the 9° of which Ibn Ezra speaks in the preface to the translation of the \emph{Reasons of the tables} of al-Khwārizmī\footnote{See Steinschneider, \emph{Zur Geschichte der Übersetzungen aus dem Indischen ins Arabische}, in Zeitschr. d. deutschen morgenländ. Gesellsch., vol. 24, p. 355.}. In the longitudes of the planets too we find this difference of 8$\tfrac{1}{2}$ to 9$\tfrac{1}{2}$° again, except for Mars, where the difference is 11$\tfrac{1}{2}$°; there are probably copyists' errors here. These differences come from the neglect of the \emph{precession} of the vernal point: indeed Kh., unlike Batt., took as the basis for determining the mean course of the sun the \emph{sidereal} year and not the \emph{tropical} year; for when we determine from Tables 4 and 5 the time the sun needs to run through a whole circle (= 360°), we find:\par
\begin{center}\small\begin{tabular}{rrrr}
in 12 · 30 = & 360 days & = 12 · (29° 34′ 5″) = & 354° 49′ \\
in & 5 days & = & 4° 55′ 41″ \\
in & 6 hours & = & 14′ 47″ \\
in & 13 min. & = & 32″ \\
so in & 365$^{d}$ 6$^{h}$ 13$^{m}$ & = & 360° 0′ 0″ \\
\end{tabular}\end{center}
\par Kh., then, took the length of the sidereal year as 365$^{d}$ 6$^{h}$ 13$^{m}$; in the \emph{Sūrya-Siddhānta}\footnote{E. Burgess, \emph{Translation of the S. S.}, in the Journal of the American Oriental Society, Vol. VI, 1860, p. 161.} we find for the length of the sidereal year 365$^{d}$ 6$^{h}$ 12.6$^{m}$, which agrees roughly with our computation\footnote{See below, chapters 21 and 22.}. — Kh. does not mention the tropical year at all in his work, and it seems indeed that the Indians too, up to and including Brahmagupta, did not know it. What is said in the \emph{Sūrya-Siddhānta}\footnote{Transl. Burgess, l. c., pp. 244–249.} about the precession is, according to Burgess and Thibaut\footnote{\emph{Astronomie, Astrologie u. Mathematik}, in the Grundriss der Indo-Arischen Philologie u. Altertumskunde, vol. 3, part 9, pp. 45–46.}, a later interpolation; on the other hand Āryabhaṭa, at the end of the fifth century AD, did teach the \emph{trepidation} of the fixed stars, or of the vernal point.\par
\par As examples of the use of the tables we give the computation of the mean place of the sun and of the moon for noon of \emph{15 Ṣafar} of the year 540 AH.\par
\begin{center}\small\begin{tabular}{llcr}
\emph{1. Sun}: & longitude at the end of the 510th year & = & 38° 35′ 45″ \\
 & increase in 29 years & = & 49° 2′ 50″ \\
 & increase in 1 month & = & 29° 34′ 5″ \\
 & increase in 14 days & = & 13° 47′ 54″ \\
 & sum & = & 131° 0′ 34″ \\
\end{tabular}\end{center}
\par So the mean longitude of the sun at the time named was 131° 0′ 34″.\par
\begin{center}\small\begin{tabular}{llcr}
\emph{2. Moon}: & longitude at the end of the 510th year & = & 40° 56′ 38″ \\
 & increase in 29 years & = & 53° 24′ 16″ \\
 & increase in 1 month & = & 35° 17′ 26″ \\
 & increase in 14 days & = & 184° 28′ 8″ \\
 & sum: & = & 314° 6′ 28″ \\
\end{tabular}\end{center}
\par So the mean longitude of the moon at the time named was 314° 6′ 28″.\par
\par \Pg{43}At the end of the rule the text remarks how one is to proceed when the place for which one wants to determine the mean longitude of a planet lies east or west of \emph{Arin}.\par
\par At the end of the chapter there is an addition, which in C and O stands only in the margin, but in M in the text, and which speaks of the relation of degrees and miles: it is stated that on the earth 66$\tfrac{2}{3}$ miles go to one degree, so that the circumference of the earth is 24000 miles; at the beginning it is remarked that according to the Chaldeans 4000 camel's paces make a mile.\par
\par This mile, of which 66$\tfrac{2}{3}$ are said to go to a degree, is the mile of 4000 cubits (\emph{dhirāʿ}), the cubit of 24 fingers (\emph{aṣbaʿ} or \emph{iṣbaʿ}), which the Arab geographers and astronomers commonly used in their data; according to Nallino\footnote{\emph{Il valore metrico del grado di meridiano secondo i geografi arabi}, in the Cosmos di Guido Cora, Vol. XI, 1892–93, fascic. I–IV, and \emph{Estratto}, Turin, 1893.} such a mile was c. 1973 m. After the measurement of a degree under al-Maʾmūn the better value of 56$\tfrac{2}{3}$ miles = 1° came into use beside this one.\par
\par M. Chasles\footnote{\emph{Recherches sur l’astronomie ancienne}, in Comptes-rendus de l’Académie des Sciences, t. XXIII, Paris, 1846, p. 851.} concludes from the above passage that the \emph{Chaldeans} (Babylonians) already knew the nearly correct circumference of the earth. We do not want to dispute this absolutely; but, first, the statement that a mile of c. 2000 metres takes 4000 camel's paces is already somewhat suspicious, and, secondly, by \emph{Caldei} the \emph{Arabs} are most probably meant here, for we hold this passage to be an interpolation of the translator or of copyists: on p. 11 of the text the word \emph{Caldei} occurs once more, in a passage that is undoubtedly an interpolation, and there it certainly designates the Arabs; and Ms. O has in fact “arabice” instead of “a Caldeis”.\par
\UnitHead{Chapter 8. On finding the (true) place of the sun}{Chapter 8. On finding the (true) place of the sun}{K04-21}
\par\begin{center}(Text pp. 9–10, Tables 21–26.)\end{center}
\par To determine the true place of the sun from the mean one, one must, according to the precept of this chapter, subtract from the mean longitude of the sun the longitude of its \emph{apogee}, which in these tables is taken as 77° 55′\footnote{When the numbers are expressed by Arabic letters, this number can easily be confused with 77° 15′, and the value would then agree with that of the later redaction of the \emph{Sūrya-Siddhānta} and of the Arabic \emph{Sindhind} (according to the testimony of Ibn Hibintā) (see Batt., ed. Nallino, I, 218).}; this gives the so-called \emph{anomaly} of the sun. With this as the argument one looks up in Tables 21–26 the so-called \emph{equation} of the sun; this equation must be subtracted from the mean longitude of the sun, or added to it, according as the anomaly is less or greater than 180°; the result gives the \emph{true longitude} of the sun.\par
\par \emph{Example}: In the preceding chapter we found the mean longitude of the sun for noon of \emph{15 Ṣafar} of the year 540 AH = 131° 0′ 34″. If we subtract from it the \Pg{44}\emph{apogee} = 77° 55′, we obtain the \emph{anomaly} = 53° 5′ 34″; with this argument one enters Table 22 and finds there, by interpolation, 1° 46′ for the \emph{equation} of the sun; this equation, subtracted from the mean longitude, gives 129° 14′ 34″ as the \emph{true place} of the sun, that is, at that moment the sun was in 9° 14′ 34″ of Leo.\par
\par The explanation of how one is to proceed with any minutes of the anomaly is corrupt in our text (pp. 9–10: \emph{Quod si ab argumento partes aliquae restiterint} . . ., but if some parts remain over from the argument ...); it is of course the same computation of proportional parts to be carried out as with our logarithmic-trigonometric tables.\footnote{See also Batt., ed. Nallino, I, 72, ll. 9–13.}\par
\par We give below the geometrical explanation of the above procedure for determining the true longitude of the sun according to Ptol.-Batt., namely according to the eccentric hypothesis, that is, that the sun moves in a circle eccentric to the ecliptic, not on an epicycle. Ptol., moreover, already showed that one arrives at exactly the same result with the help of an epicycle.\par
\par Let \emph{E} be the centre of the ecliptic, thus the centre of the earth, and \emph{M} the centre of the eccentric circle in which the sun is supposed to move uniformly\footnoteT{Here the print has a figure, which is not reproduced: two circles of the same radius, the eccentric with the centre \emph{M}, on which lie \emph{A} at the left end of a dashed horizontal diameter, \emph{F} at the upper left, \emph{G} and \emph{S} at the lower left and \emph{P} at the right end of the diameter, and the circle of the ecliptic with the centre \emph{E}, which passes through \emph{P}′ at the right end of the same dashed line; solid lines \emph{FM}, \emph{SM} produced to \emph{H}, \emph{HE} and \emph{SE}; a dashed line \emph{GM} parallel to \emph{SE}; the angle \emph{{\greekfont α}} is marked at \emph{S}, between \emph{SM} and \emph{SE}, and at \emph{M}, between \emph{MG} and \emph{MS}.}. \emph{A} is then the apogee and \emph{P} the perigee of the sun; let the sun be at the point \emph{S}, and let \emph{F} be the vernal point. Now {\SignFont\upshape ∢} \emph{FMS} is the mean longitude of the sun; from it one subtracts the longitude of the apogee, that is, {\SignFont\upshape ∢} \emph{FMA}, and then has {\SignFont\upshape ∢} \emph{AMS}, the so-called \emph{mean anomaly}, usually called \emph{argument} in our tables. But we want {\SignFont\upshape ∢} \emph{AES}, under which the arc \emph{AS}, which the sun has run through from the apogee, appears as seen from the earth (\emph{E}); we obtain this angle if we subtract from the mean anomaly {\SignFont\upshape ∢} \emph{{\greekfont α}}, the so-called \emph{equation} of the sun. {\SignFont\upshape ∢} \emph{AES}, or \emph{AMG} if \emph{MG} {\SignFont\upshape ∥} \emph{ES}, is now called the \emph{true anomaly}; if one again adds to it {\SignFont\upshape ∢} \emph{AMF}, that is, the longitude of the apogee, one has {\SignFont\upshape ∢} \emph{FMG} or the arc \emph{FG}, that is, the \emph{true longitude} of the sun. — In our text the equation \emph{{\greekfont α}} is subtracted directly from the mean longitude of the sun, which of course comes to the same thing.\par
\par How one computes the equation for each arc \emph{AS}, or for each mean anomaly, I will not set out in detail here; I only remark that Ptol.\footnote{\emph{Almagest} III, 5 (ed. Heiberg, I, 240–255; German translation by K. Manitius, I, 173–180).} and Batt.\footnote{Ed. Nallino, I, 48, 219–220.} compute it by means of the Pythagorean theorem and trigonometric relations from the triangles \emph{HME} and \emph{HES}, whereas al-Bīrūnī\footnote{See \emph{Biblioth. mathem.} 11$_{3}$, 1910, pp. 45 and 71.} first computes \emph{ES} from the triangle \emph{EMS} with the help of the cosine theorem, and then {\SignFont\upshape ∢} \emph{{\greekfont α}} from the right-angled triangle \emph{HES}.\par
\par \Pg{45}The values of the \emph{equation} of the sun, moreover, do not agree with those of Batt. (ed. Nallino, II, 78–83), which are computed by the formula $\mathrm{tg\ aeq.} = \dfrac{e\sin m}{1+e\cos m}$; already for 1° of mean anomaly the difference is 16″. This can only come from a different value of the eccentricity \emph{e} of the sun's orbit. Batt. takes it (thus \emph{ME} in the preceding figure) as 2$^{p}$ 4$\tfrac{3}{4}$′ (when \emph{MA} is put = 60$^{p}$); in Kh. there is no value for the eccentricity of the sun; from the maximum of the equation (2° 14′) at 90° we find \emph{e} = 2$^{p}$ 20′, and in other tests that we have made we found it = 2$^{p}$ 22′; the \emph{Sūrya-Siddhānta}, according to Burgess\footnote{l. c., p. 220.}, has a maximum equation of 2° 10′ 31″, which corresponds to an eccentricity \emph{e} = 2$^{p}$ 17′; Ptol. takes an eccentricity of 2$^{p}$ 29′ 30″, but sometimes rounds it up to 2$^{p}$ 30′. We see, then, that Kh. follows neither the \emph{Sūrya-Siddhānta} nor Ptol.; perhaps he took his value from the works of Brahmagupta, or from the Persian tables (see above, p. 32).\par
\UnitHead{Chapter 9. On finding the (true) place of the moon}{Chapter 9. On finding the (true) place of the moon}{K05-6}
\par\begin{center}(Text p. 10, Tables 21–26.)\end{center}
\par In Tables 6–8 one finds, besides the mean longitude of the moon, its anomaly too, in an analogous way; the column that serves for this is headed \emph{Argumentum lunae} in our tables. Kh.'s determination is very simple compared with Batt.'s: the former takes the equation of the moon from a single column, whereas the latter needs four of them, since he takes into account the inequalities of the moon's motion, the so-called \emph{evectio} and \emph{prosneusis}.\footnote{For Batt.'s determination compare the rules of the text in Nallino, I, 50–54 and 75, and the explanations of them in II, 223–226.} Let us determine, by our rules and tables, the true place of the moon for noon of 15 Ṣafar of the year 540 AH.\par
\begin{center}\small\begin{tabular}{llcr}
\emph{Example}: & the \emph{anomaly} at the end of the 510th year is & = & 61° 52′ 40″ \\
 & increase for 29 years & = & 348° 56′ 57″ \\
 & increase for the first month, Muḥarram & = & 31° 56′ 59″ \\
 & increase for 14 days & = & 182° 54′ 36″ \\
 & sum: &  & 625° 41′ 12″ \\
 &  &  & − 360° \\
 & \emph{anomaly} & = & 265° 41′ 12″ \\
\end{tabular}\end{center}
\par For this one finds in Table 24 the \emph{equation} 4° 55′ 7″\footnote{For 265° one finds 4° 54′ 52″, and for the 41′ a further 15″ results.}; it is added to the mean longitude, which we found above to be 314° 6′ 28″, because the anomaly is greater than 180°; this gives 319° 1′ 35″ for the \emph{true place} of the moon, that is, at the given time the moon was in 19° 1′ 35″ of Aquarius.\par
\par \Pg{46}By the tables of Batt. one finds 330° 41′, for they give the mean longitude at the end of the 510th year as 50° 44′ 20″ instead of the 40° 56′ 38″ of our tables, and the equation is found to be 6° 19′ 57″ instead of 4° 55′ 7″.\par
\par How Kh.'s tables of the moon are computed has not become clear to me. The greatest equation amounts to 4° 56′ at 90° of mean anomaly, which corresponds to a radius of the epicycle of c. 5$^{p}$ 10′ instead of the 5$^{p}$ 15′ of Batt.; but if one computes the equation for other anomalies (\emph{m}) by the formula\par
\par\smallskip\[\mathrm{tg\ aeq.\ anom.} = \dfrac{e\sin m}{1+e\cos m},\]\par\smallskip
\par in which \emph{e} is then to be put = $\dfrac{5^{p}\,10'}{60^{p}}$ = 0.08611, one finds no agreement; one or both of the inequalities named above must therefore surely have been taken into account after all, but in what way we could not find out.\par
\par In the headings of the tables and of the columns we find for the Arabic \emph{taʿdīl} = \emph{equation} now the Arabic word in the form \emph{Tadil}, now the Latin \emph{Examinatio}; \emph{Aequatio} does not occur; instead of this word N has \emph{Coaequatio}. The word \emph{elheza}, also written \emph{elheca}, is the Arabic \emph{al-khāṣṣa} (also written \emph{al-ḥāṣṣa}) = the (motion) peculiar (to the planet); correct would be \emph{al-ḥaraka al-khāṣṣa} = the motion peculiar (to it), that is, the so-called \emph{anomaly}, that is, the arc that the moon or the planet has run through on the epicycle from its apogee up to a given moment. In our work this word is translated, especially in Ms. C, by \emph{Argumentum}; it is confused with \emph{al-ḥiṣṣa} = share, lot, that is, the argument with which one enters a table.\par
\UnitHead{Chapter 10. Determination of the (true) place of the three upper planets}{Chapter 10. Determination of the (true) place of the three upper planets}{K05-16}
\par\begin{center}(Text pp. 10–11, Tables 27–44.)\end{center}
\par Let us first consider the \emph{apogees} of the planets and of the sun; for them we find, in the text pp. 14–15 and in Tables 27–44, the following values:\par
\begin{center}\small\begin{tabular}{cccccc}
\emph{Saturn} & \emph{Jupiter} & \emph{Mars} & \emph{Venus} & \emph{Mercury} & \emph{Sun} \\
244° 55′ & 172° 32′ & 128° 24′ & 81° 15′ & 224° 54′ & 77° 55′. \\
\end{tabular}\end{center}
\par These numbers, even when the \emph{precession} is taken into account, agree neither with those of Ptol. nor with those of Batt., nor with the values in the \emph{Sūrya-Siddhānta} and in the \emph{Pañcasiddhāntikā}.\footnote{See Burgess, l. c., pp. 216 and 427; also \emph{The Panchasiddhāntikā of Varāha Mihira}, ed. by G. Thibaut, etc., Benares 1889, p. 93 of the English translation.} In particular the values for Jupiter, Mercury and Venus differ considerably from one another in the works named. The value for the apogee of the sun agrees, if we read 15′ for 55′, as we have already remarked above (p. 43), with the value that the Arabic \emph{Sindhind} gave according to the testimony of Ibn Hibintā; and the value for Jupiter likewise with that of the “tested tables” according to Ibn Yūnus\footnote{Batt., ed. Nallino, I, 218 and 239–242.}. The value for Mercury differs most from those of the other tables, in which it does not go beyond 202°; but the \emph{Sūrya-Siddhānta} is to be excepted from this, in which \Pg{47}the apogee of Mercury is given as 220° 28′, which is still c. 4$\tfrac{1}{2}$° less than in Kh. The agreement of the apogee of Jupiter in Kh. and in the “tested tables” strengthens us in the view already expressed above (p. 32), that certain data of Kh.'s tables may have been taken from the “tested tables” of the astronomers of al-Maʾmūn. Kh., to be sure, departs again from these tables (if at least one can rely on the reports of Ibn Yūnus about the “tested tables”) in not taking the longitude of the apogee of Venus as equal to that of the apogee of the sun.\par
\par For a better understanding of these questions I give below a short account of the motion of an upper planet according to Ptol. and Batt.; I follow Nallino as closely as possible.\footnote{Ibid., II, 238–241.}\par
\par \Pg{48}It is assumed that the planet moves on an \emph{epicycle} whose centre\footnote{Since German and Latin designations will often alternate in what follows, I always write “Centrum” instead of “Zentrum”, for the sake of uniformity. [\emph{Translator:} In the English, centre.]} \emph{H} moves with unequal angular velocity, both with respect to the point \emph{M} and with respect to the point \emph{E}, on the circle with the centre \emph{M} that is eccentric with respect to the earth \emph{E}, the so-called \emph{deferent}. But the centre \emph{H} moves uniformly on a third, fictitious circle (I leave it out of the drawing) with the centre \emph{E′} (\emph{E′M} = \emph{ME}), the so-called \emph{equant}; or better, the motion of the point \emph{H} appears uniform as seen from this point \emph{E′}. Both the planet on the epicycle and the centre of the epicycle on the deferent move from west to east. The two uppermost planets (Jupiter and Saturn) complete their circuit on the epicycle in about one year, Mars, on the other hand, in c. 2 years (780 days); the centres of their epicycles complete their circuit on the deferent in different times, which agree fairly well with the present sidereal periods.\footnoteT{The print has: «The upper planets complete their circuit on the \emph{epicycle} in a sidereal year, and the centres of their epicycles their circuit on the \emph{deferent} in different times.» The wording given here is that of Suter's addendum to the corrections, a slip printed after p. 236.} \emph{EM} is the eccentricity of the orbit of the planet, \emph{EE′} twice that. The \emph{true apogee} of the epicycle is, for every position of the point \emph{H}, the point at which \emph{EH} produced meets the ecliptic (in our figure, then, \emph{R})\footnoteT{Here the print has a figure, which is not reproduced: a vertical line \emph{AB}, the diameter of a large outer circle, the ecliptic \emph{AFB}, with \emph{F} at its upper right and \emph{R} and \emph{Q} at its upper left; on the line lie, from the top, \emph{A}, \emph{G}, \emph{E}′, \emph{M}, \emph{E}, \emph{K} and \emph{B}; the inner circle, the deferent, passes through \emph{G} at its top and \emph{K} at its bottom; the small circle of the epicycle about \emph{H} stands on the deferent at the upper left; the line from \emph{Q} through \emph{H} and \emph{E}′ runs on to \emph{S} and \emph{N}; the lines from \emph{R} through \emph{H} to \emph{E}, and \emph{HM}; \emph{MS} and \emph{EN} are the perpendiculars from \emph{M} and \emph{E} to \emph{HE}′ produced.}, the \emph{mean apogee} the point at which \emph{E′H} produced meets the ecliptic (in our figure, then, \emph{Q}). The two points coincide when the centre of the epicycle is at the apogee \emph{G} or at the perigee \emph{K} of the deferent. Let \emph{AFB} be the ecliptic, and let the apogee \emph{G} have the longitude \emph{FA} in the ecliptic; the centre \emph{H} of the epicycle has, by its mean motion about the centre \emph{E′} of the equant, the longitude \emph{FQ}, but by its true motion, as seen from the centre of the earth \emph{E}, the longitude \emph{FR}; the difference between the two, the mean and the true longitude, that is, the arc \emph{RQ}, is called by Batt. the \emph{equation of the anomaly and of the centre}, or just the \emph{equation of the centre}; it is also the difference between the \emph{mean} and the \emph{true apogee} of the epicycle. — This equation is measured by {\SignFont\upshape ∢} \emph{E′HE}; if one draws \emph{MS} and \emph{EN} perpendicular to \emph{QHE′} produced, one has from the right-angled triangle \emph{ENH}:\par
\par\begin{center}$\sin E'HE = \dfrac{EN}{EH}$, or $\mathrm{tg}\,E'HE = \dfrac{EN}{HN}$.\end{center}
\par Now let \emph{m} = arc \emph{AQ}, that is, the difference between the mean longitude of the planet and that of the apogee of the eccentric circle, an arc that is measured by {\SignFont\upshape ∢} \emph{AE′Q} and was called by the Arabs the \emph{centrum} (\emph{al-markaz}), because it gives the position of the centre of the epicycle with respect to the apogee of the eccentric\footnoteT{The print has «because it is thought of as measured from the centre of the \emph{equant}»; the wording given here is that of Suter's additions and corrections, p. 235, no. 2.}; further, let \emph{e} be the eccentricity \emph{EM} = \emph{E′M} and \emph{{\greekfont ρ}} = \emph{HM} = the radius of the deferent; one then easily finds:\par
\par\begin{center}$\mathrm{tg\ aeq.\ centri} = \mathrm{tg}\,E'HE = \dfrac{2e\sin m}{\sqrt{\rho^{2}-\left(\dfrac{e\sin m}{R}\right)^{2}}+\dfrac{e\cos m}{R}}$,\footnote{\emph{R} here means 60$^{p}$, since the trigonometric functions of the Arabs always refer to a circle with the radius 60; in the second formula the radius of the circle is taken = 1 instead of = 60.}\end{center}
\par or, putting \emph{R} = 1 and $\dfrac{e}{\rho} = e_{1}$:\par
\par\begin{center}$\mathrm{tg}\,E'HE = \dfrac{2e_{1}\sin m}{\sqrt{1-(e_{1}\sin m)^{2}}+e_{1}\cos m}$.\footnote{If one wanted to determine {\SignFont\upshape ∢} \emph{E′HE} with the help of the sine, as Batt. very probably did, the formula would become considerably more complicated (see Batt., ed. Nallino, II, 240).}\end{center}
\par The \emph{equation of the anomaly and of the centre}, determined with the help of these formulas, gives the values that stand in column III of Batt.'s tables of the equations of the planets (ed. Nallino II, 108–137). But another correction must still be made to the longitude of the planets, one that depends on the position of the planet on the epicycle. Let the planet be at \emph{P} (see the following figure\footnoteT{The second figure of the print, not reproduced either, stands on p. 49: as on p. 47, the circle of the ecliptic with \emph{A} at its top, \emph{F} at its upper right, \emph{B} at its bottom and \emph{R}, \emph{Q}, \emph{T} on its upper left; the circle of the deferent with \emph{G} at its top and \emph{K} at its bottom; the vertical line \emph{AB} through \emph{G}, \emph{E}′, \emph{M}, \emph{E} and \emph{K}; the small circle of the epicycle about \emph{H}, on the deferent at the left, with the planet \emph{P} on it; the lines \emph{E}′\emph{HQ}, \emph{EHR} and \emph{EPT} run through the epicycle to the ecliptic; the radius \emph{HP}, and \emph{PS} perpendicular to \emph{HR}.}); the arc \emph{RT}, or {\SignFont\upshape ∢} \emph{RET}, is sought; if one draws \emph{PS} perpendicular to \emph{ER}, then\par
\par\begin{center}$\sin RET = \dfrac{PS}{EP} = \dfrac{PS}{\text{geocentric dist. of the planet}}$,\end{center}
\par but \emph{PS} = \emph{PH} · sin \emph{RHP} = \emph{r} sin \emph{{\greekfont ψ}}, where \emph{r} is the radius of the epicycle, expressed in decimal parts of the radius of the deferent or of the equant (\emph{MH} = 1), and \emph{{\greekfont ψ}} is the angle (arc) that the planet has described from the apogee of the epicycle to the place \emph{P}, that is, the \emph{anomaly}. But the distance \emph{EP} of the planet from the earth is $=\sqrt{(EH+HS)^{2}+PS^{2}} = \sqrt{(\text{Dist. centri epic.}+r\cos\psi)^{2}+(r\sin\psi)^{2}}$, so that\par
\par\begin{center}$\sin RET = \dfrac{r\sin\psi}{\sqrt{(\text{Dist. centri epic.}+r\cos\psi)^{2}+(r\sin\psi)^{2}}}$.\end{center}
\par This angle is called the \emph{equation of the anomaly}; it is the difference between the true longitude of the centre \emph{H} of the epicycle and the true longitude of the planet in the ecliptic. Expressed by the tangent, one finds more simply:\par
\par\begin{center}\Pg{49}$\mathrm{tg}\,RET = \dfrac{PS}{ES} = \dfrac{PS}{EH+HS} = \dfrac{r\sin\psi}{\text{Dist. centri epic.}+r\cos\psi}$.\end{center}
\par This \emph{equation of the anomaly} changes, as one sees, with the variable distance of the centre of the epicycle from the earth; Batt. has given in col. VI the equation of the anomaly for a mean distance; cols. V and VII serve to determine the corrections for the deviation of the other distances from the mean one; we cannot go into details here (see the example below).\par
\par As for Kh.'s tables, he leaves out in the first place cols. V and VII of Batt., by which the equation of the anomaly, which depends on the geocentric distance of the centre of the epicycle, is corrected. On the other hand Kh. has a column headed \emph{Sublimatio Saturni examinata}, or better \emph{Apogeum Saturni aequatum}, whose values vary between 8$^{s}$ 4° 55′ (= 244° 55′ = the real apogee of Saturn) and 8$^{s}$ 2° 3′ (= 242° 3′). This, then, is surely the point to which the statements of some Arab authors about Kh.'s tables refer. I have already mentioned above (p. 33) that Ibn al-Ādamī reported that Kh. computed his equations (of the planets) by the method of the Persians (thus probably by the tables of the \emph{Shāh}); I must now add here that Abū Naṣr Manṣūr b. ʿAlī b. ʿIrāq (c. 950–1000)\footnote{See Suter, in \emph{Abhandlgn. zur Gesch. der mathem. Wissensch.} X, 81.} addressed to al-Bīrūnī a treatise “on the reasons for the \emph{halving} of the \emph{equation} among the astronomers who followed the Sindhind”.\footnote{See Albērūnī, \emph{Chronologie orientalischer Völker}, ed. Sachau, Leipzig, 1878, p. XXXXVII.} What, then, is the case with this \emph{halving} of the \emph{equation}? M. Chasles already touched on this question in the paper cited above (p. VII\footnoteT{The print has «p. I»; Suter's list of misprints corrects it to p. VII.} of the introduction), but his account is in certain points not quite clear, and not quite correct either. He says, p. 849: \emph{Former un apogée fictif, c’est comme si l’on donnait une équation fictive à l’apogée réel. Et en effet, j’ai trouvé depuis cette équation ou la manière de la former dans d’autres tables que j’ai reconnues être imitées du Sindhind et dans lesquelles ces traces manifestes d’une astronomie étrangère avaient échappé à l’attention des astronomes historiens. L’équation en question est précisément}\footnoteT{The print has «précisement».} \emph{la moitié de l’équation du centre, c’est ainsi que l’auteur la désigne. C’est la prosthaphérèse de longitude de Ptolemée} (to form a fictitious apogee is as if one gave a fictitious equation to the real apogee. And indeed I have since found this equation, or the way of forming it, in other tables that I have recognized to be imitated from the Sindhind, and in which these manifest traces of a foreign astronomy had escaped the attention of the historians of astronomy. The equation in question is precisely half the equation of the centre; that is how the author designates it. It is the prosthaphaeresis of longitude of Ptolemy). — It is indeed true that Kh. corrects the \emph{apogee} too, and not only the \emph{anomaly} and the \emph{centrum} of the planet; it is also possible that this correction of the apogee is meant as a substitute for Ptol.'s \Pg{50}\emph{equant}, which the Indians, or the Persians, had dropped. But Chasles is mistaken, or makes a slip of the pen, when he says: \emph{L’équation en question est précisément la moitié de l’équation du centre} (the equation in question is precisely half the equation of the centre); we have found that in Kh. the equation in question (that of the apogee) is \emph{half the equation of the anomaly}. The matter stands thus.\par
\par The equation of the anomaly (in our tables \emph{Examinatio argumenti}) is halved; the one half is subtracted from the \emph{apogee} as its equation (correction), and so gives the corrected apogee, which is entered in the second main column (see, for example, Table 27 for Saturn), headed \emph{Sublimatio examinata}; the other half is probably used to correct the \emph{equation} of the \emph{centre}, but in what way we have not found out. We have made several tests, and in them have found this half now equal to the equation of the centre, now much greater, and a third time considerably smaller. I give two of these tests here, comparing the numbers in Kh. with the corresponding ones in Batt., which differ little from those of Ptol.\par
\par For the comparison we take: 1) the data of the two tables that belong to the argument 70° of Saturn (in Kh. Table 29, in Batt. II, 110):\par
\begin{center}\small\begin{tabular}{lccc}
\emph{Batt.} &  & \emph{Aequat. (anomal. et) centri} & \emph{Aequat. anomaliae} \\
 &  & 6° 2′ & 5° 37′ \\
\emph{Kh.} & \emph{Sublim. examin.} & \emph{Examin. centri} & \emph{Examin. argumenti (anomaliae)} \\
 & 242° 20′ & 8° 1′ & 5° 10′ \\
\end{tabular}\end{center}
\par If we halve the \emph{equation of the anomaly} in Kh., that is, 5° 10′, we obtain 2° 35′; this, subtracted from the real apogee of Saturn, which is given on p. 14 of the text as 244° 55′, gives the \emph{Sublimatio examinata} = 242° 20′; this half, added to the equation of the centre in Batt., = 6° 2′, gives 8° 37′, which differs by 36′ from the equation of the centre in Kh., namely 8° 1′.\par
\par 2) the data of the two tables that belong to the argument 70° of Jupiter (in Kh. Table 35, in Batt. II, 116):\par
\begin{center}\small\begin{tabular}{lccc}
\emph{Batt.} &  & \emph{Aequat. (anomal. et) centri} & \emph{Aequat. anomaliae} \\
 &  & 4° 52′ & 9° 36′ \\
\emph{Kh.} & \emph{Sublim. exam.} & \emph{Examin. centri} & \emph{Examin. argumenti (anomaliae)} \\
 & 167° 50′ & 4° 44′ & 9° 25′ \\
\end{tabular}\end{center}
\par If one halves the \emph{equation of the anomaly} in Kh., that is, 9° 25′, one obtains 4° 42′; this, subtracted from the real apogee of Jupiter, which is given on p. 14 of the text as 172° 32′, gives the \emph{Sublimatio examinata} = 167° 50′; here this half is only 2′ less than the equation of the centre in Kh. (see the additions).\par
\par \Pg{51}The reasons for this procedure, as already said, have not become clear to us; we have found out only so much, that one arrives at about the same true place of the planet whether one proceeds by Batt.'s method or by Kh.'s, provided that with the latter one takes the \emph{precession} into account, which is just what he did not do. But in our work of commentary, and especially in this chapter on the true places of the planets, we have very much regretted that the “Reasons of the tables of Kh.” by al-Bīrūnī\footnote{Or Aḥmad b. al-Muthannā; see \emph{Biblioth. mathem.}, 3rd series, XII, p. 282.}, which still exist in the Hebrew translation of Ibn Ezra, have not yet been published.\par
\par Let us see further what the \emph{Sūrya-Siddhānta} contains about determining the true place of the planets. Its rules\footnote{\emph{Transl.} Burgess, pp. 211–216.} are considerably more complicated than those of Kh., from which we may probably conclude that the Persian tables of the \emph{Shāh} had already simplified the Indian method. I quote here only verses 43 and 44 (of book II) in the English translation of E. Burgess\footnote{Ibid., pp. 214–215.}:\par
\par 43. \emph{The process of correction for the apsis is the only one required for the sun and moon; for Mars and the other planets are prescribed that for the conjunction, that for the apsis, again that for the apsis and that for the conjunction — four in succession.}\par
\par 44. \emph{To the mean place of the planet apply half the equation of the conjunction, likewise half the equation of the apsis; to the mean place of the planet apply the whole equation of the apsis and also that of the conjunction.}\par
\par Without a commentary one could hardly do much with these rules; we see from them only so much, that a \emph{halving} of the \emph{equation} is involved here too. Burgess comments on these verses in the following way:\par
\par \emph{The phraseology of the text is not entirely explicit, and would bear, if taken alone, a different interpretation from that which the commentary puts upon it, and which the rules to be given later show to be its true meaning; this is as follows: first calculate from the mean place of the planet the equation of the conjunction, and apply the half of it to the mean place; from the position thus obtained calculate the equation of the apsis, and apply half of it to the longitude as already once equated; from this result find once more the equation of the apsis and apply it to the original mean place of the planet; and finally calculate from and apply to this last place the whole equation of the conjunction.}\par
\par \Pg{52}Burgess then gives examples of the computation of the true places of the five planets by these rules. We have spent a long time comparing these computations with those carried out by Kh.'s rules, but have not been able to find out the true connection. In the \emph{Sūrya-Siddhānta} an equation is halved twice: first the \emph{equation of the conjunction} (in our terms the \emph{equation of the centre}) is halved, and half of this equation is added to the mean longitude of the planet; then the \emph{equation of the apsis} (of the epicycle; in our terms the \emph{equation of the anomaly}) is halved, and the half is again added to the new, corrected longitude of the planet; then the \emph{equation of the anomaly} is determined anew from this twice corrected longitude, and it is added to the original mean longitude; finally the \emph{equation of the centre} is determined from this result and added to the result just obtained, and so one has the \emph{true longitude} of the planet. By these rules Burgess determined the true place of Saturn for 1 Jan. 1860 for the meridian of Washington, and found it c. 11° greater than the mean longitude; we have found it, by Kh.'s rules and with the same data as Burgess used, only c. 3° greater. This is so great a difference that one must almost suppose that Burgess did not understand or did not apply the rules of the \emph{Sūrya-Siddhānta} correctly. In any case, however, the rules of the Persian tables of the \emph{Shāh}, if Kh. really followed them, must, as has already been indicated, have departed considerably from those of the \emph{Sūrya-Siddhānta}: of the \emph{two} halvings of equations that Burgess found in the Indian tables, in our view only \emph{one} has remained in Kh.\par
\par Below we give, as an example of the determination of the \emph{true place} of an upper planet, that of \emph{Saturn}, both by the tables of Batt. and by those of Kh.\par
\par \emph{Computation of the true place of Saturn} for the beginning (noon) of 15 Ṣafar of the year 540 AH.\par
\par\begin{center}1. By al-Battānī.\end{center}
\begin{center}\small\begin{tabular}{llcl}
Mean longitude of the sun & (by the tables II, 19–23) & = & 140° 46′ \\
Mean longitude of Saturn & (by the tables II, 24–28) & = & 36° 13′ (396° 13′) \\
 & anomaly of Saturn: &  & 104° 33′ \\
Mean longitude of Saturn &  & = & 396° 13′ \\
Apogee (in the year 540)\footnotemark{} &  & = & 248° 30′ \\
Centrum &  & = & 147° 43′ \\
Equation of the centre &  & = & 3° 37′ 40″\footnotemark{} (by Table II, 112) \\
Centrum aequatum &  & = & 144° 5′ 20″ \\
Anomaly &  & = & 104° 33′ \\
Equation of the centre &  & = & 3° 37′ 40″ \\
Anomalia aequata &  & = & 108° 10′ 40″ \\
\end{tabular}\end{center}
\addtocounter{footnote}{-2}\stepcounter{footnote}\footnotetext{Here the precession from 266 to 540 AH has been taken into account; the apogee is therefore greater by ($\frac{274}{68}$)° = 4° 2′ than Batt. gives it (244° 28′).}\stepcounter{footnote}\footnotetext{If the centrum is greater than 180°, the equation of the centre must be added to it, but subtracted from the anomaly. [\emph{Translator:} The print calls this note at both entries of the equation of the centre.]}
\par To the anomalia aequata of 108° 10′ 40″ there correspond (in Table II, 111), in columns VI and VII, the numbers 6° 6′ and 0° 25′, and in column IV the number 20′; 0° 25′ must now be reduced in the ratio 20′ {\SignFont\upshape ∶} 60′, which gives 8′ 20″\footnote{If the differences that correspond to the centrum aequatum (here 144° 5′) in column IV are decreasing, one must take the numbers that correspond to the anomalia aequata in cols. VI and V.}; this is added to the number of column VI, that is, to 6° 6′, which gives 6° 14′ 20″; this, added to the centrum aequatum\footnote{To be subtracted if the anomalia aequata is greater than 180°.}, gives:\par
\Pg{53}\begin{center}\small\begin{tabular}{lcl}
Centrum aequat. & = & 144° 5′ 20″ \\
Aequat. anomal. corr. & = & 6° 14′ 20″ \\
 &  & 150° 19′ 40″ \\
apogee & = & 248° 30′ \\
added: & = & 398° 50′ \\
 &  & −360° \\
 & = & 38° 50′ = \emph{true longitude} \\
\end{tabular}\end{center}
\par that is, at the given time Saturn was in 8° 50′ of Taurus.\par
\par\begin{center}2. By al-Khwārizmī.\end{center}
\begin{center}\small\begin{tabular}{lcl}
Mean longitude of the sun (by Tables 4 and 5, see above p. 42) & = & 131° 1′\footnotemark{} \\
Mean longitude of Saturn (by Tables 9 and 10) & = & 26° 33′ (386° 33′) \\
anomaly (Argumentum) & = & 104° 28′ \\
Mean longitude of Saturn & = & 386° 33′ \\
Sublimatio (Apogeum) definita (for 104° 28′) & = & 242° 5′\footnotemark{} \\
centrum & = & 144° 28′ \\
To this belongs (by Table 31) the Examinatio centri & = & 4° 56′ 30″\footnotemark{} \\
subtracted: & = & 139° 31′ 30″ \\
\end{tabular}\end{center}
\addtocounter{footnote}{-3}\stepcounter{footnote}\footnotetext{We leave out the seconds when they are under 30, and round them up to 1′ when they are over 30.}\stepcounter{footnote}\footnotetext{By the rules of the text, when the anomaly is greater than 180°, one must add the examin. argumenti (\emph{elheza}) to the sublim. examin., and subtract the sum from the mean longitude of the planet.}\stepcounter{footnote}\footnotetext{If the centrum is greater than 180°, the examin. centri must be added to the centrum, but subtracted from the anomaly. [\emph{Translator:} The print calls this note at both entries of the examinatio centri.]}
\par This is the centrum aequatum or examinatum.\par
\begin{center}\small\begin{tabular}{lcl}
anomaly & = & 104° 28′ \\
Examinatio centri & = & 4° 56′ 30″ \\
added: &  & 109° 24′ 30″ \\
\end{tabular}\end{center}
\par This is the argumentum (anomalia) examinatum. For it the examinatio argumenti (\emph{elheza}), that is, the equation of the anomaly, is looked up in Table 30; it is = 5° 35′ 35″, and it is to be added to the centrum examinatum = 139° 31′ 30″\footnote{When the argumentum (anomalia) examinatum is less than 180°; it is, on the other hand, to be subtracted from the centrum examinatum when the argumentum examinatum is greater than 180°.}; this gives 145° 7′, a quantity that Kh. calls the centrum ultimum. If one again adds to it the sublimatio examinata = 242° 5′, one has the \emph{true longitude} of Saturn for the given time, that is, = 387° 12′ = 27° 12′; that is, Saturn was in 27° 12′ of Aries.\par
\par \Pg{54}The error of c. 11° with respect to the determination of the true place by Batt. comes in the first place from the smaller mean longitude of Saturn in Kh., which, because Kh. neglects the precession, is c. 10° smaller than in Batt.; but in the second place also from deviating values in the equations of the centre and of the anomaly. If one takes the precession into account in Kh. too, but in other respects proceeds by his method, one obtains for the true place of Saturn c. 43′ less than by Batt.'s method.\par
\UnitHead{Chapter 11. Determination of the (true) place of the two lower planets}{Chapter 11. Determination of the (true) place of the two lower planets}{K08-8}
\par\begin{center}(Text p. 12, Tables 45–56.)\end{center}
\par Kh. remarks that the true place of Venus and Mercury is determined in the same way as for the upper planets, but that the mean longitude of these two planets is the same as that of the sun, and that in the tables for the mean course of these planets (Tables 15 and 16) the anomaly is given instead of the mean longitude\footnote{This is what the sentence, somewhat unclear in its first part, means to say: \emph{quod eorum in libro scriptum est argumentum atque eorundem et solis unum idemque est elwacat} (that what is written for them in the book is the argument, and that the \emph{elwacat} of them and of the sun is one and the same).}. Batt.\footnote{Ed. Nallino, I, 113.} expresses himself in a similar way.\par
\par Here too we give an example, for the planet Venus, and for the theory and the computation of the equations, which depart somewhat from those of the upper planets (especially for Mercury), we refer to Batt.\footnote{Ibid., II, 243–244.}\par
\par \emph{Determination of the true place of Venus} for noon of 15 Ṣafar of the year 540 AH.\par
\begin{center}\small\begin{tabular}{lcl}
Anomaly of Venus at the end of the 510th year (by Tables 15 and 16) &  & 232° 45′ \\
increase in 29 years &  & 216° 15′ \\
increase in 1 month &  & 18° 30′ \\
increase in 14 days &  & 8° 38′ \\
sum: &  & 476° 8′ \\
 &  & −360° \\
anomaly & = & 116° 8′ \\
Mean longitude, the same as that of the sun & = & 131° 1′ (see pp. 42 and 53) \\
Sublimatio (apog.) exam. for 116° 8′ (by Table 48) & = & 59° 9′ \\
centrum & = & 71° 52′ \\
\end{tabular}\end{center}
\begin{center}\small\begin{tabular}{lrlr}
Centrum & 71° 52′ & Anomaly & 116° 8′ \\
Examin. centri (by Table 47) & 2° 7′ & Exam. centri & 2° 7′ \\
Centrum aequat. & 69° 45′ & Argument. exam. & 118° 15′ \\
\end{tabular}\end{center}
\par \Pg{55}For this argum. exam. one finds in Table 48, as the equation of the anomaly (examinatio elheza), 44° 43$\tfrac{1}{2}$′; this is to be added to the centrum aequat., which gives 114° 28$\tfrac{1}{2}$′, the centrum ultimum; adding to it the sublim. examin. (= 59° 9′) gives 173° 37$\tfrac{1}{2}$′. This was the \emph{true longitude} of Venus at the given time; that is, the planet was in 23° 37$\tfrac{1}{2}$′ of Virgo.\par
\par By the tables of Batt. one finds the true longitude = 181° 17′. That the difference here is not 10°, as with Saturn, but only about 7° 40′, comes from the mean anomaly of Venus being c. 5° greater in Kh. than in Batt.\par
\UnitHead{Chapter 12. Determination of the longitude of the nodes of the moon}{Chapter 12. Determination of the longitude of the nodes of the moon}{K08-17}
\par\begin{center}(Text p. 12, Tables 19 and 20.)\end{center}
\par As in Batt., only the ascending node (the dragon's head) is considered, because the longitude of the descending node (the dragon's tail) differs from it by 180°. The mean longitude is computed, like that of the sun and of the moon, from Tables 19 and 20 for any given moment; the result is then subtracted from 360°, because the motion of the nodes is retrograde, and the remainder is the true longitude of the node. I refrain from giving an example of this.\par
\UnitHead{Chapter 13. On the stations and the forward and backward motion of the planets}{Chapter 13. On the stations and the forward and backward motion of the planets}{K08-20}
\par\begin{center}(Text p. 12, Tables 27–56.)\end{center}
\par For these determinations Kh. does not set up tables of their own, as Batt. does (II, 138 and 139), but he has a special column for them in the tables of the equations of the planets, headed \emph{Statio prima}; the column \emph{Statio secunda} of Batt., which contains only the complements of the numbers of the \emph{Statio prima} to 360°, Kh. leaves out.\par
\par The rule is the following: for any given moment one determines, as in computing the true place of the planet, its \emph{centrum aequat.} (\emph{examinat.}), looks up the \emph{statio prima} for this argument and subtracts it from 360°, and so has the \emph{statio secunda}; then one also determines the \emph{anomalia} (\emph{elheza}) \emph{examinata} of the planet at this moment. If it is equal to one of the two stations, the planet is stationary; if it lies between the two, that is, is greater than the first and less than the second, the planet is retrograde; in the other case (that is, when it is less than the first or greater than the second), it is moving forward.\footnote{The second sentence in Batt., \emph{si contra maior secunda statione et minor prima, directus} (if, on the contrary, greater than the second station and less than the first, direct), has a small error, for so put the case is impossible; it must simply read \emph{aut} instead of \emph{et}. The error is already in the Arabic text, where \emph{aw} (or) should stand instead of \emph{wa} (and). (Ed. Nallino, I, 114 and III, 173.)}\par
\par The rule in Batt. (I, 114) runs exactly the same, except that he leaves out (probably because he thought it obvious) the case where the anomal. exam. is equal to one of the two stations.\par
\par \emph{Example}: How did Saturn stand in this respect at noon of 15 Ṣafar of the year 540 AH?\par
\par We found above (p. 53) for Saturn:\par
\begin{center}\small\begin{tabular}{cc}
\emph{Centrum aequatum} & \emph{Anomalia aequata} \\
139° 31$\tfrac{1}{2}$′ & 109° 24$\tfrac{1}{2}$′ \\
\end{tabular}\end{center}
\Pg{56}\begin{center}\small\begin{tabular}{lr}
To 139°\footnotemark{} there belongs in Table 31 the \emph{statio prima} & 115° 10′ \\
so the \emph{statio secunda} is & = 244° 50′ \\
\end{tabular}\end{center}
\addtocounter{footnote}{-1}\stepcounter{footnote}\footnotetext{I take only the whole degrees.}
\par Now the \emph{anomal. aeq.} is less than the first station; on that day, then, Saturn was \emph{moving forward} (\emph{direct}). The rule can probably be understood without further commentary.\par
\UnitHead{Chapter 14. On the time of the direct and the retrograde motion}{Chapter 14. On the time of the direct and the retrograde motion}{K09-2}
\par\begin{center}(Text pp. 12–13, Tables 27–56.)\end{center}
\par If one wants to know the time during which a planet is direct or retrograde, one must proceed as follows, according to Kh.:\par
\par a) Let the planet be \emph{retrograde}. If one wants to know how long it has already been retrograde, one takes the difference between the \emph{statio prima} and the \emph{anomal. aeq.} and determines how much time the planet needed to run through this difference; but if one wants to know how long it will still remain retrograde, one takes the difference between the \emph{statio secunda} and the \emph{anomal. aeq.} and determines how much time the planet needs to run through this difference. To find these times one must know the daily motion of the anomaly of the planet; Batt. gives it in his account of this determination. Kh. does not speak of these daily motions; he says simply: \emph{Deinde per argumentum} (that is, here, = logical inference, procedure) \emph{ad hoc negotium proprium, quanto tempore illud transeat, coniectabitur} (then by a reasoning proper to this business it will be inferred in how much time it passes through it). Someone not familiar with the matter would hardly find out what is meant by this; here the clear rules of Batt. put us on the right path: one must divide the said differences between the stations and the anomal. aeq. by the daily motion of the anomaly of the planet\footnote{These numbers for the daily motion of the planets are found in the following tables (under the heading of the column \emph{Heza} (Anomalia or Argumentum)): Table 9, Saturn = 57′ 8″; Table 11, Jupiter = 54′ 9″; Table 13, Mars = 27′ 42″; Table 16, Venus = 37′ 0″; Table 18, Mercury = 3° 6′ 24″.}, and one then obtains the number of days for which the planet has already been retrograde, or will still remain retrograde.\par
\par b) Let the planet be \emph{direct}. For this case the rules are incomplete both in Batt. and in Kh., and in the latter unintelligible besides. Two cases must be distinguished here, namely whether the anomal. aeq. is less than the 1st station or greater than the 2nd; in both cases one can ask again how long the planet will still remain direct, and how long it has already been direct. Batt. considers only the first case and asks only the first question; Kh. considers both cases, but asks only one question in each. The complete rule would be the following.\par
\par \Pg{57}\emph{Case 1}: The \emph{anomal. aeq.} is \emph{less} than the 1st station. If one wants to know how long the planet will still remain direct, one takes the difference between the two and divides it by the daily motion of the anomaly of the planet\footnote{This case is considered by Batt. at the end of chapter 46 (ed. Nallino, I, 114) and by Kh. at the end of his rule (p. 13): \emph{Quod si non procul a statione prima fuerit} . . . . (but if it is not far from the first station ...).}. Our example above for the planet Saturn belongs to this case: it is, as we have found, direct; the difference between statio I and the anomal. aeq. is 5° 45$\tfrac{1}{2}$′; so it remains direct for another $\dfrac{5^{\circ}\,45\tfrac{1}{2}'}{57'}$ = c. 6 days, then becomes stationary and after that retrograde. — If one wants to know how long the planet has already been direct, one adds the anomal. aeq. to the 1st station\footnote{Strictly: (360° − Stat. II) + Anomal. aeq.} and divides the sum by 57′; in our case this gives $\dfrac{224^{\circ}\,34\tfrac{1}{2}'}{57'}$ = 236 days; in all, then, it was direct for 242 days.\par
\par \emph{Case 2}: The \emph{anomal. aeq.} is \emph{greater} than the 2nd station. If one wants to know how long the planet will still remain direct, one subtracts the anomal. aeq. from 360°, adds the 1st station to the remainder, and divides the result by the daily motion of the anomaly of the planet. If one wants to know how long the planet has already been direct, one takes the difference between the anomal. aeq. and the 2nd station and divides it by the daily motion of the anomaly of the planet\footnote{Batt. does not consider this case at all, and Kh. only the second question of this case, in the passage (p. 13) that begins \emph{Si vero de progrediente scire desideras} (but if you wish to know about a planet moving forward). The text of chapter 46 of Batt. in the edition of 1645, moreover, differs in some places from Nallino's translation and from the Arabic text; it is probable that the latter has omissions that the Arabic text which Plato of Tivoli had before him did not have.}.\par
\par All this can easily be recognized from the following figure\footnoteT{Here the print has a figure, which is not reproduced: the circle of the epicycle about \emph{H}, with \emph{A} at the left end and \emph{P} at the right end of a horizontal line through \emph{H}, which runs on to \emph{E} far to the right; on the circle \emph{K} at the upper left, \emph{C} at the upper right, \emph{B} at the lower right, \emph{G} at the lower left and \emph{D} at the bottom; the lines \emph{ECK}, \emph{EBG} and the tangent \emph{ED}; \emph{F} on the chord \emph{GB}, joined to \emph{H} by a line perpendicular to the chord; small arrows at \emph{A} and between \emph{P} and \emph{C} show the direction of the planet on the epicycle.}, if we add a few remarks to it.\par
\par Let \emph{E} be the earth and \emph{H} the centre of the epicycle; \emph{B} and \emph{C} are the stations I and II, at which the retrograde motion begins and ends respectively. According to the theory of Ptol., which Batt. too follows, the retrograde motion does not begin at the point of contact \emph{D} of the tangent from \emph{E} to the epicycle, but at a point \emph{B} so placed that \emph{FB} {\SignFont\upshape ∶} \emph{BE} is equal to the ratio of the speed of the centre of the epicycle to that of the planet on the epicycle; if the former ratio (\emph{FB} {\SignFont\upshape ∶} \emph{BE}) is greater than the latter, the planet is retrograde, and this takes place on the arc \emph{BC}; on the arc \emph{CAB} it is direct (the short time of standing still being included). But the stations I and II are not constant: they change with the position of the centre of the epicycle \emph{H} on the \emph{deferent}, that is, with the so-called \emph{centrum} of the planet; one must therefore enter the tables of the direct and the retrograde motion with the \emph{centrum} (\emph{aequat.}). On the computation of the stations for any centrum from 0° to 180°, compare Batt.\footnote{Ed. Nallino, II, 245–247.} and Ptol.\footnote{Ed. Heiberg, II, 450–494; translation by K. Manitius, II, 267–302.}\par
\par \Pg{58}The values of the \emph{statio prima} in Kh. agree more with those of Batt. than with those of Ptol., though differences of 1′ to 2′ sometimes occur between Kh. and Batt. The greater deviations in Ptol. come from his taking the \emph{mean centrum} as the argument, whereas the two Arab astronomers take the \emph{centrum aequatum}. The agreement between Batt. and Kh. is certainly not accidental; Batt. undoubtedly knew Kh.'s tables, and probably used them too, even though he never cites them.\par
\UnitHead{Chapter 15. On the declination of the sun}{Chapter 15. On the declination of the sun}{K09-12}
\par\begin{center}(Text p. 13, Tables 21–26.)\end{center}
\par There is little to remark about these tables. \emph{Porta} in the title is the translation of the Arabic \emph{bāb} = gate, chapter. The column of Tables 21–26 that contains the declination of the sun is headed \emph{Obliquatio solis} = inclination of the sun (Arab. \emph{mayl al-shams}); then we find added \emph{accedentis} (for \emph{ascendentis}) from 0° to 90° (or 360°–270°), and \emph{recedentis} (for \emph{descendentis}) from 90° to 180° (270°–180°). Kh. gives the \emph{greatest declination}, thus the \emph{obliquity of the ecliptic}, as 23° 51′. The \emph{Sūrya-Siddhānta}\footnote{Ed. Burgess, p. 201.} has 24°, the authors of the “tested tables” give it as 23° 33′, Ptol. has the value 23° 51′ 20″, Batt. 23° 35′; Kh., then, comes nearest to Ptol.: he merely leaves out the 20″ (see above, p. 33). It should be noted, moreover, that a number of Arab astronomers likewise give the Ptolemaic value as 23° 51′, leaving out the seconds. Kh.'s values differ by at most 2″ from those computed by the formula sin \emph{{\greekfont δ}} = sin \emph{l} · sin \emph{e} (where \emph{l} denotes the degrees of the ecliptic counted from the vernal point, thus the longitude of the sun, and \emph{e} the obliquity of the ecliptic). These values of the declination are a fine testimony to the accuracy with which Kh. proceeded in composing his tables.\par
\UnitHead{Chapter 16. On the latitude of the moon}{Chapter 16. On the latitude of the moon}{K09-15}
\par\begin{center}(Text p. 14, Tables 21–26 and 19–20.)\end{center}
\par The latitude of the moon is determined by the same formula as the declination of the sun, that is, \emph{sin latit. lunae} = \emph{sin of the argument of latitude} × \emph{sin of the inclination of the orbit of the moon.} By the argument of latitude one understands the distance of the true place of the moon in the ecliptic from the northern node (the dragon's head). This \emph{argument of latitude} is found in the following way.\par
\par To the true place of the moon one adds the mean longitude of the dragon's head [and subtracts 360° from this sum, when it exceeds this quantity]\footnote{What stands in square brackets is missing in Kh.}.\par
\par Batt. has the same rule, but then also the following one, which Kh. leaves out: one subtracts the true place of the dragon's head from the true place of the moon, and so has the argument of latitude\footnote{Ed. Nallino, I, 76.}.\par
\par \Pg{59}The latter rule also occurs among the Indians\footnote{\emph{Sūrya-Siddhānta}, ed. Burgess, p. 228.}, and in this form it holds for the three upper planets too; it is worth noting that Kh. does not give the rule in the Indian form (at least in that of the \emph{Sūrya-Siddhānta}).\par
\par As the greatest latitude (thus the inclination of the orbit of the moon to the ecliptic) Kh. has, at the argument 90°, the value 4° 30′, while Batt., like Ptol., has 5°. Here Kh. follows the \emph{Sūrya-Siddhānta} (or the tables of the \emph{Shāh}); at least the same value is found in the former\footnote{Ibid., pp. 188–189.}, but it may also have passed into the later \emph{Siddhāntas}.\par
\par \emph{Example}: The latitude of the moon is to be determined for noon of 15 Ṣafar of the year 540 AH.\par
\begin{center}\small\begin{tabular}{rcl}
\multicolumn{3}{c}{\emph{True place of the moon (see above p. 45):}} \\
319° 2′\footnotemark{} &  &  \\
279° 41′ & = & mean longitude of the dragon's head \\
598° 43′ &  &  \\
−360° &  &  \\
238° 43′ & = & argument of latitude (by the 1st rule) \\
319° 2′ & = & true place of the moon \\
80° 19′ & = & corrected longitude (= true place) of the dragon's head \\
238° 43′ & = & argument of latitude (by the 2nd rule) \\
\end{tabular}\end{center}
\addtocounter{footnote}{-1}\stepcounter{footnote}\footnotetext{I round the seconds down or up to minutes. [\emph{Translator:} The print calls this note at 319° 2′ and again at 92° 39′.]}
\begin{center}\small\begin{tabular}{lr}
\multicolumn{2}{c}{\emph{Mean longitude of the dragon's head (by Tables 19 and 20):}} \\
at the end of the 510th year & 92° 39′ \\
for 29 years & 184° 42′ \\
for 1 month & 1° 35′ \\
for 14 days & 45′ \\
sum: & 279° 41′ \\
\multicolumn{2}{c}{360° − 279° 41′ = 80° 19′ = the corrected longitude, since the dragon's head is retrograde.} \\
\end{tabular}\end{center}
\par With this argument one enters Table 25, and finds there:\par
\begin{center}\small\begin{tabular}{lcl}
in the last column, for 238° (7$^{s}$ 28°), the latitudo lunae & = & 3° 48′ 57″ \\
difference for 43′ & = & 1′ 46″ \\
\emph{latitude} of the moon at the given time: &  & 3° 50′ 43″ \\
\end{tabular}\end{center}
\par According to the remarks that Batt. adds to his rules, but Kh. leaves out, the latitude of the moon in this case is \emph{southern} and \emph{increasing}, that is, the moon is going still further towards the south, which is also indicated by the heading of the last column, \emph{Latitudo lunae recedentis} (latitude of the receding moon).\par
\par The headings of the first two columns of Tables 21–26, \emph{Computatio minuenda et augmentum examinis et numeri} — \emph{Computatio augenda et diminutio examinis et numeri} (computation to be subtracted, and increase of the equation and of the number — computation to be added, and decrease of the equation and of the number), are not quite clear, but they probably mean: if the argument is less than 180°, which is the case in the 1st column, the equation found (\emph{computatio}) must be subtracted from the mean longitude; if it is greater than 180°, which is the case in the 2nd column, the equation must be added; further, when the argument increases from 0° to 180°, the equation (\emph{examen}) increases too, but when the argument grows from 180° to 360°, the equation decreases again.\par
\UnitHead{Chapter 17. On the latitude of the planets}{Chapter 17. On the latitude of the planets}{K10-9}\Pg{60}
\par\begin{center}(Text p. 14, Tables 27–56.)\end{center}
\par The Latin title of this chapter is incorrect: in its first paragraph the chapter deals with the three upper planets, in the second with the two lower ones. The theory of the change of the latitude of the planets is the most difficult and obscure part (book 13) of Ptol., so difficult and obscure that Batt. renounced going into this theory more closely; he gives only the rules by which the latitude of the planets can be found from the tables\footnote{Batt., ed. Nallino, I, 115–116; II, 140–141.}.\par
\par But this obscurity has been cleared up in an excellent way by Nallino, after Delambre and others had failed at it, and he has thereby brought the theory of Ptol. nearer to our understanding; he has also tried to give it a mathematical foundation, in which he has likewise succeeded finely. We cannot go into these foundations more closely here, but must refer the reader to Nallino's work\footnote{Ibid., II, 247–255.}; we shall only mention briefly the main factors on which the latitude of the planets in their various positions depends. The latitude of the \emph{three upper} planets depends, first, on the constant inclination of the plane of the deferent to the ecliptic, and secondly on the variable inclination of the epicycle to the plane of the deferent. The latitude of the \emph{two lower} planets depends, first, on the inclination of the plane of the deferent to the ecliptic, which Ptol. takes to be variable for the lower planets; secondly, on the variable inclination of the diameter apogee–perigee of the epicycle to the plane of the deferent; and thirdly, on the variable inclination to the plane of the ecliptic of the diameter perpendicular to the one just named (for the three upper planets this diameter is always parallel to the plane of the ecliptic). — The only thing we have to object to in Nallino's accounts is that he has not tried to make Ptol.'s unclear figures more vivid by adding perspective figures.\par
\par If we now compare the rules that Kh. has given for determining the latitude of the planets with those of Batt., we see at once that the former must rest on other theories than the latter; and it is quite understandable that the difficulty of the Ptolemaic account deterred his successors from following the path of the great master and compelled them to use other, simpler methods.\par
\par \Pg{61}If we now look at the rule of the \emph{Sūrya-Siddhānta}\footnote{\emph{Transl.} Burgess, \emph{Journal of the American Orient. Soc.}, Vol. VI, p. 228.}, we find no direct agreement with Kh.'s here either; but that both proceeded from similar points of view is nevertheless evident. The rule of the \emph{Sūrya-Siddhānta} says: the sine of the arc (the difference of the true longitude of the planet and the longitude of the node), multiplied by the maximum latitude and divided by the radius (for this position of the planet), gives the latitude for the position (time) concerned. Burgess says in his commentary\footnote{Ibid., pp. 229–230.} on this rule that it results from the composition of two proportions, of which the one holds only for the moon, while the other is added for the planets, namely:\par
\begin{center}\small\begin{tabular}{llcl}
\emph{1st prop.} & \emph{r} {\SignFont\upshape ∶} sin dist. from the node & = & greatest latitude {\SignFont\upshape ∶} latitude at the dist. \emph{r}\footnotemark{} \\
\emph{2nd prop.} & var. \emph{r} {\SignFont\upshape ∶} \emph{r} & = & latitude at the dist. \emph{r} {\SignFont\upshape ∶} latitude at the true dist. \\
\end{tabular}\end{center}
\addtocounter{footnote}{-1}\stepcounter{footnote}\footnotetext{Strictly, the sine should also be taken of the last two terms of the proportion; but since these inclinations are only small in the orbits of the planets, one can put the arcs instead of the sines without great error.}
\par From the two proportions it follows, by multiplication:\par
\par\begin{center}var. \emph{r} {\SignFont\upshape ∶} sin dist. from the node = greatest latitude {\SignFont\upshape ∶} latitude at the true dist.,\end{center}
\par so: latitude at the true dist. = $\dfrac{\text{sin dist. from the node} \times \text{greatest latitude}}{\text{var. } r}$.\par
\par As for the first proportion, it follows from the rule of the four quantities in spherical trigonometry, sin \emph{AB} {\SignFont\upshape ∶} sin \emph{AB′} = sin \emph{BC} {\SignFont\upshape ∶} sin \emph{B′C′}, which passes into the formula\par
\par\begin{center}$\sin\alpha = \dfrac{\sin B'C'}{\sin AB'}$\end{center}
\par when \emph{AB} becomes = 90°\footnoteT{Here the print has a figure, which is not reproduced: a spherical triangle with \emph{B} at the upper left, \emph{C} at the lower left and \emph{A} at the right; the arc \emph{BC} on the left, the arc \emph{CC}′\emph{A} below and the arc \emph{BB}′\emph{A} above, curving down to \emph{A}; the arc \emph{B}′\emph{C}′ from \emph{B}′ on the upper arc down to \emph{C}′ on the lower; the angle \emph{{\greekfont α}} is marked at \emph{A}.}. Applied to our case, \emph{A} is one of the nodes (the points of intersection of the deferent with the ecliptic), \emph{{\greekfont α}} = the inclination of the deferent to the ecliptic, \emph{BC} = the maximum latitude, when \emph{AB} = 90°, \emph{B′C′} = the latitude of the planet at its position \emph{B′} on the deferent, and \emph{AB′} = the difference of the true longitude of the planet and the longitude of the node.\par
\par The second proportion is in any case only approximately correct; we can derive\footnoteT{The verb is missing in the print.} it from the following figure\footnoteT{Here the print has a figure, which is not reproduced: two large ellipses, the circles of the ecliptic and of the inclined deferent seen obliquely, crossing at two points; a straight line through \emph{E} from the top of the outer ellipse to the lower crossing point; the small ellipse of the epicycle about \emph{H} at the upper right, with \emph{P} on it; the lines \emph{EP}, \emph{MH} and \emph{EA}, the horizontal line \emph{EMB}, and the vertical lines \emph{PA} and \emph{HB}.}:\par
\par Let \emph{P} be the position of the planet on the epicycle, \emph{A} its projection on the plane of the ecliptic, and \emph{B} the projection of the centre of the epicycle on the same plane; then the angles \emph{HMB} and \emph{PEA} are what in the \emph{Sūrya-Siddhānta} are called the latitude of the planet at the distance \emph{r} (the mean distance \emph{MH}) and the latitude of the planet at the distance \emph{r′} (the variable distance \emph{EP}) respectively. Because of the small inclination of the plane of the epicycle to the ecliptic, the Indians took the two perpendiculars \emph{PA} and \emph{HB} as equal; one then has:\par
\par\begin{center}$\sin HMB = \dfrac{HB}{MH} = \dfrac{HB}{r}$; $\sin PEA = \dfrac{PA}{EP} = \dfrac{PA}{r'}$\end{center}
\par so that $\dfrac{\sin HMB}{\sin PEA} = \dfrac{r'}{r}$, \Pg{62}or, as was remarked above, putting the arcs (angles) instead of the sines:\par
\par\begin{center}$\dfrac{\text{latitude at the mean dist. } r}{\text{latitude at the true dist. } r'} = \dfrac{r'}{r}$,\end{center}
\par and this is the second proportion of the \emph{Sūrya-Siddhānta}.\par
\par Now the values of the latitude of the planets computed according to Batt. by no means agree with those computed according to the \emph{Sūrya-Siddhānta}; but the latter do not agree with those computed according to Kh. either. We have already remarked above that Kh.'s rule shows no direct agreement with that of the \emph{Sūrya-Siddhānta}; we shall now try to show how far they do agree.\par
\par We have found that the numbers of the last column of Tables 27–56, which bears the heading \emph{Latitudo secunda}, are obtained by the formula\par
\par\begin{center}sin \emph{lat. sec.} = sin \emph{argum.} × sin \emph{lat. maxim.},\end{center}
\par where the \emph{argumentum} = the true longitude of the planet less the longitude of its node = the distance of the planet from the node, and the \emph{lat. max.} (at the mean distance) = the \emph{lat. secunda} for 90° (for Saturn = 5° according to Kh.); this \emph{lat. secunda}, then, agrees with the numerator of the value for the latitude at the true distance (see p. 61)\footnote{If one puts “\emph{sin} greatest latitude” instead of “greatest latitude”.}. Now Kh.'s rule says that one must divide the \emph{lat. secunda} by the \emph{lat. prima} to obtain the latitude sought: not by the \emph{lat. prima} that corresponds to the preceding argument (the distance of the planet from the node), but by the one that corresponds to the argument \emph{anomalia aequata} (\emph{elheza examinatum} (!) in the text). If one compares this with the formula given in the \emph{Sūrya-Siddhānta} for the latitude at the true distance, this \emph{lat. prima}, which corresponds to the \emph{anomalia aequata}, would then have to be equal to the variable radius \emph{r′}. The tables of the planets (nos. 27–56), however, do not contain this \emph{r′} in absolute magnitude, but give its ratio to the mean radius \emph{r}; in other words, the numbers of the column \emph{lat. prima} represent the values of \emph{r′} when \emph{r} is taken = 1.\par
\par How this \emph{r′} was determined we could not find out; but in any case it is wrong that the column of the \emph{lat. prima} is headed with degrees and minutes: \emph{r′}, or rather $\dfrac{r'}{r}$, is a pure number, and so instead of \emph{degrees} and \emph{minutes} it should read \emph{partes} and \emph{minutae}; one then also obtains degrees and minutes again for the quotient of the \emph{lat. sec.} by the \emph{lat. prima}.\par
\par We shall now give an \emph{example} that follows Kh.'s rules exactly. Let the \emph{latitude} of \emph{Saturn} be determined for the beginning of 15 Ṣafar of the year 540 AH.\par
\par The \emph{anomalia aequata} (\emph{elheza examinata}) for this moment is = 109° 24$\tfrac{1}{2}$′ (see p. 53); to it belongs, by Table 30, the latitudo prima 1° 15′, or rather 1$^{p}$ 15′ = 1$\tfrac{1}{4}$; the true place of the planet for this moment is = 27° 12′ (see p. 53), and the longitude of the node = 103° 12′ (see p. 64); the two added give 130° 24′; to this argument there corresponds in Table 31 the latitudo secunda 3° 48′ 27″\footnoteT{The print has «3° 48″ 27′», with the signs of the minutes and the seconds exchanged.}; we now have:\par
\par\begin{center}3° 48$\tfrac{1}{2}$′ {\SignFont\upshape ∶} 1$^{p}$ 15′ = 228$\tfrac{1}{2}$′ {\SignFont\upshape ∶} 1$\tfrac{1}{4}$ = 183′ = 3° 3′ = \emph{the latitude sought.}\end{center}
\par \Pg{63}The computation by Batt.'s rules gives the value 2° 34′; the computation with Batt.'s data, but by Kh.'s rules, gives 2° 28′.\par
\par Whether the latitude is \emph{northern} or \emph{southern} our text does not say at all, while Batt. (I, 115–116) speaks clearly about it: by his rules it is \emph{southern} for Saturn in our example; further, the planet is \emph{descendens}, that is, going from the ecliptic towards the south, which can also be read directly in our Table 31: the column from which we took the \emph{lat. sec.} is headed \emph{latitudo Saturni recedentis} (the latitude of Saturn receding). — By the rule of our text the argument for the \emph{lat. sec.} is, it seems, counted from the ascending node through the perigee of the deferent of the planet; counted from the same node through the apogee it would be = 360° − 130° 24′ = 229° 36′, a value from which it would be easier to see that the latitude is southern. — The rule given here for determining the \emph{lat. sec.} is, moreover, the same as the one given for determining the latitude of the moon (see above, pp. 58–59).\par
\par If one compares our rule with those in Ptol. and Batt., which agree fairly closely, one sees that they behave in exactly opposite ways as regards the arguments for the latitude and the proportional numbers: Ptol. determines the latitude number with the \emph{anomalia aequata} as argument, and the proportional number (by which the latitude found must be multiplied) with the true longitude of the planet increased by 50° (for Saturn)\footnote{Batt. takes the centrum instead of the true longitude.} as argument; Kh., on the other hand, determines the latitude number with the distance of the planet from the ascending node (that is, indirectly with the true longitude), but the proportional number (\emph{latit. prima}) with the \emph{anomalia aequata} as argument. These different conceptions then naturally entail a quite different arrangement of the tables as well; we shall mention only one main difference: for the upper planets Ptol. and Batt. have the maximum of the latitude at the argument 180° and the minimum at 0° (360°); Kh., on the other hand, has the maximum of the \emph{latitudo secunda} at 90° (270°) and the minimum at 0° and 180°.\par
\par For the \emph{lower} planets the theory of the determination of the \emph{latitude}, and the computation with the help of the tables, are considerably more complicated still in Ptol. and Batt. than for the upper ones; I may be excused from going into it here, and for details I refer the reader to Nallino's edition of Batt.\footnote{I, 115–116, and II, 249–253.}\par
\par As an \emph{example} we give the computation of the \emph{latitude} of \emph{Venus} for the beginning of 15 Ṣafar 540 AH by Kh.'s rules and tables.\par
\par The \emph{anomalia} (\emph{argumentum} in the text) \emph{aequata} of Venus for this moment, computed from Tables 15–17, is = 118° 15′ (see above, p. 54); to it belongs in Table 48 the \emph{latit. prima} = 1$^{p}$ 15′.\par
\begin{center}\small\begin{tabular}{lcl}
anomalia aequata & = & 118° 15′ \\
longitude of the node & = & 59° 27′ \\
Mean longitude of Venus, the same as that of the sun & = & 131° 1′ \\
sum: &  & 308° 43′ \\
\end{tabular}\end{center}
\par to which there belongs in Table 46 the \emph{latitud. sec.} 3° 54′.\par
\par \Pg{64}3° 54′ {\SignFont\upshape ∶} 1$^{p}$ 15′ = 234′ {\SignFont\upshape ∶} 1.25 = 187′ = 3° 7′ = the \emph{latitude} of \emph{Venus} at the given moment. By the tables of Batt. one finds the latitude = 57′; this great difference suggests that Kh.'s rule is not right. Indeed it is unclear to us why the \emph{anomalia aequata} must be added as well to the sum of the mean longitude and the longitude of the node; if one leaves it out, the argument for the lat. sec. comes out as 190° 28′, and to it belongs in Table 50 the lat. 54$\tfrac{1}{2}$′; the \emph{latitude} is then =\par
\par\begin{center}54.5′ {\SignFont\upshape ∶} 1.25 = c. 44′.\end{center}
\par Here too our text says nothing about the kind of latitude; by \emph{Batt.'s} rules it is \emph{northern}.\par
\par With this we leave this difficult chapter on the latitude of the planets. In our case the question was made harder still by the fact that the Indian astronomical works are not in agreement on this point either: the rule of the \emph{Sūrya-Siddhānta}, which we quoted above, does not agree with that of the \emph{Pañcasiddhāntikā}\footnote{Ed. G. Thibaut, pp. 95 and 96 of the English translation.}; there, for Saturn and Mercury, the sine of the distance from the node is multiplied by the constant $\tfrac{9}{8}$ (if the commentator is right), whereas in the \emph{Sūrya-Siddhānta} it is multiplied by the variable quantity $\dfrac{\text{greatest latitude}}{r'}$.\par
\UnitHead{Chapter 18. On the apogee and the perigee of the planets}{Chapter 18. On the apogee and the perigee of the planets}{K11-23}
\par\begin{center}(Text pp. 14–15.)\end{center}
\par These technical terms are rendered by the translator by \emph{sublimatio} (O and C sometimes have \emph{sullimatio}) and \emph{submissio}. We have already given above (p. 46) the apogees of the planets and of the sun according to Kh., and have also spoken of their deviations from the values of Batt. and of the \emph{Sūrya-Siddhānta}; we add here only that, with regard to the moon, the remark is made that its apogee is movable, and that this motion amounts to 6′ 40″ 18‴ in a day and 16″ 42‴ 27$^{\mathit{IV}}$ in an hour. This motion, which is the cause of the anomalistic month, we do not find mentioned in Batt.; according to the more recent determinations the line of the apsides needs 3231.466$^{d}$ (= c. 9$^{a}$) for a full revolution, which makes a motion of 6′ 41″ in a day; if in the above numbers one compares the hourly motion with the daily one, the latter, if the former is right, would have to be 6′ 40″ 58‴; so 58‴ should probably be read instead of 18‴, a confusion that is easily possible in numbers expressed in Arabic letters; and then Kh.'s data agree excellently with the modern values.\par
\UnitHead{Chapter 19. On the nodes of the planets}{Chapter 19. On the nodes of the planets}{K11-26}
\par\begin{center}(Text p. 15.)\end{center}
\par The longitudes of the dragon's head for the individual planets are:\par
\begin{center}\small\begin{tabular}{ccccc}
\emph{Saturn} & \emph{Jupiter} & \emph{Mars} & \emph{Venus} & \emph{Mercury} \\
103° 12′ & 82° 1′ & 21° 54′ (44′) & 59° 27′ & 21° 10′.\footnotemarkT{} \\
\end{tabular}\end{center}
\addtocounter{footnoteT}{-1}\stepcounter{footnoteT}\footnotetextT{The print has «21° 10.», without the sign of the minutes.}
\par \Pg{65}For the moon he remarks again that its nodes are movable, and that their motion amounts to 3′ 11″ 48‴ a day and 7″ 57‴ an hour; if one multiplies the latter number by 24, one obtains 3′ 10″ 48‴, and so it must read instead of 3′ 11″ 48‴; then it also agrees with Table 20, which, like Batt., has 3′ 11″ for the daily motion of the dragon's head and 8″ for the hourly. That this motion of the nodes is opposite to that of the apsides is not mentioned by Kh. According to the more recent determinations the line of the nodes needs 6798.3355$^{d}$ for a full revolution, which makes a motion of 3′ 10″ 38‴ in a day, so again a very fine agreement. The tables of Batt.\footnote{Ed. Nallino, II, 19–22.} agree with those of Kh. only up to 4 days; from there on the deviations begin, and they grow to 1′ 14″ in one year and to 34′ 46″ in 30 years.\par
\UnitHead{Chapters 20 and 21. On the daily motion of the moon and of the sun}{Chapters 20 and 21. On the daily motion of the moon and of the sun}{K12-1}
\par\begin{center}(Text pp. 15–16.)\end{center}
\par These two chapters do not stand in the right place here; they should immediately follow chapter 7, “On the mean motion of the planets”, or better be united with it.\par
\par The mean daily motion of the moon is given as 13° 10′ 34″ 52‴ 48$^{\mathit{IV}}$, the hourly as 32′ 56″ 47‴ 52$^{\mathit{IV}}$; the last two numbers are wrong: it must read 27‴ 12$^{\mathit{IV}}$.\par
\par The mean daily motion of the sun is given as 59′ 8″, the hourly as 2′ 27″ 50‴ 25$^{\mathit{IV}}$; but the division of 59′ 8″ by 24 gives only 2′ 27″ 50‴, so ten thirds have probably dropped out, or been left out, after 59′ 8″.\par
\par \Pg{66}If we take as the basis for the motion of the sun the above mean hourly motion of 2′ 27″ 50‴ 25$^{\mathit{IV}}$, we obtain for the mean motion in 365 days 359° 44′ 40″ 50‴. To this mean motion there corresponds a length of the year of 365$^{d}$ 6$^{h}$ 13$^{m}$ 2$^{s}$, about the same quantity as we found above (p. 42), which shows us again that Kh. took the sidereal year as the basis of his tables of the sun and the planets. But this length of the year does not quite correspond to the values found in the Indian \emph{Siddhāntas}\footnote{Batt., ed. Nallino, I, 206.}: it is c. 25$^{s}$ greater than that given by the \emph{Sūrya-Siddhānta} and the later \emph{Ārya-Siddhānta}, and 53$^{s}$ greater than that given by the \emph{Brahma-Siddhānta}. If, conversely, we start from the length of the year of these \emph{Siddhāntas}, we arrive, for the \emph{Sūrya-Siddhānta} (365$^{d}$ 6$^{h}$ 12$^{m}$ 36.5$^{s}$), at a mean daily motion of 59′ 8″ 10‴ 10$^{\mathit{IV}}$, and for the \emph{Brahma-Siddhānta} (365$^{d}$ 6$^{h}$ 12$^{m}$ 9$^{s}$) at one of 59′ 8″ 10‴ 21$^{\mathit{IV}}$ 22$^{\mathit{V}}$\footnoteT{Recomputed from the year of 365$^{d}$ 6$^{h}$ 12$^{m}$ 9$^{s}$, the daily motion is 59′ 8″ 10‴ 21$^{\mathit{IV}}$ 33$^{\mathit{V}}$ and the hourly motion 2′ 27″ 50‴ 25.9$^{\mathit{IV}}$; the two hourly motions of the text are truncated, not rounded.}; to the former value there corresponds an hourly motion of 2′ 27″ 50‴ 25.4$^{\mathit{IV}}$, to the latter one of 2′ 27″ 50‴ 25.8$^{\mathit{IV}}$. Since Kh. gives the hourly motion as 2′ 27″ 50‴ 25$^{\mathit{IV}}$, the probability is somewhat greater that he followed the \emph{Sūrya-Siddhānta} and the later \emph{Ārya-Siddhānta}; but this question can by no means be decided definitively from the numbers given\footnote{See also p. 103. [\emph{Translator:} The print has «Kap. 38» (chapter 38); Suter's list of misprints corrects it to p. 103.]}.\par
\par In a note that only Ms. C contains, the data of Ptol. for the length of the year and for the monthly, daily and hourly mean motion of the sun are given correctly, and so is the path of the sun in a year (of 365 whole days) and in 18 years\footnote{This statement is faulty in the degrees: instead of 1350 it must read 355.}. As for the length of the year, it is given as 365$^{d}$ 14′ 48″, that is, 365 + $\tfrac{14}{60}$ + $\tfrac{48}{3600}$ days; converted into hours, minutes and seconds this makes 365$^{d}$ 5$^{h}$ 55$^{m}$ 12$^{s}$, that is, the length of Ptol.'s tropical year. — It is curious that nowhere in the whole of Kh.'s work is there a statement of the length of the sidereal year on which it is based.\par
\UnitHead{Chapter 22. On the appearance of the moon on the evening of the 29th day (of the month)}{Chapter 22. On the appearance of the moon on the evening of the 29th day (of the month)}{K12-8}
\par\begin{center}(Text pp. 16–17, Table 57$^{a}$.)\end{center}
\par This chapter is one of the most difficult of Arabic astronomy. Al-Bīrūnī says about it in his \emph{Chronology}\footnote{Ed. Sachau, Arabic text p. 196, English translation p. 176.}: “As for the appearance of the new moon, its exact determination is lengthy and difficult (to explain), and it needs complicated computations and many tables; let it suffice to refer to what is said on this subject in the tables of al-Battānī and of Ḥabash the Reckoner”.\par
\par Batt. treats this question in chapter 41\footnote{Ed. Nallino, I, 85–92.}, and G. Schiaparelli and Nallino have commented on it excellently\footnote{Ibid., I, 265–272.}. I cannot go into this subject more closely here; let it suffice to remark that Kh.'s rule and his table must probably rest on a different basis from Batt.'s account. I bring out only one point that can support this assertion. According to Batt., the first appearance of the lunar crescent depends not only on the difference of the true longitudes of the sun and the moon, but also on the latitude of the moon; he gives (p. 87) the precept that one must add the square of the difference of the true longitudes of the sun and the moon to the square of the latitude of the moon, and one then obtains approximately the square of the distance of the moon from the sun; the correctness of this precept follows at once from the figure beside\footnoteT{Here the print has a figure, which is not reproduced: an arc, concave downward, labelled «Ekliptik» (ecliptic), with \emph{C} and \emph{S} on it; \emph{M} above \emph{C}, joined to \emph{C} by a short vertical line and to \emph{S} by a straight line.}:\par
\par\begin{center}\emph{MS} {\SignFont\upshape ∽} $\sqrt{CS^{2}+MC^{2}}$.\end{center}
\par \Pg{67}But Kh., who does take the latitude of the moon into account too, says (p. 16): “If the latitude of the moon is northern, it is to be added to the difference of the true longitudes of the sun and the moon; if it is southern, it is to be subtracted from that difference.” — How Table 57$^{a}$ is to be used to determine the time of the first visibility of the lunar crescent will best become clear from an example.\par
\par \emph{Example:} We are to examine whether the lunar crescent became visible on the evening of 29 \emph{Shaʿbān}\footnote{My reason for choosing this day is that the determination of the visibility of the lunar crescent is really of importance only for this day, for with the end of 29 \emph{Shaʿbān} (at 6 o'clock in the evening) the month of fasting, \emph{Ramaḍān}, begins. The fast, however, does not begin at 6 o'clock, but at the moment when the lunar crescent first becomes visible. This can happen some time before or after 6 o'clock. Since the weather does not always permit the appearance of the lunar crescent to be observed, computation and the tables must be called in to help.} of the year 540 AH.\par
\begin{center}\small\begin{tabular}{lr}
\multicolumn{2}{c}{\emph{\emph{Mean longitude of the sun:}}} \\
at the end of the 510th year & = 38° 35′ 45″ \\
increase for 29 years & = 49° 2′ 50″ \\
increase for 7 months & = 204° 1′ 11″ \\
increase for 29 days\footnotemarkT{} & = 28° 34′ 57″ \\
increase for 6 hours & = 14′ 47″ \\
 & 320° 29′ 30″ \\
apogee of the sun: & 77° 55′ \\
\emph{anomaly:} & 242° 34′ 30″ \\
\multicolumn{2}{c}{For this one finds in Table 24 the equation 1° 57′ 35″, which, added to the mean longitude, gives in round numbers 322° 27′ for the \emph{true longitude} of the sun.} \\
\end{tabular}\end{center}
\addtocounter{footnoteT}{-1}\stepcounter{footnoteT}\footnotetextT{Counted as on p. 42, where noon of 15 Ṣafar is 1 month and 14 days after the end of the year, 7 months, 29 days and 6 hours lead to 6 in the evening of 30 Shaʿbān; the evening that ends 29 Shaʿbān would be 7 months, 28 days and 6 hours.}
\begin{center}\small\begin{tabular}{lr}
\multicolumn{2}{c}{\emph{\emph{Mean longitude of the moon:}}} \\
at the end of the 510th year & = 40° 56′ 38″ \\
increase for 29 years & = 53° 24′ 16″ \\
increase for 7 months & = 207° 30′ 20″ \\
increase for 29 days & = 22° 6′ 52″ \\
increase for 6 hours & = 3° 17′ 39″ \\
 & 327° 15′ 45″ \\
\multicolumn{2}{c}{\emph{\emph{Anomaly of the moon:}}} \\
at the end of the 510th year & = 61° 52′ 40″ \\
increase for 29 years & = 348° 56′ 57″ \\
increase for 7 months & = 184° 27′ 14″ \\
increase for 29 days & = 18° 53′ 5″ \\
increase for 6 hours & = 3° 15′ 58″ \\
 & 617° 25′ 54″ \\
 & −360° \\
\emph{anomaly:} & 257° 25′ 54″ \\
\multicolumn{2}{c}{For this one finds in Table 24 the equation 4° 48′ 29″, which, added to the mean longitude, gives in round numbers 332° 4′ for the \emph{true longitude} of the moon.} \\
\end{tabular}\end{center}
\par Difference of the two true longitudes = 9° 37′. We now have to determine the \emph{latitude} of the moon for this time:\par
\begin{center}\small\begin{tabular}{lr}
\multicolumn{2}{c}{True longitude of the moon = 332° 4′.} \\
\multicolumn{2}{c}{\emph{Longitude of the dragon's head (by Tables 19 and 20):}} \\
at the end of the 510th year & = 92° 38′ 42″ \\
increase for 29 years & = 184° 41′ 57″ \\
increase for 7 months & = 10° 58′ 16″ \\
increase for 29 days & = 1° 32′ 14″ \\
increase for 6 hours & = 48″ \\
 & 289° 51′ 57″ \\
\end{tabular}\end{center}
\par \Pg{68}True longitude less the longitude of the dragon's head = 42° 12′\footnoteT{Here Suter takes the sum of Tables 19 and 20 (289° 51′ 57″) itself away from the true longitude of the moon. By his own rule (pp. 55 and 59) the node stands at 360° − 289° 51′ 57″ = 70° 8′, and the argument of latitude is 332° 4′ − 70° 8′ = 261° 56′; the latitude is then about 4° 27′ south (4° 30′ × sin 81° 56′), the distantia definita by the rule of the text about 9° 37′ − 4° 27′ = 5° 10′, less than the 9° 21′ of Table 57$^{a}$, and the crescent would not have been visible that evening.}. This is the \emph{argument} for the \emph{latitude}; to it belongs in Table 22 the \emph{latitude} 3° 1′ 22″; this latitude, added to the difference of the longitudes found before (9° 37′\footnoteT{The print has «9° 37″», with the sign of the seconds.}), since the latitude of the moon is northern, gives in round numbers 12° 38′; this quantity is called the \emph{distantia definita}. One now enters Table 57$^{a}$ with the true longitude of the moon, in the following way: the true longitude is expressed in signs, degrees and minutes, so in our case it is = 11$^{s}$ 2° 4′, and the moon therefore stands in the 2nd degree (of the first face)\footnote{A sign of the zodiac is divided into three decans (\emph{facies}), each of which comprises 10°.} of Pisces. In our tables the number 9° 21′ corresponds to the column \emph{Pisces} in the row X; but 9° 19′ should rather be taken, since the place of the moon is nearer to 30° of Aquarius than to 10° of Pisces; this, however, makes no difference here. This number (9° 21′) is less than the distance found above; the moon is thus farther from the sun than 9° 21′, and so on this evening the lunar crescent must be visible at 6 o'clock. If the number found in the tables were greater than the distance, the lunar crescent would not become visible on the evening of 29 \emph{Shaʿbān}, and the fast would then begin only on the following evening, that is, on the evening of 1 \emph{Ramaḍān}.\par
\par The computation given holds for the meridian of \emph{Arin}. In the last paragraph Kh. adds the remark that for a place other than Arin the difference in time must be taken into account and the longitudes of the sun and the moon corrected accordingly.\par
\par The middle paragraph states how one is to proceed if one wants to know whether or not the lunar crescent will be visible on the morning of the 27th day of the month (thus one day before the conjunction); the procedure is of course analogous to the preceding one.\par
\par How the numbers of Table 57$^{a}$ were found it was impossible for us to establish. Kh.'s rules also differ considerably from those of the \emph{Sūrya-Siddhānta}\footnote{\emph{Transl.} Burgess, pp. 365 and 370.} and of the \emph{Pañcasiddhāntikā}\footnote{Ed. Thibaut, pp. 38–39 of the translation.}; in these two astronomical works one also finds the statement, missing in Kh., that the lunar crescent becomes visible in the west when its distance from the sun amounts to 12 or more equatorial degrees.\par
\par Table 57$^{b}$ does not belong to this chapter; I have put it here to make use of the little room that 57$^{a}$ leaves free. As the reader will recognize at once, it is a multiplication table for sexagesimal parts, and so needs no further explanation; nor is it mentioned anywhere in the text.\par
\UnitHead{Chapter 23. On finding the sine from the arc and conversely}{Chapter 23. On finding the sine from the arc and conversely}{K13-5}
\par\begin{center}(Text pp. 17–18, Tables 58 and 58$^{a}$.)\end{center}
\par On this trigonometrical chapter we have a few remarks to add to what Bjørnbo has already published about the table of sines and cotangents of this work.\footnote{Al-Chwarizmi’s \emph{Trigonometriske Tavler.} \emph{Af Festskrift til} H. G. Zeuthen, København, 1909.}\par
\par \Pg{69}Kh. calls the sine of an angle \emph{elgeib elmustewi} = \emph{al-jayb al-mustawī}, and Adelard renders this by \emph{elgeib planum} (he takes most Arabic words as neuters); \emph{planus} is indeed one meaning of \emph{mustawī}, but \emph{rectus} = straight, or set straight upright, would be better here. Adelard, then, did not yet know (about the year 1126) the translation of \emph{jayb} by the Latin \emph{sinus}; on the other hand we see from the Madrid Ms. N that Robert of Chester, who probably carried out the better “ordering” of Adelard's translation between about 1143 and 1150, knew the word \emph{sinus}. We do not know when Gerard of Cremona, who mostly did not date his translations, translated the astronomy of Jābir b. Aflaḥ, in which the word \emph{sinus} occurs, and so we cannot decide whether he or Robert of Chester was the first to use this word; but the probability is greater for the latter.\par
\par The \emph{versed sine} Adelard calls \emph{elgeib elmankūs} (or \emph{maʿkūs}) = the reversed sine; but he translates \emph{mankūs} wrongly by \emph{diminutum}, confusing it with \emph{manqūṣ}, which means “diminished”, “reduced”, and is pronounced almost the same. This may well be a proof that at that time one learned a language not so much from books as rather in spoken intercourse with people, and therefore in translating relied more on the spoken word than on the written or read one; other wrong spellings of Arabic words may well come from this too.\par
\par Adelard's wrong translations, \emph{mustawī} by \emph{planus} and \emph{mankūs} by \emph{diminutus}, Robert of Chester did not correct; he too calls the sine \emph{sinus planus} and the versed sine \emph{sinus diminutus}.\par
\par We must also point out in particular that in the table of sines, in the four columns headed \emph{Semitae numerorum}, Kh. gives the angles of all four quadrants that have the same sine (apart from its sign); thus in the 6th row of numbers:\par
\begin{center}\small\begin{tabular}{cccccccc|l}
\emph{Sig.} & \emph{Gr.} & \emph{Sig.} & \emph{Gr.} & \emph{Sig.} & \emph{Gr.} & \emph{Sig.} & \emph{Gr.} &  \\
0 & 6 & 5 & 24 & 6 & 6 & 11 & 24 & that is: \\
\multicolumn{2}{c}{6°} & \multicolumn{2}{c}{174°} & \multicolumn{2}{c}{186°} & \multicolumn{2}{c}{354°} &  \\
\end{tabular}\end{center}
\par The sine is given in \emph{partes} (the radius = 60 of these), or, as the text has it, \emph{partitiones}, \emph{dakaicae} = Arab. \emph{daqāʾiq} = minutes, and \emph{theniae} (with the article \emph{eltheniae}) = Arab. \emph{thāniya} = seconds\footnote{See above, p. 36.}. Thus the sine of 6° = 6$^{p}$ 16′ 18″ = 6 + $\tfrac{16}{60}$ + $\tfrac{18}{60^{2}}$ = 6.27166, which gives, with the radius put = 1, sin 6° = 0.10453 (correct value = 0.104528).\par
\par Where the tables deviate from those of Batt., we have taken the reading of the manuscripts into the table when the error was only 1″, but the correct value of Batt. when the error was more than 1″, contrary to the procedure of Bjørnbo in the paper named above, who everywhere took the values of C and O, even when they were wrong.\par
\par \Pg{70}We are of Bjørnbo's opinion that Kh. took his table of sines from the Ptolemaic table of chords by halving the angles and the chords, for he himself mentions the \emph{Almagest} at the end of the chapter; perhaps the rather frequent deviations of 1″ also come from this, one rounding up and the other down in halving odd numbers of seconds\footnoteT{Here the print has a figure across the page, which is not reproduced; it stands at the end of chapter 23 in Ms. N only, and Suter explains it on p. 71: a circle divided into the twelve signs, numbered I to XII; the vertical and the horizontal diameter and the chords that join the points of division, two horizontal chords above and below the horizontal diameter and two vertical chords on either side of the vertical one, forming a tall and a wide rectangle; along the upper half of the vertical diameter «Sinus diminuti» (versed sines), below the horizontal diameter «Sinus plani» (sines); above the horizontal diameter, which small strokes divide into six parts, six times «XX»; below it, in the segments between the circle and the vertical chords and between the chords, from the left: «Gr. VIII M. II S. XIX», «Gr. XXI M. LVII S. XLI», «Gr. XXX», «Gr. XXX», «Gr. XXI M. LVII S. XLI», «Gr. VIII M. II S. XIX».}. It would not, to be sure, be impossible that the Indian scholar who came to the court of al-Manṣūr about the year 750 brought with him in his \emph{Siddhānta} a table of sines proceeding from degree to degree, which was then reduced from the Indian radius 3438 to the radius 60. But so far we know of no work earlier than the \emph{Siddhānta-Śiromaṇi} of Bhāskara in which the computation of the sine of 1° is taught. It must not go unmentioned, moreover, that a probably somewhat older contemporary of Kh., Yaʿqūb b. Ṭāriq, composed two works, entitled “on the division of the \emph{kardaja}” (that is, of the sine or arc of 3° 45′) and “astronomical tables composed after the \emph{Sindhind}, from degree to degree”.\footnote{See H. Suter, \emph{Die Mathematiker und Astronomen der Araber}, etc., p. 4.} From this one must, if \Pg{71}these titles are exact, draw the conclusion that tables of sines proceeding from degree to degree already existed before Muḥ. b. Mūsā al-Kh.\par
\par As for the rules for finding the sine and the versed sine of an angle and conversely, I shall hardly need to add anything to the text. Strange is the remark in paragraph 4, “\emph{ut in sequenti tabula subnotatur}” (as is noted below in the following table), as if another table were needed to find the angle from the sine. We point out here that A. Bjørnbo was of the opinion\footnote{See the \emph{Festskrift til} H. G. Zeuthen cited above, p. 14, and v. Braunmühl, \emph{Vorlesungen über Geschichte der Trigonometrie}, Leipzig 1900, I, p. 84.} that Kh.'s work might also have contained a table of arc sines, such as occurs in Abū l-Ḥasan ʿAlī of Morocco, who calls it the “table of Kh.”.\par
\par I add here a figure\footnoteT{The figure stands on p. 70; see the note there.} that stands at the end of this chapter, but only in Ms. N, and that certainly did not belong to Adelard's autograph.\par
\par The numbers all round, from I to XII, mean the twelve signs (\emph{signa}) of 30° each; the six XX above the horizontal diameter each designate a third of the radius, thus 20 partes; below the horizontal diameter P. (= partes) should stand instead of Gr. (= gradus); one can then read from the figure: sin I = sin 30° = sin V = sin 150° = sin VII = sin 210° = sin XI = sin 330° = XXX P.; sin (I + II) = sin 60° = sin 120° = sin 240° = sin 300° = XXX P. + XXI P. LVII M. XLI S. = LI P. LVII M. XLI S.; sin (I + II + III) = sin 90° = sin 270° = LI P. LVII M. XLI S. + VIII P. II M. XIX S. = LX P. — It should also be noted that in this figure the lines of the sines are horizontal; the translation of \emph{jayb mustawī} by \emph{sinus planus} would then be easier to understand.\par
\UnitHead{Chapter 24. On finding the latitude of any place}{Chapter 24. On finding the latitude of any place}{K14-4}
\par\begin{center}(Text p. 18.)\end{center}
\par First it is stated that, when the sun stands at the beginning of Aries or of Libra (that is, on the equator), its meridian altitude (\emph{H}) is to be subtracted from 90° to obtain the geographical latitude of the place (\emph{{\greekfont φ}} = 90° − \emph{H}).\par
\par If the sun stands in another sign, one subtracts its declination from its meridian altitude when it is northern, or adds it to the meridian altitude when it is southern, and subtracts the result from 90° (\emph{{\greekfont φ}} = 90° − (\emph{H} {\SignFont\upshape ∓} \emph{{\greekfont δ}})).\par
\par Then the method is given of determining with the astrolabe the altitude of the pole, which is equal to the geographical latitude of the place.\par
\par Finally the method is treated of determining the altitude of the pole, that is, the geographical latitude, from the greatest and the least altitude of a circumpolar star ($\varphi = h' + \dfrac{h''-h'}{2} = \dfrac{h'+h''}{2}$).\par
\par The same determinations are also found in Batt.\footnote{Ed. Nallino, I, 15, 29–30.}\par
\UnitHead{Chapter 25. The ascensions of the signs of the zodiac in the right sphere}{Chapter 25. The ascensions of the signs of the zodiac in the right sphere}{K14-11}\Pg{72}
\par\begin{center}(Text pp. 18–19, Tables 59–59$^{b}$.)\end{center}
\par This chapter with the tables that belong to it teaches how to compute the \emph{right ascension} (\emph{a}) of every degree of the ecliptic in the right sphere (that is, at the equator, or, as our tables have it, \emph{secundum terram Arin}, for the land of Arin)\footnote{\emph{Arin} (\emph{Azin}, today \emph{Ujjain}, see above p. 33) was placed on the equator by the older Arab astronomers and regarded as the middle of the inhabited half of the earth.}. This is done by the same formula as in Batt., namely\par
\par\begin{center}$\sin a = \dfrac{\sin\delta\cdot\cos\varepsilon}{\cos\delta\cdot\sin\varepsilon}\cdot r$.\end{center}
\par Since Batt., like Kh., does know the tangents and cotangents, but has tables of them only for the radius (gnomon) 12, they do not use the simpler formula sin \emph{a} = tg \emph{{\greekfont δ}} · cot \emph{{\greekfont ε}}, which corresponds to the formula $\sin a = \dfrac{\mathrm{tg}\,b}{\mathrm{tg}\,B}$ of the right-angled spherical triangle\footnoteT{Here the print has a figure, which is not reproduced: a right-angled spherical triangle with \emph{B} at the left, \emph{C} at the lower right and \emph{A} at the upper right; the arc \emph{c} from \emph{B} to \emph{A} above, the arc \emph{a} from \emph{B} to \emph{C} below, the arc \emph{b} from \emph{A} down to \emph{C}, and the right angle at \emph{C}.}. In our case the angle \emph{B} is the obliquity (\emph{{\greekfont ε}}) of the ecliptic, \emph{c} a given arc of the ecliptic, \emph{b} the declination (\emph{{\greekfont δ}}) of the point \emph{A} of the ecliptic, and \emph{a} the corresponding arc of the equator, that is, the \emph{right ascension} of the point \emph{A}; the above formula therefore reads in our case:\par
\par\begin{center}$\sin a = \dfrac{\mathrm{tg}\,\delta}{\mathrm{tg}\,\varepsilon} = \mathrm{tg}\,\delta\cdot\cot\varepsilon = \dfrac{\sin\delta\cdot\cos\varepsilon}{\cos\delta\cdot\sin\varepsilon}$.\end{center}
\par The Arab mathematicians still had to multiply this expression by \emph{r} = 60$^{p}$, because their trigonometric functions were computed for this radius. We give an \emph{example}.\par
\par What right ascension corresponds to the arc 25° of Aries? For this case we must first look up in Table 21 the declination (obliquatio) \emph{{\greekfont δ}} of the sun for the 25th degree of Aries; it is 9° 50′ 21″; Kh. takes the obliquity of the ecliptic as 23° 51′; the computation by the formula sin \emph{a} = tg \emph{{\greekfont δ}} · cot \emph{{\greekfont ε}} gives (computed with five-place logarithms) the value 23° 5′ 50″ for \emph{a}; our Table 59 has, at 25° of longitude (in the section \emph{Horoscopus arietis}), the number 113° 5′ 49″; since Kh., like Batt. in the corresponding tables\footnote{Ed. Nallino, II, 61–64.}, counts the right ascension from 0° of Capricorn instead of from 0° of Aries, one must subtract 3 signs = 90°, and so obtains \emph{a} = 23° 5′ 49″.\footnote{In six tests we have not met an error greater than 1″–2″.} — Batt.'s values naturally differ from Kh.'s, since he takes the obliquity of the ecliptic = 23° 35′; Batt. also contents himself with degrees and minutes: in our example he has a right ascension of 113° 8′ (= 23° 8′).\par
\par \Pg{73}That some Arab astronomers count the right ascension from 0° of Capricorn has its reason in this, that this right ascension, increased by 90°, then gives precisely the degree of the equator that rises at the same time as the degree of the ecliptic in question culminates; and this was important to know for astrological purposes. In our example, then, the point of the equator with the right ascension 113° 5′ 49″ rises at the same time as the point of the ecliptic with the longitude 25° culminates.\footnote{We should say: the point of the ecliptic with the longitude 25° rises at the same time as the point of the equator with the right ascension 23° 5′ 49″.}\par
\par Kh. does not yet use the expression “sine of the complement” (that is, cosine), as Batt. does; he says: “subtract the arc from 90° and look up the sine of the difference.” The concept of a special function, then, was not yet as far developed in Kh. as in Batt.\par
\par The explanation of the tables agrees almost throughout with that in Batt.\footnote{Ed. Nallino, I, 13–14.} We still have to mention some errors in our tables. The right ascensions of the first three degrees of Cancer (Table 59$^{a}$) do not agree in the seconds with those of Capricorn, although they should be the same increased by 180°; our computation even gives deviations in the minutes in part. These may also come from errors in the declinations, for at longitudes that differ by not more than 1°–2° from 90° or 270°, a difference of 1″ in the declination causes one of 1$\tfrac{1}{2}$′ to 2$\tfrac{1}{2}$′ in the right ascension: if, for example, one takes the declination at 89° or 271° as 23° 50′ 46″ instead of the 23° 50′ 47″ of the table, the right ascension of 1° of Capricorn becomes = 1° 5′ 55″ instead of 1° 3′ 30″. By our computations the following values result for the right ascensions of the first three degrees of Capricorn and of Cancer, if one takes the declinations as they are given in Tables 23 and 24:\par
\begin{center}\small\begin{tabular}{rr}
\emph{Capricorn:} & \emph{Cancer:} \\
1° 3′ 30″ & 181° 3′ 30″ \\
2° 10′ 34″ & 182° 10′ 34″ \\
3° 16′ 48″ & 183° 16′ 48″ \\
\end{tabular}\end{center}
\par Of these values only the one for 3° of Cancer agrees, to 1″, with the one given in Table 59$^{a}$.\par
\UnitHead{Chapters 26 and 26$^{a}$. With how many degrees (of the equator) each sign of the zodiac rises at any place}{Chapters 26 and 26a. With how many degrees (of the equator) each sign of the zodiac rises at any place}{K14-24}
\par\begin{center}(Text pp. 20 and 21, Tables 59–59$^{b}$.)\end{center}
\par For places that have a given latitude, and so do not lie on the equator, our work has no special tables for determining the ascensions: first, probably, because the number of tables would grow very large if one wanted to take many latitudes into account, and secondly perhaps also because these so-called oblique ascensions (\emph{ascensiones obliquae}) can be derived from the right ones in a fairly simple way. Batt. has such tables for the latitudes \Pg{74}of \emph{Mecca}, \emph{Baghdad}, \emph{Ḥarrān} and \emph{al-Raqqa}, and then also generally for all places whose day lengths amount to 13$^{h}$ to 16$^{h}$, at intervals of 15$^{m}$.\footnote{Ibid., II, 65–71.}\par
\par Batt. gives two ways of determining the oblique ascensions for any place\footnote{Ibid., II, 27 and 187–189.}; the second, which Kh. too follows in part, is the following:\par
\par\begin{center}\emph{Ascensio obliqua} for the latitude \emph{{\greekfont φ}} = \emph{ascensio recta} {\SignFont\upshape ∓} arc sin (tg \emph{{\greekfont φ}} · tg \emph{{\greekfont δ}}),\end{center}
\par where \emph{{\greekfont δ}} is the declination of the point of the ecliptic whose \emph{ascensio obliqua} is sought. This formula follows from the following figure\footnoteT{Here the print has a figure, which is not reproduced: a circle, the meridian, with \emph{Z} at the top, \emph{P} at the upper right, \emph{A} and \emph{E} at the upper left, \emph{H} and \emph{T} at the ends of a horizontal dashed diameter and \emph{K} and \emph{Q} at the lower right; dashed lines from the centre to \emph{Z}, \emph{P}, \emph{A}, \emph{E}, \emph{K} and \emph{Q}; within the circle the arcs of the horizon from \emph{H} through \emph{D} and \emph{L} to \emph{T}, of the equator from \emph{A} through the vernal point, \emph{D} and \emph{N} to \emph{Q}, of the ecliptic from \emph{E} through the vernal point and \emph{L} to \emph{K}, and of the circle of declination from \emph{P} through \emph{L} to \emph{N}; the angle at \emph{D} is marked.}, from which Ptol. too derives it by means of the theorem of Menelaus\footnote{Ed. Heiberg, I, 121 ff.; translation by K. Manitius, I, 83–85.}: let \emph{HT} be the horizon, \emph{AQ} the equator, \emph{EK} the ecliptic, {\SignFont\upshape ♈} the vernal point, \emph{P} the north pole of the equator and \emph{PN} the circle of declination of the rising point \emph{L} of the ecliptic; then {\SignFont\upshape ♈}\emph{N} is the right ascension and {\SignFont\upshape ♈}\emph{D} the oblique ascension of the point \emph{L} of the ecliptic, and \emph{DN} is the so-called \emph{ascensional difference}. Now\par
\par\begin{center}{\SignFont\upshape ♈}\emph{D} = {\SignFont\upshape ♈}\emph{N} − \emph{DN},\end{center}
\par but \emph{DN} follows from the right-angled spherical triangle \emph{DNL} in the following way: \emph{LN} is the declination \emph{{\greekfont δ}} of the point \emph{L}, and {\SignFont\upshape ∠} \emph{D} = arc \emph{AH} = 90° − \emph{{\greekfont φ}}; by the formula given above ($\sin a = \dfrac{\mathrm{tg}\,b}{\mathrm{tg}\,B}$), then,\par
\par\begin{center}$\sin DN = \dfrac{\mathrm{tg}\,LN}{\mathrm{tg}\,D} = \dfrac{\mathrm{tg}\,\delta}{\mathrm{tg}\,(90-\varphi)} = \mathrm{tg}\,\delta\cdot\mathrm{tg}\,\varphi$,\end{center}
\par so that {\SignFont\upshape ♈}\emph{D} = {\SignFont\upshape ♈}\emph{N} − arc sin (tg \emph{{\greekfont δ}} · tg \emph{{\greekfont φ}}), or: \emph{ascensio obliqua} = \emph{ascensio recta} − arc sin (tg \emph{{\greekfont δ}} · tg \emph{{\greekfont φ}}). If the declination \emph{{\greekfont δ}} of the point \emph{L} is southern, the arc whose sine is tg \emph{{\greekfont δ}} · tg \emph{{\greekfont φ}} must be added to the \emph{ascensio recta}.\par
\par \Pg{75}Batt. again uses the formula in the form $\sin DN = \dfrac{\sin\varphi\cdot\sin\delta}{\cos\varphi\cdot\cos\delta}$, whereas Kh., curiously, determines tg \emph{{\greekfont φ}} with the help of the gnomon, or takes it from the table of shadows (the table of cotangents)\footnote{The shadows, especially the equinoctial shadow (= tg \emph{{\greekfont φ}}), occur already in the \emph{Sūrya-Siddhānta}, though it has no table of them: each time the sine is divided by the cosine, or the shadow is determined separately for each case with the help of the gnomon.}; but he does not multiply this tg \emph{{\greekfont φ}} by tg \emph{{\greekfont δ}}, but, in the various cases, by certain factors whose origin is not really clear and which also give only approximate results. Kh. first gives the rule for determining the \emph{ascensio obliqua} for the whole sign of Aries; he says: multiply the noon shadow of the sun, when it is in the 1st degree of Aries (really the 0th), by 114 and divide the product by 115, and you have (in degrees) what must be subtracted from, or added to, the right ascension to obtain the \emph{ascensio obliqua}. We give below the computation for \emph{al-Raqqa} (geogr. latitude = 36°) by this rule:\par
\par\begin{center}cot 54° = tg 36° = 8$^{p}$ 44′ (by Table 60); 8$^{p}$ 44′ · $\tfrac{114}{115}$ = 8$^{p}$ 39′;\end{center}
\par Kh. puts this = 8° 39′ and subtracts it from the right ascension of 30° of Aries, which by Table 59 is = 117° 50′ − 90° = 27° 50′; this gives the \emph{ascensio obliqua} of the sign of Aries for \emph{al-Raqqa} = 19° 11′. By the tables of Batt. one obtains 19° 21′; Ptol. (l. c.) has 19° 12′, since he takes about the same obliquity of the ecliptic as Kh. If one computes \emph{DN} by Batt.'s formula, but with Kh.'s values, one obtains 8° 37$\tfrac{1}{2}$′; the factor $\tfrac{114}{115}$ is therefore a little too large.\par
\par If one were to compute sin \emph{DN} entirely with Kh.'s table of cotangents, one would first find in Table 21 the declination of 30° of Aries = 11° 39′ 50″ = 11° 40′, to which belongs in Table 60 the tangent 2$^{p}$ 28′, and one would then have sin \emph{DN} = 8$^{p}$ 44′ · 2$^{p}$ 28′; but since the tangents are computed for the gnomon 12 and the sines for the radius 60, this must be divided by 144 and multiplied by 60, which gives:\par
\par\begin{center}$\sin DN = \dfrac{8^{p}\,44'\cdot 2^{p}\,28'\cdot 60}{144} = 8^{p}\,44'\cdot\dfrac{148}{144} = 8^{p}\,59'$;\end{center}
\par to this sine there belongs in Table 58 the angle 8° 37′. To obtain the sine, then, one must multiply 8$^{p}$ 44′ by the factor $\tfrac{148}{144}$; to obtain the angle at once, Kh. multiplies 8$^{p}$ 44′ by the factor $\tfrac{114}{115}$; whether and how this was obtained from the former can hardly be found out; in our view it was found empirically.\par
\par If one wants the ascensional difference for only a certain part of the sign, one simply takes the corresponding part of the factor, for example, for 15°, only $\tfrac{57}{115}$; this then gives the value 4° 19$\tfrac{1}{2}$′ for the arc \emph{DN}, and so the value 9° 26$\tfrac{1}{2}$′ for the \emph{ascensio obliqua} of 15° of Aries (Batt. has 9° 28′).\par
\par For the sign of Taurus Kh. has the factor $\tfrac{113}{116}$; there is certainly an error here: the ascensional difference becomes too great; a better value would be $\tfrac{103}{126}$. — For the sign of Gemini Kh. gives the factor $\tfrac{1}{3}$, which is fairly good: with this factor one obtains for the \emph{ascensio obliqua} of this sign the value 29° 21′, while Batt. has 29° 17′. — In the \emph{Sūrya-Siddhānta} we find nothing of these factors, but in verses 60 and 61 of chapter 2\footnote{\emph{Transl.} Burgess, pp. 232–233. The same rule is also found in the \emph{Pañcasiddhāntikā}, ed. G. Thibaut, p. 27 of the translation.} the rule\par
\par\begin{center}$\sin\text{ascens.-diff.} = \dfrac{\sin\delta\cdot\mathrm{tg}\,\varphi\cdot 60}{12\cdot\cos\delta}$,\end{center}
\par which agrees with Batt.'s formula given above.\par
\par \Pg{76}The same factor that is used for Aries also gives the ascensional difference for Pisces, Virgo and Libra; the factor for Taurus also that for Aquarius, Leo and Scorpio; and the factor for Gemini also that for Capricorn, Cancer and Sagittarius. For Virgo, Libra, Leo, Scorpio, Cancer and Sagittarius, however, the ascensional difference must be added to the right ascension.\par
\par In Ms. O there is (in our text p. 21) a second rule for determining the \emph{ascensio obliqua}, which I have designated chapter 26$^{a}$, since it belongs to chapter 26 but is inserted at a later place. It agrees with the first in saying that one first determines tg \emph{{\greekfont φ}} (the equinoctial shadow); but then it does not multiply this by the factors named above, but refers to a table \emph{differentiarum} (of differences), which is missing, however, in O as well as in C and N, and in which one would find for any degree of the ecliptic the quantity by which tg \emph{{\greekfont φ}} must be multiplied to obtain the sine of the ascensional difference; this table must therefore have contained the values of tg \emph{{\greekfont δ}} · $\tfrac{60}{144}$ for all the degrees of the ecliptic, and it would in any case have given more exact values than one obtains with the factors above.\par
\UnitHead{Chapter 27. Determination of the size of the seasonal hour\footnote{As is well known, the Greeks, and then the Arabs too, distinguished seasonal and equinoctial hours; the former were the twelfth part of the day arc, and so unequal according to the geographical latitude and the season; the latter are our present equal hours.} of any day for any place}{Chapter 27. Determination of the size of the seasonal hour of any day for any place}{K15-20}
\par\begin{center}(Text pp. 20–21.)\end{center}
\par For this our work has the short rule: “Divide the day arc of the sun by 12, and you have what is sought.” This, to be sure, is only the length of a seasonal hour in degrees; to convert it into time one has, as is well known, the relations 15° = 1$^{h}$, 15′ = 1$^{m}$, 15″ = 1$^{s}$. But how the day arc is to be determined our tables give no information.\footnote{See Batt., ed. Nallino, I, 28–29.}\par
\UnitHead{Chapters 28–28$^{b}$. How the two shadows are determined from the altitude of the sun, and conversely\footnote{These chapters, and the table, are missing in N.}}{Chapters 28–28b. How the two shadows are determined from the altitude of the sun, and conversely}{K15-23}
\par\begin{center}(Text pp. 21–22, Table 60.)\end{center}
\par Bjørnbo too has already spoken about this chapter in the paper cited above (p. 68). The radius (\emph{gnomon}) is thus taken = 12$^{p}$; as is well known, it was only Ḥabash\footnote{See \emph{Abhandlgn. z. Gesch. d. mathem. Wissensch.} X, 1900, p. 12.}, a contemporary of Batt., who computed the cotangents for the radius 60, and this probably after the appearance of Batt.'s tables, since otherwise Batt. would surely have used these tables.\par
\par \Pg{77}Kh. calls the cotangent \emph{al-ẓill al-mustawī} = the plane or straight shadow, and the tangent \emph{al-ẓill al-mankūs} (or \emph{maʿkūs}) = the reversed shadow. Better are the expressions that Batt. uses, namely \emph{al-ẓill al-mabsūṭ} = the spread-out, plane shadow, and \emph{al-ẓill al-muntaṣib} or \emph{al-qāʾim} = the erect, vertical shadow; they are also called the \emph{second} and the \emph{first} shadow respectively.\par
\par Like Batt.\footnote{Ed. Nallino, I, 22–23 and 180–182.}, Kh. too teaches how to compute the tangent and the cotangent from the sine and the cosine, and so uses the formulas\par
\par\begin{center}$\cot\alpha = \dfrac{\cos\alpha}{\sin\alpha}\cdot 12$ and $\mathrm{tg}\,\alpha = \dfrac{\sin\alpha}{\cos\alpha}\cdot 12$.\end{center}
\par For determining the altitude of the sun \emph{{\greekfont α}} from the cotangent the following formula is given:\par
\par\begin{center}$\sin(90-\alpha) = \dfrac{\cot\alpha\cdot 60}{\sqrt{12^{2}+\cot^{2}\alpha}}$,\end{center}
\par that is, our present formula $\cos\alpha = \dfrac{\cot\alpha}{\sqrt{1+\cot^{2}\alpha}}$.\par
\par Batt. also gives the formulas $\sin\alpha = \dfrac{1}{\sqrt{1+\cot^{2}\alpha}}$, $\sin\alpha = \dfrac{\mathrm{tg}\,\alpha}{\sqrt{1+\mathrm{tg}^{2}\alpha}}$, $\cos\alpha = \dfrac{1}{\sqrt{1+\mathrm{tg}^{2}\alpha}}$, which are missing in Kh.\par
\par At the end Kh., as Batt. does too, refers to the table of cotangents, from which one can determine the shadow directly from the altitude and conversely.\par
\par We find, then, the first table of cotangents in Kh., and not only in Batt. or Ḥabash. But, contrary to Bjørnbo\footnote{\emph{Al-Chwarizmi’s Trigonometriske Tavler}, pp. 16–17.}, we refrain from drawing further conclusions from this; for it is by no means impossible, indeed not even improbable, that this table of cotangents was not in Kh.'s original work but was taken into his recension only by Maslama, who, after all, also composed an abridgement of Batt.'s tables: the rules of the text are very similar to those of Batt., only somewhat more briefly put. — As for the numerical values that we have taken into the table, the same holds here as we said above (p. 69) in discussing the table of sines.\par
\par \Pg{78}We come here to a statement of Bjørnbo's in the paper cited; he says, p. 17: “Abū l-Wafāʾ finally took the last and greatest step. He threw both the Greek division into sixty and the Indian division into twelve overboard, put the radius = 1, made, without changing the names, the plane shadow into what we understand by cotangent and the erect shadow into what we understand by tangent, and with this concluded the history of the development of the two functions.” In our view this goes somewhat too far. Even though Abū l-Wafāʾ says in one place of his Almagest: “from this it is clear that, if one takes the radius = 1, the ratio of the sine to the cosine is equal to the tangent, and the ratio of the cosine to the sine equal to the cotangent”\footnote{See Carra de Vaux, \emph{L’Almageste d’Aboul-Wéfa}, in the \emph{Journal asiatique}, Série 8, T. XIX, p. 420.}, it is nevertheless as good as certain that neither he himself nor any of his successors carried out this simplification in practice. Carra de Vaux, to be sure, says that Abū l-Wafāʾ set up a table of the tangent and the cotangent both for the radius 1 and for the radius 60; but this is not quite established, and the tables have unfortunately been lost; and even if such tables did once exist, they were certainly never used again afterwards.\par
\UnitHead{Chapter 29. On \emph{elbuht} and how to find it}{Chapter 29. On elbuht and how to find it}{K16-10}
\par\begin{center}(Text p. 22, Tables 61–66.)\end{center}
\par By \emph{al-buht} the Arab and especially the Persian\footnote{Ulugh Beg uses the expression in his tables.} astronomers understood the true, or unequal, motion of the sun, the moon and the planets in a day; the word comes from the Indian \emph{bhukti}, which literally means “enjoyment, pleasure”\footnote{\emph{Sūrya-Siddhānta}, transl. Burgess, p. 158: “\emph{When a planet has} enjoyed \emph{the whole troop of asterisms, it has made a complete revolution.}”}, and which became \emph{buht} in Persian and was so taken over by the Arabs\footnote{Batt. does not know the word, or does not use it; see I, XLIII and II, 324.}.\par
\par As is well known, the motion of the planets is non-uniform, faster at the perigee than at the apogee; we therefore find the greatest true motion at 180° of anomaly and the least at 0°. Our Tables 61–66 are therefore easy to understand: the column headed \emph{semitae numerorum} contains the anomalies from 0° to 180° (or 360° to 180°), the second column the true motion of the sun in an hour, the third column that of the moon in an hour\footnote{For the five planets the corresponding tables are missing, as they are in Batt.; for the sun and the moon the latter gives them only for anomalies at intervals of 6°; see II, 88.}. For the sun the minimum of the motion is 2′ 22″ and the maximum 2′ 34″; for the moon, on the other hand, 30′ 12″ and 35′ 40″ respectively; Batt. (l. c.) has the values 2′ 23″ and 2′ 33″ for the sun, 30′ 18″ and 36′ 4″ for the moon.\par
\UnitHead{Chapters 30 and 30$^{a}$. On the apparent diameters of the sun and the moon}{Chapters 30 and 30a. On the apparent diameters of the sun and the moon}{K16-14}
\par\begin{center}(Text pp. 22–23, Tables 61–66.)\end{center}
\par \Pg{79}The apparent diameters of the sun and the moon Kh. calls \emph{quantitas rotae solis} and \emph{rotae lunae} (the size of the wheel of the sun, of the moon). These diameters naturally change with the distance of these heavenly bodies from the earth. While Batt. contents himself with giving the maximum, the mean and the minimum size of these diameters\footnote{Ed. Nallino, I, 58.}, Kh. takes the trouble to determine these sizes for every degree of the anomaly from 0° to 180°. In columns 4 and 5 of Tables 61–66 he gives the apparent semidiameters (\emph{dimidium rotae solis vel lunae}) of the sun and the moon for the various anomalies. Here again constant factors occur, by which he multiplies the true \emph{daily} motions of the sun and the moon to obtain the apparent diameters: the factor for the sun is = 33′ = $\tfrac{33}{60}$ = $\tfrac{11}{20}$, that for the moon = 2′ 25″\footnote{The text has, wrongly, 2′ 16″.} = 2$\tfrac{25}{60}$ {\SignFont\upshape ∶} 60 = $\tfrac{145}{3600}$ = $\tfrac{29}{720}$. Batt.\footnote{Ibid., I, 97.} gives a similar rule for computing the diameter of the moon: he says that one must multiply the true \emph{hourly} motion of the moon by 5$\tfrac{7}{8}$ and divide by 6 (that is, multiply by $\tfrac{47}{48}$) to obtain the apparent diameter of the moon; if we reduce Kh.'s factor to the \emph{hourly} true motion, it becomes = $\tfrac{29}{30}$, and the difference of the two factors is = $\tfrac{1}{80}$. Let us give an \emph{example} each for the sun and for the moon.\par
\par The daily true motion of the sun for 90° of anomaly is, by Table 63, = 2′ 28″ · 24; this, multiplied by $\tfrac{11}{20}$, gives 32′ 33$\tfrac{1}{2}$″ for the diameter of the sun, and so 16′ 16$\tfrac{3}{4}$″ for the semidiameter; the table has 16′ 17″\footnote{For other values of the anomaly, to be sure, the computation does not always agree so well.}.\par
\par The daily true motion of the moon for 90° of anomaly is, by the same table, 32′ 56″ · 24; this, multiplied by $\tfrac{29}{720}$, gives 31′ 50″ for the diameter of the moon, and so 15′ 55″ for the semidiameter; the table has 15′ 58″.\par
\par One will not always find the correct value of the daily true motion of the sun and the moon, moreover, by multiplying the hourly one by 24, for the latter values are rounded in the tables partly down and partly up, and multiplication by 24 makes the error that many times greater. The text therefore also remarks that one finds the daily motion by determining the true place of the heavenly body for noon of a given day and likewise for noon of the following day: the difference of the two places is the \emph{buht} of this heavenly body.\par
\par The rules of Kh. named above for determining the apparent diameters of the sun and the moon seem to have been taken from the Indian works. The \emph{Sūrya-Siddhānta}, to be sure, does not contain them in this form, and its values for the mean diameters of the sun and the moon (32′ 25″ and 32′ respectively)\footnote{Transl. Burgess, pp. 269–272.} do not quite agree with those of Kh. (32′ 34″ and 31′ 50″), though the differences are small. But al-Bīrūnī has a passage in his work on India\footnote{\emph{Alberuni's India}, English edition by E. C. Sachau, London 1888 and 1910, II, 79 (Arabic text, London 1887, p. 241).} which proves to us that Kh. relied on Indian sources on this point: he says that the method of computing the apparent diameters of the sun and the moon contained in the \emph{Khaṇḍakhādyaka}\footnote{A work that Brahmagupta wrote after his \emph{Siddhānta}, built on somewhat different principles.} of Brahmagupta is the same as the one found in the tables of Kh. Indeed the Arabs knew early a translation of this work under the name \emph{al-Arkand}, which was later improved by al-Bīrūnī himself, as he states in the list of his writings. There is thus probably no doubt that Kh. knew and used this work of Brahmagupta too. (See further below, the end of chapter 33.)\par
\par \Pg{80}Interesting is the Latin word \emph{stupor}, by which Adelard renders the Persian \emph{buht}; we can probably translate it best by “stopping”, “halting” (in the motion), for when the heavenly body goes now faster, now slower, it stops, as it were, now and then in its motion\footnote{Compare Ovid's \emph{stupuit Ixionis orbis} (the wheel of Ixion stood still), or also \emph{stupuerunt verba palato} (the words stuck to the palate) (Georges' Latin dictionary).}; for the planets, to be sure, this expression fits better than for the sun and the moon.\par
\par I shall speak of the meaning of the 6th and 7th columns of Tables 61–66, which are headed \emph{primum dimidium rotae draconis} and \emph{secundum dimidium rotae draconis} (first and second half of the wheel of the dragon), at the end of chapter 33.\par
\UnitHead{The tables of the equation of days (the equation of time)}{The tables of the equation of days (the equation of time)}{K17-2}
\par\begin{center}(Tables 67 and 68.)\end{center}
\par The text of the present work seems to be defective, especially towards the end: for these tables (67 and 68), for example, there is no explanatory chapter in the text. The tables correspond to the columns of Batt.\footnote{Ed. Nallino, II, pp. 61–64.} headed \emph{Aequatio nychthemerōn}, except that Batt.'s tables begin with the 1st degree of Capricorn, Kh.'s with the 1st degree of Aries. A further difference is that Batt. gives the \emph{equation of time} (that is, the difference between \emph{true} and \emph{mean} time) in degrees and minutes, Kh., on the other hand, in minutes and seconds of time. But they agree in taking the equation of time as zero when true time is furthest behind mean time; at all other times, then, true time must be ahead of mean time, that is, the \emph{equation of time} that the table gives must always be \emph{subtracted}\footnote{See for this Batt., ed. Nallino, I, pp. 221–222 (after Schiaparelli).}.\par
\par \emph{Example:} The equation of time is \emph{zero} when the sun stands in the 22nd degree of Aquarius (Batt. has zero at the 18th and 19th degree, and 0° 1′ at the 22nd); at this position of the sun, then, the Arab astronomers took true time to be equal to mean time. When the sun stands in the 9th to 10th degree of Scorpio, the equation of time reaches its \emph{maximum}, 34$^{m}$ 28$^{s}$, that is, true time is ahead of mean time by 34$^{m}$ 28$^{s}$ (Batt. has 7° 54′ = 31$^{m}$ 36$^{s}$). If we compare these values with the present ones, the Arab equation of time zero corresponds to our equation of time of + 14$^{m}$ 31$^{s}$ on 12 Feb., and the Arab equation of time of 34$^{m}$ 28$^{s}$ to our equation of time of − 16$^{m}$ 18$^{s}$ on 18 Nov.; the difference of the two is 30$^{m}$ 49$^{s}$, so that Batt.'s value is nearer the truth than Kh.'s.\par
\par \Pg{81}Whom Kh. followed in this question we cannot decide; the maximum of the equation of time in Kh. agrees better with that in Ptol. (33$^{m}$ 30$^{s}$) than with that in Batt. In the \emph{Sūrya-Siddhānta} we have found no values of the equation of time\footnote{See \emph{Sūrya-Siddhānta}, transl. Burgess, pp. 220–221.}. — These tables of the equation of time are mentioned at the end of the following chapter, where the conversion of equinoctial hours into unequal hours and conversely is dealt with.\par
\UnitHead{Chapter 31. On the conjunction and the opposition of the sun and the moon}{Chapter 31. On the conjunction and the opposition of the sun and the moon}{K17-7}
\par\begin{center}(Text pp. 23–25, Tables 69–72.)\end{center}
\par With the help of these tables, according to the rules of the text, the times of the conjunctions and oppositions of the sun and the moon are to be determined. On p. 25 of the text we find the remark: \emph{Haec autem omnia pro situ Cordubae dicta sunt} (all this, however, has been said for the situation of Córdoba). This proves that Maslama al-Majrīṭī either added these tables of his own accord or recomputed them from those of Kh. for the meridian of Córdoba. Tables arranged quite similarly, but differing somewhat in their numerical values, are found in Batt.\footnote{Ed. Nallino, II, 306–307.}; but they contain only the \emph{anni collecti}, and the \emph{anni plani} (\emph{singuli}) and the \emph{menses} are missing; Nallino has rightly counted them among the spurious tables. For the meridian of \emph{al-Raqqa} Batt. set up tables of the syzygies in Roman years\footnote{Ibid., II, 84–86.}, which are arranged on the same principle as ours, except that the \emph{anni collecti} proceed by 24 years instead of 30; and they begin with 1 \emph{Ādhār} 879 of the Seleucid era = 1 March 568 AD.\par
\par The determination of a \emph{conjunction} or an \emph{opposition} of the sun and the moon rests, of course, in the first place on the computation of the true longitudes of these heavenly bodies for the moment in question. If one finds these longitudes equal, the two bodies are then in conjunction; if they differ by 180°, they are in opposition. If they do not quite agree, one can infer from the difference of the longitudes and from the true hourly motions of the two bodies the time that is still needed until they come into true conjunction or opposition. The justification of this correction is set out by Batt.\footnote{Ibid., I, 93–94; see also the explanations of Schiaparelli, I, 273–274.} in a rather complicated way, but he also gives a simpler account\footnote{Ibid., I, 94, ll. 26–34.}, which agrees with that of Ptol.\footnote{Ed. Heiberg, I, 472–475; translation by K. Manitius, I, 342–349.} and which Kh. too follows. — Both the construction of the tables and the computation of the conjunctions and oppositions with their help are far simpler in the Muslim reckoning than in the Seleucid or the Julian, because the conjunctions fall on the last or the second-to-last day of the month, and the oppositions always in its middle.\par
\par \Pg{82}The principle on which the arrangement of the tables, and especially of the 3rd column, headed \emph{Dies et horae} or only \emph{Horae}, rests is to determine the time needed from a given moment, or from a given position of the heavenly bodies (\emph{radix}), to the conjunction sought, that is, until the two bodies have the same longitude. Thus in Table 69, for example, in the column headed \emph{Dies et horae}, the number 29$^{d}$ 0$^{h}$ 14$^{m}$ 46$^{s}$, which corresponds to the number 31 in the column headed \emph{Anni collecti}, means that after 30 whole Arab years have elapsed\footnote{The 31 means 0 \emph{Muḥarram} of the year 31 AH.} a further 29$^{d}$ 0$^{h}$ 14$^{m}$ 46$^{s}$ are needed until the next conjunction of the sun and the moon; and correspondingly in the other tables for the single years and the months: for example, after the end of the 1st year after the 30th, a further 8$^{h}$ 48$^{m}$ 7$^{s}$\footnoteT{So Table 71 prints it; the value of the table is 8$^{h}$ 48$^{m}$ 27$^{s}$ (see the critical notes of the edition of the tables), and Suter quotes the misprint.} are needed, by Table 71, until the next conjunction. — The 4th column gives the mean longitude of the sun and the moon for the moment in question, the 5th column contains the anomaly of the moon, and the 6th the argument of latitude of the moon for that moment. The first column contains the \emph{signa} (or \emph{notae}) of the years or months concerned; we need neither it nor the last (6th) column for determining the conjunctions and oppositions.\par
\par We give below an \emph{example} of the computation of a conjunction by the rules of the text, although here, perhaps more than anywhere else, it would have been necessary to give an intelligible German translation of this chapter, which is badly translated, or has at least come down to us in very bad Latin. We choose for it the moment of the eclipse of the sun that Batt. cites in chapter 30\footnote{Ed. Nallino, I, 56–57.}, and whose elements Schiaparelli has recomputed from Batt.'s tables\footnote{Ibid., I, 228–229.}.\par
\par At what time did the \emph{conjunction} of the sun and the moon (the \emph{new moon}) take place at the end of the month \emph{Muḥarram} of the year 288 AH (= 22/23 Jan. 901)?\par
\par From Tables 69–72 one finds:\par
\begin{center}\footnotesize\begin{tabular}{lrrrr}
 & \emph{Dies et horae} & \begin{tabular}[b]{@{}c@{}}\emph{Longit. media}\\\emph{solis et lunae}\end{tabular} & \begin{tabular}[b]{@{}c@{}}\emph{Anomalia}\\\emph{lunae}\end{tabular} & \begin{tabular}[b]{@{}c@{}}\emph{Argumentum}\\\emph{latitud. lunae}\end{tabular} \\
at the end of the year 270 & 29$^{d}$ 2$^{h}$ 4$^{m}$ 51$^{s}$ & 123° 49′ 42″ & 252° 31′ 18″ & 30° 19′ 50″ \\
for 17 years & 0 5 43 50 & 177 30 35 & 226 38 40 & 136 48 10 \\
for Muḥarram & 0 0 0 0 & 0 0 0 & 0 0 0 & 0 0 0 \\
for 0 Muḥ. 288: & 29$^{d}$ 7$^{h}$ 48$^{m}$ 41$^{s}$ & 301° 20′ 17″ & 119° 9′ 58″\footnotemark{} & 167° 8′ 0″ \\
\end{tabular}\end{center}
\addtocounter{footnote}{-1}\stepcounter{footnote}\footnotetext{Strictly 479° 9′ 58″; but 360° = a whole revolution is subtracted from it.}
\par The sun and the moon, then, have the same mean longitude, namely 301° 20′ 17″, on 29 Muḥ. of the year 288 AH at 7$^{h}$ 49$^{m}$ in the evening\footnote{The days of the tables begin at noon, at the culmination of the sun. I have rounded the seconds up to 1 min.}. One must now determine the true places of the sun and the moon for this time by chapters 8 and 9 (Tables 21–26):\par
\begin{center}\small\begin{tabular}{lrlr}
 & \emph{Mean longitude of the sun:} &  & \emph{Mean longitude of the moon:} \\
 & 301° 20′ 17″ &  & 301° 20′ 17″ \\
Equation of the sun &  & Equation of the moon &  \\
by Table 25\footnotemark{} & 1° 30′ 39″ & by Table 24\footnotemark{} & 4° 15′ 45″ \\
 & 302° 50′ 56″ &  & 297° 4′ 32″ \\
\end{tabular}\end{center}
\addtocounter{footnote}{-2}\stepcounter{footnote}\footnotetext{This equation corresponds to the argument 223° 25′ 17″, which one obtains by subtracting the longitude of the apogee (77° 55′) from the mean longitude of the sun.}\stepcounter{footnote}\footnotetext{This equation corresponds to the argument 119° 9′ 58″ (119° 10′), namely the anomaly of the moon.}
\par \Pg{83}The moon, then, is 302° 50′ 56″ − 297° 4′ 32″ = 5° 46′ 24″ behind the sun, that is, the sun \emph{praecessit} (has gone ahead), as the text puts it; the true conjunction therefore takes place only after the time found above, and, if the sun were regarded as standing still, by as many hours as the true hourly motion of the moon is contained in the difference of longitude 5° 46′ 24″. But since the longitude of the sun changes too, a correction is added to this difference, namely its 12th part = 28′ 52″, and the difference of longitude is now 6° 15′ 16″; now 6° 15′ 16″ {\SignFont\upshape ∶} 34′ 20″\footnote{This is, by Table 64, the true hourly motion of the moon for the argument (anomaly) 119° 10′.} = 10$^{h}$ 55$^{m}$ 48$^{s}$; the \emph{true conjunction}, then, took place on 29 Muḥ. 288 AH, 18$^{h}$ 44$^{m}$ 29$^{s}$ after the culmination of the sun, that is, on 23 January 901 at 6$^{h}$ 44$^{m}$ in the morning\footnote{Schram's tables (\emph{Kalendariograph. u. chronolog. Tafeln}, Leipzig, 1908) give 23 Jan., 6$^{h}$ 43$^{m}$ in the morning for \emph{Greenwich}, and so 6$^{h}$ 23$^{m}$ for \emph{Córdoba}, which lies c. 5° further west; the error would thus be c. 21$^{m}$; but Schram remarks that his tables give the syzygies to an accuracy of about half an hour, so that the error might reduce to zero. For \emph{Antioch} Schiaparelli finds the time of the true conjunction = 9$^{h}$ 20$^{m}$ in the morning; since \emph{Antioch} lies c. 41° east of \emph{Córdoba}, the time for Antioch by our tables would come to 9$^{h}$ 29$^{m}$, and so differ by only 9$^{m}$ from Batt.'s tables.}.\par
\par To obtain the places of the sun and the moon at the true conjunction, one must, according to the precepts given, add 6° 15′ 16″ to the true place of the moon, but to that of the sun only the 12th part of the difference of longitude (or the 13th part of 6° 15′ 16″), that is, 28′ 52″; and one must then of course obtain the same quantities:\par
\begin{center}\small\begin{tabular}{lrlr}
Place of the sun: & 302° 50′ 56″ & Place of the moon: & 297° 4′ 32″ \\
 & + 28′ 52″ &  & + 6° 15′ 16″ \\
 & 303° 19′ 48″ &  & 303° 19′ 48″ \\
\end{tabular}\end{center}
\par This, then, is the \emph{true longitude} of the sun and the moon at the time of their true conjunction, on 29 Muḥ. 288 at 6$^{h}$ 44$^{m}$ in the morning.\par
\par \Pg{84}The reason why $\tfrac{1}{12}$ of the difference of longitude must be added to the longitude of the sun, and $\tfrac{13}{12}$ of it to the longitude of the moon (in the case where the moon is ahead of the sun they must be subtracted), is surely that the mean hourly motions of the sun and the moon are nearly as 1 {\SignFont\upshape ∶} 13\footnote{Mean hourly motion of the moon 32′ 56″, of the sun 2′ 28″; the division gives 13.3.}. Let \emph{K} be the place of the true conjunction, \emph{S} the place of the sun and \emph{M} that of the moon, after the correction by their equations\footnoteT{Here the print has a small figure, which is not reproduced: a horizontal line with short upright strokes at its two ends and one between them, nearer the right end, with \emph{M}, \emph{S} and \emph{K} above the left end, the inner stroke and the right end; it is not drawn to scale (\emph{MS} about three times \emph{SK}, where the text makes it twelve times).}; then 13 parts fall to \emph{MK} and one part to \emph{SK}, and so 12 parts to \emph{MS}. One therefore divides the difference of longitude \emph{MS} by 12 and so has one part; adding this to the longitude of the sun, and thirteen times it to the longitude of the moon, one has the same true longitude of the two bodies at the time of the conjunction; and \emph{MK} divided by the true hourly motion of the moon gives the number of hours that the two bodies still have to run from the positions \emph{M} and \emph{S} to the conjunction \emph{K}. Since the ratio of the hourly motions of the two bodies is not constant, this procedure does not give quite exact results, and Batt. therefore remarks that his more lengthy procedure is the more correct one.\par
\par Two further precepts are now given, which are of no consequence for determining the conjunction, but are needed for the later computation of an eclipse that may occur. \emph{First} it is remarked (p. 24) that the equation of the moon is also to be subtracted from (or added to) the argument of latitude; in our case this gives 167° 8′ − 4° 15′ 45″ = 162° 52′ 15″; and to this quantity the corrected difference of longitude (6° 15′ 16″) is to be added, which gives 169° 7′ 31″ as the corrected argument of latitude. \emph{Secondly} it is added (pp. 24–25) that one must also correct the anomaly of the moon for the time correction found, 10$^{h}$ 55$^{m}$ 48$^{s}$; this is done by multiplying this quantity by 32′ 40″, that is, by the mean hourly motion of the anomaly of the moon; taking the decimal fraction 10.93 instead of 10$^{h}$ 55$^{m}$ 48$^{s}$, one obtains 5° 57′ 3″, and this, added to the anomaly (119° 10′)\footnote{This correction is to be subtracted when the moon is ahead of the sun.}, gives 125° 7′ = the corrected anomaly (in the text: \emph{argumentum ratum}).\par
\par At the end it is also stated how one must correct the time of the conjunction when the place of observation lies west or east of Córdoba.\footnote{The text is faulty here: \emph{oriente} and \emph{occidente} must be exchanged.} Then, for the case that one wants to convert equinoctial hours into unequal hours or conversely, it is pointed out that one must use the tables of the equation of time (67–68) for this.\par
\par We also give an \emph{example} of the computation of an \emph{opposition}.\par
\par At what time about the middle of the month \emph{Rabīʿ II} (which counts 29 days) of the year 146 AH did the \emph{opposition} of the sun and the moon (the \emph{full moon}) take place?\footnote{1 Rabīʿ II of the year 146 corresponds to 17 June 763.}\par
\par From Tables 70–72 one finds:\par
\begin{center}\footnotesize\begin{tabular}{lrrrr}
 & \emph{Dies et horae} & \begin{tabular}[b]{@{}c@{}}\emph{Longit. media}\\\emph{solis et lunae}\end{tabular} & \begin{tabular}[b]{@{}c@{}}\emph{Anomalia}\\\emph{lunae}\end{tabular} & \begin{tabular}[b]{@{}c@{}}\emph{Argumentum}\\\emph{latitud. lunae}\end{tabular} \\
at the end of the year 120 & 14$^{d}$ 6$^{h}$ 34$^{m}$ 9$^{s}$ & 279° 28′ 28″ & 29° 12′ 6″ & 67° 12′ 57″ \\
for 25 years & 4 11 21 & 91 37 55 & 185 3 54 & 201 10 50 \\
for 4 months & $\underline{1}$ 14 12 9\footnotemark{} & 87 18 58 & 77 27 4 & 92 0 42 \\
for 0 Rabīʿ II, 146: & 14$^{d}$ 0$^{h}$ 57$^{m}$ 39$^{s}$ & 458° 25′ 21″ & 291° 43′ 4″ & 360° 24′ 29″ \\
 &  & = 98° 25′ 21″ &  & = 0° 24′ 29″ \\
\end{tabular}\end{center}
\addtocounter{footnote}{-1}\stepcounter{footnote}\footnotetext{$\underline{1}$ means that this day is to be taken negatively, that is, subtracted from the sum of the other days and hours.}
\par One now determines the true longitudes of the sun and the moon by Tables 21–26:\par
\Pg{85}\begin{center}\small\begin{tabular}{lrlr}
 & \emph{Mean longitude of the sun:} &  & \emph{Mean longitude of the moon:\footnotemark{}} \\
 & 98° 25′ 21″ &  & 98° 25′ 21″ \\
Equation of the sun &  & Equation of the moon &  \\
by Table 21: & 0° 45′ 44″ & by Table 23: & 4° 36′ 36″ \\
 & 97° 39′ 37″ &  & 103° 1′ 57″ \\
\end{tabular}\end{center}
\addtocounter{footnote}{-1}\stepcounter{footnote}\footnotetext{For the further computations one keeps this longitude of the moon, which agrees with the mean longitude of the sun, although at the opposition it is really 180° greater (or smaller).}
\par The moon, then, is 5° 22′ 20″ ahead of the sun, and so the opposition has already taken place. One divides the difference of the longitudes by 12, which gives 26′ 52″, and adds this to 5° 22′ 20″, which gives 5° 49′ 12″; this quantity, divided by the true hourly motion of the moon, which by Table 63 is = 31′ 50″, gives 10.97$^{h}$ = 10$^{h}$ 58$^{m}$; and this difference of time, subtracted from the time at which the mean opposition took place, gives, in round numbers, 13$^{d}$ 14$^{h}$ as the time of the true opposition; that is, this \emph{opposition} took place on 13 Rabīʿ II, 146, 14$^{h}$ after 12 o'clock noon, that is, in our (civil) reckoning (with the day beginning at midnight), on 30 June 763 at 2 o'clock in the morning\footnote{Schram's tables give 30 June 763, 1$^{h}$ 41$^{m}$ in the morning for Greenwich, and so 1$^{h}$ 21$^{m}$ for Córdoba; the difference here would thus amount to c. 40$^{m}$; by Schram's remark quoted above it might reduce to 10$^{m}$.}.\par
\par To obtain the true longitudes of the sun and the moon at the opposition, one must subtract from the longitude of the sun found above the 12th part of the difference of the longitudes, that is, 26′ 52″, but from the longitude of the moon thirteen times that, that is, 5° 49′ 12″; one then obtains:\par
\begin{center}\small\begin{tabular}{lrlr}
Longitude of the sun: & 97° 39′ 37″ & Longitude of the moon: & 103° 1′ 57″ \\
 & − 26′ 52″ &  & − 5° 49′ 12″ \\
 & 97° 12′ 45″ &  & 97° 12′ 45″ \\
\end{tabular}\end{center}
\par This number gives the true longitude of the sun at the time of the opposition; that of the moon is therefore 180° greater, that is, = 277° 12′ 45″.\par
\par According to the precepts of the text, the argument of latitude of the moon and its anomaly must now still be corrected.\par
\begin{center}\small\begin{tabular}{lr}
argument of latitude & = 360° 24′ 29″ \\
 & − 5° 49′ 12″ \\
 & 354° 35′ 17″ \\
to this is added the equation of the moon: & 4° 36′ 36″ \\
corrected argument of latitude: & 359° 11′ 53″ \\
\end{tabular}\end{center}
\begin{center}\small\begin{tabular}{lr}
anomaly of the moon & = 291° 43′ 4″ \\
correction corresponding to the time correction of 10$^{h}$ 58$^{m}$\footnotemark{} & 5° 58′ 21″ \\
corrected anomaly\footnotemark{}: & 285° 44′ 43″ \\
\end{tabular}\end{center}
\addtocounter{footnote}{-2}\stepcounter{footnote}\footnotetext{As above, one must multiply the time correction by 32′ 40″, by the rule (p. 24).}\stepcounter{footnote}\footnotetext{Called \emph{anomalia aequata} by Batt.}
\par \Pg{86}From these examples one sees that Kh. (or al-Majrīṭī) on the whole followed Ptol. in determining the syzygies, as Batt. did too. The determinations of the \emph{Sūrya-Siddhānta} are incomplete and obscure in this respect\footnote{Transl. Burgess, p. 275.}.\par
\UnitHead{Chapter 33. How the eclipse of the moon is determined\footnote{The title of the text speaks of the sun and the moon, but in this chapter only the eclipses of the moon are in fact treated; the eclipses of the sun follow in chapter 35, and I have therefore changed the title in this sense.}}{Chapter 33. How the eclipse of the moon is determined}{K19-1}
\par\begin{center}(Text pp. 26–28, Tables 73–76.)\end{center}
\par Since chapter 32 appears to have been inserted at this place\footnote{It stands at this place only in Ms. C; in N only on fol. 56$^{r}$, that is, in its right place; in O it is missing altogether.}, as it does not fit at all the subject treated in chapters 31 and 33, namely the computation of the syzygies and of the eclipses, we take chapter 33 directly here.\par
\par Batt.'s rules for determining the time and the magnitude of eclipses are much more complicated than Kh.'s; he deals with them in chapters 43 and 44\footnote{Ed. Nallino, I, 96–113, and 275–282.}. Kh.'s tables, on the other hand, are worked out somewhat further: Batt. has left to computation various things that can be read off from the tables in Kh.\par
\par To compute an \emph{eclipse of the moon}, one first determines and notes: the time of the opposition, the true places of the sun and the moon, and the corrected argument of latitude. For our example of an \emph{opposition} given above (pp. 84–85) these quantities are the following:\par
\begin{center}\small\begin{tabular}{lll}
\emph{Time of the opposition:} & \begin{tabular}[b]{@{}c@{}}\emph{True place of the sun}\\\emph{(or of the moon):}\end{tabular} & \begin{tabular}[b]{@{}c@{}}\emph{Corrected argument}\\\emph{of latitude:}\end{tabular} \\
13 Rabīʿ II, 146, 14$^{h}$ & 97° 12′ 45″ (277° 12′ 45″) & 359° 11′ 53″ \\
\end{tabular}\end{center}
\par We must now first examine whether the moon is near the earth or far from it, for on this depends the size of the earth's shadow, and with it also the duration and the magnitude of the eclipse. According to a remark at the end, which stands only in Mss. C and N\footnote{In N the wording of C has been corrected by a later hand, see p. 28, note 1; in my account I follow Ms. C, whose wording I regard as the right one.}, the moon is \emph{in spatio longinquo} (at the far distance) when its longitude amounts to 89°–91°, and \emph{in spatio propinquo} (at the near distance) when it amounts to 269°–271°. In our case it amounts to c. 277° 13′, which is no longer taken as the near distance, at least as regards our computation; in this case, then, according to the precept, we must enter with the argument of latitude (359° 12′)\footnote{Batt.'s tables (II, 90) have the latitude itself in the 1st column.} both the tables \emph{pro spatio propinquo} and those \emph{pro spatio longinquo}. We now find in Table 76, for this argument of latitude (by interpolation), the following data:\par
\Pg{87}\begin{center}\small\begin{tabular}{lccc}
 & \emph{Digits of the eclipse:} & \emph{Min. and sec. of the casus:} & \emph{Min. and sec. of the half mora:} \\
 & 20$^{\text{dig.}}$ 14′ & 34′ 45″ & 27′ 8″ \\
\multicolumn{4}{c}{in Table 74 the following:} \\
 & 19$^{\text{dig.}}$ 16′ & 29′ 39″ & 21′ 2″ \\
differences: & 0$^{\text{dig.}}$ 58′ & 5′ 6″ & 6′ 6″ \\
\end{tabular}\end{center}
\par These are the differences of the three data at the near and the far distance. What, then, do these data mean? Our text tells us nothing about it; but Batt. instructs us in chapter 43.\footnote{Ed. Nallino, I, 97, l. 25 – 98, l. 5.} He, like Kh., understands by the \emph{digits} of the eclipse, in a \emph{partial} as well as in a \emph{total} eclipse of the moon, the distance from the edge of the earth's shadow to the edge of the disc of the moon that is convex towards it, in parts of which 12 go to the diameter of the moon.\par
\par If \emph{S} is the centre of the earth's shadow, which must lie in the ecliptic, and \emph{M} the centre of the moon (in Fig. 1 for a partial, in Fig. 2 for a total eclipse), then in Fig. 1 \emph{AM}, expressed in digits (inches), is the magnitude of the eclipse, in\par
\begin{center}\small\begin{tabular}{cc}
Fig. 1.\footnotemarkT{} & Fig. 2. \\
\end{tabular}\end{center}
\addtocounter{footnoteT}{-1}\stepcounter{footnoteT}\footnotetextT{Here the print has two figures side by side, which are not reproduced. Fig. 1: a large semicircle, the earth's shadow, on a base line with \emph{E} at its left end, \emph{S} at its middle and \emph{K} at its right end; a vertical line from \emph{S} upwards through \emph{A}, \emph{M} and \emph{C}; a small circle, the moon, about \emph{M}, which lies on the large semicircle, \emph{A} the lowest and \emph{C} the highest point of the small circle. Fig. 2: the same semicircle with \emph{E}, \emph{S} and \emph{K}; the vertical line from \emph{S} upwards through \emph{A}, \emph{M}, \emph{B} and \emph{C}, \emph{B} where it meets the semicircle and \emph{C} above it; the small circle about \emph{M} lies within the semicircle, \emph{A} its lowest point.}
\par Fig. 2 \emph{AB}. It is computed in the following way (by Batt.'s simple precept): one adds the two radii, of the shadow and of the moon (in both figures let this be \emph{SC}), subtracts from this sum the latitude of the moon (\emph{SM}), and so has the magnitude of the eclipse in minutes (in Fig. 1 \emph{AM}, in Fig. 2 \emph{CM} = \emph{AB}). One multiplies this by 12 and divides the product by the diameter of the moon in minutes, and so has the eclipse in \emph{digits}.\par
\par \Pg{88}As for the terms \emph{casus} and \emph{mora}, the former is the translation of the Arab. \emph{suqūṭ}, which we must translate here by “falling in”; by this word the Arab astronomers render the Ptolemaic \emph{{\greekfont ἔμπτωσις}} = the “falling into”; \emph{casus} therefore means the phase from the beginning of the eclipse to the beginning of totality. The second (\emph{mora}) is the translation of the Arab. \emph{mukth} (or \emph{makth}) = “staying”, which in turn translates the Ptolemaic \emph{{\greekfont μονή}}; \emph{mora} therefore means the phase of the total eclipse. — In the text of this chapter there also occurs \emph{respiratio}, as the translation of the Arab. \emph{injilāʾ} = “appearing”, which in turn corresponds to the Ptolemaic \emph{{\greekfont ἀναπλήρωσις}} = the “making full again” (of the moon); \emph{emersio} would be better than \emph{respiratio}.\par
\par For computing the minutes of the \emph{casus} and of the \emph{mora} we likewise give Batt.'s rules (I, 98, ll. 5–14), and show their correctness with the help of the following two figures. Batt. says first: from the square of the sum of the two radii (\emph{MS}$^{2}$ in Fig. 3) subtract the square of the latitude of the moon (\emph{SN}$^{2}$); then the root of the difference (thus \emph{MN}), divided by the hourly motion of the moon, gives the time from the beginning of the eclipse to the middle of the total eclipse (thus the \emph{casus} and the \emph{half mora} together). — Secondly he says: if a total eclipse takes place (and one wants to know its time alone), subtract the diameter of the moon (\emph{MT} in Fig. 4) from the sum of the two radii (\emph{ST})\footnote{He could also have said more simply: the radius of the moon (\emph{MP}) from the radius of the shadow (\emph{SP}).}, and from the square of this difference (\emph{SM}$^{2}$) subtract the square of the latitude of the moon (\emph{SN}$^{2}$); then the root of the difference (thus \emph{MN}), divided by the hourly motion of the moon, gives half the time of the total eclipse (thus the \emph{half mora}). If one subtracts this from the time obtained by the first rule, one obtains the time of the \emph{casus}.\par
\begin{center}\small\begin{tabular}{cc}
Fig. 3.\footnotemarkT{} & Fig. 4. \\
\end{tabular}\end{center}
\addtocounter{footnoteT}{-1}\stepcounter{footnoteT}\footnotetextT{Here the print has two figures side by side, which are not reproduced. Fig. 3: a large semicircle, the earth's shadow, on a base line \emph{ESK}, \emph{S} its centre; a vertical line from \emph{S} to the top of the semicircle, with \emph{N} on it; the small circle of the moon about \emph{M} at the upper right, outside the semicircle, which crosses the lower left part of the small circle; the horizontal line \emph{NM} and the line \emph{SM}. Fig. 4: the same semicircle with \emph{E}, \emph{S} and \emph{K} and the vertical line from \emph{S} with \emph{N} on it; the small circle about \emph{M} inside the semicircle, touching it at \emph{P} at the upper right; the horizontal line \emph{NM} and the line from \emph{S} through \emph{M} and \emph{P} to \emph{T}, outside the semicircle.}
\par Let us now return to the differences obtained above between the data of the near and the far distance. They now undergo a correction which comes from the change of the diameter of the moon in its motion from the apogee to the perigee, and which is determined with the help of the second-to-last column of Table 74 (or 73), headed \emph{Tabula proportionis}. One enters this column with the \emph{anomalia aequata} (the corrected anomaly), which in our case is, in round numbers, 285° 45′ (see above, p. 85); the corresponding value, called \emph{dakaicae superationum} (minutes of the excesses), is 20′ 43″; by this number, which is thus = 20$\tfrac{43}{60}$ {\SignFont\upshape ∶} 60, that is, as a decimal fraction, = 0.345, one multiplies the above three differences, and obtains the values:\par
\begin{center}\small\begin{tabular}{ccc}
0$^{\text{dig.}}$ 20′ & 1′ 46″ & 2′ 6″ \\
\end{tabular}\end{center}
\par these one adds to the corresponding values of the second row (for the far distance), and so obtains the corrected values:\par
\begin{center}\small\begin{tabular}{ccc}
\emph{Digits of the eclipse:} & \emph{Min. and sec. of the casus:} & \emph{Min. and sec. of the half mora:} \\
19$^{\text{dig.}}$ 36′ & 31′ 25″ & 23′ 8″ \\
\end{tabular}\end{center}
\par \Pg{89}Now two passages of our text contradict each other. On p. 26 it says: \emph{In supposita autem tabula tam casus quam morae tempora per dakaicas et elthenias horarum determinata sunt} (in the table below, the times both of the casus and of the mora are given in minutes and seconds of hours); but only O has the word \emph{horarum}, C and N do not. C and N now add to this chapter an addition headed \emph{Scientia examinationis lunarium defectuum} (the knowledge of the examination of the eclipses of the moon) (see p. 27). In this addition is found the precept to increase these minutes and seconds found for the casus and the mora by a twelfth, and to divide the sum by the hourly motion of the moon; one then has the time of the partial and of the total eclipse. C, then, takes the minutes and seconds found as minutes and seconds of arc, O, on the other hand, as minutes and seconds of time. After comparison with Batt.'s precepts, which are fairly identical with those of C, we must accept the account of the latter Ms. as the right one. But there is still a difference between the precepts of Ms. C and those of Batt.: the latter does not require the addition of a twelfth to the values of the casus and the mora mentioned above. We now proceed by the precepts of Ms. C and find:\par
\begin{center}\small\begin{tabular}{lrrrr}
\emph{Digits of the eclipse:} & \multicolumn{2}{c}{\emph{Min. and sec. of the casus:}} & \multicolumn{2}{c}{\emph{Min. and sec. of the half mora:}} \\
19$^{\text{dig.}}$ 36′ &  & 31′ 25″ &  & 23′ 8″ \\
 & $\tfrac{1}{12}$ of it: & 2′ 37″ & $\tfrac{1}{12}$ of it: & 1′ 55″ \\
 &  & 34′ 2″\footnotemark{} &  & 25′ 3″ \\
\end{tabular}\end{center}
\addtocounter{footnote}{-1}\stepcounter{footnote}\footnotetext{The corresponding numbers by Batt.'s precepts and tables are 31′ 41″ and 25′ 0″. [\emph{Translator:} The print calls this note after 34′ 2″ and again after 25′ 3″.]}
\par If we now divide the minutes and seconds of the casus and of the mora by 32′ 9″, that is, by the hourly motion of the moon at an anomaly of 285° 45′, we obtain 1$^{h}$ 3$^{m}$ 30$^{s}$ for the time of the \emph{casus} and 46$^{m}$ 45$^{s}$ for the time of the \emph{half mora}. We now have the following data for this eclipse:\par
\par 1) A \emph{total}, but not central, eclipse, since the number of digits is greater than 12 — 2) magnitude of the eclipse 19$^{\text{dig.}}$ 36′ — 3) duration of the partial eclipse (\emph{casus} + \emph{emersio}) = 2 (1$^{h}$ 3$^{m}$ 30$^{s}$) = 2$^{h}$ 7$^{m}$ — 4) duration of the total eclipse (\emph{mora}) = 2 (46$^{m}$ 45$^{s}$) = 1$^{h}$ 33$^{m}$ 30$^{s}$ — 5) duration of the \emph{whole} eclipse 3$^{h}$ 40$^{m}$ 30$^{s}$.\par
\par The opposition took place, as we found above (p. 85), on 13 Rabīʿ II, 146, 14$^{h}$ after 12 o'clock noon = 30 June 763, 2 o'clock in the morning; this is the time of the middle of the eclipse, and so its beginning, if we neglect the seconds, was on 30 June 763 at 0$^{h}$ 10$^{m}$ in the morning, and its end at 3$^{h}$ 50$^{m}$ in the morning. — According to Oppolzer\footnote{\emph{Kanon der Finsternisse} (Denkschriften der Wiener Akademie, mathem. Klasse, vol. LII), p. 355, no. 3050 of the eclipses of the moon.} the following data result (for Greenwich): middle of the eclipse 30 June 763, 1$^{h}$ 42$^{m}$ in the morning; magnitude = 22.3 inches; duration of the partial eclipse 2$^{h}$; duration of totality 1$^{h}$ 44$^{m}$; duration of the whole eclipse 3$^{h}$ 44$^{m}$. This is surely a fine result to obtain with the help of these tables of the Arabs, more than 1000 years old!\par
\par \Pg{90}We have still to add here something about the determination of the size of the \emph{earth's shadow} in eclipses of the moon. In the passage cited above (p. 79) al-Bīrūnī makes the further remark that the method of computing the size of the earth's shadow in eclipses of the moon is likewise the same in the \emph{Khaṇḍakhādyaka} as in Kh. Of this method, however, we find nothing at all in our work; we have only, in Tables 61–66, two columns each (the 6th and the 7th), headed \emph{Primum dimidium rotae draconis} and \emph{Secundum dimidium rotae draconis} (first and second half of the wheel of the dragon), with which we shall deal directly. That those rules for determining the size of the earth's shadow are missing in this work may have its reason either in al-Bīrūnī's having had before him another edition of Kh.'s tables than Maslama al-Majrīṭī, or, what is more probable, in the present work's having been much abridged; whether by Maslama, or by the translator, or only by later copyists, remains undecided.\par
\par As regards the 6th and 7th columns of Tables 61–66, I have found the following. The numbers of the 6th column are 5 times the corresponding values of the hourly unequal motion (\emph{buht}) of the sun, which are entered in column 2; the numbers of the 7th column are 1$\tfrac{3}{5}$ times the corresponding values of the hourly unequal motion of the moon, which are given in column 3; if one subtracts the values of column 6 from the corresponding ones of column 7, one obtains in each case the values of the radius of the earth's shadow that correspond to the anomalies of column 1. If, then, Kh. carried out these computations according to the \emph{Khaṇḍakhādyaka}, as al-Bīrūnī reports, the rule of the Indian work must have run roughly as follows: subtract 5 times the hourly \emph{buht} of the sun from 1$\tfrac{3}{5}$ times the hourly \emph{buht} of the moon, and one has the apparent radius of the earth's shadow. \emph{Example:} at an anomaly of 3$^{\text{sig.}}$ = 90° (Table 63, bottom line) the hourly \emph{buht} of the sun (\emph{motus solis in hora}, the motion of the sun in an hour) is = 2′ 28″, that of the moon = 32′ 56″; the former number multiplied by 5 gives 12′ 20″, and this is the corresponding number of the 6th column, that is, the \emph{primum dimidium rotae draconis}; the second number multiplied by 1$\tfrac{3}{5}$ gives 52′ 42″, while the corresponding number of the 7th column, that is, the \emph{secundum dimidium rotae draconis}, is = 52′ 47″\footnote{The discrepancy amounts to 5″; it is the greatest that we have found in our trials.}; the apparent semidiameter of the earth's shadow at this anomaly is therefore = 52′ 47″ − 12′ 20″ = 40′ 27″.\par
\par \Pg{91}Ptol.\footnote{Ed. Heiberg, I, 421–423; translation by Manitius, I, pp. 309–310.} and after him Batt.\footnote{Ed. Nallino, I, 59–60.} state that the radius of the earth's shadow is c. 2$\tfrac{3}{5}$ times the apparent radius of the moon; in Kh. this factor is not quite constant: according to various trials that we have made, its values vary between 2.5 and 2.6, and so differ little from Ptol.'s value. This leads us to the conjecture that Brahmagupta, in his \emph{Khaṇḍakhādyaka}, may perhaps have followed Ptol. in this question more than his Indian predecessors; at least the rules of the \emph{Sūrya-Siddhānta}\footnote{Transl. Burgess, p. 273.} and of the \emph{Pañcasiddhāntikā}\footnote{Ed. Thibaut, pp. 63–65 of the translation.} for determining the size of the earth's shadow show no great kinship with that of Kh. In the first work the true diameters of the earth, the sun and the moon and the mean and true motions of the latter two heavenly bodies are drawn into the computation, and the result gives the true diameter of the earth's shadow in \emph{yojanas}\footnote{One \emph{yojana} = 8 Arab miles = c. 16000 metres.}, which can then, to be sure, easily be converted into minutes. In the second work the rule runs: “Multiply the true distance of the moon by 36, and divide this product by the true distance of the sun multiplied by 90 and divided by 276; subtract this quotient from 36, and one has the diameter of the earth's shadow (in \emph{yojanas}).”\par
\UnitHead{Chapter 34. On the parallaxes of the moon in longitude and latitude}{Chapter 34. On the parallaxes of the moon in longitude and latitude}{K20-8}
\par\begin{center}(Text pp. 28–29, Tables 77 and 77$^{a}$.)\end{center}
\par Kh.'s rules for determining the \emph{parallaxes} of the moon in longitude and latitude are, in my view, somewhat corrupt, especially the former, or at least quite unclear. For orientation we state first that, since the determination of the parallax is used in this work only for computing the eclipses of the sun, the moon is assumed to be in the ecliptic when the parallax in longitude is determined, that is, its latitude is set equal to zero. It is further to be noted that, when the moon is in the meridian of the place of observation at the time of the conjunction, its parallax in longitude is equal to zero.\par
\par As for Batt., his treatment of this subject is a very thorough and complicated one\footnote{Ed. Nallino, I, 76–84, and 251–265.}; I cannot go into the details here, and remark only, in summary, that he needs in the first place the \emph{zenith distance z} of the moon, which he finds by the formula\par
\par\smallskip\[\cos z = \sin h = \frac{\sin EH\cdot\sin MB}{\sin EB}\,,\]\par\smallskip
\par which follows at once from the triangles \emph{EHB} and \emph{MAB}\footnoteT{Here the print has a figure at the right of the text, which is not reproduced: a circle, the meridian, with \emph{Z} at its top, \emph{P} at the upper right, \emph{E} at the upper left, \emph{H} at the left end and \emph{T} at the right end of a horizontal diameter, and \emph{K} at the lower right; dashed lines: the vertical diameter from \emph{Z}, the diameter \emph{HT}, the diameter from \emph{E} through the centre to \emph{K}, and a diameter from \emph{P} through the centre to the lower left. Solid arcs: the horizon from \emph{H} to \emph{T}, curving down through \emph{A} and \emph{B}; the ecliptic from \emph{E} through \emph{M} and \emph{B} to \emph{K}; the circle of altitude from \emph{Z} through \emph{M} to \emph{A}, marked \emph{z} between \emph{Z} and \emph{M} and \emph{h} between \emph{M} and \emph{A}; and the arc from \emph{M} to \emph{P}; the angle at \emph{M} between \emph{ME} and \emph{MZ} is marked with a small arc.} by the rule of four quantities. \emph{EH} is the altitude of the ecliptic, easily found from the declination of the sun and the geographical latitude of the place of observation; \emph{E} is the so-called \emph{medium coelum} (midheaven), \emph{EB} the distance of the \emph{medium coelum} from the \emph{ascendant B}, \emph{MB} the distance of the moon from the ascendant. Batt. then computes the so-called \emph{angles} of \emph{latitude} and \emph{longitude}, that is, the angle that the ecliptic makes with the circle of altitude passing through the moon, and its complement (thus the angles \emph{EMZ} and \emph{ZMP}, if \emph{Z} is the zenith and \emph{P} the pole of the \Pg{92}ecliptic). After he has then determined from the tables (II, 93–94), with the zenith distance as argument and with various corrections, the parallax of the moon in the circle of altitude \emph{ZA}, he multiplies this parallax both by the sine of the angle of longitude and by that of the angle of latitude; the first product gives the \emph{parallax} of the moon in \emph{longitude}, the second that in \emph{latitude}. According to Batt., then, the parallaxes of the moon in longitude and in latitude are to each other as the sines of the angles of longitude and of latitude; or, expressed otherwise, he resolves the parallax of altitude into its components in longitude and in latitude.\par
\par The rules of the \emph{Sūrya-Siddhānta} on the determination of parallax are in its 5th chapter\footnote{\emph{Transl.} Burgess, pp. 285–295.}. In principle they do not differ essentially from Batt.'s: trigonometric relations are used there too, the so-called \emph{meridian distance} and \emph{ecliptic zenith distance} are computed, and from these the parallaxes in longitude and latitude are then determined.\par
\par Kh.'s rules are much simpler; but since he only explains the use of the tables and gives no trigonometric derivations at all, it is difficult to decide by what route he found the numbers of his tables. His principle, briefly put, is the following: the \emph{parallax} in \emph{longitude} is a function of the deviation of the moon from the meridian to the east or to the west, that is, in our present-day way of speaking, a \emph{function} of the \emph{hour angle}; the \emph{parallax} in \emph{latitude} is a function of its \emph{zenith distance}. The numbers in the first column (Table 77) therefore mean two things: if it is a matter of determining the \emph{parallax} in \emph{longitude}, they denote the distance of the moon from the meridian of the place of observation, measured on the equator (or on a parallel circle), and thus the \emph{hour angle} of the moon; if it is a matter of determining the \emph{parallax} in \emph{latitude}, they mean the zenith distance of the moon, namely its so-called \emph{ecliptic zenith distance}, that is, its zenith distance when it is in the \emph{medium coelum}. — In this respect, then, Kh. seems to come nearer to the Indians than Batt. does.\par
\par \Pg{93}The text of chapter 34 requires first the determination of the \emph{ascendant} (\emph{horoscopus}) for the given time of the conjunction, that is, of the longitude of the point of the ecliptic which — to take up again our \emph{example} of a \emph{conjunction} of the sun and the moon given above (pp. 82–84) — rises on 29 Muḥarram of the year 288 AH at 18$^{h}$ 44$^{m}$ 29$^{s}$. The text says nothing at all about how this ascendant is to be determined; nor, in our view, could it be found with the tables that this work contains, for the tables of the oblique ascensions are missing, of which we ought to have at least those for the latitude of Córdoba; it could, however, be determined by computation with the help of the ascensional difference. We have therefore determined the \emph{ascendant} by Batt.'s precepts\footnote{Ed. Nallino, I, 73.} and found it = 297° 3′. Kh.'s rule for finding the \emph{medium coelum} also seems to be faulty; at the least it is unclear from what the 90° are to be subtracted, whether from the \emph{ascendant} itself, or from its \emph{right ascension}, or from its \emph{ascensio obliqua} (oblique ascension); we have therefore determined it likewise by Batt., and found it = 228° 40′. Now Kh.'s rules say that one must enter Tables 59–59$^{b}$, by which longitudes are converted into right ascensions and conversely, both with the \emph{true longitude} of the moon (= 303° 20′) and with that of the \emph{medium coelum} (= 228° 40′); doing this, one finds 305° 43′ for the right ascension of the moon and 226° 10′ for that of the medium coelum; the difference of the two, = 79° 33′, is the \emph{hour angle} of the moon at the given time.\footnote{By trigonometry we find this hour angle = 79° 36′.} For better understanding I give here the adjoining figure\footnoteT{Here the print has a figure at the right of the text, which is not reproduced: a circle, the meridian, with \emph{Z} at its top, \emph{A} and \emph{E} at its upper left, \emph{P} at its upper right, a small \emph{s} beside \emph{P} and a small arc marking the angle at \emph{P}; \emph{H} and \emph{T} at the ends of a horizontal dashed diameter, \emph{K} at the lower right and \emph{Q} at the bottom right; dashed lines from the centre to \emph{Z}, \emph{A}, \emph{E}, \emph{P}, \emph{K} and \emph{Q}. Solid arcs: the horizon from \emph{H} through \emph{L} to \emph{T}, below the diameter; the equator from \emph{A} through \emph{W}, \emph{D}, \emph{N} and \emph{R} to \emph{Q}; the ecliptic from \emph{E} through \emph{W}, \emph{L} and \emph{M} to \emph{K}; and the arcs from \emph{P} through \emph{L} to \emph{N} and from \emph{P} through \emph{M} to \emph{R}. \emph{W} stands on the diameter \emph{HT}, \emph{D} and \emph{N} below \emph{L}.}, in which, so that it might agree better with our example, the intersection of the equator and the ecliptic that lies on the rear hemisphere had to be taken as the vernal point, and so the one that lies on the front hemisphere as the autumnal point \emph{W}. Then (if \emph{HT} = the horizon, \emph{AQ} = the equator, \emph{EK} = the ecliptic, \emph{L} = the ascendant, \emph{M} = the moon) 180° + \emph{WM} is the longitude of the moon, 180° + \emph{WR} its right ascension, 180° + \emph{WL} the longitude of the ascendant, 180° + \emph{WN} its right ascension, 180° − \emph{EW} the longitude of the medium coelum, 180° − \emph{AW} its right ascension; consequently 180° + \emph{WR} − (180° − \emph{AW}) = \emph{WR} + \emph{AW} = \emph{AR} = {\SignFont\upshape ∠} \emph{APR} = \emph{s} = the hour angle of the moon\footnoteT{In the figure \emph{E} and \emph{A} stand before \emph{W}, and these relations hold for it. In the example the medium coelum, 228° 40′, and its right ascension, 226° 10′, lie 48° 40′ and 46° 10′ beyond \emph{W} in the order of the signs, so that for these numbers they are 180° + \emph{WE} and 180° + \emph{WA}. The hour angle is in either case the arc \emph{AR} of the equator between the meridian and the circle through the pole and the moon, here \emph{WR} − \emph{WA} = 125° 43′ − 46° 10′ = 79° 33′.}. — With the hour angle of the moon (79° 33′) one now enters the table of parallaxes (Table 77) and finds there c. 1$^{h}$ 33$\tfrac{1}{2}$$^{m}$ for the \emph{parallax in longitude}; this quantity is multiplied by the unequal hourly motion of the moon, which for the anomaly 125° 7′ is = 34′ 34″, and one so obtains c. 54′. — Kh., then, gives the parallax of longitude in his tables in time, Batt., on the other hand, in degrees of arc; in our case, then, the apparent place of the moon, as it shows itself to our eyes, is 54′ distant from the true place, as it has resulted from the computation and the tables; and by the precepts of the text this parallax must be added to the true place of the moon (303° 19′ 48″)\footnote{For the moon is between the ascendant and the medium coelum, or, as Batt. (I, 81 and 83) puts it more clearly, it is less than 90° distant from the ascendant. [\emph{Translator:} The print has no call for this note in the text; it is placed here, where the note gives the reason for the addition.]} to obtain the apparent place of the moon, which is therefore 304° 13′ 48″.\par
\par \Pg{94}The rule for finding the \emph{parallax} in \emph{latitude} is then a little simpler, since the questions of the \emph{ascendant} and of the \emph{medium coelum} are settled for us. It is now required to determine the \emph{obliquatio} (declination) of the \emph{medium coelum}, that is, the declination of the moon (= that of the sun) at that time; by Table 22 this is = − 19° 44′ 38″; we take the latitude of Córdoba as 37° 50′ in round numbers, and this gives for the ecliptic zenith distance \emph{EZ} = 57° 34′ 38″; our text calls this the \emph{interstitium medii coeli huius horoscopi et imminentis capitibus nostris} (the interval between the midheaven of this ascendant and the point that hangs over our heads). Now the \emph{latitude} of the moon is to be determined: its corrected argument is, by p. 84, = 169° 7$\tfrac{1}{2}$′; to it belongs, by Table 26, the latitude 50′ 53″; it is northern, and is therefore to be subtracted from the ecliptic zenith distance found just now, which gives 56° 43′ 45″; with this argument one enters Table 77 and finds there the \emph{parallax} of \emph{latitude} = 40′ 43″, and it is to be taken as \emph{southern}, because the \emph{medium coelum} too lies south of the zenith of Córdoba.\par
\par By Theon's tables (in Batt. II, 98) one finds, with a correction (see Batt. I, 83), the parallax of longitude c. 35′ and that of latitude c. 47′; roughly the same results follow from the rules that Batt. sets out in I, 79–81. We see, then, that the parallax of longitude in Kh. is considerably greater than in Batt., and that the parallax in latitude also differs, by 6′; but Batt. himself remarks that these tables of Theon cannot give quite satisfactory results.\par
\UnitHead{Chapter 35. On the eclipse of the sun}{Chapter 35. On the eclipse of the sun}{K21-6}
\par\begin{center}(Text p. 29, Table 78.)\end{center}
\par This chapter is very short and in any case incomplete. From Table 78 we see that Kh. takes the limits of the argument of latitude within which an eclipse can take place at all as follows: from 172° 49′ to 187° 30′, from 0° to 7° 11′, and from 352° 30′ to 360°; Batt. (II, 88) has the following numbers: 159° 44′ to 191° 16′, 0° to 20° 16′, and 348° 44′ to 360°.\par
\par Nothing is said in this chapter about how the tables \emph{pro spatio longinquo} and \emph{pro spatio propinquo} (for the far and for the near distance) are to be used; we assume that this must be done in a similar way as with the tables of the eclipses of the moon, and as in Batt.\footnote{Ed. Nallino, I, 111, and II, 91.}; but Table 78 lacks the \emph{tabulae proportionis} (tables of proportion) that are found in the tables of the moon; probably, however, the same tables of proportion are to be used as there (Tables 73–76).\par
\par Since the argument of latitude at our conjunction (see pp. 82–84), namely 169° 7$\tfrac{1}{2}$′, is not found in Table 78, we must conclude that no eclipse of the sun takes place for Córdoba at this conjunction. By Batt.'s tables, however, one can take place for other places, and we shall now carry out the further computation by these tables; but we take the values found so far by Kh.'s tables:\par
\begin{center}\small\begin{tabular}{ll}
latitude of the moon & 50′ 53″, \\
parallax of latitude & 40′ 43″. \\
\end{tabular}\end{center}
\par Since the \emph{latitude} is \emph{northern} but the \emph{parallax} is to be taken as \emph{southern}, we must subtract the parallax from the latitude; this gives 10′ 10″, and this is the apparent latitude of the moon at the conjunction visible to us\footnote{By Batt. one finds: apparent latitude = 58′ 25″ − 47′ = 11′ 25″; but the parallax cannot be determined exactly by Theon's tables, so that it would be possible that the two values for the apparent latitude might after all nearly agree.}. For this \emph{apparent} \Pg{95}\emph{latitude} we now find, by Batt.'s tables (II, 91), the following data:\par
\begin{center}\footnotesize\begin{tabular}{llll}
\multicolumn{2}{c}{\begin{tabular}[b]{@{}c@{}}\emph{in maxima longinquitate lunae}\\\emph{(at the greatest distance of the moon):}\end{tabular}} & \multicolumn{2}{c}{\begin{tabular}[b]{@{}c@{}}\emph{in minima longinquitate lunae}\\\emph{(at the least distance of the moon):}\end{tabular}} \\
\begin{tabular}[b]{@{}c@{}}\emph{Digiti diametri solaris}\\\emph{(digits of the diameter of the sun):}\end{tabular} & \begin{tabular}[b]{@{}c@{}}\emph{Quantitas incidentiae}\\\emph{(amount of the falling in):}\end{tabular} & \emph{Digiti diametri solaris:} & \emph{Quantitas incidentiae:} \\
7$^{\text{dig.}}$ 42′ & 29′ 15″ & 8$^{\text{dig.}}$ 48′ & 32′ 26″ \\
\end{tabular}\end{center}
\par The differences of the corresponding values are:\par
\par\begin{center}1$^{\text{dig.}}$ 6′ and 3′ 11″.\end{center}
\par Of these, by the precepts (I, 111), one must take proportional parts, namely they must be reduced in the ratio 46$\tfrac{1}{3}$ {\SignFont\upshape ∶} 60 (the number 46$\tfrac{1}{3}$, more exactly 46′ 17″, is found in the \emph{tabula correctionis} (table of correction) (II, 89) as the value corresponding to the anomaly 125° 7′ of the moon). Carrying out these corrections, one changes those differences into 51′ and 2′ 27″ respectively; these values, added to the data for the \emph{maxima longinquitas lunae} (greatest distance of the moon) (= \emph{pro spatio longinquo}), give\par
\par\begin{center}8$^{\text{dig.}}$ 33′ and 31′ 42″.\end{center}
\par The minutes of the \emph{quantitas incidentiae} (in our tables the \emph{casus}) are divided by the unequal hourly motion of the moon at the anomaly of 125° 7′, that is, by 34′ 34″; this gives c. 55$^{m}$. We now have the following results:\par
\par 1) For the time of the true conjunction we found above (p. 83): 18$^{h}$ 44$^{m}$ 29$^{s}$; by the precepts of our text, and of Batt. too (I, 105 and 110), we must subtract from this the parallax of longitude in time, that is, 1$^{h}$ 33$\tfrac{1}{2}$$^{m}$; this gives, in round numbers, 17$^{h}$ 11$^{m}$ for the time of the \emph{apparent conjunction}, that is, of the \emph{middle} of the eclipse; that is, the middle took place on 23 Jan. 901 at 5$^{h}$ 11$^{m}$ in the morning (for Córdoba).\par
\par 2) The \emph{beginning} of the eclipse took place at 5$^{h}$ 11$^{m}$ − 55$^{m}$ = 4$^{h}$ 16$^{m}$, the \emph{end} at 5$^{h}$ 11$^{m}$ + 55$^{m}$ = 6$^{h}$ 6$^{m}$.\par
\par 3) The \emph{maximum} of the eclipse amounted to 8$^{\text{dig.}}$ 33′, that is, somewhat more than $\tfrac{2}{3}$ of the diameter of the sun.\par
\par I remark expressly that these data can make no claim to accuracy, since they are computed partly by Kh.'s tables and partly by Batt.'s. It will nevertheless be of interest to compare the results with those that one obtains for this eclipse from Oppolzer's \emph{Kanon}: for Córdoba the greatest phase (the middle) of the eclipse was at 5$^{h}$ 14$^{m}$ in the morning, and the end at 6$^{h}$ 0$^{m}$. But since at Córdoba on 23 Jan. 901 the sun rose only at 7$^{h}$, the eclipse was not visible there\footnote{I owe the computation for \emph{Córdoba} by Oppolzer's \emph{Kanon} to the kindness of Mr Broger, assistant at the Zurich observatory, to whom I here express my best thanks.}. So once again a quite fine result!\par
\par \Pg{96}Let us further compare these data with Batt.'s statements about this eclipse. He says\footnote{Ed. Nallino, I, 56–57.} that the middle of the eclipse took place at \emph{Antioch} 3$^{h}$ 40$^{m}$ before the beginning of 23 Jan. 901 (the beginning being, according to the Arab astronomers, at 12 o'clock noon), and so, in our civil reckoning, on 23 Jan. at 8$^{h}$ 20$^{m}$ in the morning. Now \emph{Antioch} lies c. 41° east of \emph{Córdoba}; this makes 2$^{h}$ 44$^{m}$ in time, and this added to 5$^{h}$ 11$^{m}$ gives 7$^{h}$ 55$^{m}$ as the time of the middle of the eclipse for \emph{Antioch}; the difference from Batt.'s statement is therefore c. 25$^{m}$. The magnitude of the eclipse, says Batt., amounted to little more than half, or at \emph{al-Raqqa} to somewhat less than two thirds. These statements too differ little from those of our computation.\par
\par As we have remarked above, the two chapters 34 and 35 are very unclear, and most probably incomplete and in part wrongly rendered in the translation. We shall hardly go wrong if we assume that Adelard did not properly understand these somewhat difficult astronomical sections, and later copyists, naturally, still much less; from this resulted partly false interpretations, partly omissions, two of which I wish to mention here and to try to correct. In chapter 34 (p. 28) it says: \emph{erit hoc productum horis coniunctionis addendum} (this product is to be added to the hours of the conjunction); instead of \emph{horis coniunctionis} it must probably read \emph{loco coniunctionis} (to the place of the conjunction), or better \emph{loco vero lunae} (to the true place of the moon); for here it is not specifically a matter of a conjunction, but generally of the place of the moon as changed by the \emph{parallax}. And further it says in the same place: \emph{Quodque inde surget, horam eclipsis supra locum pro quo inquiris ex diversitate aspectus mutatam ponit} (and what results from this gives the hour of the eclipse, as changed by the parallax, for the place for which you are inquiring); instead of \emph{horam eclipsis} it must read: \emph{locum lunae apparentem in ecliptica} (the apparent place of the moon in the ecliptic)\footnote{Compare Batt., ed. Nallino, I, 83–84.}. In chapter 35 the passage from the beginning to \emph{Si enim fuerit locus coniunctionis} (for if the place of the conjunction is) is quite insufficient for determining an eclipse of the sun; but I shall not be blamed if I refrain from an attempt to restore the right text.\par
\UnitHead{Chapters 32 and 36. On the division of the twelve astrological houses}{Chapters 32 and 36. On the division of the twelve astrological houses}{K22-13}
\par\begin{center}(Text pp. 25 and 29–30, Tables 79–90.)\end{center}
\par These two chapters and the last, which follows next, deal with purely \emph{astrological} questions. Chapter 32 stands only in Mss. C and N, and in the former in the wrong place\footnote{See above, p. 86, note 3.}; it should immediately precede chapter 36, for it contains the precepts for finding the beginnings of the twelve houses on the ecliptic, chapter 36, on the other hand, the instructions for using the tables that belong to them.\par
\par \Pg{97}The Arabic \emph{taswiya} = distribution, division, arrangement has been rendered by the translator partly by \emph{aequatio}, partly by the more correct \emph{dispositio}. It is a matter of the distribution of the twelve astrological houses over the celestial sphere. Various kinds of this distribution existed, for which I refer to Batt.\footnote{Ed. Nallino, I, 246–249.}; all have in common that six great circles of the celestial sphere intersect in one diameter, by which the sphere is divided into twelve spherical lunes. Batt.'s division, which is probably also the Ptolemaic one in the \emph{Quadripartitum}, and which Kh. too follows, is the following. The six great circles pass through the poles of the equator, and are thus meridians or circles of declination; the first, passing through the \emph{ascendant}, marks the beginning of house I. Now, to obtain the point on the ecliptic through which the second circle passes (counted from the ascendant towards the east, in the order of the signs of the zodiac), with which, then, house II begins, one must divide the night arc of the ascendant by 6\footnote{According to Kh. this gives the \emph{tempora duarum horarum nocturnarum} (the times of two night hours); by \emph{tempora} equatorial degrees are meant here.}, add the result to the right ascension of the ascendant, and look up the longitude for this new right ascension in Tables 59–59$^{b}$; one then has the degree of the beginning of house II, and, increased by 180°, the degree of the beginning of house VIII, which is the nadir (\emph{naẓīr}) of that of house II. If one wants the beginning of house III (or IX), one adds the 6th part of the night arc once more to the right ascension found just now, and the sum is the right ascension of the beginning of house III (or IX); the same holds also for computing the beginning of house IV (or X). To obtain the beginning of house V (or XI), one must add to the right ascension of house IV the 6th part of the day arc of the ascendant (= $\tfrac{360^{\circ} - \text{night arc}}{6}$ = 60° − $\tfrac{\text{night arc}}{6}$), and again look up the longitude for this right ascension; this is the beginning of house V (or XI). One now proceeds in the same way for finding the beginnings of houses VI and VII; if the computation has been carried out correctly, the beginning of house VII must be the nadir of the beginning of house I, that is, of the ascendant.\par
\par One sees from this, then, that these twelve houses were of unequal size; still, houses I, II, III, VII, VIII and IX are congruent with one another, and so are houses IV, V, VI, X, XI and XII. What is peculiar is that smaller and larger houses fall both on the day arc and on the night arc. Each geographical latitude has its own size of the houses, and this changes again with the season; only at the equator are all the houses of equal size at every season\footnote{At the time of the equinox this is the case at the other latitudes too.}.\par
\par In chapter 32 it is then also stated how one can determine the size (in degrees) of a \emph{temporal hour}, that is, both of a \emph{day hour} and of a \emph{night hour}; for this one needs the tables of the oblique ascensions for the latitude concerned, but these tables, as has been remarked before, are missing in this work. To make a trial of the beginnings of the twelve houses, we have determined the size of the temporal hour for Córdoba approximately by Batt.'s tables (II, 67), and found it = 12° 4′ for the morning of 23 Jan. 901, that is, for the ascendant 297° 3′ (= 27° 3′ of Capricorn)\footnote{This value rests on the geographical latitude 38° 30′ for Córdoba; for we found afterwards, in the tables of the \emph{proiectio radiorum} (projection of the rays), that al-Majrīṭī took it to be so large.}; one can, to be sure, also compute it trigonometrically, by looking for the hour angle of the sun at its rising on the day concerned at Córdoba. The length of the double day hour is therefore 24° 8′, and the length of the double night hour = 35° 52′.\par
\par \Pg{98}By the precepts of our text (chapter 32), one has to add the length of the double night hour to the right ascension of the ascendant, which by Table 59 is = 29° 10′; this gives 65° 2′; for this right ascension one finds in the same table the longitude 3° 1′ of Pisces, and this is the beginning of house II counted from the ascendant. If one adds the amount of the double night hour once more to the right ascension 65° 2′, one obtains 100° 54′; to this belongs, in the same Table 59, the longitude 11° 54′ of Aries, and this is the beginning of house III. 35° 52′ added once more to 100° 54′ gives 136° 46′; to this belongs, in Table 59$^{a}$, the longitude 19° 19′ of Taurus, and this is the beginning of house IV. — Now one adds the length of the double day hour (= 24° 8′) to 136° 46′; this gives 160° 54′; to this belongs, in the same Table 59$^{a}$, the longitude 12° 25′ of Gemini, and this is the beginning of house V; 24° 8′ added once more to 160° 54′ gives 185° 2′; to this belongs, in the same table, the longitude 4° 36′ of Cancer, and this is the beginning of house VI; 24° 8′ added once more to 185° 2′ gives 209° 10′; to this belongs, in the same table, the longitude 27° 3′ of Cancer; this is the beginning of house VII, and it is the nadir of the ascendant, = 27° 3′ of Capricorn.\par
\par We therefore have the following values for the beginnings of the twelve houses:\par
\begin{center}\small\begin{tabular}{cccccc}
\emph{I and VII} & \emph{II and VIII} & \emph{III and IX} & \emph{IV and X} & \emph{V and XI} & \emph{VI and XII} \\
27° 3′ & 3° 1′ & 11° 54′ & 19° 19′ & 12° 25′ & 4° 36′ \\
of Capricorn & of Pisces & of Aries & of Taurus & of Gemini & of Cancer \\
or & or & or & or & or & or \\
of Cancer & of Virgo & of Libra & of Scorpio & of Sagittarius & of Capricorn \\
\end{tabular}\end{center}
\par Now in our Table 88 we find, for the ascendant 27° 3′ of Capricorn, the following values for the beginnings of the twelve houses (the signs of the zodiac being left out):\par
\begin{center}\small\begin{tabular}{cccccc}
\emph{I and VII} & \emph{II and VIII} & \emph{III and IX} & \emph{IV and X} & \emph{V and XI} & \emph{VI and XII} \\
27° 3′ & 3° 8′ & 12° 2′ & 19° 31′ & 12° 33′ & 4° 36′ \\
\end{tabular}\end{center}
\par The differences range between 0′ and 12′; they may come in the first place from a different length of the temporal hour, but errors of writing in the numbers may also be the cause.\par
\par For casting a horoscope at a birth etc. these tables were of great importance, and replaced laborious computations and observations that were not always possible: if one knew the mean and true places of the sun, the moon and the five planets from Tables 4–56, one could at once determine by Tables 79–90 in which of the twelve houses each of these heavenly bodies stood at the time concerned.\par
\UnitHead{Chapter 37. On the places of the casting of the rays}{Chapter 37. On the places of the casting of the rays}{K23-7}
\par\begin{center}(Text pp. 30–31, Tables 91–114.)\end{center}
\par \Pg{99}Of as great importance in astrology as the knowledge of the distribution of the twelve houses, or of greater importance still, was the determination of the \emph{proiectiones radiorum} (projections of the rays), that is, of those points of the ecliptic that stand 60°, 90° or 120° from a given point in the heavens, for example from a planet, or, in other words, that stand with it in \emph{sextile}, \emph{square} or \emph{trine aspect}. To these three \emph{aspects} are then added the \emph{conjunction} and the \emph{opposition} (the \emph{opposite aspect}). It is characteristic that Kh. devoted the greatest space of his tables precisely to this chapter: text and tables together take up 26 pages; this illustrates better than any other considerations or investigations the fact, known already, that the whole scientific work of antiquity\footnote{At least from Ptol. on.} and of the Middle Ages in the field of astronomy was not an end in itself, but had only to supply the means for the needs of astrology. It is well known, moreover, that even \emph{Regiomontanus} wrote a work of his own on a related astrological subject\footnote{\emph{Tabulae directionum profectionumque}, Augsburg 1490, Tübingen 1550.}, in which he likewise (Problema 31) deals with the question of determining the \emph{place} of the \emph{casting of the rays} (\emph{radiatio}).\par
\par For determining these places too there were various routes, for whose details I must again refer to Batt.\footnote{Ed. Nallino, I, 307–313.}. The simplest were those that neglected the latitude of the planets, that is, assumed them to lie in the ecliptic; the most obvious thing then, naturally, was to go 60° (or 90° or 120°) forwards or backwards in the ecliptic from the planet, and one had the two points (the left and the right) of the sextile aspect (or of the square or the trine aspect): this primitive route was taken, for example, by al-Qabīṣī (Alcabitius) and Abraham ibn ʿEzra\footnote{Ibid., I, 309.}. Others, for example the Greek astrologer Dorotheus, are said, according to Ibn Hibintā, to have measured the place of the casting of the rays not on the ecliptic but on the equator; that is, they regarded as the radius of the sextile aspect the arc of the equator that was equal to the difference: right ascension of the planet minus right ascension of [longitude of the planet ± 60°], or conversely; in this way they came to the same point of the ecliptic as those named before. — Ptol. took a somewhat more complicated route in the \emph{Quadripartitum}\footnote{Ibid., I, 311.}, but likewise neglected the latitude of the planets. — The first to take it into account seems to have been Batt., who set up the formula\par
\par\smallskip\[\mathrm{crd.}\,x = \sqrt{R\cdot\frac{\mathrm{crd.}^{2}\,60^{\circ}-\mathrm{crd.}^{2}\,b}{\cos b}}\,,\]\par\smallskip
\par where \emph{b} is the latitude of the planet and \emph{x} the arc of the ecliptic that corresponds to the hypotenuse 60° (the radius of the sextile aspect); this formula already occurs implicitly in the problem of determining the distance of two stars from their latitudes and longitudes, and can, as Nallino has shown\footnote{Ibid., I, 307–308.}, easily be transformed into the formula\par
\par\smallskip\[\cos x = \frac{\cos 60^{\circ}}{\cos b}\,,\]\par\smallskip
\par which follows at once from the right spherical triangle \emph{ABC}\footnoteT{Here the print has a figure at the right of the text, which is not reproduced: a right spherical triangle, with \emph{B} at the left, \emph{A} at the upper right and \emph{C} below \emph{A}; the arc from \emph{B} to \emph{A} is marked 60°, the short side from \emph{A} down to \emph{C} is marked \emph{b}, and the arc from \emph{B} to \emph{C} below it is marked \emph{x}; the right angle is at \emph{C}.} by the formula $\cos c = \cos a\cdot\cos b$, which Batt. did not yet know.\par
\par \Pg{100}Which route, then, did Kh. take? From what Nallino\footnote{Ibid., I, 312.} quotes from Ibn Hibintā one would have to conclude that Kh. proceeded roughly as was reported above of Dorotheus, except that he converted the right ascension of the arc of the equator back into degrees of the ecliptic, and then took this arc of the ecliptic as the radius of the sextile aspect. But the tables before us here differ essentially, both in arrangement and in content, from Kh.'s tables as Ibn Hibintā has handed them down to us, a sample of which Nallino very kindly sent me: first, those go only from 10° to 10° of the signs, ours from 5° to 5°; secondly, they do not give in each case the whole oblique ascension\footnoteT{The print has «Rectascension» (right ascension) here and in the next clause; Suter's list of misprints corrects both to «Ascensio obliqua» (oblique ascension). It leaves the footnote to this passage, whose numbers are right ascensions, as it stands.} counted from the vernal point, but only the oblique ascension of each sign by itself\footnote{That is, for example, the right ascension of Aries, that is, of 30°, is 27° 54′, that of Aries and Taurus together, that is, of 60°, is 57° 48′; hence that of Taurus alone is equal to the difference of these numbers, that is, = 29° 54′.}; the numbers for the sextile, square and trine aspects are also different. Our tables take as their basis the \emph{ascensiones obliquae} (oblique ascensions) for the latitude of \emph{Córdoba}, which Maslama gives as 38° 30′; Kh.'s tables for the \emph{proiectiones radiorum} have therefore likewise been reworked for Córdoba. But now a difficult problem arises for us: the oblique ascensions agree in each case only with the numbers of the first three columns; after that they differ rather strongly. Let us give an example of this:\par
\par Table 91, first line, has the following numbers:\par
\begin{center}\small\begin{tabular}{lcccccc}
\emph{Gradus} & \emph{Aries} & \emph{Taurus} & \emph{Gemini} & \emph{Cancer} & \emph{Leo} & \emph{Virgo} \\
\begin{tabular}[b]{@{}c@{}}\emph{aequales}\\\emph{(equal degrees)}\end{tabular} & \emph{(Libra)} & \emph{(Scorpius)} & \emph{(Sagittar.)} & \emph{(Capricorn.)} & \emph{(Aquar.)} & \emph{(Pisces)} \\
5° & 2° 58′ & 21° 50′ & 45° 20′ & 76° 0′ & 106° 32′ & 131° 9′ \\
\multicolumn{7}{c}{\begin{tabular}[b]{@{}c@{}}But the oblique ascensions of the first 5 degrees of each sign\\are, for the latitude 38° 30′:\end{tabular}} \\
 & 3° 0′ & 21° 55′ & 45° 5′ & 75° 18′ & 111° 24′ & 149° 2′ \\
\end{tabular}\end{center}
\par These deviations in the last three to four columns show themselves everywhere; where do they come from? I have not so far succeeded in explaining this, and the way in which the tables were computed is therefore still obscure to me; but from the rules of the text, very short as they are and often disfigured by the corrupt Latin translation, I have found out how one must use the tables to determine the place of the casting of the rays. I give three \emph{examples} of this:\par
\par 1. To compute the place one needs in the first place the \emph{horoscope}\footnote{\emph{Horoscopus} = \emph{ascendens} (the ascendant).}, that is, the degree of the ecliptic that rises at the given time (of the computation); let it be 55° (= 25° of Taurus); in the second place the \emph{longitude} of the planet; let it be 65° (5° of Gemini). One now first looks up in the tables the field (\emph{area}) that corresponds to the horoscope 25° of Taurus; it is the field on Table 94 headed \emph{Viginti quinque gradus tauri} (twenty-five degrees of Taurus); one then enters this field with the longitude of the planet, 5° of Gemini, that is, one looks for 5° in the first column and goes horizontally along to the column that is headed at the top Gemini / Sagittar.; here one finds \Pg{101}the number 26° 45′; to this one must add the number that stands at the head of this field at \emph{Exagonum} (sextile aspect), that is, 53° 30′\footnote{This, naturally, when one is looking for the sextile aspect; if one is looking for the square or the trine aspect, one must add the number that stands at the head of the field at \emph{Tetragonum} or \emph{Trigonum} respectively.}; this gives 80° 15′. One looks for this number in the same field among the numbers of the six lines and six columns, and finds the next higher, 80° 23′, in the 5th line, 4th column; to it belongs (it is taken as an \emph{ascensio obliqua}) the longitude 25° of Cancer; for the difference 80° 23′ − 80° 15′ = 8′ the interpolation gives the number 7′, and so the place (= the longitude) of the \emph{left sextile aspect} is = 24° 53′ of Cancer\footnote{This, then, is c. 10° less than would result if one simply counted 60° forwards on the ecliptic from the place of the planet; one would then come to 5° of Leo.}. If one wants the place of the right sextile aspect, one must look for the place of the left trine aspect and add 180° to it.\par
\par 2. \emph{Horoscope} = 317° (= 17° of Aquarius), \emph{longitude} of the planet = 138° (= 18° of Leo). In the field of Table 111 headed \emph{Quindecim gradus Aquarii} (fifteen degrees of Aquarius) one finds for 18° of Leo, with the help of interpolation, the number 14° 16′; adding to it the number that stands at the head of the field at \emph{Exagonum}, 49° 28′, gives 63° 44′; for this number one finds in the same field, with the help of interpolation, the longitude 3° 40′ of Scorpio. But since the horoscope is not 15° but 17° of Aquarius, one must also determine the sextile aspect for 20° of Aquarius; by the same procedure one finds it in the next field (Table 112) = 3° 55′ of Scorpio; the difference of the two values is therefore 15′ for 5° of the horoscope, and so 6′ for 2° (from 15° to 17°); this, added to 3° 40′, gives 3° 46′ of Scorpio for the place sought of the \emph{left sextile aspect}.\par
\par 3. \emph{Horoscope} = 335° (= 5° of Pisces), \emph{longitude} of the planet = 115° (= 25° of Cancer). In the field of Table 113 headed \emph{Quinque gradus Piscium} (five degrees of Pisces) one finds for 25° of Cancer the number 117° 13′; adding to it the number that stands at \emph{Exagonum}, 48° 19′, gives 165° 32′; this number is no longer found in this field, whose highest number is 144° 58′; one subtracts this from the number obtained, and 20° 34′ remains; for this one finds in the 1st–2nd column of the same field, with the help of interpolation, the longitude 3° 45′ of Libra as the place of the \emph{left sextile aspect}.\par
\par When one has determined the \emph{left} sextile, square and trine aspects in this way, one finds the \emph{right} ones by the rules:\par
\begin{center}\small\begin{tabular}{llll}
The right & trine aspect & is the nadir of the left & sextile aspect. \\
The right & square aspect & is the nadir of the left & square aspect. \\
The right & sextile aspect & is the nadir of the left & trine aspect. \\
\end{tabular}\end{center}
\par Thus in the 1st example the right trine aspect is = 24° 53′ of Capricorn. In the three examples given, all the cases are shown that can possibly occur in these determinations and that are taken into account in the rules of the text.\par
\par \Pg{102}Kh.'s method probably belongs to the category named by \emph{Regiomontanus} (l. c.), of which he says: \emph{Alii autem exequuntur id negotium per ascensiones quidem rectas stella meridianum tenente, per obliquas autem ascensiones regionis, si in oriente extiterit, aut per descensiones si in occidente} (others, however, carry out this business by the right ascensions when the star holds the meridian, but by the oblique ascensions of the region if it is in the east, or by the descensions if it is in the west); only, Kh. does not treat separately the case where the star is in the meridian.\par
\par One must admit that these tables of Kh.'s were very convenient for the astrologers in determining the various aspects; only, they may have left something to be desired as regards accuracy.\par
\par With this the commentary on the text and the tables, as far as these agree in the manuscripts C, O and N, is ended. But Ms. O contains two further tables at the end (fol. 158$^{r}$ and 158$^{v}$), which Bjørnbo had not taken into his copy at all; he probably overlooked the sentence \emph{Explicit ezich} (here ends the zīj), etc., which stands at the end of Table 116, and therefore took these two tables for appendages added by a copyist that did not belong to the present work. It must be admitted that this would well be possible; but I think it more probable that these tables, since they do stand in some connection with the preceding astrological tables, formed the end of Maslama's recension of Kh.'s tables, especially as the second table (as the astrological rose) is also found among Batt.'s \emph{tabulae spuriae}\footnote{Ed. Nallino, II, 299. [\emph{Translator:} The print has no call for this note in the text; it is placed here, where the note gives the page of the astrological rose among the \emph{tabulae spuriae}.]}, which all derive from Maslama. In the two following (unnumbered) final chapters I therefore give a short description and explanation of these two tables.\par
\UnitHead{On the conversion of the years of birth}{On the conversion of the years of birth}{K24-9}
\par\begin{center}(Table 115.)\end{center}
\par Table 115 teaches how to determine the place in the heavens that the sun occupies after a number of full years have elapsed, for example on the 35th birthday of a person, or, in other words, how many degrees (or how much time) it still has to run until it has arrived again at the same point (fixed star) at which it stood at the birth. Understood from the purely astronomical standpoint, one can also regard this table as a table for converting civil (Julian) years into sidereal years. Batt. has (II, 188) a similar table for converting civil years into tropical years, but its \emph{anni collecti} proceed from 20 to 20 instead of from 10 to 10 as in our table, and the excesses over the full civil years are given only in degrees of arc, whereas in our table they are also given in measure of time.\par
\par \emph{Example:} Someone was born on 4 January 962 at a certain hour; how much time must he still wait, from the same hour of 4 Jan. 997, until the sun has arrived at the same place in the heavens, that is, at the same fixed star, at which it stood at the moment of the birth?\par
\par \Pg{103}Table 115 gives for 30 years (in the column \emph{anni deceni}, the years by tens) 7$^{d}$ 18$^{h}$ 4$^{m}$, for 5 years (in the column \emph{anni singulares}, the single years) 1$^{d}$ 7$^{h}$ 1$^{m}$, together 9$^{d}$ 1$^{h}$ 5$^{m}$. Only 9$^{d}$ 1$^{h}$ 5$^{m}$ after his 35th birthday\footnoteT{Table 115 counts the excess over years of 365 days: with al-Khwārizmī's year of 365.2584375 days (p. 42), 30 years exceed 30 × 365 days by 7$^{d}$ 18$^{h}$ 4$^{m}$, and 5 years exceed 5 × 365 days by 1$^{d}$ 7$^{h}$ 1$^{m}$. From 4 January 962 to 4 January 997 the Julian calendar counts 35 × 365 days and 9 leap days; the 9$^{d}$ 1$^{h}$ 5$^{m}$ thus run from 26 December 996, the end of 35 years of 365 days, and the sun returns to its place of the birth 1$^{h}$ 5$^{m}$ after the same hour of 4 January 997. The «civil (Julian) years» of p. 102 are likewise years of 365 days.}, then, does the sun stand at the same fixed star as at his birth. This time, however, undergoes a further correction through the determination of the \emph{true} place of the sun, for the place found is only the \emph{mean} one; but I shall not go into this further, and refer for it to Batt. I, 148. The second half of the table gives the path still to be run by the sun in degrees, minutes and seconds. — These \emph{revolutiones} or \emph{conversiones annorum nativitatum} (revolutions or conversions of the years of the nativities), and likewise the \emph{revolutiones annorum mundi} (or \emph{civitatum}) (revolutions of the years of the world, or of the cities), played a great part in the astrology of antiquity and of the Middle Ages\footnote{Compare Batt., ed. Nallino, I, 304–305.}.\par
\par It is still to be mentioned here that in this table the excess of the sidereal year over the civil year is taken as 6$^{h}$ 12$^{m}$, or, if we convert the 93° 2′ 15″ of the last column into time, more exactly as 6$^{h}$ 12$^{m}$ 9$^{s}$, whereas by chapters 20 and 21 6$^{h}$ 13$^{m}$ 2$^{s}$ would have been expected. The value of 6$^{h}$ 12$^{m}$ 9$^{s}$ corresponds, according to the statements of various Arab astronomers, to the length of the sidereal year in the Arabic \emph{Sindhind}; Maslama al-Majrīṭī too, then, seems to have adopted this value, and so not to have followed Kh. in this.\footnote{See above, pp. 42 and 65.}\par
\UnitHead{On the domicilia, exaltationes, decani, etc. (the houses, exaltations, decans, etc.) of the planets}{On the domicilia, exaltationes, decani, etc. (the houses, exaltations, decans, etc.) of the planets}{K24-15}
\par\begin{center}(Table 116.)\end{center}
\par This table is on the whole identical with the \emph{astrological rose} in Batt., as I have already mentioned above; but in single points, which are still to be named, the two differ from each other.\par
\par The 1st column\footnote{For convenience of citation I have provided the columns of the upper half of the table with numbers, which are missing in the Ms.} contains the twelve signs of the zodiac, which in the rose named stand outside the outermost circle, which is accordingly divided into twelve sectors. In the 2nd column stand the lords (\emph{domini}) of the corresponding houses (\emph{domicilia}) of the 1st column; it would therefore rather have to be headed \emph{domini} instead of \emph{domicilia}; thus the Ram (\emph{Aries}) is the house of \emph{Mars}, or \emph{Mars} is the lord of the house of the \emph{Ram}, \emph{Venus} that of the \emph{Bull} (\emph{Taurus}), and so on. The sun and the moon have only \emph{one} house each, but the five planets \emph{two} each; thus besides the \emph{Ram} the \emph{Scorpion} too is the house of \emph{Mars}, and besides the \emph{Bull} the \emph{Balance} (\emph{Libra}) too is the house of \emph{Venus}. — These \emph{domini} (Arab. \emph{ṣāḥib} or \emph{rabb}) stand in Batt.'s rose in the second innermost circle.\par
\par In the 4th column, headed \emph{alienatio} (estrangement), stand the planets that are in the so-called \emph{detriment} (\emph{detrimentum}, Arab. \emph{wabāl}) in the corresponding sign of the zodiac; a planet is in its detriment in the house that is 180° distant from its domicilium, or that is the nadir of its domicilium; it is then, as it were, estranged from this, hence the name \emph{alienatio}, which I have not so far found in any other medieval astrologer. Thus \emph{Venus} is in her \emph{detriment} in the \emph{Ram} and the \emph{Scorpion}, because her houses are the Balance and the Bull, \emph{Mars} in the \emph{Bull} and in the \emph{Balance} in \Pg{104}his \emph{detriment}, because his houses are the Scorpion and the Ram. — These \emph{alienationes} are missing in Batt.'s rose.\par
\par Columns 3 and 5 contain the so-called \emph{exaltations} (\emph{dignitates}, \emph{exaltationes}, Arab. \emph{sharaf}, pl. \emph{ashrāf}) and \emph{dejections} (\emph{deiectiones}, \emph{casus}, Arab. \emph{hubūṭ}), or rather the planets that are in exaltation or in dejection (fall) in the corresponding sign of the 1st column. Thus, for example, \emph{Jupiter} is in \emph{exaltation} in the \emph{Crab} (\emph{Cancer}), that is, in this sign he has the highest power, the greatest influence; on the other hand he is in \emph{dejection} or in \emph{fall} in the \emph{Goat} (\emph{Capricornus}), the sign that stands 180° from the Crab; in this sign he has the least influence, or none at all. \emph{Saturn} is in \emph{exaltation} in the \emph{Balance}, but in \emph{fall} in the \emph{Ram}. In the Lion and the Water-bearer (\emph{Aquarius}) no planet stands in exaltation and none in fall, in the Bull none in fall, in the Scorpion none in exaltation. Besides the sun, the moon and the five planets, the nodes of the moon, the dragon's head and the dragon's tail, also have their exaltationes and casus. To the 3rd column corresponds the innermost circle of Batt.'s rose; the 5th column is missing in it. In the rose there is also given the degree of the sign in which the planet reaches the maximum of its power (for example Venus in the 27th degree of the Fishes); these degrees are missing in our table. The headings of these columns too, \emph{regnum} (kingship) for \emph{dignitas} or \emph{exaltatio}, and \emph{servitus} (servitude) for \emph{casus} or \emph{deiectio}, we meet here for the first time.\par
\par Columns 6 and 7 contain the \emph{lords} or \emph{rulers} (here called \emph{Judices}, judges) of the so-called \emph{triplicities} (here called \emph{trigonalia}), that is, of the sets of three signs of the zodiac that form the corners of an equilateral triangle. Thus, for example, the \emph{Ram}, the \emph{Lion} and the \emph{Archer} form a \emph{triplicity} (that is, they stand in trine aspect to one another), and their lords (\emph{Judices}, \emph{domini}) are the \emph{sun}, \emph{Jupiter} and \emph{Saturn} by \emph{day}, and \emph{Jupiter}, the \emph{sun} and \emph{Saturn} by \emph{night}. Strictly, only the first two are the lords; the third, Saturn, is only their (somewhat subordinate) companion. These \emph{Judices} (\emph{domini}) of the triplicities are found in the astrological rose in the second outermost ring, and there too for day and for night.\par
\par The lower half of the table contains on the left the so-called \emph{facies} (faces) (\emph{decani}, Arab. \emph{wajh}, pl. \emph{wujūh}) and their lords (\emph{domini}). For each sign of the zodiac is divided into three equal parts of 10° each, and to each of these parts a planet (the sun and the moon included) is assigned as lord. This distribution was not carried out on the same principle by all the peoples who have occupied themselves with astrology; the distribution given here is the Egyptian-Greek one\footnote{Ptolemy excepted, who does not know the \emph{facies}.}, which has been adopted by most of the Arab astrologers: to the first decan of the Ram its lord, Mars, is assigned, to the second the sun, to the third Venus, to the first decan of the Bull Mercury, and so on, always in the old order of these heavenly bodies: Saturn, Jupiter, Mars, sun, Venus, Mercury, moon. — The astrological rose contains the facies in the 3rd (middle) ring, distributed in the same way as in our table.\par
\par \Pg{105}On the right in the lower half of the table the \emph{bounds} (\emph{termini}, here called \emph{fines}, Arab. \emph{ḥadd}, pl. \emph{ḥudūd}) are entered. Each sign of the zodiac is also divided into five parts, in general unequal, and to each of them one of the five planets (the sun and the moon drop out here) is assigned as lord. Thus to the first 6° of the Ram \emph{Jupiter} is assigned, to the next 6° \emph{Venus}, to the following 8° \emph{Mercury}, to the further 5° \emph{Mars} and to the last 5° \emph{Saturn}. Here too various systems of distribution existed; the one followed here differs from that of Batt.'s rose (in the outermost ring), and seems to be the system that Nallino\footnote{Batt. II, 309.} ascribes to the Egyptians, which was also adopted by most of the Arab astrologers. The title of this part of the table is \emph{Fines secundum Medos} (the bounds according to the Medes); the author, then, ascribes this method of distribution to the Medes; does he mean by this the Babylonians or the Persians? Batt. in his tables of kings (II, 2) calls a Darius, the predecessor of Cyrus, a Mede (\emph{Medus}), and so does al-Bīrūnī in his Chronology.\footnote{Ed. Sachau, English translation, p. 101.}\par
\par For further details concerning these astrological notions I refer, besides the Latin editions of Greek and Arab astrologers such as Ptol., Albumasar, Alcabitius, Albohazen, etc., also to the following writings: Batt., ed. Nallino, II, 309–314; Loth, \emph{Al-Kindi als Astrolog}, in \emph{Morgenländische Forschungen}, Festschrift f. H. L. Fleischer, Leipzig 1875; Dieterici, \emph{Die Propädeutik der Araber im 10. Jahrh.}, Berlin 1865, pp. 46–85; Boll, \emph{Sphaera}, Leipzig 1903; Bouché-Leclercq, \emph{L'astrologie grecque}, Paris 1899; Uhlemann, \emph{Grundzüge der Astronomie und Astrologie der Alten, besonders der Aegypter}, Leipzig 1857.\par
\par \Pg{106}If at the end we try to survey what the study of this work has yielded us about the sources that Kh. followed, we must admit that the result is not a decided one, and is certain only in a few points. It is certain that in determining the \emph{mean places} of the planets he adopted the \emph{Indo-Persian} procedure, that is, took the sidereal year as his basis; further, that in the size of the \emph{obliquity} of the \emph{ecliptic} he followed Ptol. It is very probable\footnote{See above, p. 32.} that in computing the \emph{equations} of the \emph{planets} he went the way of the \emph{Persians}, that is, of the \emph{Tables} of the \emph{Shāh}, which in turn are based on the \emph{Siddhāntas} of the \emph{Indians}; on which of them cannot be said with certainty. \emph{Persian} traces we notice, further, in the tables of the \emph{true} or \emph{unequal motion} of the sun and the moon, as the retention of the expression \emph{buht}, which Batt. does not know, proves. In the tables of the \emph{conjunctions} and \emph{oppositions} (\emph{syzygies}) Kh. (or Maslama?) probably followed mainly Ptol.; in the computation of the \emph{eclipses} and of the \emph{parallax of the moon} we again find an approach to the \emph{Indians}. That Kh. may have kept more to the indirect Persian sources than to the direct Indian ones is probable, and also easily explained, since he was indeed a Persian by birth; it is also possible that in the first edition of his tables he leaned still more on the Indian \emph{Siddhāntas}, and in the second more on the \emph{Tables} of the \emph{Shāh}, on Ptol., and perhaps also on the \emph{tested tables} of al-Maʾmūn's astronomers.\par
\par Our knowledge of the changes that Maslama al-Majrīṭī made in Kh.'s tables is likewise deficient in certain points. We can assert with certainty only that Maslama, in the tables for determining the \emph{mean places} of the planets, replaced the \emph{Persian era} of Yazdegerd by the \emph{Muslim} one, and that he reworked for the meridian and the latitude of \emph{Córdoba} the tables of the \emph{syzygies}, the \emph{eclipses} and the \emph{parallax of the moon} (if these were present at all in Kh.'s original work), and likewise the tables of the \emph{distribution} of the twelve \emph{astrological houses} and those of the \emph{proiectiones radiorum}. As for the latter tables, he made, according to what we know from Ibn Hibintā (see above, p. 100), some further changes in them. Added by Maslama himself are perhaps (this too is not certain) some of Tables 1 to 3$^{a}$ (see above, pp. 40–41), Table 60 of the \emph{cotangents}, Tables 67–68 of the \emph{equation of time}, and the two last Tables 115 and 116, on the \emph{conversion} of the \emph{years of birth} and on the \emph{domicilia}, \emph{exaltationes}, etc.\par
\par To make the judgement of what Kh.'s work may originally have contained, and what not, more difficult still, there is finally the further fact that the translator (or translators), perhaps also later revisers and arrangers of the translation (Robert of Chester), or copyists too, made additions to the text and also to some of the tables, and that these additions are not always as easy to recognize as such as is the case with some passages of the text that were discussed above, pp. 35 and 37.\par
\par But in spite of this manifold work of collaboration and alteration on these tables, they nevertheless keep, even in the form in which they lie before us, a value beside those of Batt. that is not to be underrated: they very probably form the basis of later Western Arabic tables, and so perhaps also of certain tables of the Christian Middle Ages; direct proofs of this view, to be sure, we cannot yet bring forward, but they may well still be found; perhaps it is precisely the tables of Roger of Hereford mentioned in the introduction that rest on the present ones. They show us further, very clearly, the great progress that Arabic astronomy made from Kh. to Batt.: in the latter, trigonometric formulas and mathematical proofs for the problems; in the former, only bare, short rules without any explanation, and very few traces of trigonometric derivation. If, then, on the one hand the greatness of the work of the most important Arab astronomer stands out all the more brilliantly through the publication of these tables, on the other hand they prove in turn that even at that time one knew how to produce, with quite elementary mathematical means, astronomical tables that, as we have seen, were able to satisfy no small demands of accuracy.\par
\PartHead{Additions and corrections}{B}\Pg{233}
\par \textbf{1.} To chapter 10 of my commentary (the determination of the true place of the three upper planets) I have some additions and corrections to give, to which my attention was drawn by a communication by letter from the late\footnoteT{The print has the abbreviation «w.», read as \emph{weiland} (the late): Pechüle died in May 1914, as the foreword (p. VI) says.} Mr. C. F. Pechüle, astronomer at the royal observatory in Copenhagen.\par
\par Nallino (ed. Batt., II, 241) was the first to point out that the values of the \emph{prosthaphaeresis} of longitude (the equation of the \emph{centre}) that Ptol. gives in the 3rd column of his tables for the true longitude of the planets\footnote{Edition of Heiberg, II, 436–444; translation by Manitius, II, 261–265.} do not correspond to the formula (p. 48)\par
\par\begin{center}$\mathrm{tg\ aeq.\ centri} = \dfrac{2e_{1}\sin m}{\sqrt{1-(e_{1}\sin m)^{2}}+e_{1}\cos m}$\end{center}
\par but to the simpler one\par
\par\begin{center}$\mathrm{tg\ aeq.\ centri} = \dfrac{EN}{HN} = \dfrac{2e_{1}\sin m}{1+2e_{1}\cos m}$,\end{center}
\par which follows from the figure on p. 47 of the commentary if one takes the centre \emph{H} of the epicycle to be on the \emph{equant} (which we have left out), and then sets the radius \emph{HE′} of the equant = 1. But since this is only approximately right, the centre \emph{H} moving not on the equant but on the \emph{deferent}, Ptol. corrects the values found by this formula by the differences of the 4th column; these corrected values Batt. then took into the 3rd column of his tables of the equations of the planets (ed. Nallino, II, 108–135).\par
\par \Pg{234}Now, since the values of the equation of the centre in Kh. differ essentially from those in Ptol.-Batt., I drew the conclusion (p. 50) that one half of the halved equation of the anomaly might have been used in Kh. to correct the equation of the centre. But this seems in fact not to be the case: Mr. Pechüle has informed me that the values of the equation of the centre in Kh. were probably computed by the formula $C_{\alpha} = C_{90}\cdot\sin\alpha$, where $C_{\alpha}$ means the equation of the centre for the argument \emph{{\greekfont α}}, and $C_{90}$ that for the argument 90°, that is, the greatest equation (= 2\emph{e}). Better, to be sure, would be in my view: $\sin C_{\alpha} = 2e\cdot\sin\alpha$; for very small angles one can, of course, put the arcs instead of the sines, but here it is after all a matter of angles of 8°–9°. If one computes the equations of the centre by the formula named, one finds the greatest deviation between the values so computed and those of Kh.'s tables to be 3$\tfrac{1}{2}$′–4′. — If Pechüle's view is right, and it really seems to me to be so, Kh. would then, in the formula (see the figure on p. 47)\par
\par\begin{center}$\sin E'HE = \sin\mathrm{aeq.\ centri} = \dfrac{EN}{HE}$,\end{center}
\par take \emph{HE} = 1 (or = 60$^{p}$); then, since $EN = 2e_{1}\sin m$,\par
\par\begin{center}$\sin\mathrm{aeq.\ centri} = 2e_{1}\sin m$.\end{center}
\par By this formula one can also determine the eccentricities of the orbits of the planets, which are recorded nowhere in Kh., neither in the text nor in the tables. For Saturn, for example, Kh. has the greatest equation 8° 36′ at 90° of centre, so that $\sin 8^{\circ}36' = 2e_{1}\sin 90^{\circ} = 2e_{1}$; this gives $2e_{1}$ = 0.1495 = 8$^{p}$ 58′; Ptol. has 6$^{p}$ 50′. The eccentricities of the orbits of Jupiter, Mars and Venus in Kh., computed in this way, agree fairly well with those in Ptol.-Batt.; in those for Saturn and Mercury, on the other hand, greater deviations show themselves, for Saturn more than 2$^{p}$.\par
\par As for the computation of the equation of the \emph{anomaly}, I gave on p. 49 the formula\par
\par\begin{center}$\mathrm{tg}\,RET = \mathrm{tg\ anom.} = \dfrac{r\sin\psi}{\text{Dist. centri epic.}+r\cos\psi}$,\end{center}
\par where \emph{r} means the radius of the epicycle expressed in decimal parts of the radius of the deferent, and \emph{{\greekfont ψ}} the anomaly. Mr. Pechüle informed me that in Ptol. the values of the equations of the anomaly (in col. 6 of his tables for the true longitude of the planets) correspond almost exactly to the formula\par
\par\begin{center}$\mathrm{tg}\,B_{\alpha} = \dfrac{\mathrm{tg}\,B_{90}\cdot\sin\alpha}{1+\mathrm{tg}\,B_{90}\cdot\cos\alpha}$,\end{center}
\par where $B_{\alpha}$ means the equation of the anomaly for the argument (anomaly) \emph{{\greekfont α}}. I informed him that Ptol.'s values would agree with the formula still better if one put \emph{sin aeq. anom. max.} (the sine of the greatest equation of the anomaly) instead of tg $B_{90}$, thus, for example, for Saturn sin $B_{96}$ instead of tg $B_{90}$; for Ptol. has the maximum of the anomaly for Saturn at 96° (Kh. likewise at 95°–98°), because it is not at 90° but at 96° of anomaly that the line \emph{EP} (see the figure on p. 49) touches the epicycle, so that the triangle \emph{EHP} then becomes right-angled at \emph{P}, and sin \emph{HEP} becomes = $\tfrac{HP}{HE}$, thus sin \emph{HEP} = \emph{HP} = \emph{r} if \emph{HE}, that is, the radius of the deferent, is set = 1; the formula would then run:\par
\par\begin{center}$\mathrm{tg}\,B_{\alpha} = \dfrac{r\cdot\sin\alpha}{1+r\cdot\cos\alpha}$,\end{center}
\par \Pg{235}thus quite corresponding to the formula for the equation of the centre if in place of $2e_{1}$, that is, of the eccentricity of the orbit of the planet, the radius \emph{r} of the epicycle is put. — Mr. Pechüle has now also tested the values of the equation of the anomaly in Kh.'s tables against this formula, and found on the whole fair agreement (greatest deviations 4′–5′); only for Mars did differences of up to 21′ show themselves (at an anomaly of 160°–170°). Kh. seems, then, in the formula given above for tg \emph{RET}, to have taken the distance of the centre of the epicycle from the earth as constant, that is, = 1 (or 60$^{p}$), for then it changes into\par
\par\begin{center}$\mathrm{tg\ aeq.\ anom.} = \dfrac{r\sin\psi}{1+r\cdot\cos\psi}$.\end{center}
\par Only for Mars, which has a fairly large eccentricity, does he seem to have taken the variability of this distance somewhat into account; but in what way can hardly be found out.\par
\par With these additions to chapter 10, however, the reasons why Kh. diminished the mean apogee of the planets by half the equation of the anomaly are by no means revealed.\par
\par \textbf{2.} On p. 48, line 10 from the top, instead of “because it is thought of as measured from the centre of the equant” it must read: “because it gives the position of the centre of the epicycle with respect to the apogee of the eccentric”.\par
\par \textbf{3.} On p. 46, etc., I have given the astronomical work of the Indian Varāhamihira in the form \emph{Panchasiddhāntikā}, which Thibaut used after the English manner of transcription; by the manner of transcription now generally adopted by the orientalists it should be written \emph{Pañcasiddhāntikā}.\footnoteT{This translation writes \emph{Pañcasiddhāntikā} throughout, except in the title of Thibaut's edition.}\par
\par \textbf{4.} On p. 42 (chapter 7) I computed the mean motion of the sun in 360 days by multiplying the mean motion in 30 days by 12; this is not quite right, for if the 5″ in the mean motion in 30 days (= 29° 34′ 5″) still had a fraction that was thrown away, the error so arising would become that many times greater through the multiplication by 12; we therefore give another determination below:\par
\begin{center}\small\begin{tabular}{lllrlcrl}
Mean motion & in & 1 year of & 354 & days & = & 348° 54′ 13″ & (see T. 4) \\
— & — &  & 11 & — & = & 10° 50′ 30″ & (see T. 5) \\
— & — &  & 6 & hours & = & 14′ 47″ & (see T. 5) \\
— & — &  & 12 & minutes & = & 30″ & (see T. 5) \\
\midrule
Mean motion & in & \multicolumn{3}{c}{365$^{d}$ 6$^{h}$ 12$^{m}$} & = & 360° 0′ 0″ &  \\
\end{tabular}\end{center}
\par \Pg{236}that is, Kh. took the length of the sidereal year as 365$^{d}$ 6$^{h}$ 12$^{m}$. This would agree with what al-Bīrūnī reports to us in his Chronology (ed. Sachau, Arabic text p. 11, English transl. pp. 12–13) about the Persian solar year; in another place of the same work (Arabic text p. 52, English transl. pp. 61–62), to be sure, he then adds a further 9$^{s}$ to this length of the year, and so obtains the value that the \emph{Brahma-Siddhānta} and the Arabic \emph{Sindhind} also have. It seems, then, that in this point too Kh. had probably followed the \emph{Persian} tables of the \emph{Shāh} (compare also Batt., ed. Nallino, I, 206–209). — This remark is also to be taken into account in my statements on pp. 65 and 103.\par
\par \textbf{5.} On Table 92, in the notes, \emph{S... elmirech} should probably read \emph{Secundum Emirethium} (according to Emirethius) (compare also Table 7), and refer to the astronomical tables of Eumathius (that is, probably Ammonius), which were improved by al-Zarqālī; of this improvement there is a Ms. in Munich (no. 853 of Aumer's catalogue), whose end runs: “Finished is the \emph{Kānōn} of Eumathius, improved by Abū Isḥāq al-Naqqāsh, known by the name al-Zarqālah, in the middle of \emph{Dhū l-Ḥijja} 655 (end of 1257).” — Ammonius lived at Alexandria at the end of the 5th century AD (compare also Batt., ed. Nallino, I, XXXV, n. 5). — I have also thought of \emph{al-mirrīkh} = Mars, and likewise of \emph{al-mizāj} = temperament, the astrological nature of a heavenly body (compare Batt., ed. Nallino, II, 352–353), but neither gives any sense in that place (Table 92).\par
\par \textbf{6.} When this work was already half printed, I received from Prof. C. H. Haskins in Cambridge (Mass.) the information that in Ms. 283 of \emph{Corpus Christi Coll.} in Oxford there is a copy of Adelard's translation of Kh.'s\footnoteT{The print has «Kb.’s».} tables, which, however, is probably rather incomplete. He remarks further that, as far as he could establish, this text agrees rather with that of the Oxford Ms. \emph{Auct. F. I. 9} than with that of the Mss. of Chartres and Madrid. It would certainly be of interest to draw on this Ms. for comparison even now; but I must leave this work to such scholars as have easier access to this Ms. than I, that is, in the first place to the English ones.\par
\par \textbf{7.} In chapter 34 it was forgotten to draw the reader's attention to the fact that in Tables 77 and 77$^{a}$ the parallax of latitude is given in parts of the circle, not in parts of time like that of longitude. — In Table 77$^{a}$ the heading \emph{Longitudo et Latitudo} (longitude and latitude) is misleading, for the values entered in these columns for the arguments 91–150 are only parallaxes of longitude.\par
\UnitHead{Addendum to the corrections}{Addendum to the corrections}{B01-30}
\par On p. 47, l. 14 from below, instead of “The upper planets complete . . . . times” it should read: The two uppermost planets (Jupiter and Saturn) complete their circuit on the epicycle in about one year, Mars, on the other hand, in c. 2 years (780 days); the centres of their epicycles complete their circuit on the deferent in different times, which agree fairly well with the present sidereal periods.\footnoteT{A slip printed on one side, without a page number, tipped in after p. 236. The translation of p. 47 gives the new wording.}\par
\UnitHead{Misprints}{Misprints}{B01-32}\Pg{237}
\par\begin{center}(In giving the line numbers the footnotes are not counted.)\footnoteT{The corrections to the commentary are made in the translation of the commentary; the copy transcribed has «167°» on p. 50 already. The corrections to the Latin text and the tables concern the edition of the Latin, which is not translated here.}\end{center}
\begin{center}\small\begin{tabular}{rr>{\raggedright\arraybackslash}p{0.75\textwidth}}
Page & 3, & note 32: read kiteb hea (librum canonum N) for librum hea (canonum N). \\
— & 6, & l. 10 from below: read $^{26}$) for $^{21}$). \\
— & 17, & note 8 must be inserted between CO and alius: \emph{sed supra} planum \emph{scrips.} C elmustewi, \emph{supra} diminutum elmaruz (but above planum C has written elmustewi, above diminutum elmaruz). \\
— & 22, & l. 9 from the top: read $^{11a}$) for $^{10a}$). \\
— & 26, & l. 1 from below: read $^{39}$) and $^{40}$) for $^{38}$) and $^{39}$), and likewise in the notes (last line) $^{40}$) for $^{39}$). \\
— & —, & l. 8 from below: read tabulam for tabulum. \\
— & 31, & l. 15 from the top (beginning): read tetragoni for trigoni. \\
— & 49, & l. 11 from below: read p. VII for p. I. \\
— & 50, & l. 7 from below: read 167° for 67°. \\
— & 62, & l. 2 from below: read 48′27″ for 48″27′. \\
— & 66, & note 1: read p. 103 for Kap. 38 (chapter 38). \\
— & 92, & l. 19 from below: read so for zo. \\
— & 100, & ll. 9 and 10 from the top: read Ascensio obliqua for Rectascension. \\
— & 120, & (Tab. 9): in the section headed “Heza Saturni supradictum”, at 60 (Dies), in the column of the “Signa” read 1 for 0. \\
— & 131 & (Tab. 20) n. 7: read in for tn. \\
— & 145 & (Tab. 34) n. 1: read Jauzegarhu for Jauzezarhu. \\
\end{tabular}\end{center}
\PartHead{Glossary and index of names and subjects}{G}\Pg{239}
\UnitHead{Glossary\footnoteT{The Arabic forms keep Suter's transcription, explained in the introduction (p. XXV); his \emph{dj} is the \emph{j} and his \emph{ḳ} the \emph{q} of the transliteration used elsewhere in this translation. The meanings are translated into English.}}{Glossary}{G01-1}\Pg{241}
\par\begin{center}I. Arabic (and some Persian) words\end{center}
\par\begin{center}\textbf{A.}\end{center}
\Entry{\emph{Aakrab} see \emph{Akrab}.}
\Entry{\emph{Aalemet} = \emph{ʿalāmāt} (pl. of \emph{ʿalāma}) \emph{Notulae}, signs, marks, signatures, here the numbers of the days of the week (1 . . . 7 = Sunday to Saturday), T. 69.}
\Entry{\emph{Aard} see \emph{Ard}.}
\Entry{\emph{Abaad} = \emph{abʿad}, farther away (\emph{longinquior}), T. 73, 74.}
\Entry{\emph{Achir} = \emph{ākhir}, (the) last, second, T. 2, 3 etc.}
\Entry{\emph{Acofra} = \emph{aṣ-ṣufr} (\emph{Ta’rīkh aṣ-ṣufr}), the Spanish era, also called simply the era (in the text \emph{Era}), 3, 35, T. 1.}
\Entry{\emph{Acrab} see \emph{Akrab}.}
\Entry{\emph{Aded} and \emph{Adet} = \emph{ʿadad}, number, T. 21 etc.; perhaps it sometimes also means the plural \emph{aʿdād}.}
\Entry{\emph{Adel} (?) T. 10.}
\Entry{\emph{Adhel} = \emph{al-dhel}, see \emph{Dhel}.}
\Entry{\emph{Aguze} = \emph{adjzā’} (pl. of \emph{djuz}), parts, T. 73.}
\Entry{\emph{Ahemar} = \emph{aḥmar}, red, here also the name of the planet Mars, besides \emph{Mirrīkh}, T. 39–42.}
\Entry{\emph{Ajuze} see \emph{Aguze}.}
\Entry{\emph{ʿAḳd}, node (of the orbits of the planets); \emph{ʿaḳd al-shamālī}, the northern node (of the orbit of the moon) 41.}
\Entry{\emph{Akrab}, 1) = \emph{aḳrab}, nearer (\emph{propinquior}) T. 75, 76; 2) = \emph{ʿaḳrab}, scorpion (the 8th sign of the zodiac) T. 59$^{b}$.}
\Entry{\emph{Anāhīd}, Persian name of the planet Venus, XXII, T. 45.}
\Entry{\emph{Arafa} (?) 15 (N. 20).}
\Entry{\emph{Ard} = \emph{ʿarḍ}, latitude of a heavenly body, or geographical latitude of a place, 19 (N. 1), T. 22, 79.}
\Entry{\emph{Arhik} T. 15 (N. 7) see \emph{Achir}.}
\Entry{\emph{Arin} = \emph{Uzain} = \emph{Udjain}, VII, XVIII, 1, 8, 19, 33, etc.}
\Entry{\emph{Artifa} = \emph{irtifāʿ}, altitude (of a heavenly body), 18, 21, 22, etc. T. 60.}
\Entry{\emph{Aṣbaʿ}, finger, inch, an Arabic measure of length, 43.}
\Entry{\emph{Asceir} see \emph{Sceir}.}
\Entry{\emph{Ascowa} see \emph{Shoa}.}
\Entry{\emph{Ashoa} see \emph{Shoa}.}
\Entry{\emph{Asmir} = \emph{as-sair}, or \emph{al-masīr}, the course, the path (of a heavenly body), XXI, T. 9.}
\Entry{\emph{Athofen} = \emph{aṭ-ṭūfān}, the Flood (\emph{diluvium}) 5, 35, T. 1.}
\Entry{\emph{Aub}, corrupted from \emph{Aug}; see this.}
\Entry{\emph{Audel}, probably corrupted from \emph{ʿarḍ el}, latitude of the . . .\footnoteT{Suter writes «Breite des . . ., oder Breite der . . .», with the German article of the genitive in its two forms.} T. 10.}
\Entry{\emph{Aug} = \emph{audj} (\emph{awdj}), \emph{apogeum} = the greatest distance of a planet from the earth, 9, 10, 14; \emph{aug el-musteri muadala} = the corrected apogee of Jupiter, T. 36; see also \emph{Muadala}.}
\Entry{\emph{Autarid} see \emph{Otari[d]}.}
\Entry{\emph{Avvel} see \emph{Awel}.}
\Entry{\emph{Avtarid} see \emph{Otari[d]}.}
\Entry{\emph{Awel} = \emph{awwel}, (the) first, T. 2, 3, etc.}
\Entry{\emph{Azabaa} = \emph{aṣābiʿ} (pl. of \emph{aṣbaʿ}), fingers (\emph{digiti}), the measure for the diameter of the sun and the moon, also for the earth's shadow in eclipses of the moon, T. 73.}
\Entry{\emph{Azofar} see \emph{Acofra}.}
\par\begin{center}\textbf{B.}\end{center}
\Entry{\emph{Bāb}, gate, chapter, 58.}
\Entry{\emph{Baheram} and \emph{Bahram} = \emph{Bahrām}, Persian name for the planet Mars, XXII, T. 39–42.}
\Entry{\emph{Batiae}, probably the feminine of \emph{baṭī’} = \emph{baṭī’a}, slow, T. 24.}
\Entry{\emph{Bei ed halt} see the following.}
\Entry{\emph{Beied hakht} = \emph{Bēdukht}, Persian name for the planet Venus, XXII, T. 45.}
\Entry{\Pg{242}\emph{Bergis} = \emph{Birdjis}, Persian name for the planet Jupiter, T. 33, 36.}
\Entry{\emph{Bi}, joined with the article \emph{bil}, Arabic preposition, in, on, with, for example \emph{bilfelek}, in (or on) the sphere (see \emph{felek}).}
\Entry{\emph{Bhukti} (Sanskrit) enjoyment, pleasure, 78.}
\Entry{\emph{Bood} = \emph{buʿd}, distance, T. 73.}
\Entry{\emph{Buht} (Persian), in astronomy the true (unequal) daily motion of the sun, the moon and the five planets, 22, 78, 90.}
\Entry{\emph{Buiut} = \emph{buyūt} (pl. of \emph{bait}), houses, especially the twelve astrological houses, 25, T. 79.}
\Entry{\emph{Burug} = \emph{burūdj} (pl. of \emph{burdj}), castles, houses, especially the twelve signs (houses) of the zodiac, 11 (N. 40), 18 (N. 23), T. 79.}
\Entry{\emph{Buruga} see \emph{Burug}.}
\Entry{\emph{Burus} see \emph{Burug}.}
\par\begin{center}\textbf{C.}\end{center}
\Entry{\emph{Caat} = \emph{sāʿāt} (pl. of \emph{sāʿa}), hours, perhaps also the singular, T. 5.}
\Entry{\emph{Caratan} = \emph{saraṭān}, crab, here the sign of the Crab in the zodiac, T. 59$^{a}$.}
\Entry{\emph{Care} = \emph{sarīʿ}, fast, T. 24.}
\Entry{\emph{Carne} (?) T. 9.}
\Entry{\emph{Carneli} (?) T. 13.}
\Entry{\emph{Catar} = \emph{saṭar} or \emph{saṭr}, line, row, column, XIV; pl. \emph{cotor} = \emph{suṭūr}; \emph{catar eladed} = \emph{saṭar al-ʿadad}, column of the number, T. 21, 22, 25–30, 32, etc.}
\Entry{\emph{Catur} see \emph{Catar}.}
\Entry{\emph{Cauz} = \emph{ḳaus} (\emph{ḳaws}), arc (of a circle), also the sign of the Archer in the zodiac, T. 59$^{b}$.}
\Entry{\emph{Cavsc} see \emph{Cauz}.}
\Entry{\emph{Ciicuth} see \emph{Zukut}.}
\Entry{\emph{Ciikuf} see \emph{Zukut}.}
\Entry{\emph{Cotoreladet} = \emph{suṭūr al-ʿadad}, columns of the number, 27 (N. 16); see also \emph{Catar}.}
\Entry{\emph{Cukut} see \emph{Zukut}.}
\Entry{\emph{Cumbula} = \emph{sumbula}, ear of corn, here the sign of the Virgin in the zodiac, T. 59$^{b}$.}
\Entry{\emph{Cuzuf} see \emph{Kucuf}.}
\par\begin{center}\textbf{D.}\end{center}
\Entry{\emph{Dakaicae} = \emph{daḳā’iḳ} (pl. of \emph{daḳīḳa}) minutes (the 60th part of the degree or of the hour) XXIII, XXV; 3, 4, 36, etc.}
\Entry{\emph{Dakaik} see \emph{Dakaicae}.}
\Entry{\emph{Dakeca} see \emph{Dakaicae}.}
\Entry{\emph{Darega} = \emph{daradja} (pl. \emph{daradj}) degree (the 360th part of the circumference of the circle) 11 (N. 40).}
\Entry{\emph{Dareis}, corrupted from \emph{daregis}, dat. plur. of \emph{darega}; see this.}
\Entry{\emph{Dehi} = \emph{dalw} (\emph{dalu}), large jug, bucket, the sign of the Water-bearer (\emph{Aquarius}) in the zodiac, T. 59.}
\Entry{\emph{Delidhe} see \emph{Thelitha}.}
\Entry{\emph{Dhel} = \emph{ẓill} shadow, then the trigonometric function of the \emph{cotangent}, T. 60.}
\Entry{\emph{Dhirāʿ}, ell (a measure of length) 43.}
\Entry{\emph{Djebr}, restoration; \emph{al-djebr wa’l-muḳābala} = algebra, VIII; see also \emph{Muḳābala}.}
\Entry{\emph{Dulheia} and \emph{Dulheiga} = \emph{Dhū’l-ḥidjdja}, the twelfth month of the Muslims, the month of the pilgrimage, XIV, XV, 2, T. 2, 3, etc.}
\Entry{\emph{Dulkada} = \emph{Dhū’l-ḳaʿda}, the eleventh month of the Muslims, XIV, XV, 2, T. 2, 3, etc.}
\Entry{\emph{Dulkarnain} = \emph{Dhū’l-ḳarnain}, the two-horned (\emph{bicornus}), the Arabs' epithet for Alexander the Great, 3, 35.}
\Entry{\emph{Dulkeda} and \emph{Dulkidda} see \emph{Dulkada}.}
\par\begin{center}\textbf{E.}\end{center}
\Entry{\emph{Ecumbula} see \emph{Cumbula}.}
\Entry{\emph{Ekucufet} see \emph{Kucufet}.}
\Entry{\emph{Eneid} see \emph{Anāhīd}.}
\Entry{\emph{Etdeiard} see \emph{Jazdeiard}.}
\Entry{\emph{Ezdeiard} see \emph{Jazdeiard}.}
\Entry{\emph{Ezed} = \emph{asad}, lion, the sign of the Lion in the zodiac, T. 59$^{a}$.}
\Entry{\emph{Ezeig} see \emph{Zich}.}
\Entry{\emph{Ezich} see \emph{Zich}.}
\par\begin{center}\textbf{F.}\end{center}
\Entry{\emph{Felek} = \emph{falak}, sphere, vault of heaven, circular path of a heavenly body about the earth, 7, 18, 25, etc.; \emph{felek elmustakim} see \emph{Mustakim}; \emph{felek eltedewir} see \emph{Tedewir}.}
\Entry{\emph{Feler} see \emph{Felek}.}
\Entry{\emph{Fi}, joined with the article \emph{fil} (= \emph{fi’l}), Arabic preposition, in, on, for example \emph{fiaard} = \emph{fī ʿarḍ}, in latitude (with a following genitive) T. 79 (here \emph{Cordubae} has dropped out after \emph{fiaard}) — \emph{filcaat} and \emph{filacaat} = \emph{fī’l-sāʿa}, in the hour, or \emph{fī’l-sāʿāt}, in the hours, T. 5 — \emph{filcinun} = \emph{fī’l-sinīn}, in the years, T. 69 — \emph{fil dakaik elcaat} = \emph{fī daḳā’iḳ al-sāʿa}, in the minutes of the hour, T. 5 — \emph{filiaum} = \emph{fī’l-yaum}, in the day, T. 5 — \emph{filshohor} = \emph{fī’l-shuhūr}, in the months, T. 9 — \emph{fil zahat} see \emph{filcaat}.}
\Entry{\emph{Fudhul} = \emph{fuḍūl} (pl. of \emph{faḍl}), excesses, differences, 27 (N. 19), T. 73.}
\par\begin{center}\Pg{243}\textbf{G.}\end{center}
\Entry{\emph{Gauze} see \emph{Jauze}.}
\Entry{\emph{Gauzehar} = \emph{djauzahar} (Persian), the points where the orbits of the planets intersect the ecliptic, especially the nodes of the moon (see also \emph{ʿAḳd}) 15, 41, T. 19, 20.}
\Entry{\emph{Gavzehar} see \emph{Gauzehar}.}
\Entry{\emph{Gedi} = \emph{djady} (or \emph{djadi}), he-goat, kid, the sign of the Goat in the zodiac, T. 59.}
\Entry{\emph{Gedval} = \emph{djadwal}, table, tabulation (astronomical) 28, T. 60–66, 69–71.}
\Entry{\emph{Geib} = \emph{djaib}, the trigonometric function of the \emph{sine}, 17–19, 21, 22, 69, 71, T. 58, 58$^{a}$.}
\Entry{\emph{Genubi} = \emph{djanūbī}, southern, 13, 16, 18, 29 (N. 1, 6, 14, 15); \emph{genubia} = \emph{djanūbīya} (fem.). 29 (N. 1, 6, 15).}
\Entry{\emph{Genubium} see the preceding.}
\Entry{\emph{Gidher} = \emph{djidhr}, root, especially square root, 22.}
\Entry{\emph{Gumedi} see \emph{Jumedi}.}
\par\begin{center}\textbf{H.}\end{center}
\Entry{\emph{Haazem} (?) XXII, 20 (N. 3).}
\Entry{\emph{Hacib} = \emph{ḥāsib}, calculator, XXII, 20 (N. 3).}
\Entry{\emph{Hacil} = \emph{ḥāṣil}, outcome, result (for example of a computation) 11.}
\Entry{\emph{Ḥadd} (pl. \emph{ḥudūd}), bounds, parts of the signs of the zodiac (Latin \emph{termini}, \emph{fines}) 105.}
\Entry{\emph{Hadhid} = \emph{ḥaḍīḍ}, lowest part, in astronomy the \emph{perigeum} of a planet, rendered in this work by \emph{submissio}, 14 (N. 28).}
\Entry{\emph{Hamel} = \emph{ḥamal}, ram, the sign of the Ram (\emph{Aries}) in the zodiac, T. 59.}
\Entry{\emph{Harket} = \emph{ḥaraka}, motion; \emph{harket elscems} = \emph{ḥarakat al-shams}, motion of the sun, T. 61.}
\Entry{\emph{Hea} = \emph{hai’a}, outward appearance (of the heavens); \emph{ʿilm al-hai’a}, the science of the heavens, astronomy: see also \emph{Kiteb}.}
\Entry{\emph{Heca} see \emph{Heza}.}
\Entry{\emph{Helhacib} see \emph{Hacib}.}
\Entry{\emph{Hemar} see \emph{Ahemar}.}
\Entry{\emph{Hemeli} see \emph{Shemeli}.}
\Entry{\emph{Heza} = \emph{ḥāṣṣa} and \emph{khāṣṣa}, the motion peculiar to the planet on the epicycle, the so-called \emph{anomaly}, rendered in this work by \emph{argumentum}, XIV, XXI, 9–14 (and notes), 28 (N. 1) 46, 53; T. 6–9, 11, 13, 15–18, etc.}
\Entry{\emph{Higera} = \emph{Hidjra}, Muhammad's flight from Mecca, the beginning of the Muslim era, 1, 3, 33, etc. T. 1.}
\Entry{\emph{Ḥiṣṣa}, share, lot, the argument with which one enters a table, 46.}
\Entry{\emph{Hormi} and \emph{Horom}, corrupted from \emph{Hormuzd} or \emph{Hormuz} = \emph{Ormuzd}, Persian name for the planet Jupiter, XXII, T. 33, 36.}
\Entry{\emph{Hout} = \emph{ḥūt}, fish, the sign of the Fishes in the zodiac, T. 59.}
\Entry{\emph{Hubūṭ}, fall, falling, in astrology the dejection, the fall (\emph{casus}) of a planet, 104.}
\par\begin{center}\textbf{I.}\end{center}
\Entry{\emph{Iblīs}, Satan, the devil, XXI.}
\Entry{\emph{Idhnascer} = \emph{ithnā ʿashar}, twelve, T. 79.}
\Entry{\emph{Igstimaat} = \emph{idjtimāʿāt} (pl. of \emph{idjtimāʿ}) see \emph{Istima}.}
\Entry{\emph{Ihtilef} = \emph{ikhtilāf}, difference T. 73; \emph{elihtilef lilmanshar} = \emph{al-ikhtilāf lil-manẓar}, the difference of the aspect (of a heavenly body), that is, its \emph{parallax} 28 (N. 3).}
\Entry{\emph{Ihtilief}, corrupted from \emph{Ihtilef}; see this.}
\Entry{\emph{Ingile} = \emph{indjilā’}, disappearance, especially the going out or the end of an eclipse, the phase from the totality to the end (Latin \emph{emersio}), 26, 87.}
\Entry{\emph{Ingine} see \emph{Ingile}.}
\Entry{\emph{Iṣbaʿ} see \emph{Aṣbaʿ}.}
\Entry{\emph{Ishemz} corrupted from \emph{eshemz} = \emph{el-shems}; see \emph{Shems}.}
\Entry{\emph{Isticbelet} = \emph{istiḳbālāt} (pl. of \emph{istiḳbāl}) see \emph{Istikbel}.}
\Entry{\emph{Istikbel} = \emph{istiḳbāl}, opposition (of two heavenly bodies), especially of the sun and the moon, 23, T. 70.}
\Entry{\emph{Istima} = \emph{idjtimāʿ}, conjunction (of two heavenly bodies), especially of the sun and the moon, 23, T. 69.}
\Entry{\emph{Istimaz} see \emph{Istima}.}
\par\begin{center}\textbf{J.}\end{center}
\Entry{\emph{Jauze} = \emph{djauzā’} (or \emph{djawzā’}), the sign of the Twins in the zodiac, T. 59$^{a}$; also the constellation of Orion.}
\Entry{\emph{Jauzehar} see \emph{Gauzehar}.}
\Entry{\emph{Jazdeiard} = \emph{Yezdegird} III, the last king of the Sasanians, XV, XVI, T. 1.}
\Entry{\emph{Jenubi} see \emph{Genubi}.}
\Entry{\emph{Jumedi elachir} = \emph{Djumādā al-ākhira}, \emph{Djumādā II.}, the sixth month of the Muslim year, XIV, XV, T. 2, 3, etc.}
\Entry{\emph{Jumedi elule} = \emph{Djumādā al-ūlā}, \emph{Djumādā I}, the fifth month of the Muslim year, XIV, XV, T. 2, 3, etc.}
\par\begin{center}\Pg{244}\textbf{K.}\end{center}
\Entry{\emph{Kamar} = \emph{ḳamar}, moon, T. 2$^{a}$, 6–8, 21–26, etc.}
\Entry{\emph{Kamarie} = \emph{ḳamarīya}, fem. of \emph{ḳamarī}, relating to the moon, T. 73, 74.}
\Entry{\emph{Kamene}, corrupted from \emph{Kamarie}.}
\Entry{\emph{Karnain} see \emph{Dulkarnain}.}
\Entry{\emph{Kawekib} (also \emph{Kavvekib}) = \emph{kawākib} (pl. of \emph{kaukab} or \emph{kawkab}), the heavenly bodies, especially the seven planets, T. 55 (N.).}
\Entry{\emph{Keiwen} = \emph{Kaiwān}, Persian name for the planet Saturn, XXII, T. 27–29.}
\Entry{\emph{Ketib} = \emph{Kātib}, scribe, then a name for the planet Mercury used mainly in Spain and North Africa, T. 51, 53.}
\Entry{\emph{Kiteb} = \emph{Kitāb}, book; \emph{Kiteb hea} = \emph{kitāb al-hai’a}, book of astronomy, 3 (N. 32); perhaps the astronomy of \emph{Djābir b. Aflaḥ} is meant by this, which bears the title \emph{kitāb al-hai’a} in the Ms. of the Escorial.}
\Entry{\emph{Kucuf} and \emph{Kuzuf} = \emph{kusūf}, eclipse, darkening (of the sun and the moon) 27, T. 73, 74.}
\Entry{\emph{Kucufet} = \emph{kusūfāt} (pl. of \emph{kusūf}) see \emph{Kucuf}.}
\Entry{\emph{Kuscus}, probably = \emph{ḳaus}, arc, 9; see also \emph{Cauz}.}
\Entry{\emph{Kuszuf} see \emph{Kucuf}.}
\par\begin{center}\textbf{L.}\end{center}
\Entry{\emph{Li}, joined with the article \emph{lil}, Arabic preposition, to, for, because of, with respect to; it also denotes simply the dative; also a conjunction: that, so that. For example \emph{libood elabaad} = \emph{lil-buʿd al-abʿad}, for the greatest distance (\emph{pro spatio longinquo}) T. 73 — \emph{liburug elhamel} = \emph{liburdj al-ḥamal}, for the sign of the Ram, T. 79 — \emph{lieztekim} = \emph{liyastaḳīm}, that he may determine, regulate . . . XXI, T. 9 — \emph{lilkawekib} = \emph{lil-kawākib}, for the heavenly bodies, or with respect to the heavenly bodies, T. 91 — \emph{lilmanshar} = \emph{lil-manẓar}, with respect to the aspect (the place of appearance) 28; see also \emph{Manshar} — \emph{liwacat alcuzuf} = \emph{liwasaṭ al-kusūf}, for the middle of the eclipse, 27; see also \emph{Wacat}.}
\par\begin{center}\textbf{M.}\end{center}
\Entry{\emph{Mabsūṭ}, spread out (\emph{planus}), 36, 77; see also \emph{Mobzota}.}
\Entry{\emph{Magesti} (also \emph{Megesti}) = \emph{Madjisṭī}, the Almagest (the great Syntaxis) of Ptolemy, 2, 18, etc.}
\Entry{\emph{Magmoa} = \emph{madjmūʿ}, sum, totality; also a participle = collected (fem. \emph{madjmūʿa}) T. 15 (N. 1); \emph{filcinun elmagmoa} = \emph{fī’l-sinīn al-madjmūʿa}, in the years united (by 30) (\emph{in annis collectis}) T. 69.}
\Entry{\emph{Mahwar} = \emph{miḥwar}, axis, the axis of the earth, here its pole, 18.}
\Entry{\emph{Maimoa} see \emph{Magmoa}.}
\Entry{\emph{Makuz} = \emph{maʿkūs}, reversed, 22 (N. 6), 69; see also \emph{Mankuz}.}
\Entry{\emph{Mandhar} = \emph{manẓar}, aspect, place of appearance, 28 (N. 3), T. 77; see also \emph{Ihtilef}.}
\Entry{\emph{Mankut} see \emph{Mankuz}.}
\Entry{\emph{Mankuz}, 1) = \emph{mankūs}, reversed, 17, 22, 69, 77; 2) = \emph{manḳūṣ}, diminished (\emph{diminutus}), 17, 69.}
\Entry{\emph{Manshar} see \emph{Mandhar}.}
\Entry{\emph{Marech} = \emph{Mirrīkh}, the planet Mars, T. 13, 14, 39, 41–44.}
\Entry{\emph{Mareh} see \emph{Marech}.}
\Entry{\emph{Markez} = \emph{markaz}, middle point, centre, 48; \emph{Tadil elmarkez} = \emph{Taʿdīl al-markaz}, equation of the centre, T. 54.}
\Entry{\emph{Maruz}, corrupted from \emph{Mankuz}; see this.}
\Entry{\emph{Masīr}, course, path, XXI; \emph{al-masīr al-wasaṭ}, the mean course, the mean motion, 41.}
\Entry{\emph{Matale} = \emph{maṭāliʿ} (pl. of \emph{maṭlaʿ}), ascensions of the heavenly bodies (\emph{ascensiones rectae et obliquae}), rendered in this work by \emph{horoscopi}, 18 (N. 23), T. 59–59$^{b}$, 91 (N. 1), 95 (N. 1), 97 (N. 4).}
\Entry{\emph{Matarech} see \emph{Matareh}.}
\Entry{\emph{Matareh} = \emph{maṭāriḥ} (pl. of \emph{maṭraḥ}) see \emph{Matrak}.}
\Entry{\emph{Matrak} = \emph{maṭraḥ}, place of the projection; \emph{matrak elshoa} = \emph{maṭraḥ al-shuʿāʿ}, place of the casting of the rays, 31 (N. 6), T. 91; see also \emph{Shoa}.}
\Entry{\emph{Maudho} = \emph{mauḍūʿ}\footnoteT{The print has «maudūʿ» here, without the dot under the \emph{d}.}, place, position; \emph{maudho Jauzegarhu} = \emph{mauḍūʿ djauzaharhu}, place of its (ascending) node, T. 34.}
\Entry{\emph{Medek} = \emph{mudaḳḳ}, finely crushed, divided into very small parts, minute, 7 (N. 20).}
\Entry{\emph{Medte} = \emph{madda} (or \emph{mudda}), extension, duration, T. 13 (N. 17).}
\Entry{\emph{Megil} see \emph{Meil}.}
\Entry{\emph{Meil} = \emph{mail}, inclination, the declination of every point of the ecliptic to the equator; 13, 14, 58, T. 21 (N. 4); \emph{al-mail al-kullī}, the whole inclination, that is, the obliquity of the ecliptic.}
\Entry{\emph{Mezen} = \emph{mīzān}, balance, the sign of the Balance in the zodiac, T. 59$^{b}$.}
\Entry{\emph{Mirech}, T. 92, see addition 5.}
\Entry{\emph{Mirrīkh}, the planet Mars, 236.}
\Entry{\Pg{245}\emph{Mizādj}, temperament, the astrological nature of a heavenly body; see addition 5.}
\Entry{\emph{Mobzota} = \emph{mabsūṭa} (fem. of \emph{mabsūṭ}), spread out, plane (\emph{planus}), T. 15 (N. 2).}
\Entry{\emph{Muadala} and \emph{Muaddala} = \emph{muʿaddala} (fem. of \emph{muʿaddal}), equated, corrected, 13 (N. 2), T. 36, 54.}
\Entry{\emph{Muʿāmalāt} (pl. of \emph{muʿāmala}), commercial or business computations, IX.}
\Entry{\emph{Muchre} = \emph{mu’akhkhara}, conclusion, end, here the eastern and the western end of the inhabited half of the earth, 9.}
\Entry{\emph{Muharram} = \emph{Muḥarram}, the first month of the Muslim year, XIV, XV, 1, 2, 34, 37, etc. T. 2, 3, etc.}
\Entry{\emph{Muhit} = \emph{muḥīṭ}, circumference, periphery, 18.}
\Entry{\emph{Muht} see \emph{Mukdh}.}
\Entry{\emph{Mukaam} = \emph{muḳām}, stay, station (of the planets), 12, 13 (N. 8, 12, 14), T. 54.}
\Entry{\emph{Muḳābala}, confrontation, balancing, VIII; see also \emph{Djebr}.}
\Entry{\emph{Mukatil} = \emph{muḳātil}, fighter, murderer, epithet of the planet Saturn, probably used only in Spain and North Africa, T. 27–29.}
\Entry{\emph{Mukdh} = \emph{mukth} (or \emph{makth}), staying, stay (\emph{mora}), the time during which an eclipse is total, 26, 27, 87.}
\Entry{\emph{Mult} see \emph{Muhit}.}
\Entry{\emph{Munegim} = \emph{munadjdjim}, astronomer, astrologer, XXII, 20 (N. 3).}
\Entry{\emph{Mustakim} = \emph{mustaḳīm}, straight, right; \emph{elfelek elmustakim} = \emph{al-falak al-mustaḳīm}, the right sphere, that is, the celestial sphere as it appears at the equator, 18, 20, 25, etc. T. 59–59$^{b}$.}
\Entry{\emph{Musteri} = \emph{Mushtarī}, the planet Jupiter, T. 11, 12, 33–38.}
\Entry{\emph{Mustewi} = \emph{mustawī}, level, straight, 69, 77; fem. \emph{mustewia} = \emph{mustawiya}, 21.}
\Entry{\emph{Muth} see \emph{Mukdh}.}
\par\begin{center}\textbf{N.}\end{center}
\Entry{\emph{Nadair} see the following.}
\Entry{\emph{Nadir} = \emph{naẓīr}, the point of the heavens opposite the zenith (\emph{samt al-ra’s}), then in general the opposite point of any point in the heavens, XX, 97, 98, 101.}
\Entry{\emph{Nāhīd} see \emph{Anāhīd}.}
\Entry{\emph{Nuzf} = \emph{nuṣf} and \emph{niṣf}, half (\emph{dimidium}), T. 73.}
\par\begin{center}\textbf{O.}\end{center}
\Entry{\emph{Otari[d]} = \emph{ʿUṭārid}, the planet Mercury, XXI, T. 17, 18, 51–56.}
\Entry{\emph{Ozari} see the preceding.}
\par\begin{center}\textbf{R.}\end{center}
\Entry{\emph{Rabb}, lord, master, 103.}
\Entry{\emph{Raba} see \emph{Rabia}.}
\Entry{\emph{Rabe} = pl. of \emph{Raba} or \emph{Rabia}; see this.}
\Entry{\emph{Rabe elachir} = \emph{Rabīʿ al-ākhir}, the fourth Muslim month, XIV, XV, 34, 84, 86, 89, T. 2, 3, etc.}
\Entry{\emph{Rabe elawel} = \emph{Rabīʿ al-awwal}, the third Muslim month, XIV, XV, 34, T. 2, 3, etc.}
\Entry{\emph{Rabia} = \emph{rābiʿ}, fem. \emph{rābiʿa}, (the) fourth; in computing with sexagesimal fractions the \emph{partes quartae}, that is, the multiples of $\tfrac{1}{60^{4}}$, 15 (N. 26).}
\Entry{\emph{Race} = \emph{ra’s}, head; \emph{race} is probably meant to be the ablative of \emph{rac}; \emph{Race elhamel} = \emph{ra’s al-ḥamal}, head, beginning (of the sign) of the Ram, 20.}
\Entry{\emph{Rageb} see \emph{Regeb}.}
\Entry{\emph{Ramadham} = \emph{Ramaḍān}, the ninth Muslim month, the month of fasting, XIV, XV, 34, 37, 67, 68, T. 2, 3, etc.}
\Entry{\emph{Raz} see \emph{Race}.}
\Entry{\emph{Regeb} = \emph{Radjab}, the seventh Muslim month, XIV, XV, 34, 38, 39, T. 2, 3, etc.}
\par\begin{center}\textbf{S.}\end{center}
\Entry{\emph{Saaben} see \emph{Shabin}.}
\Entry{\emph{Safar} see \emph{Zafar}.}
\Entry{\emph{Ṣāḥib}, lord, master of the house, owner, in astrology the \emph{dominus} of a sign, or of a part of a sign, 103.}
\Entry{\emph{Sana kabīsa}, leap year, 34.}
\Entry{\emph{Scabin} see \emph{Shabin}.}
\Entry{\emph{Scawal} see \emph{Shawel}.}
\Entry{\emph{Sceir} = \emph{sair}, going, course, path, XXI.}
\Entry{\emph{Sewil} see \emph{Shawel}.}
\Entry{\emph{Shabin} = \emph{Shaʿbān}, the eighth Muslim month, XIV, XV, 34, 67, 68, T. 2, 3, etc.}
\Entry{\emph{Sharaf} (pl. \emph{ashrāf}), height, elevation, distinction; in medieval astrology \emph{exaltatio}, XXI, 104.}
\Entry{\emph{Shawel} (also \emph{Shavvel}) = \emph{Shawwāl}, the tenth Muslim month, XIV, XV, 34, 39, T. 2, 3, etc.}
\Entry{\Pg{246}\emph{Shemeli} = \emph{shamālī}, northern, 13, 16, 18, 29 (N. 5, 12, 13).}
\Entry{\emph{Shemelia} = \emph{shamālīya}, fem. of the preceding, 29 (N. 5, 13).}
\Entry{\emph{Shems} = \emph{Shams}, sun, 58, T. 21–26.}
\Entry{\emph{Shoa} = \emph{shuʿāʿ} (pl. \emph{shuʿuʿ} and \emph{shuʿāʿāt}), ray of a heavenly body, for example of the sun, 31 (N. 6), T. 91; \emph{matrak elshoa} see \emph{Matrak}.}
\Entry{\emph{Shohor} = \emph{shuhūr}, pl. of \emph{shahr}, month, T. 9; see also \emph{filshohor} (under \emph{fi}).}
\Entry{\emph{Ṣifr}, empty, zero, 35.}
\Entry{\emph{Sohor} see \emph{Shohor}.}
\Entry{\emph{Ṣūra}, form, figure; \emph{ṣūrat al-arḍ}, picture (figure) of the earth, IX.}
\Entry{\emph{Szukuf} see \emph{Zukut}.}
\Entry{\emph{Szukut} see \emph{Zukut}.}
\par\begin{center}\textbf{T.}\end{center}
\Entry{\emph{Ṭabaḳāt} (pl. of \emph{ṭabaḳa}), classes; \emph{ṭabaḳāt al-umam}, the classes of the nations, IX; see also \emph{Umam}.}
\Entry{\emph{Tabeb} = \emph{ṭabīb}, scholar, physician, XXII, 20.}
\Entry{\emph{Tadil} = \emph{taʿdīl}, equation (of the sun, the moon and the planets), rendered in the Latin text by \emph{aequatio}, \emph{coaequatio} and \emph{examinatio}, 9, 10, 46 etc. T. 21–56.}
\Entry{\emph{Tarbea} = \emph{tarbīʿ}, quartile or square aspect (astrol.) T. 91.}
\Entry{\emph{Ta’rīkh}, reckoning of time, chronicle, date, IX; \emph{Ta’rīkh al-ṣufr}, the Spanish era, 35; see also \emph{Acofra}.}
\Entry{\emph{Tedewir} = \emph{tadwīr}, turning (in a circle); \emph{felek eltedewir} = \emph{falak al-tadwīr}, the epicycle, T. 9 (N. 20).}
\Entry{\emph{Tethlith} = \emph{tathlīth}, trine or trigonal aspect (astrol.) T. 91.}
\Entry{\emph{Tezdiz} = \emph{tasdīs}, sextile or hexagonal aspect (astrol.) T. 91.}
\Entry{\emph{Tezweit} see \emph{Tezwiet}.}
\Entry{\emph{Tezwiet} = \emph{taswiya}, equalizing, arrangement, distribution, 25, 96, T. 79.}
\Entry{\emph{Thaur}, bull, the sign of the Bull (\emph{Taurus}) in the zodiac, T. 59$^{a}$.}
\Entry{\emph{Thedewir} see \emph{Tedewir}.}
\Entry{\emph{Thelide} and \emph{Thelidhe} = pl. of \emph{Thelidha}; see the following.}
\Entry{\emph{Thelitha} = \emph{thālitha}, fem. of \emph{thālith}, (the) third; in computing with sexagesimal fractions the \emph{partes tertiae}, that is, the multiples of $\tfrac{1}{60^{3}}$, 15 (N.).}
\Entry{\emph{Theni} = \emph{thānī}, (the) second, 12.}
\Entry{\emph{Thenia} = \emph{thāniya}, fem. of the preceding, second (of arc or time); in computing with sexagesimal fractions the \emph{partes secundae}, that is, the multiples of $\tfrac{1}{60^{2}}$, XXIII, XXV, 11, 14, 15, 36, 69, T. 4, 5, etc.}
\Entry{\emph{Thuid} see \emph{Thuld}.}
\Entry{\emph{Thuld} = \emph{thulth}, a third, 9.}
\Entry{\emph{Thur} see \emph{Thaur}.}
\par\begin{center}\textbf{U.}\end{center}
\Entry{\emph{Udjain} see \emph{Arin}.}
\Entry{\emph{Ule} = \emph{ūlā}, fem. of \emph{awwal}, (the) first, T. 2, 3, etc.}
\Entry{\emph{Umam} (pl. of \emph{umma}), peoples, nations.}
\Entry{\emph{Uzain}, the Arabic form of the Indian \emph{Udjain}; see \emph{Arin}.}
\par\begin{center}\textbf{W.}\end{center}
\Entry{\emph{Wa}, Arabic preposition, and; for example \emph{wakhod} = \emph{wakhudh}, and take, XXII, 27 (N. 18); \emph{wahezatu} = \emph{wakhāṣṣathu}, and its anomaly, T. 6–8; see also \emph{Heza}.}
\Entry{\emph{Wabāl}, detriment, 161.\footnoteT{So the print; \emph{wabāl} is explained on p. 103, and p. 161 is Table 50.}}
\Entry{\emph{Wacat} = \emph{wasaṭ} (or \emph{wasṭ}), middle, mean, mean motion (of the sun, the moon and the planets), rendered in the text and in the tables mostly by \emph{medialitas}, 7–12 etc., 41, T. 4–20, etc.; \emph{wacat alkamar} = \emph{wasaṭ al-ḳamar}, the mean course (the mean motion) of the moon, T. 6–8.}
\Entry{\emph{Wadjh}, pl. \emph{wudjūh}, front side, face, manner, etc.; in astrology the third part (10°) of a sign of the zodiac (Latin \emph{decanus}, also \emph{facies}) 104.}
\Entry{\emph{Wakhod mebibelche} XXII, 27 (N. 18).}
\Entry{\emph{Wazat} see \emph{Wacat}.}
\par\begin{center}\textbf{Y.}\end{center}
\Entry{\emph{Yle} see \emph{Ule}.}
\par\begin{center}\textbf{Z.}\end{center}
\Entry{\emph{Zafar} = \emph{Ṣafar}, the second Muslim month, XIV, XV, 2, 34, 42, 43, 45 etc. T. 2, 3, 4 etc.}
\Entry{\emph{Zahat} see \emph{Caat}.}
\Entry{\emph{Zare} see \emph{Care}.}
\Entry{\emph{Zenia} see \emph{Thenia}.}
\Entry{\emph{Zich} = \emph{zīdj} (Persian \emph{zīg}), a work of astronomical tables, VIII, 1, 32, etc. T. 116; \emph{zīdj al-shāh} or \emph{al-shahriyār}, the tables of the Shāh (of the king) = the Persian tables, 32, \Pg{247}33, 52, etc.; see also the index of names and subjects, under “Tafeln” (tables).}
\Entry{\emph{Zīg} see \emph{Zich}.}
\Entry{\emph{Ẓill}, shadow, the trigonometric function of the \emph{tangent} and the \emph{cotangent}, 77; see also \emph{Dhel}.}
\Entry{\emph{Zohal} = \emph{Zuḥal}, the planet Saturn, T. 9, 10, 27–32.}
\Entry{\emph{Zohara} = \emph{Zuhara}, the planet Venus, XXII, T. 45, 47, 49, 50.}
\Entry{\emph{Zohari} = \emph{zuharī}, belonging to Venus, T. 15, 16.}
\Entry{\emph{Zukut} = \emph{suḳūṭ}, (the) falling, falling in; in astronomy the phase of an eclipse from its beginning to the totality, rendered in this work by \emph{casus}, elsewhere by \emph{immersio}, 26, 27, 87, T. 73, 78.}
\Entry{\emph{Zukuz} see \emph{Zukut}.}
\par\begin{center}II. Egyptian (Coptic) names of the months (Tab. 2$^{b}$).\end{center}
\par (The Coptic names of the months are taken in these tables from the Upper Egyptian (Theban) dialect; in parentheses I have also added the forms of the Lower Egyptian (Memphitic) dialect.)\par
\begin{center}\small\begin{tabular}{rlll}
1. & Tut = thowt (thōwt) & Greek & {\greekfont θώθ} \\
2. & Beba = paape (paopi) & — & {\greekfont φαωφί} \\
3. & Hettur = hatōr (athōr) & — & {\greekfont ἀθύρ} \\
4. & Keihec = kiakh (khoiak) & — & {\greekfont χοιάκ} \\
5. & Tuba = tōbe (tōbi) & — & {\greekfont τυβί} \\
6. & Amushir = $^{e}$mshir (mekhir) & — & {\greekfont μεχίρ} \\
7. & Burimihez = parmhāt (phamenōth) & — & {\greekfont φαμενώθ} \\
8. & Burmuda = parmūte (pharmūthi) & — & {\greekfont φαρμουθί} \\
9. & Besenz = pashons (pakhōn) & — & {\greekfont παχών} \\
10. & Buba = paōne (paōni) & — & {\greekfont παϋνί} \\
11. & Abib = epep (epēp) & — & {\greekfont ἐπιφί} \\
12. & Muzre = mesōrē (mesōrē) & — & {\greekfont μεσορή} \\
\end{tabular}\end{center}
\par\begin{center}III. Persian names of the months (Tab. 2$^{a}$).\end{center}
\par\begin{center}(\emph{māh} = month)\end{center}
\begin{center}\small\begin{tabular}{rll}
1. & Affrordinme & = Farwardīn-māh. \\
2. & Azdihestme & = Ardībihist-māh.\footnotemarkT{} \\
3. & Chordezme & = Khurdādh-māh. \\
4. & Tiierme & = Tīr-māh. \\
5. & Merdezme & = Murdādh-māh. \\
6. & Shaharirme & = Shahrīwar-māh. \\
7. & Meherime & = Mihr-māh. \\
8. & Abenme & = Ābān-māh. \\
9. & Adarme & = Ādhar-māh. \\
10. & Dime & = Dai-māh. \\
11. & Bahamanme & = Bahman-māh. \\
12. & Izfindarme & = Isfandārmadh-māh. \\
\end{tabular}\end{center}
\addtocounter{footnoteT}{-1}\stepcounter{footnoteT}\footnotetextT{So the print; the name of the second Persian month is Ardībihisht.}
\par\begin{center}\Pg{248}IV. Syrian (Seleucid) names of the months (Tab. 3$^{a}$).\end{center}
\par\begin{center}(\emph{Nomina mensium Romanorum}, the names of the months of the Romans)\end{center}
\begin{center}\small\begin{tabular}{rllll}
1. & Tisrin elawel & = Teshrīn q$^{e}$dīm (I.) & = October & (of the Julian year). \\
2. & Tisrin elachir & = Teshrīn ḥ$^{e}$rai (II.) & = November & ( — — ). \\
3. & Keinun elawel & = Kānūn q$^{e}$dīm (I.) & = December & ( — — ). \\
4. & Keinun elachir & = Kānūn ḥ$^{e}$rai (II.) & = Januarius & ( — — ). \\
5. & Sabat & = Sh$^{e}$bāṭ & = Februarius & ( — — ). \\
6. & Adar & = Ādār & = Martius & ( — — ). \\
7. & Nizan & = Nīsān & = Aprilis & ( — — ). \\
8. & Iiar (Iyar) & = Iyār & = Maius & ( — — ). \\
9. & Heziren & = Ḥ$^{e}$zīrān & = Junius & ( — — ). \\
10. & Temuz & = Tāmūz & = Julius & ( — — ). \\
11. & Ab & = Āb & = Augustus & ( — — ). \\
12. & Elul & = Īlūl & = September & ( — — ). \\
\end{tabular}\end{center}
\UnitHead{Index of names and subjects\footnoteT{The entries keep the order of the German index. German headwords are translated, with the German in square brackets; the Latin, Greek and Arabic headwords and the names stand as printed (in the rest of this translation Athelhard is Adelard and Björnbo Bjørnbo). The numbers are Suter's pages, given in the margin of this translation; T. is a table, N. a footnote.}}{Index of names and subjects}{I01-0}\Pg{249}
\par\begin{center}\textbf{A.}\end{center}
\Entry{Abacus X.}
\Entry{ʿAbbās b. Saʿīd al-Djauharī 32.}
\Entry{Abraham b. ʿEzrā 41, 51, 99.}
\Entry{Abū’l-Ḥasan ʿAlī of Morocco 71.}
\Entry{Abū Isḥāḳ al-Naḳḳāsh see al-Zarḳālī.}
\Entry{Abū Naṣr Manṣūr see Manṣūr b. ʿAlī b. ʿIrāḳ.}
\Entry{Abū’l-Wefā 77, 78.}
\Entry{Adam 35, T. 1.}
\Entry{Adelhard see Athelhard.}
\Entry{Adiectiones XVIII.}
\Entry{Equant [Aequans] 47, 48, 50, 233.}
\Entry{Aequatio see Examinatio.}
\Entry{— domorum, see Houses, astrological [Häuser, astrolog.].}
\Entry{— nychthemerōn 80.}
\Entry{Equinoctial hours (equal) [Aequatorial-Stunden] 24, 76, 81, 84.}
\Entry{Aera Arabum (or of the \emph{Hidjra}) VII, IX, 1–3, 33–39, 106; — Augusti (or Alexandrina) XVI; — Coptorum XV; — Diocletiani XV; — Domini (Christian, Julian) XVII, 3, 4, 35, 36, 102; — Nabonassaris XVI; — Philippi Aridaei XVI; — Romanorum (Seleucid) 2–4, 6, 7, 34–40, 81; — Spanish [spanische] (\emph{acofra}) XVIII, 3, 35; — of Yezdegird VII, VIII, XV, 33, 37, 106, T. 1.}
\Entry{Aḥmed b. al-Muthannā 51.}
\Entry{Albohazen 105.}
\Entry{Albumasar 105.}
\Entry{Alcabitius see Ḳabīṣī.}
\Entry{Algebra (al-Khwārizmī's) VIII, XII.}
\Entry{Alienatio 103, 104, T. 116.}
\Entry{Alexander, the Great, 5\footnoteT{So the print; Alexander is named on p. 3 of the text, not on p. 5.}, 34, 35, 38, T. 1.}
\Entry{Almagest, the, of Ptolemy XI, 2, 16, 18, 33, 44, 70, etc.}
\Entry{Altitudo see Altitude [Höhe].}
\Entry{Ammonius see Emirethius.}
\Entry{{\greekfont ἀναπλήρωσις} = the making full again (of the moon, or of the sun in eclipses) 88.}
\Entry{Andersen, D. VI.}
\Entry{Anni collecti XV, XVII, 4–7, 23, 36, etc. T. 2, 3, 4, 6, etc.}
\Entry{— plani (vel expansi, singuli, residui) XV, 4–7, 23, 36, etc. T. 2, 3, 4, 6, etc.}
\Entry{— vagi (of the Egyptians) 38.}
\Entry{Annus Arabum, description of it, 2, 34; its conversion into Roman years, 6–7, 38–40, T. 3, 3$^{a}$.}
\Entry{— bissextus or bissextilis (leap year) XVII, XVIII, 6–7, 34, 35, T. 3$^{a}$.}
\Entry{— Persarum, see Aera, of Yezdegird.}
\Entry{— Romanorum vel Alexandri, description of it, 2–4, 34–36; finding it, 6–7, 38–40, T. 1, 3.}
\Entry{Anomalia see Argumentum.}
\Entry{Antioch [Antiochia] 83 (N. 2), 95, 96.}
\Entry{Apogeum see Sublimatio.}
\Entry{Line of apsides [Apsidenlinie] 65.}
\Entry{Argumentum (Anomalia, \emph{Heza}) XIV, 9–12, 43–46, 48–55, 67, 82–85, 93, T. 6–9, 11, 13, 15–18, 27–56, 69–72.}
\Entry{— (examinatum, or aequatum) 11–13, 14, 52–57, 62–64.}
\Entry{— (the quantity with which one enters a table) 19, 20, 55, 92.}
\Entry{— (latitudinis lunae) 23, 24, 26, 27, 29, 58, 59, 68, 82, 84–86, 94, T. 69–76, 78.}
\Entry{— (procedure, logical inference) 13, 56.}
\Entry{— (ratum) 25, 84.}
\Entry{Arin VII, XVIII, 1, 8, 9, 17–19, 33, 41, 43, 68, 72.}
\Entry{Arkand, the book al-, see Khaṇḍakhādyaka.}
\Entry{Āryabhaṭṭa (or Āryabhaṭa) 42.}
\Entry{Ārya-Siddhānta see Siddhāntas.}
\Entry{Ascendens see Horoscopus.}
\Entry{Ascensional difference [Ascensional-Differenz] 74–76, 92.}
\Entry{Ascensiones obliquae XII, 20, 21, 73–76, 92, 93, 100–102.}
\Entry{\Pg{250}Ascensiones rectae (= right ascension) 17–20, 72–76, 93, 99, 100, 102.}
\Entry{Asín, Miguel, XI, XVII.}
\Entry{Astrolabe [Astrolabium] IX, X, XX, 18, 71.}
\Entry{Ascendant [Aszendent] see Ascendens.}
\Entry{Athelhard of Bath VI–X, XII, XIII, XXIII (N. 2), 1, 37, 69, 71, 80, 96, 236, T. 116.}
\Entry{Aug see Sublimatio.}
\Entry{Augustus (emperor) XVI, XVII.}
\Entry{Azin see Arin.}
\par\begin{center}\textbf{B.}\end{center}
\Entry{Baghdād 74.}
\Entry{al-Battānī, Muḥ. b. Djābir b. Sinān (and his astronomical tables, ed. Nallino) IX, XIII, XV, XVI, XVIII, XXI, XXIII, XXIV, 32, 34, 35, 38–50, 52, 54–60, 62–66, 69, 71–97, 99, 100, 102–106, 233, 234, 236.}
\Entry{Bernelinus X.}
\Entry{Bernhelmus see the preceding.}
\Entry{Besthorn, R. O., V, XXV.}
\Entry{Motion, mean (of the sun) [Bewegung, mittlere] XIV, XVIII, 7–9, 23, 33, 36, 41–44, 65, 67, 82, 84, 85, 105, 106, 235, T. 4, 5.}
\Entry{— mean (of the moon) XIV, XVIII, 7, 8, 10, 23, 33, 36, 41–43, 45, 46, 65, 67, 82, 84, 85, 105, 106, T. 6, 7, 8.}
\Entry{— mean (of the upper planets) XIV, XVIII, 7, 8, 10–12, 36, 41–44, 51–54, 63, 64, 105, 106, T. 9–14.}
\Entry{— mean (of the dragon's head) XIV, XIX, 7, 8, 12, T. 19, 20.}
\Entry{— mean (of the anomaly of the lower planets) 54, T. 15–18.}
\Entry{— true (that is, unequal, of the moon) 22, 23, 28, 78, 79, 81, 83, 105, T. 61–66.}
\Entry{— true (of the sun) 22, 23, 78, 79, 81, 83, 105, T. 61–66.}
\Entry{Bhāskara, 70.}
\Entry{Billeter, G. VI.}
\Entry{al-Bīrūnī, Muḥ. b. Aḥmed, 37, 38, 44, 49, 51, 66, 79, 89, 90, 105, 235.}
\Entry{Björnbo, A. V, VII, X, XI, XIII, XXIV, 68, 69–71, 76, 77, 102.}
\Entry{Boll, F. 105.}
\Entry{Boncompagni, B. VIII.}
\Entry{Bouché-Leclercq, A. 105.}
\Entry{Brahmagupta 32, 42, 45, 79, 90.}
\Entry{von Braunmühl, A. 71.}
\Entry{Latitude, of the moon [Breite] 14, 27, 29, 58, 59, 66–68, 86, 88, 94, T. 21–26.}
\Entry{— of the planets, 14, 60–64, T. 27–56.\footnoteT{The print has «S. 27—56»; the latitudes of the planets stand in Tables 27–56.}}
\Entry{Latitude, apparent, of the moon (produced by the parallax) [Breite, scheinbare] 94.}
\Entry{— geographical, of a place (its determination) 18, 28, 71.}
\Entry{Brevis circulus see Epicycle [Epicyklus].}
\Entry{Broger, M. 95 (N. 1).}
\Entry{Bukhtanaṣṣar see Nabonassar.}
\Entry{Burgess, E. 42, 45, 46, 51, 52, 58–61, 68, 75, 78–80, 86, 90, 92.}
\par\begin{center}\textbf{C.}\end{center}
\Entry{Caldei 9, 11, 43.}
\Entry{Caput draconis see Nodes (of the moon) [Knotenpunkte].}
\Entry{Carra de Vaux 78.}
\Entry{Casae see Houses [Häuser].}
\Entry{Casus 26, 27, 87–89, 95, see also Dejection [Tiefstand].}
\Entry{Cauda draconis see Nodes (of the moon) [Knotenpunkte].}
\Entry{Celada (cardinal) XI.}
\Entry{Centrum (astron.) 11, 48–54, 57, 58, see also Examinatio centri.}
\Entry{— examinatum vel aequatum 11, 12, 52–55, 57.}
\Entry{— ultimum 11, 53, 54.}
\Entry{Chaldei see Caldei.}
\Entry{Chasles, M. VII, X, 43, 49, 50.}
\Entry{Cheikho, L. IX (N. 3).}
\Entry{Cifra see Zero [Null].}
\Entry{Circulus aequinoctialis 18–21, 25.}
\Entry{— excentricus 26, 44, 47.}
\Entry{Circumpolar stars [Circumpolarsterne] 71.}
\Entry{Coaequatio see Examinatio.}
\Entry{Computatio XIV, 59, T. 21–26.}
\Entry{Computus (Gerlandi) X.}
\Entry{Conversio annorum see Conversion [Umwandlung].}
\Entry{Corduba (Cordova) VII, XII, 25, 33, 81, 83–85, 92–97, 100, 106.}
\Entry{Cotangent [Cotangens] see Umbra.}
\Entry{Curtze, M. XXIII.}
\par\begin{center}\textbf{D.}\end{center}
\Entry{Darius Medus 105.}
\Entry{Decani see Facies.}
\Entry{Defectus see Eclipse [Verfinsterung].}
\Entry{Deferent [Deferens] 16, 17\footnoteT{So the print; the deferent is introduced on pp. 47–48, the 16th and 17th pages of the commentary; pp. 16–17, of the Latin text, do not name it.}, 57, 60, 61, 233, 234.}
\Entry{Deiectio see Dejection [Tiefstand].}
\Entry{Declination (obliquatio) of the sun [Deklination] 13, 18, 19, 58, 91, T. 21–26.}
\Entry{— of any point of the ecliptic 74, 75, see also Ecliptic [Ekliptik].}
\Entry{Delambre 60.}
\Entry{Detrimentum see Alienatio.}
\Entry{Dieterici, F. 105.}
\Entry{Digiti = parts of the gnomon 20–22, T. 60.}
\Entry{\Pg{251}Digiti = parts of the diameter of the moon or of the sun 26, 27, 87–89, 95.}
\Entry{Dignitates see Exaltations [Erhöhungen].}
\Entry{Diluvium 3, 35, T. 1.}
\Entry{Diocletianus (emperor) XV, T. 1.}
\Entry{Dispositio (casarum duodecim) see Houses, astrological [Häuser, astrol.].}
\Entry{Distantia definita 16, 68.}
\Entry{Distance, geocentric, of the centre of the epicycle [Distanz] 48, 49.}
\Entry{Diversitas aspectus vel aspectuum see Parallax [Parallaxe].}
\Entry{— respectus, see the same.}
\Entry{Djābir b. Aflaḥ 69.}
\Entry{Domicilia 103, 106, T. 116.}
\Entry{Domini 31, 103, 104, T. 116.}
\Entry{Dorotheus 99, 100.}
\Entry{Dozy, R. 36.}
\Entry{Dragon's head [Drachenkopf] see Nodes (of the moon) [Knotenpunkte].}
\Entry{Dragon's tail [Drachenschwanz] see the same.}
\Entry{Dracones see the same.}
\Entry{Diameter (rota), apparent (of the moon) [Durchmesser] 23, 78–79, 90, T. 61–66.}
\Entry{— apparent (of the sun) 22–23, 78–79, T. 61–66.}
\Entry{— of the earth's shadow 90, 91, T. 61–66.}
\par\begin{center}\textbf{E.}\end{center}
\Entry{Eclipsis see Eclipse [Verfinsterung].}
\Entry{Ecliptic [Ekliptik] 44, 47, 58–61, 63, 72–74, 76, 96–101; its obliquity 19, 33, 58, 72, 105.}
\Entry{Ecliptic zenith distance [Ekliptik-Zenitdistanz] 92, 93.}
\Entry{Emersio see Respiratio.}
\Entry{Emirethius (?), perhaps Eumathius (probably = Ammonius), whose canon (astronomical tables) is said to have been improved by al-Zarḳālī (compare addition 5) T. 7, 92.}
\Entry{{\greekfont ἔμπτωσις} = the falling in 87.}
\Entry{Epicycle [Epicyclus] XX, XXI, 44, 47–49, 51, 57, 60, 61, 234; its radius for the moon 46, for the planets 48.}
\Entry{Era see Aera.}
\Entry{Earth's shadow [Erdschatten] see Umbra terrae.}
\Entry{Exaltations [Erhöhungen] 103, 104, 106, T. 116.}
\Entry{Eumathius see Emirethius.}
\Entry{Evectio 45.}
\Entry{Exagonum see Sextile aspect [Sextilschein].}
\Entry{Exaltatio see Exaltation [Erhöhung].}
\Entry{Examen see the following.}
\Entry{Examinatio (equation) of the sun, the moon and the planets XIV, 9, 10, 24, 43–46, 82, 83, T. 21–26.}
\Entry{— argumenti (= anomaliae, \emph{elheza}) 10, 11, 48–54, 233–235, T. 27–56.}
\Entry{Examinatio centri 11, 12, 48–54, 233, 234, T. 27–56.}
\Entry{Eccentricity, of the orbits of the planets [Exzentrizität] 47, 48, 234.}
\Entry{— of the orbit of the sun 45.}
\par\begin{center}\textbf{F.}\end{center}
\Entry{Facies (= decanus), the third part (10°) of a sign of the zodiac, 16, 68, 103, 104, T. 116.}
\Entry{Fall see Dejection [Tiefstand].}
\Entry{Feria = day of the week 1, 5–7, 36, 37.}
\Entry{Fines see Bounds [Grenzen].}
\Entry{Eclipses [Finsternisse] see Eclipse [Verfinsterung].}
\Entry{Fleischer, H. L. 105.}
\par\begin{center}\textbf{G.}\end{center}
\Entry{Opposite aspect [Gegenschein] 99.}
\Entry{Geib see Sinus.}
\Entry{Gerbert (Pope Sylvester II) X.}
\Entry{Gerhard of Cremona VIII, 69.}
\Entry{Gerlandus X.}
\Entry{Ghoarizmi see Khwārizmī.}
\Entry{Equation of the anomaly and of the centre (in Battānī) [Gleichung] 48, 50.}
\Entry{— of the apogee 50.}
\Entry{— of the apsis 51.}
\Entry{— of the conjunction 51.}
\Entry{Equations (of the sun, of the moon, of the anomaly, of the centre) [Gleichungen] see Examinatio.}
\Entry{— (of the planets) 33, 49, 55, 105; see also Examinatio (argumenti and centri).}
\Entry{Gnomon 72, 74–76.}
\Entry{Goebel, J. XXII (N. 1).}
\Entry{Gonzalez XVII (N.).}
\Entry{Greenwich 83, 85.}
\Entry{Bounds (unequal parts of the signs of the zodiac) [Grenzen] 104, 105, T. 116.}
\Entry{Reasons of the tables of Khwārizmī [Gründe der Tafeln des Khwārizmī] 41, 51.}
\par\begin{center}\textbf{H.}\end{center}
\Entry{Ḥabash, the Calculator, 66, 76, 77.}
\Entry{Halving of the equation (of the planets) [Halbierung der Gleichung] 49–52.}
\Entry{Ḥarrān 74.}
\Entry{Haskins, C. H. VII, X, 236.}
\Entry{Houses, the 12 astrological [Häuser], 25; their distribution (aequatio, dispositio) 29, 30, 96–98, 106, T. 79–90.}
\Entry{Heiberg, J. L. XIII, XXIV, XXV, 44, 57, 74, 81, 90, 233.}
\Entry{Henry, C. XXIII.}
\Entry{Hermes Trismegistos XXII.}
\Entry{Hermannus Contractus X.}
\Entry{— Dalmata XIII, XVIII (N.).}
\Entry{\Pg{252}Heza see Argumentum.}
\Entry{Altitude of the sun (artifa solis) [Höhe] 21–22, 76, 77.}
\Entry{Horae aequales see Equinoctial hours [Aequatorial-Stunden]; — inaequales see Temporal hours [Zeitstunden].}
\Entry{— diurnae see Day hours [Tagesstunden].}
\Entry{— nocturnae see Night hours [Nachtstunden].}
\Entry{Horoscopi signorum (= the risings of the signs of the zodiac) see Signum, and the glossary, under Matale.}
\Entry{Horoscopus (Ascendens), the point of the ecliptic that rises at a given time, 25, 28–31, 91–94, 97, 98, 100, 101, T. 79.}
\Entry{Houtsma, M. Th. XXV.}
\par\begin{center}\textbf{I.}\end{center}
\Entry{Ibn Abī Uṣaibiʿa VII.}
\Entry{— al-Adamī 33, 49.}
\Entry{— ʿEzrā see Abraham b. ʿEzrā.}
\Entry{— Hibintā 43, 46, 99, 100, 106.}
\Entry{— al-Ḳifṭī VII, 33.}
\Entry{— Ṣāʿid see Ḳāḍī Ṣāʿid.}
\Entry{— Yūnus 47.}
\Entry{Ibrāhīm al-Fazārī VIII, 32, 33.}
\Entry{Ideler, L. XVIII, 36.}
\Entry{Immersio see Casus.}
\Entry{Intercalation (of the leap day) XV, 34.}
\Entry{Interstitium 24, 29, 93.}
\Entry{Isidore of Seville [Isidor von Sevilla] XVIII.}
\par\begin{center}\textbf{J.}\end{center}
\Entry{Year, Arab [Jahr], see Annus Arabum.}
\Entry{— Roman (Seleucid) see Annus Romanorum.}
\Entry{— sidereal (star year) 42, 65, 66, 102, 105, 235.}
\Entry{— tropical 42, 66, 102.}
\Entry{Johannes de Sacrobosco XXIII (N. 1).}
\Entry{— of Seville XI.}
\Entry{Jordanus Nemorarius XI.}
\Entry{Judices (= Domini) 104, T. 116.}
\par\begin{center}\textbf{K.}\end{center}
\Entry{al-Ḳabīṣī 99, 105.}
\Entry{Ḳāḍī Ṣāʿid VII, IX, 33.}
\Entry{Kardadja, its division 70.}
\Entry{Karpinski, L. C. XII.}
\Entry{Kaurezmi see Khwārizmī.}
\Entry{Khaṇḍakhādyaka 79, 90.}
\Entry{al-Khwārizmī, Muḥ. b. Mūsā VII–X, XII, XIII, XV, XXI, XXII, XXIV, 1, 32–35, 38–42, 45–47, 49–60, 62–66, 68–75, 77–81, 86, 87, 90–97, 99–103, 105, 106, 233–236, T. 116.}
\Entry{al-Kindī 105.}
\Entry{Line of the nodes, of the moon [Knotenlinie] 65.}
\Entry{Nodes, of the moon [Knotenpunkte] 12, 14, 15, 27, 41, 55, 65, 67, 68.}
\Entry{— of the planets 15, 64.}
\Entry{Conjunction and opposition (syzygies) of the sun and the moon [Konjunktion und Opposition] XIX, 23–29, 33, 81–86, 89, 92, 94–96, 106, T. 69–72.}
\Entry{— apparent, of the sun and the moon 94, 95.}
\Entry{Circle, eccentric [Kreis, exzentrischer] see Circulus excentricus.}
\par\begin{center}\textbf{L.}\end{center}
\Entry{Longitude, mean [Länge, mittlere] see Motion, mean [Bewegung, mittlere].}
\Entry{Longitude, true, of the moon [Länge, wahre] 10, 14, 16, 33, 45, 46, 58, 59, 66–68, 81–86, 93, T. 6–8, 21–26.}
\Entry{— — of the sun 9, 10, 16, 33, 43–45, 66, 67, 81–86, T. 4, 5, 21–26.}
\Entry{— — of the upper planets 10, 11, 33, 46–54, 60, 62, T. 9–14, 27–44.}
\Entry{— — of the lower planets 12, 33, 54, 55, 60, 62, T. 15–18, 45–56.}
\Entry{— — of the dragon's head 12, 55, 58, 59, 67, 68, T. 19, 20.}
\Entry{Latitudo lunae, etc., see Latitude [Breite].}
\Entry{— regionum (habitationum) see Latitude, geographical [Breite, geogr.].}
\Entry{Liber Embadorum X.}
\Entry{Lippert, J. 33.}
\Entry{London, meridian of, XVIII.}
\Entry{Longinquitas maxima lunae see Spatium longinquum.}
\Entry{— minima lunae see Spatium propinquum.}
\Entry{Loth, O. 105.}
\par\begin{center}\textbf{M.}\end{center}
\Entry{al-Madjrīṭī see Maslama.}
\Entry{al-Ma’mūn, caliph VIII, 32, 43, 47, 106.}
\Entry{Manitius, K. 44, 57, 74, 81, 90, 233.}
\Entry{al-Manṣūr, caliph 70.}
\Entry{— b. ʿAlī b. ʿIrāḳ 49.}
\Entry{Maslama b. Aḥmed al-Madjrīṭī VII–X, XIII, XV, XXI, XXII, 32–35, 39, 40, 77, 81, 86, 90, 100, 102, 103, 105, 106.}
\Entry{Medialitas see Motion, mean [Bewegung, mittlere].}
\Entry{Mediocritas see the same.}
\Entry{Medium coelum 25, 28, 29, 91–94.}
\Entry{Medius cursus see Motion, mean [Bewegung, mittlere].}
\Entry{Megistophiel XXII.}
\Entry{Mecca [Mekka] 74.}
\Entry{Menelaus, theorem of, 74.}
\Entry{Mephistopheles XXII.}
\Entry{Meridian distance [Meridian-Distanz] 92.}
\Entry{\Pg{253}Month, anomalistic [Monat] 64.}
\Entry{Months, Arab [Monate], their beginning 2, 5, 36, 37, T. 2; their names XIV, XV, 2, 34, T. 2, 3, 4, etc. and the glossary.}
\Entry{— Egyptian or Coptic 37, 38, T. 2$^{b}$; their names T. 2$^{b}$ and the glossary.}
\Entry{— Persian 37, T. 2$^{a}$; their names T. 2$^{a}$ and the glossary.}
\Entry{— Roman (Seleucid), their beginning 6, 7, 38–40, T. 3$^{a}$; their names T. 3$^{a}$ and the glossary.}
\Entry{Moon [Mond] see Motion [Bewegung], Latitude [Breite], Longitude [Länge], Sublimatio, Diameter [Durchmesser], Examinatio, Eclipse [Verfinsterung]; its appearing see New moon [Neumond].}
\Entry{Orbit of the moon, its obliquity to the ecliptic [Mondbahn] 59.}
\Entry{Nodes of the moon, their longitude [Mondknoten], see Longitude of the dragon's head [Länge].}
\Entry{{\greekfont μονή} = the remaining, staying, see Mora.}
\Entry{Mora 26, 27, 87–89.}
\Entry{Motus verus vel inaequalis see Motion, true [Bewegung, wahre].}
\Entry{Muḥammed b. Mūsā al-Khwārizmī see al-Khwārizmī.}
\Entry{Muḥammed b. Mūsā b. Shākir VIII.}
\Entry{Multiplication table for sexagesimal parts [Multiplikationstafel] 68, T. 57$^{b}$.}
\par\begin{center}\textbf{N.}\end{center}
\Entry{Nabonassar XVI, 35, T. 1.}
\Entry{Nabukodonossor see the preceding.}
\Entry{Night arc, of the ascendant [Nachtbogen] 97.}
\Entry{Night hours [Nachtstunden] 20, 25, 97, 98.}
\Entry{Nallino, C. A. IX, XIII, XV, XVI, XVIII (N.), XXI–XXIII (N.), XXIV, 32, 34, 38, 43–48, 54–58, 60, 63, 65, 66, 71–73, 76–78, 80–82, 86, 87, 90–92, 94–96, 99, 100, 102, 103, 105, 233, 236; see also al-Battānī.}
\Entry{New moon, its appearing [Neumond], 16, 17, 66–68, T. 57$^{a}$.}
\Entry{Notae, or Notulae (also Signa) 5, 34, 37, 39, 82.}
\Entry{Zero (\emph{cifra}) [Null] 8, 10, 36; its sign XXIII, 8.}
\par\begin{center}\textbf{O.}\end{center}
\Entry{Obliquatio see Declination [Deklination].}
\Entry{Ocreat, N. XXIII (N. 2).}
\Entry{Ophiel XXII.}
\Entry{Oppolzer, Th. 89, 95.}
\Entry{Opposition see Conjunction [Konjunktion].}
\Entry{Place of the casting of the rays [Ort der Strahlenwerfung] see Radiatio stellarum.}
\Entry{Places, mean, of the planets [Örter] see Motion, mean [Bewegung, mittlere].}
\Entry{— true, of the planets see Longitude, true [Länge, wahre].}
\Entry{Ortus iuxta circulum aequinoctialem see Ascensiones rectae.}
\par\begin{center}\textbf{P.}\end{center}
\Entry{Panchasiddhāntikā (= the book of the five Siddhāntas) 46, 64, 68, 75, 90, 235.}
\Entry{Parallax of the moon in longitude [Parallaxe] 28, 29, 91–93, 95, 96, 106, T. 77, 77$^{a}$.}
\Entry{— — in latitude 28, 29, 93, 94, 96, 106, T. 77, 77$^{a}$.}
\Entry{Pechüle, C. F. VI, 233–235.}
\Entry{Perigeum see Submissio.}
\Entry{Petrus de Dacia XXIII (N. 1).\footnoteT{The print has «(N. 2)»; Petrus de Dacia is cited in note 1 of p. XXIII, Ocreatus in note 2.}}
\Entry{Philippus Aridaeus XVI, T. 1.}
\Entry{Planets [Planeten] see Motion [Bewegung], Latitude [Breite], Longitude [Länge], Sublimatio, Examinatio, Eccentricity [Exzentrizität], Stations [Stationen].}
\Entry{Planisphaerium, of Ptolemy IX, XIII.}
\Entry{Plato of Tivoli 57.}
\Entry{Altitude of the pole [Polhöhe] see Latitude, geographical [Breite, geogr.].}
\Entry{Precession of the vernal point [Praezession] 42, 46, 50, 52, 54.}
\Entry{Proiectio lucis stellarum see Radiatio stellarum.}
\Entry{— radiorum see the same.}
\Entry{Prosneusis 45.}
\Entry{Prosthaphaeresis of longitude 49, 233.}
\Entry{Ptolemy [Ptolemäus] IX, XIII, XXIV, 2, 16, 18, 33, 35, 38, 44–47, 49, 50, 57–60, 63, 66, 70, 74, 80, 81, 86, 87, 90, 96, 99, 104–106, 233, 234.}
\Entry{Puncta (= Minuta) 36.}
\par\begin{center}\textbf{Q.}\end{center}
\Entry{Square aspect (tetragonum) [Quadratschein] 30, 31, 98, 99, 101, T. 91–114.}
\Entry{Quadripartitum 96, 99.}
\Entry{Quantitas incidentiae see Casus.}
\par\begin{center}\textbf{R.}\end{center}
\Entry{Radiatio stellarum (exagonalis, tetragonalis, trigonalis) 30, 31, 98–102, 106, T. 91–114.}
\Entry{— dextra 31, 101.}
\Entry{— sinistra 31, 101.}
\Entry{Radius of the earth's shadow [Radius des Erdschattens] see Diameter [Durchmesser].}
\Entry{Radix XVI, 7, 33, 36, 41, 81, T. 3, 4, 6, 9, etc.}
\Entry{al-Raḳḳa 41, 74, 75, 81, 96.}
\Entry{Direct motion of the planets [Rechtläufigkeit], see Progression [Vorwärtsschreiten].}
\Entry{Right ascension [Rectascension] see Ascensiones rectae.}
\Entry{Rule of the four quantities [Regel der vier Grössen] 61.}
\Entry{Rulers [Regenten] see Judices.}
\Entry{Regiomontanus, Joh. 99, 101.}
\Entry{Regnum see Exaltation [Erhöhung].}
\Entry{Respiratio 26, 28, 87, 89.}
\Entry{Revolutiones annorum see Conversion [Umwandlung].}
\Entry{Robert of Chester [Robert von Chester] VIII, XI–XIII, XVIII, 1, 69, 106.}
\Entry{\Pg{254}Robertus Ketenensis or Retinensis see Robert of Chester.}
\Entry{Rodbertus Cestrensis see Robert of Chester.}
\Entry{Roger of Hereford [Roger von Hereford] XII, 106.}
\Entry{Rose, astrological [Rose, astrologische] 102–105.}
\Entry{Rota = circumferentia, peripheria, sphaera 4, 7; see also Diameter [Durchmesser].}
\Entry{— draconis see Diameter of the earth's shadow [Durchmesser].}
\Entry{Rudolf of Bruges [Rudolf von Brügge] XIII, XVIII.}
\Entry{Retrogradation [Rückläufigkeit] see the following.}
\Entry{Retrogression, of the planets [Rückwärtsschreiten] 12, 55, 56; its time 12, 13, 56–58.}
\Entry{Rumor (?) 31.}
\par\begin{center}\textbf{S.}\end{center}
\Entry{Sachau, E. 37, 38, 49, 66, 79, 105, 235.}
\Entry{Savasorda X.}
\Entry{Detriment [Schaden] see Alienatio.}
\Entry{Leap years [Schaltjahre] of the Arabs 34, 37; of the Julian reckoning see Annus bissextus.}
\Entry{Shadow [Schatten] see Umbra.}
\Entry{Schiaparelli, Giov. XXIV, 66, 80–83.}
\Entry{Obliquity of the ecliptic [Schiefe] see Ecliptic [Ekliptik].}
\Entry{— of the orbit of the moon see Orbit of the moon [Mondbahn].}
\Entry{Schram, R. 83, 85 (N. 2).}
\Entry{Semita numerorum (column of the numbers) 69, 78, T. 21, 22 etc., 58, 58$^{a}$, 61–66.}
\Entry{Sened b. ʿAlī 32.}
\Entry{Servitus see Dejection [Tiefstand].}
\Entry{Sexagesimal parts [Sexagesimalteile] see Multiplication table [Multiplikationstafel].}
\Entry{Sextile aspect (exagonum) [Sextilschein] 30, 31, 98–101, T. 91–114.}
\Entry{Seybold, C. F. XXI.}
\Entry{Siddhāntas VIII, 32, 59, 65, 70, 105, 106; Ārya-S. 65, 66; Brahma-S. 32, 65, 235; Śiromani-S. 70; Sūrya-S. 32, 33, 42, 43, 45, 46, 51, 52, 58–60, 62, 64, 65, 68, 74, 75, 78–80, 86, 90, 92.}
\Entry{Signa see Notae.}
\Entry{Signum, sign of the zodiac, 14, 15, 41, 71; their risings in the right sphere 18–20, 72, 73, T. 59, 59$^{a}$; see also Ascensiones rectae; in the oblique sphere 20, 21, 73–76; see also Ascensiones obliquae.}
\Entry{Sindhind 32, 33, 43, 46, 49, 70, 103, 235.}
\Entry{Sinus, finding it, 17, 18, 68–71, T. 58 and 58$^{a}$.}
\Entry{— of the complement (cosinus) 73, 77.}
\Entry{— diminutus (versus) 17, 69, 71.}
\Entry{— planus (rectus) 17, 21, 22, 69, 71, 77.}
\Entry{Sun [Sonne] see Motion [Bewegung], Longitude [Länge], Sublimatio, Declination [Deklination], Diameter [Durchmesser], Examinatio, Eclipse [Verfinsterung].}
\Entry{Spatium longinquum (remotum) 26–28, 86, 87, 94, 95.}
\Entry{— propinquum 26–28, 86, 87, 94, 95.}
\Entry{Stations, of the planets [Stationen] 12, 13, 55–58, T. 27–56.}
\Entry{Steinschneider, M. XIII, 41.}
\Entry{Star year [Sternjahr] see Year, sidereal [Jahr, siderisches].}
\Entry{Standing still of the planets [Stillstand], see Stations [Stationen].}
\Entry{Hours, unequal [Stunden, ungleiche], see Temporal hours [Zeitstunden].}
\Entry{Hour angle [Stundenwinkel] 92, 93.}
\Entry{Stupor 22–24, 80.}
\Entry{Sublimatio (apogeum) of the moon 14, 15, 64, 88; of the planets 14, 15, 46–50, 52, 53, 60, 63, 64, 235; of the sun 9, 14, 33, 43, 44, 46, 47, 67, 82.}
\Entry{— examinata vel definita 10, 11, 49, 50, 53, 54, T. 27–56.}
\Entry{Submissio (perigeum) 14, 15, 47, 63, 64, 88.}
\Entry{Sullimatio see Sublimatio.}
\Entry{Syzygies [Syzygien] see Conjunction and opposition [Konjunktion und Opposition].}
\par\begin{center}\textbf{T.}\end{center}
\Entry{Tabula (vel pagina) differentiarum 21, 76.}
\Entry{— correctionis 95.}
\Entry{Tabulae proportionis 26, 27, 88, 94, T. 73–75.}
\Entry{— spuriae (al-Battānī's) XV, 40, 102.}
\Entry{Tables (astron.) [Tafeln]: the Alfonsine 35; those of Battānī see al-Battānī; the tested 32, 46, 47, 58; the Persian (of the Shāh) 32, 45, 49, 51, 52, 105, 106, 235; those composed after the Sindhind VIII, 70; Theon's 94.}
\Entry{Day arc of the ascendant [Tagebogen] 97.}
\Entry{Day hours [Tagesstunden] 20, 25, 76, 97, 98.}
\Entry{Tangent [Tangens] see Umbra.}
\Entry{Tempora = equatorial degrees 97.}
\Entry{Termini see Bounds [Grenzen].}
\Entry{Tetragonum see Square aspect [Quadratschein].}
\Entry{Thābit b. Ḳurra XI.}
\Entry{Thebit see Thābit.}
\Entry{Thibaut, G. 42, 46, 64, 68, 75, 90, 235.}
\Entry{Thoth (Egyptian deity) XXII.}
\Entry{Dejection [Tiefstand] 104, T. 116.}
\Entry{Sign of the zodiac [Tierkreiszeichen] see Signum.}
\Entry{Trepidation of the fixed stars 42.}
\Entry{Trigonalia see Triplicities [Triplizitäten].}
\Entry{Trine aspect [Trigonalschein] 30, 31, 98, 99, 101, T. 91–114.}
\Entry{Trigonum see the preceding.}
\Entry{Triplicities [Triplizitäten] 104.}
\par\begin{center}\textbf{U.}\end{center}
\Entry{Udjain see Arin.}
\Entry{Uhlemann 105.}
\Entry{\Pg{255}Ujjayinī see Arin.}
\Entry{Ulūgh Beg 78.}
\Entry{Umbra meridiana et aequinoctialis 20, 21, 74, 76.}
\Entry{— mankūs (versa) = tangent 22, 76–78.}
\Entry{— mustewia (plana) = cotangent 21, 22, 68, 76–78, 106, T. 60.}
\Entry{— terrae 26, 86–91; see also Diameter of the earth's shadow [Durchmesser].}
\Entry{Conversion (conversio, revolutio) of the years of birth (annorum nativitatum) [Umwandlung] 102, 103, 106, T. 115.}
\Entry{— of the years of the world (annorum mundi) 103.}
\par\begin{center}\textbf{V.}\end{center}
\Entry{Varāha-Mihira 46, 235.}
\Entry{Eclipse of the moon [Verfinsterung] 26–28, 86–90, T. 73–76.}
\Entry{— of the sun 28, 29, 94–96, T. 78.}
\Entry{Distribution of the 12 astrological houses [Verteilung], see Houses [Häuser].}
\Entry{Full moon [Vollmond] 84.}
\Entry{Progression of the planets [Vorwärtsschreiten] 12, 55, 56; its time 12, 13, 56–58.}
\par\begin{center}\textbf{W.}\end{center}
\Entry{Washington, meridian of, 52.}
\Entry{Angles of latitude and longitude [Winkel] 91, 92.}
\Entry{Wüstenfeld, F. XIII.}
\par\begin{center}\textbf{Y.}\end{center}
\Entry{Yaḥyā b. Abī Manṣūr 32.}
\Entry{Yaʿḳūb b. Ṭāriḳ 32, 70.}
\Entry{Yezdegird III, the last king of the Sasanians, VII, VIII, XV, 33, 35, 37, 106, T. 1, see also Aera.}
\Entry{Ysidorus see Isidore [Isidor].}
\Entry{Yojana 91.}
\par\begin{center}\textbf{Z.}\end{center}
\Entry{al-Zarḳālī 236.}
\Entry{Time, mean [Zeit, mittlere] 80.}
\Entry{— true 80.}
\Entry{Equation of time [Zeitgleichung] 25, 80, 84, 106, T. 67–68.}
\Entry{Temporal hours (unequal) [Zeitstunden] 20, 21, 25, 76, 81, 84, 97, 98; see also Night hours [Nachtstunden] and Day hours [Tagesstunden].}
\Entry{Zenith distance (of the moon) [Zenitdistanz] 92.}
\Entry{Centre [Zentrum] see Centrum.}
\Entry{Zeuthen, H. G. VII, 68, 71.}
\Entry{Zodiacus 29; see also Ecliptic [Ekliptik].}
\end{document}
